\documentclass[11pt,a4paper,oneside,openany]{book}
\usepackage[margin=23mm,top=24mm,bottom=25mm,headheight=15pt,headsep=8mm]{geometry}
\usepackage{fontspec}
\setmainfont{Latin Modern Roman}
\setsansfont{Latin Modern Sans}
\setmonofont[Scale=0.83]{DejaVu Sans Mono}
\usepackage{amsmath,amssymb,mathtools}
\usepackage{unicode-math}
\setmathfont{Latin Modern Math}
\usepackage{microtype,xcolor,graphicx,adjustbox,booktabs,tabularx,array,enumitem}
\usepackage[most]{tcolorbox}
\usepackage{titlesec,fancyhdr,needspace,xurl}
\usepackage[unicode,colorlinks=true,linkcolor=atlasgreen,urlcolor=atlasgreen,bookmarksnumbered=true,pdfborder={0 0 0}]{hyperref}
\definecolor{atlasgreen}{HTML}{183F36}
\definecolor{atlaspale}{HTML}{F4F6F1}
\definecolor{atlasgold}{HTML}{946B30}
\definecolor{atlasgray}{HTML}{5A6260}
\hypersetup{pdftitle={Polynomial Arithmetic Atlas: Complete Handbook},pdfauthor={Polynomial Arithmetic Atlas},pdfsubject={Irreducibility, Galois groups, methods, examples, and an annotated bibliography},pdfkeywords={polynomials, irreducibility, Galois theory, Newton polygons, research bibliography}}
\setlength{\parindent}{0pt}
\setlength{\parskip}{3pt}
\linespread{1.06}
\setcounter{tocdepth}{0}
\setcounter{secnumdepth}{1}
\setlist{topsep=4pt,itemsep=3pt,parsep=1pt,leftmargin=1.6em}
\frenchspacing
\emergencystretch=3em
\widowpenalty=10000\clubpenalty=10000
\newcommand{\atlasdisplay}[1]{\par\begingroup\centering\vspace{3pt}\noindent\adjustbox{max width=\linewidth}{\(\displaystyle #1\)}\par\vspace{5pt}\endgroup}
\providecommand{\lt}{<}\providecommand{\gt}{>}
\newtcolorbox{atlasblock}[1][standard]{enhanced,breakable,colback=atlaspale,colframe=atlasgreen,coltext=black,boxrule=0pt,borderline west={1pt}{0pt}{atlasgreen},arc=0pt,outer arc=0pt,left=9pt,right=9pt,top=8pt,bottom=6pt,before skip=9pt,after skip=9pt,standard/.style={},boundary/.style={colback=atlasgold!6,borderline west={1pt}{0pt}{atlasgold}},#1}
\newenvironment{sources}{\par\smallskip\begingroup\footnotesize\color{atlasgray}\leftskip=10pt\setlength{\parskip}{2pt}}{\par\endgroup\smallskip}
\titleformat{\chapter}[display]{\sffamily\bfseries\raggedright\color{atlasgreen}}{\fontsize{12}{15}\selectfont CHAPTER \thechapter}{8pt}{\fontsize{25}{30}\selectfont}
\titlespacing*{\chapter}{0pt}{-6pt}{20pt}
\titleformat{\section}{\sffamily\large\bfseries\color{atlasgreen}}{\thesection}{0.7em}{}
\titlespacing*{\section}{0pt}{16pt}{7pt}
\titleformat{\subsection}{\sffamily\normalsize\bfseries\color{atlasgreen}}{}{0em}{}
\pagestyle{fancy}\fancyhf{}
\fancyhead[L]{\sffamily\scriptsize\color{atlasgray}\nouppercase{\leftmark}}
\fancyhead[R]{\sffamily\small\thepage}
\fancyfoot[L]{\sffamily\scriptsize\color{atlasgray}POLYNOMIAL ARITHMETIC ATLAS}
\fancyfoot[R]{\sffamily\scriptsize\color{atlasgray}12 NOVEMBER 2026}
\renewcommand{\headrulewidth}{0.3pt}
\renewcommand{\chaptermark}[1]{\markboth{#1}{}}
\fancypagestyle{plain}{\fancyhf{}\fancyfoot[R]{\sffamily\small\thepage}\renewcommand{\headrulewidth}{0pt}}
\begin{document}
\begin{titlepage}\thispagestyle{empty}
\vspace*{14mm}
{\sffamily\large\color{atlasgold}A CONNECTED RESEARCH HANDBOOK\par}
\vspace{14mm}
{\sffamily\bfseries\fontsize{38}{43}\selectfont\color{atlasgreen}Polynomial\\Arithmetic\\Atlas\par}
\vspace{11mm}
{\fontsize{18}{24}\selectfont Irreducibility, Galois groups,\\and the methods between them\par}
\vspace{10mm}\textcolor{atlasgold}{\rule{35mm}{1.2pt}}\par\vspace{8mm}
{\large Precise statements. Proof mechanisms. Exact examples.\\Connected literature, with explicit boundaries.\par}
\vfill
{\sffamily\large 37 explanatory articles\quad /\quad 232 annotated sources\par}
\vspace{4mm}
{\color{atlasgray}Complete offline edition / Literature checked through 12 November 2026\par}
\vspace{7mm}
{\small This handbook records the current atlas. It includes the complete article text, section-level citations, research questions, reading routes, and annotated reference catalogue. It is a curated collection, not a claim to contain every publication.\par}
\end{titlepage}
\frontmatter
\tableofcontents
\chapter{How to use this handbook}
The organizing question is not only whether a polynomial is irreducible or which group it has, but \emph{why} a particular method answers that question. Start with the ground field and the exact object, then separate factorization, transitivity, cycle evidence, and identification of the complete group.

This edition contains all 37 articles and all 232 annotated reference records in the current website source. Mathematics, qualifications, and source-specific research status have been retained. Section citations carry a stable reference key, precise source location where available, and a clickable source link; the annotated catalogue prints the complete URLs with publication details, DOI links, and verification links. Internal links and the PDF outline support offline navigation.

\section*{Read the status together with the statement}
\begin{description}
\item[Classical or published result.] Read the field, parameter range, exceptional cases, and quantifiers alongside the theorem.
\item[Preprint boundary.] A preprint is labelled as such. An abstract-level verification is not silently presented as a line-by-line proof audit.
\item[Open in a cited source.] This records a question's status in the named source; it is not an unsupported assertion that all subsequent literature has been exhausted.
\item[Research direction.] This is a useful question suggested by the connections, not necessarily a new conjecture or a verified unsolved problem.
\item[Exact example.] A computation supports the specified polynomial. A finite sample cannot establish a uniform theorem or an asymptotic statement.
\end{description}

\section*{Scope and editions}
The literature-review date is 12 November 2026. This handbook is a new, complete snapshot of the expanded atlas. The earlier 188-source, 37-page bibliography remains a separate fixed edition; it does not contain the full current explanations or all subsequent reference additions. Linked papers and books remain with their authors and publishers. The catalogue stores metadata, annotations, and links rather than reproducing their full text.

\section*{A reliable order of work}
\begin{enumerate}
\item Specify the field, degree, separability, coefficient assumptions, and any factor removed by normalization.
\item Establish irreducibility or list the possible factor degrees.
\item Obtain permutations or structural upper bounds without confusing local evidence with a complete group determination.
\item Use a stated group theorem, resolvent, or certificate to finish; use the discriminant for parity only within its proper scope.
\item Save the precise theorem, source location, prime, identity, or polygon that makes the conclusion reproducible.
\end{enumerate}
\chapter{Reading routes and reference books}
\section*{From a polynomial to its Galois group}
Prove irreducibility, obtain permutation information, and finish a complete group determination.\par
\textbf{Background:} A first course in field theory\par
\begin{enumerate}
\item \textbf{\hyperlink{article:foundations}{Fields, factors, and the action on roots}}. Fix the ground field, separate irreducibility from separability, and translate factorization into a permutation action.\par
\item \textbf{\hyperlink{article:modular-irreducibility}{Irreducibility by reduction modulo primes}}. Use one finite-field factorization, or intersect the possible degrees supplied by several primes. Understand why the method can fail for an irreducible polynomial.\par
\item \textbf{\hyperlink{article:frobenius}{Frobenius: from factor degrees to permutations}}. A factorization over a finite field reveals a cycle type in a global Galois group. Combine these local witnesses with a separate group-theoretic argument.\par
\item \textbf{\hyperlink{article:permutation-groups}{The group-theoretic step: transitivity, alternating groups, and symmetric groups}}. Local arithmetic supplies permutations. Block systems, transpositions, and long cycles explain when those permutations force a large group.\par
\item \textbf{\hyperlink{article:discriminants}{Discriminants, parity, and ramification}}. The discriminant detects collisions of roots and the sign character of a Galois group. Its square class is powerful, but it does not measure the size of the group.\par
{\footnotesize References: \hyperlink{ref:Milne2022}{Milne2022}, \hyperlink{ref:Cox2012}{Cox2012}, \hyperlink{ref:ConradIrredQ}{ConradIrredQ}, \hyperlink{ref:ConradModP}{ConradModP}, \hyperlink{ref:ConradSnAn}{ConradSnAn}}\par
\item \textbf{Checkpoint}. Give a proof of \(\operatorname{Gal}(x^5-x-1/\mathbb Q)=S_5\) and state exactly what each prime contributes.\par
\end{enumerate}
\section*{From valuations to large groups}
Learn how coefficient divisibility restricts factor degrees and produces long prime cycles.\par
\textbf{Background:} Readers comfortable with valuations and finite permutation groups\par
\begin{enumerate}
\item \textbf{\hyperlink{article:eisenstein-dumas}{Schönemann-Eisenstein and one-edge Newton polygons}}. A divisibility pattern can force every factor to have the full degree. Translation and Newton slopes extend the elementary criterion.\par
\item \textbf{\hyperlink{article:newton-polygons}{Newton polygons, factor degrees, and large cycles}}. Read p-adic root valuations from coefficients, intersect the permitted factor degrees, and use suitable slope denominators to force a large Galois group.\par
\item \textbf{\hyperlink{article:permutation-groups}{The group-theoretic step: transitivity, alternating groups, and symmetric groups}}. Local arithmetic supplies permutations. Block systems, transpositions, and long cycles explain when those permutations force a large group.\par
\item \textbf{\hyperlink{article:taylor}{Taylor truncations: two exact benchmarks}}. Exponential truncations have alternating or symmetric groups, while geometric truncations are cyclotomic. Their contrast identifies which arithmetic features a general truncation theory must track.\par
\item \textbf{\hyperlink{article:laguerre}{Generalized Laguerre polynomials: two parameter regimes}}. Factorial products connect local slopes to large Galois groups. Fixing the parameter and increasing the degree is a different problem from fixing the degree and varying the parameter.\par
{\footnotesize References: \hyperlink{ref:Serre1979}{Serre1979}, \hyperlink{ref:dumas1906}{dumas1906}, \hyperlink{ref:Coleman1987}{Coleman1987}, \hyperlink{ref:Hajir2005}{Hajir2005}, \hyperlink{ref:FilasetaLam2002}{FilasetaLam2002}}\par
\item \textbf{Checkpoint}. Separate the irreducibility input, the Newton-polygon input, Jordan's theorem, and the discriminant step.\par
\end{enumerate}
\section*{Study a structured family}
Identify forced factors and symmetries before choosing the expected Galois group.\par
\textbf{Background:} Researchers moving from examples to a family theorem\par
\begin{enumerate}
\item \textbf{\hyperlink{article:reciprocity}{Reciprocal symmetry and the group on root pairs}}. Pass from a reciprocal polynomial to a polynomial of half the degree. The remaining quadratic extensions explain both irreducibility and the correct ambient Galois group.\par
\item \textbf{\hyperlink{article:sparse-gaps}{Sparse polynomials, factor reversal, and large gaps}}. Separate persistent reciprocal factors from the nonreciprocal part, then use coefficient sparsity to constrain possible factorizations.\par
\item \textbf{\hyperlink{article:trinomials}{Trinomials: uniform families and explicit classifications}}. The polynomials x\ensuremath{{}^n}\ensuremath{-}x\ensuremath{-}1 realize S\ensuremath{{}_n} in every degree. Prime reductions give short certificates, while inertia explains the uniform theorem.\par
\item \textbf{\hyperlink{article:eulerian}{Eulerian polynomials: descents, shifts, and inverse pairs}}. Descent statistics lead to reciprocal polynomials. A prime-index Eisenstein argument and an exact quartic example show how combinatorics, local arithmetic, and Galois theory fit together.\par
\item \textbf{\hyperlink{article:fekete}{Fekete polynomials: character sums and reciprocal arithmetic}}. Quadratic characters impose exact rational factors and reciprocal symmetry. Trace polynomials connect their irreducibility and Galois groups to special values and square classes.\par
\item \textbf{\hyperlink{article:orthogonal-families}{Hermite, Bessel, and Legendre: shared tools, different answers}}. Classical polynomial families share differential and recurrence identities, but their arithmetic factorization problems differ. Parity reduction and coefficient products expose the correct questions.\par
{\footnotesize References: \hyperlink{ref:Schinzel2000}{Schinzel2000}, \hyperlink{ref:Prasolov2004}{Prasolov2004}, \hyperlink{ref:sawinshustermanstoll2020}{sawinshustermanstoll2020}, \hyperlink{ref:MinacNguyenTan2023}{MinacNguyenTan2023}, \hyperlink{ref:CullinanHajir2014}{CullinanHajir2014}}\par
\item \textbf{Checkpoint}. Write the normalization, the natural ambient permutation group, a proved benchmark, and a source-scoped unresolved question.\par
\end{enumerate}
\section*{Read probabilistic results without changing the model}
Keep the asymptotic parameter, coefficient law, conditioning, and analytic hypotheses visible.\par
\textbf{Background:} Researchers reading counting and random-polynomial papers\par
\begin{enumerate}
\item \textbf{\hyperlink{article:height-counting}{Counting exceptional Galois groups by height}}. At fixed degree, almost all integer polynomials have full symmetric group. The sharp exceptional order, and its change under reciprocity, reveal which arithmetic obstructions dominate.\par
\item \textbf{\hyperlink{article:random-models}{Random polynomials: models and proof strategies}}. The coefficient law and the limiting parameter determine the problem. A guide to irreducibility, alternating-group containment, and the additional step needed for a full symmetric group.\par
\item \textbf{\hyperlink{article:littlewood}{Littlewood polynomials: signs, parity, and special degrees}}. Independent signs create strong congruence constraints. Exact examples show both their usefulness and why the all-degree probabilistic problem needs more than reduction modulo two.\par
\item \textbf{\hyperlink{article:random-reciprocal}{Random reciprocal polynomials of growing degree}}. Independent free coefficients generate paired roots and a signed-permutation group. The trace polynomial connects this structured model to ordinary factorization problems.\par
\item \textbf{\hyperlink{article:random-matrices}{Characteristic polynomials of random matrices}}. Matrix entries can be independent while characteristic-polynomial coefficients are strongly dependent. Local rank estimates replace the independent-coefficient tools.\par
{\footnotesize References: \hyperlink{ref:Bhargava2025}{Bhargava2025}, \hyperlink{ref:BreuillardVarju2019}{BreuillardVarju2019}, \hyperlink{ref:BarySorokerKoukoulopoulosKozma2023}{BarySorokerKoukoulopoulosKozma2023}, \hyperlink{ref:BarySorokerHokkenKozmaPoonen2025}{BarySorokerHokkenKozmaPoonen2025}, \hyperlink{ref:Eberhard2022}{Eberhard2022}}\par
\item \textbf{Checkpoint}. Explain why a fixed-degree height theorem does not imply a growing-degree fair-sign theorem.\par
\end{enumerate}
\section*{From generic families to arithmetic specializations}
Connect generic monodromy, Hilbert irreducibility, and specialization while tracking the quantifiers.\par
\textbf{Background:} Graduate readers interested in arithmetic geometry and inverse Galois theory\par
\begin{enumerate}
\item \textbf{\hyperlink{article:critical-points-monodromy}{Critical points and monodromy: a geometric route to the symmetric group}}. Derivatives and critical values determine the generic group; the cubic example separates that group from four different rational fibers.\par
{\footnotesize References: \hyperlink{ref:Serre2007Topics}{Serre2007Topics}}\par
\item \textbf{\hyperlink{article:specialization}{Specialization and Hilbert irreducibility}}. Pass from a polynomial over a rational function field to arithmetic fibers, and understand how the exceptional parameters are organized.\par
\item \textbf{\hyperlink{article:inverse-galois}{Inverse Galois theory and generic polynomials}}. Construct a polynomial with a prescribed group, then distinguish one realization from regular families and universal parametrizations.\par
\item \textbf{\hyperlink{article:composition-iteration}{Composition, iteration, and arboreal groups}}. Use extension towers and rooted trees to study repeated polynomial composition. Separate stability from maximality of the Galois action.\par
\item \textbf{\hyperlink{article:laguerre}{Generalized Laguerre polynomials: two parameter regimes}}. Factorial products connect local slopes to large Galois groups. Fixing the parameter and increasing the degree is a different problem from fixing the degree and varying the parameter.\par
{\footnotesize References: \hyperlink{ref:Serre2007Topics}{Serre2007Topics}, \hyperlink{ref:FriedJarden2023}{FriedJarden2023}, \hyperlink{ref:HajirWong2006}{HajirWong2006}, \hyperlink{ref:JensenLedetYui2002}{JensenLedetYui2002}, \hyperlink{ref:Jones2013Arboreal}{Jones2013Arboreal}}\par
\item \textbf{Checkpoint}. Distinguish infinitely many good specializations, all but finitely many, density one, and an infinite tower.\par
\end{enumerate}
\section*{Turn a computation into a certificate}
Understand the algebra behind factorization and group algorithms and record what has been proved.\par
\textbf{Background:} Readers using computer algebra in proofs\par
\begin{enumerate}
\item \textbf{\hyperlink{article:finite-field-irreducibility}{Irreducibility over finite fields}}. Turn Frobenius powers into exact irreducibility certificates, count irreducibles, and distinguish irreducible from primitive polynomials.\par
\item \textbf{\hyperlink{article:factorization-algorithms}{Exact factorization: finite fields, Hensel lifting, and lattices}}. Factorization algorithms turn structural ideas into certificates: Frobenius and linear algebra over finite fields, then p-adic lifting and lattice reduction over the rationals.\par
\item \textbf{\hyperlink{article:resolvents}{Resolvents: turn group containment into factorization}}. A resolvent records how a group acts on selected expressions in the roots. Rational factors can detect a hidden partition or force descent to a smaller subgroup.\par
\item \textbf{\hyperlink{article:galois-computation}{Computing a Galois group with a proof certificate}}. A reliable computation combines exact arithmetic, subgroup constraints, and certification. Save enough evidence that the conclusion survives a change of software.\par
{\footnotesize References: \hyperlink{ref:LidlNiederreiter1997}{LidlNiederreiter1997}, \hyperlink{ref:GathenGerhard2013}{GathenGerhard2013}, \hyperlink{ref:LLL1982}{LLL1982}, \hyperlink{ref:FiekerKlueners2014}{FiekerKlueners2014}, \hyperlink{ref:MagmaGaloisDocumentation}{MagmaGaloisDocumentation}, \hyperlink{ref:PARIDocumentation}{PARIDocumentation}}\par
\item \textbf{Checkpoint}. Record the exact input, software version, field, factorization or resolvent witness, and certification step.\par
\end{enumerate}
\section*{The book shelf}
\needspace{5\baselineskip}\textbf{\hyperlink{ref:Milne2022}{Fields and Galois Theory}} / Milne, James S.\par
Field extensions, separability, and the fundamental theorem of Galois theory; Chapter 4 gives explicit Galois-group computations. The frozen course notes v5.10 and the 2022 book have different pagination. Cite the version actually used.\par
{\small Free entry point}\par\medskip
\needspace{5\baselineskip}\textbf{\hyperlink{ref:Cox2012}{Galois Theory}} / Cox, David A.\par
Work through explicit polynomials, resolvents, radicals, and permutation actions alongside the general theory.\par
{\small Examples and classical constructions}\par\medskip
\needspace{5\baselineskip}\textbf{\hyperlink{ref:Prasolov2004}{Polynomials}} / Prasolov, Victor V.\par
Read the chapters on irreducible polynomials, special forms, root properties, and Galois theory together. The linked electronic edition carries a 2009 release date; the English book is ©2004.\par
{\small Connections across polynomial methods}\par\medskip
\needspace{5\baselineskip}\textbf{\hyperlink{ref:Schinzel2000}{Polynomials with Special Regard to Reducibility}} / Schinzel, Andrzej\par
Consult for structured, sparse, and multivariable polynomials and for the historical trail behind criteria. The publisher states that results valid only over finite fields, local fields, or Q lie outside this monograph's scope. Pair it with the field-specific references in this guide.\par
{\small Central research reference on reducibility}\par\medskip
\needspace{5\baselineskip}\textbf{\hyperlink{ref:Serre1979}{Local Fields}} / Serre, Jean-Pierre\par
Learn valuations, ramification, and local extensions before using Newton polygons to produce Galois elements.\par
{\small Local arithmetic behind the methods}\par\medskip
\needspace{5\baselineskip}\textbf{\hyperlink{ref:DixonMortimer1996}{Permutation Groups}} / Dixon, John D. and Mortimer, Brian\par
Use transitivity, blocks, primitivity, and cycle criteria to turn arithmetic evidence into group determination.\par
{\small The permutation-group step}\par\medskip
\needspace{5\baselineskip}\textbf{\hyperlink{ref:LidlNiederreiter1997}{Finite Fields}} / Lidl, Rudolf and Niederreiter, Harald\par
Irreducible and primitive polynomials, extension fields, counting, constructions, and factorization.\par
{\small Finite-field foundation}\par\medskip
\needspace{5\baselineskip}\textbf{\hyperlink{ref:GathenGerhard2013}{Modern Computer Algebra}} / von zur Gathen, Joachim and Gerhard, Jürgen\par
Polynomial arithmetic, finite-field algorithms, Hensel lifting, lattice reduction, and factorization over the rationals.\par
{\small Algorithmic proofs and complexity}\par\medskip
\needspace{5\baselineskip}\textbf{\hyperlink{ref:Cohen1993}{A Course in Computational Algebraic Number Theory}} / Cohen, Henri\par
Connect defining polynomials to number fields, integral bases, discriminants, and explicit arithmetic.\par
{\small Practical computational number theory}\par\medskip
\needspace{5\baselineskip}\textbf{\hyperlink{ref:Serre2007Topics}{Topics in Galois Theory}} / Serre, Jean-Pierre\par
Hilbert irreducibility, thin sets, regular realizations, and sieve arguments. The electronic edition has copyright 2007; the physical second edition is also catalogued as 2008.\par
{\small Compact route to specialization and inverse Galois theory}\par\medskip
\needspace{5\baselineskip}\textbf{\hyperlink{ref:FriedJarden2023}{Field Arithmetic}} / Fried, Michael D. and Jarden, Moshe\par
Hilbertian fields, PAC fields, absolute Galois groups, and specialization in the fourth edition.\par
{\small Deep field arithmetic}\par\medskip
\needspace{5\baselineskip}\textbf{\hyperlink{ref:MalleMatzat2018}{Inverse Galois Theory}} / Malle, Gunter and Matzat, B. Heinrich\par
Regular realizations, rigidity, embedding problems, and explicit construction methods.\par
{\small Inverse Galois research reference}\par\medskip
\chapter{Choose a method}
\section*{Is a given polynomial irreducible over Q?}
\textbf{Available information.} A useful factorization modulo a small prime\par
\textbf{Method.} \hyperlink{article:modular-irreducibility}{Degree-preserving reduction and intersection of possible factor degrees}.\par
\textbf{Conclusion.} An exact irreducibility certificate when the remaining proper degrees are eliminated\par
\textbf{Boundary.} An irreducible polynomial need not have an irreducible reduction at any prime.\par\medskip
\section*{Can coefficient divisibility prove irreducibility?}
\textbf{Available information.} One prime divides all lower coefficients, or the valuations have a simple pattern\par
\textbf{Method.} \hyperlink{article:eisenstein-dumas}{Schönemann-Eisenstein and Dumas/Newton polygons}.\par
\textbf{Conclusion.} Irreducibility over a local field, hence over Q, or restrictions on factor degrees\par
\textbf{Boundary.} A criterion is sufficient; failure of its hypotheses is not evidence of reducibility.\par\medskip
\section*{Can several local tests work together?}
\textbf{Available information.} Different primes exclude different factor degrees\par
\textbf{Method.} \hyperlink{article:newton-polygons}{Intersect Newton-polygon degree constraints}.\par
\textbf{Conclusion.} A global irreducibility proof even when each individual prime permits factors\par
\textbf{Boundary.} The same polygons may still leave common possible degrees.\par\medskip
\section*{Can root sizes rule out factors?}
\textbf{Available information.} A monic integer polynomial satisfying Perron's coefficient inequality, or a root bound at a prime-valued argument\par
\textbf{Method.} \hyperlink{article:archimedean-criteria}{Perron's criterion and archimedean factor bounds}.\par
\textbf{Conclusion.} Irreducibility from root location and integral constant terms or factor evaluations\par
\textbf{Boundary.} Perron's test needs monicity and a nonzero constant term; the displayed bounds are strict.\par\medskip
\section*{Can integer evaluations rule out every factor?}
\textbf{Available information.} Many distinct integer arguments with small nonzero values, or at least 2n+1 unit or prime values for a primitive degree-n polynomial\par
\textbf{Method.} \hyperlink{article:integer-values}{Schur unit-value arguments, Pólya interpolation, and prime-value counting}.\par
\textbf{Conclusion.} An exact irreducibility certificate from restrictions on the values of hypothetical integral factors\par
\textbf{Boundary.} Use the stated degree-dependent threshold and distinct arguments; one prime value alone is insufficient.\par\medskip
\section*{Is a composition or binomial irreducible?}
\textbf{Available information.} An irreducible outer polynomial F, a nonconstant inner polynomial g, or a binomial in characteristic zero\par
\textbf{Method.} \hyperlink{article:capelli-composition}{Capelli root-field reduction and the exact binomial power criterion}.\par
\textbf{Conclusion.} A lower-degree irreducibility question over a root field; for binomials, necessary and sufficient power exclusions\par
\textbf{Boundary.} Irreducibility of both inputs is insufficient. Binomials with exponent divisible by 4 require the exceptional negative-four-times-a-fourth-power exclusion.\par\medskip
\section*{What happens when a sparse polynomial has a large gap?}
\textbf{Available information.} A family x\textasciicircum{}N c(x\textasciicircum{}\{-1\})+d(x) with fixed end pieces\par
\textbf{Method.} \hyperlink{article:sparse-gaps}{Nonreciprocal-part theorems and Capelli exclusions}.\par
\textbf{Conclusion.} Irreducibility of a specified nonreciprocal part for sufficiently large N\par
\textbf{Boundary.} Reciprocal factors and Capellian exceptions cannot be discarded silently.\par\medskip
\section*{Which permutations occur in the Galois group?}
\textbf{Available information.} Squarefree factorization at a prime of good reduction\par
\textbf{Method.} \hyperlink{article:frobenius}{Frobenius cycle types}.\par
\textbf{Conclusion.} An element of the specified cycle type\par
\textbf{Boundary.} Bad primes and unseen cycle types require separate treatment.\par\medskip
\section*{When does cycle evidence force A\_\allowbreak{}n or S\_\allowbreak{}n?}
\textbf{Available information.} Transitivity, primitivity, and an appropriate prime cycle or transposition\par
\textbf{Method.} \hyperlink{article:permutation-groups}{Permutation-group recognition}.\par
\textbf{Conclusion.} A proved group containment or equality\par
\textbf{Boundary.} A transitive subgroup containing a transposition need not be \(S_n\).\par\medskip
\section*{How can A\_\allowbreak{}n be distinguished from S\_\allowbreak{}n?}
\textbf{Available information.} Containment of \(A_n\) has already been proved\par
\textbf{Method.} \hyperlink{article:discriminants}{Discriminant square class or an odd permutation}.\par
\textbf{Conclusion.} The remaining parity distinction\par
\textbf{Boundary.} A nonsquare discriminant alone does not prove \(S_n\).\par\medskip
\section*{Does a root symmetry force a smaller ambient group?}
\textbf{Available information.} Reciprocal, even, or compositional structure\par
\textbf{Method.} \hyperlink{article:reciprocity}{Invariant blocks and quotient variables}.\par
\textbf{Conclusion.} A wreath-product or other structural upper bound\par
\textbf{Boundary.} A generic \(S_n\) expectation can be incompatible with the root symmetry.\par\medskip
\section*{How can remaining candidate groups be separated?}
\textbf{Available information.} Several subgroups fit the observed Frobenius types\par
\textbf{Method.} \hyperlink{article:resolvents}{Resolvents, relative invariants, and subgroup descent}.\par
\textbf{Conclusion.} Additional exact constraints or a certified group\par
\textbf{Boundary.} Numerical approximation and a plausible candidate are not by themselves a proof.\par\medskip
\section*{Do generic results survive parameter specialization?}
\textbf{Available information.} An irreducible family over Q(t), or over K(t) with K Hilbertian\par
\textbf{Method.} \hyperlink{article:specialization}{Hilbert irreducibility and specialized geometric criteria}.\par
\textbf{Conclusion.} Many good specializations under the stated hypotheses\par
\textbf{Boundary.} Infinitely many, density one, and all but finitely many are different conclusions.\par\medskip
\section*{Are all iterates irreducible?}
\textbf{Available information.} A polynomial f and the tower of its iterates\par
\textbf{Method.} \hyperlink{article:composition-iteration}{Capelli-type steps, critical orbits, and arboreal groups}.\par
\textbf{Conclusion.} Finite-level or all-level conclusions only under the relevant hypotheses\par
\textbf{Boundary.} Specializing each finite level does not select one specialization valid at every level.\par\medskip
\section*{Can an irreducible polynomial over F\_\allowbreak{}q be constructed or certified?}
\textbf{Available information.} Specified degree and possibly prescribed coefficients\par
\textbf{Method.} \hyperlink{article:finite-field-irreducibility}{Frobenius congruences, gcd tests, and finite-field constructions}.\par
\textbf{Conclusion.} A finite-field irreducibility certificate or an existence theorem\par
\textbf{Boundary.} Primitive is stronger than irreducible; an irreducible degree-n polynomial has cyclic splitting-field group.\par\medskip
\section*{Is a multivariate polynomial absolutely irreducible?}
\textbf{Available information.} Sparse exponent support or an effective specialization problem\par
\textbf{Method.} \hyperlink{article:multivariate-irreducibility}{Newton polytopes and absolute-irreducibility criteria}.\par
\textbf{Conclusion.} Irreducibility after extending the coefficient field to an algebraic closure\par
\textbf{Boundary.} Exclude monomial factors; this differs from ordinary irreducibility over the ground field.\par\medskip
\section*{What does a random-polynomial theorem actually apply to?}
\textbf{Available information.} A coefficient law and a specified limiting parameter\par
\textbf{Method.} \hyperlink{article:random-models}{Local mixing, sieve methods, and factor exclusion}.\par
\textbf{Conclusion.} A probability estimate for the exact model in the theorem\par
\textbf{Boundary.} Track conditioning, \(A_n\) versus \(S_n\), and any RH or GRH hypothesis.\par\medskip
\section*{What is the generic Galois group of F(X)\ensuremath{-}T?}
\textbf{Available information.} Simple critical points with pairwise distinct critical values in characteristic zero\par
\textbf{Method.} \hyperlink{article:critical-points-monodromy}{Morse criterion: finite branch inertia generates transpositions}.\par
\textbf{Conclusion.} Geometric and arithmetic group \(S_n\); a regular splitting field\par
\textbf{Boundary.} A full generic group does not determine every rational fiber. Check specialization separately.\par\medskip
\mainmatter
\part{Foundations}
\chapter{Fields, factors, and the action on roots}\label{article:foundations}\hypertarget{article:foundations}{}\color{black}
{\sffamily\small\color{atlasgold}foundation / Start here / general fields / Classical theory}\par\medskip
Fix the ground field, separate irreducibility from separability, and translate factorization into a permutation action.\par\medskip
\textbf{Questions answered}\begin{itemize}
\item What does irreducibility depend on?
\item How do irreducible factors appear in a Galois group?
\item Why does proving irreducibility not determine the group?
\end{itemize}
\needspace{8\baselineskip}\section{Start with the field}
\color{black} A nonconstant polynomial \(f\in K[x]\) is irreducible over a field \(K\) if a factorization \(f=gh\) in \(K[x]\) forces one factor to be constant. The field is part of the question: \(x^2-2\) is irreducible over \(\mathbb Q\), but splits over \(\mathbb Q(\sqrt2)\). Over an algebraic closure, every polynomial of degree greater than one factors. Thus absolute irreducibility is a substantive higher-degree question for polynomials in several variables, not for a single-variable polynomial.\par

\color{black} For \(f\in\mathbb Z[x]\), first remove the content, the greatest common divisor of its coefficients. Gauss's lemma says that a primitive polynomial has a factorization of positive degrees over \(\mathbb Q\) exactly when it has one over \(\mathbb Z\). This is why congruences, divisibility, and integral coefficient bounds can answer a rational question. A constant factor such as the \(2\) in \(2x+2\) is a unit over \(\mathbb Q\), but is not a unit over \(\mathbb Z\).\par
\begin{sources}
\textbf{Sources for this section}\par\smallskip
\noindent\hyperlink{ref:Lang2002}{[Lang2002]}. Polynomial rings, content, and irreducibility over fields and unique factorization domains. \href{https://link.springer.com/book/10.1007/978-1-4613-0041-0}{Source}.\par
\noindent\hyperlink{ref:ConradIrredQ}{[ConradIrredQ]} / §2. A proof of Gauss's lemma and its use in rational irreducibility tests. \href{https://kconrad.math.uconn.edu/blurbs/ringtheory/irredtestsoverQ.pdf}{Source}.\par
\end{sources}\color{black}
\needspace{8\baselineskip}\section{Keep repeated roots separate}
\color{black} Here \emph{separable} means that all roots in an algebraic closure are distinct. For nonzero \(f\), this is equivalent to \(\gcd(f,f')=1\). Every irreducible polynomial over a field of characteristic zero is separable. In positive characteristic, the qualifications matter: \(x^p-t\in\mathbb F_p(t)[x]\) is irreducible, yet has only one distinct root in its splitting field.\par

\color{black} Before using a permutation group on \(n\) roots, verify separability. Before studying an intended polynomial family, remove known factors only when the normalization is stated explicitly. The group of the original polynomial and the group acting on a selected factor need not be the same permutation representation.\par
\begin{sources}
\textbf{Sources for this section}\par\smallskip
\noindent\hyperlink{ref:Milne2022}{[Milne2022]} / Chapter 2, Propositions 2.21 and 2.24. Separable polynomials, the derivative criterion, and perfect fields; numbering refers to the free version 5.10. \href{https://www.jmilne.org/math/Books/index.html}{Author's book page}.\par
\end{sources}\color{black}
\needspace{8\baselineskip}\section{The bridge: factors are orbits}

\begin{atlasblock}[standard]
\textbf{\sffamily Theorem  /  factor-orbit correspondence}\par\smallskip
\color{black} Let \(f\in K[x]\) be separable of degree \(n\), let \(L\) be its splitting field, and let \(G=\operatorname{Gal}(L/K)\). The action of \(G\) on the roots is faithful. Its orbits are exactly the root sets of the irreducible factors of \(f\) over \(K\). In particular, \(f\) is irreducible if and only if \(G\le S_n\) is transitive.\par
\end{atlasblock}\color{black}


\begin{atlasblock}[standard]
\textbf{\sffamily Proof mechanism}\par\smallskip
\color{black} An automorphism fixing \(K\) carries a root to another root of its minimal polynomial. Conversely, the isomorphism between the simple extensions generated by two roots of the same irreducible polynomial extends to an automorphism of the normal splitting field. These are exactly the roots in one orbit. Faithfulness follows because the roots generate \(L\).\par
\end{atlasblock}\color{black}

\color{black} This gives a useful order of work: determine possible factor degrees, establish transitivity, then obtain enough additional permutations or invariants to identify the group. A transitive action is a starting point, not a claim that every permutation occurs.\par
\begin{sources}
\textbf{Sources for this section}\par\smallskip
\noindent\hyperlink{ref:Milne2022}{[Milne2022]} / Proposition 4.5. Irreducibility is equivalent to transitivity for a separable polynomial. \href{https://www.jmilne.org/math/Books/index.html}{Author's book page}.\par
\end{sources}\color{black}
\needspace{8\baselineskip}\section{Worked example: one cubic, two extensions}

\begin{atlasblock}[standard]
\textbf{\sffamily Example  /  \(x^3-2\)}\par\smallskip
\color{black} Eisenstein at \(2\) proves irreducibility over \(\mathbb Q\). Write \(\alpha=\sqrt[3]{2}\) and let \(\zeta_3\) be a primitive cube root of unity. Its splitting field is \(L=\mathbb Q(\alpha,\zeta_3)\). The real field \(\mathbb Q(\alpha)\) has degree \(3\), and adjoining the nonreal element \(\zeta_3\) has degree \(2\). Hence \([L:\mathbb Q]=6\), so the faithful subgroup of \(S_3\) is all of \(S_3\).\par
\color{black} Over \(\mathbb Q(\zeta_3)\), the same splitting field has degree \(3\), and the group is cyclic of order \(3\). Over \(\mathbb Q(\alpha)\), the polynomial already has a linear factor and the splitting-field group has order \(2\). The polynomial expression alone does not specify its arithmetic.\par
\end{atlasblock}\color{black}
\begin{sources}
\textbf{Sources for this section}\par\smallskip
\noindent\hyperlink{ref:Cox2012}{[Cox2012]}. Field extensions, splitting fields, and explicit low-degree Galois theory; the displayed degree calculation is supplied here. \href{https://onlinelibrary.wiley.com/doi/book/10.1002/9781118218457}{Source}.\par
\end{sources}\color{black}
\needspace{8\baselineskip}\section{Four questions that should remain distinct}

\begin{itemize}
\item\leavevmode\color{black} \textbf{Factorization:} is \(f\) irreducible, or what are its irreducible factor degrees?
\item\leavevmode\color{black} \textbf{Permutation group:} which subgroup of \(S_n\) acts on the roots?
\item\leavevmode\color{black} \textbf{Field realization:} which extension of \(K\) is the splitting field, including its ramification?
\item\leavevmode\color{black} \textbf{Uniform family statement:} does the conclusion hold for every parameter, outside an explicit exceptional set, or with a limiting probability?
\end{itemize}

\color{black} The word \emph{primitive} also has several meanings: relatively prime integer coefficients; a permutation action with no nontrivial block system; or a finite-field element generating the multiplicative group. None of these definitions implies the others. Throughout the atlas, hypotheses belong beside the claim they support.\par
\begin{sources}
\textbf{Sources for this section}\par\smallskip
\noindent\hyperlink{ref:Schinzel2000}{[Schinzel2000]}. The research reference for reducibility questions over different fields and within special polynomial families. \href{https://doi.org/10.1017/CBO9780511542916}{Source}.\par
\noindent\hyperlink{ref:DixonMortimer1996}{[DixonMortimer1996]}. Transitive and primitive permutation actions. \href{https://link.springer.com/book/10.1007/978-1-4612-0731-3}{Source}.\par
\end{sources}\color{black}
\needspace{11\baselineskip}\section*{Questions and research boundaries}
\needspace{7\baselineskip}\subsection*{Does irreducibility force the full symmetric group?}
{\small\sffamily\color{atlasgold}Settled in stated scope / Separable polynomials over Q}\par\smallskip
\color{black} No. The polynomial \(x^4+1=\Phi_8(x)\) is irreducible, but its splitting field is \(\mathbb Q(\zeta_8)\) and its group is \((\mathbb Z/8\mathbb Z)^\times\cong C_2\times C_2\), acting transitively on four roots. The modular and reciprocity entries explain two consequences of this example.\par
\begin{sources}
\textbf{Sources for this section}\par\smallskip
\noindent\hyperlink{ref:Milne2022}{[Milne2022]} / Chapter 5. Cyclotomic extensions and their automorphism groups. \href{https://www.jmilne.org/math/Books/index.html}{Author's book page}.\par
\end{sources}\color{black}
\needspace{6\baselineskip}\par\medskip\begingroup\footnotesize\setlength{\parskip}{1pt}\textbf{\sffamily\color{atlasgreen}Connections}\par
\color{black}\textbf{Applies to.} \hyperlink{article:modular-irreducibility}{Irreducibility by reduction modulo primes}. Reduction modulo primes converts integer factorizations into restrictions on rational factor degrees.\par
\color{black}\textbf{Applies to.} \hyperlink{article:frobenius}{Frobenius: from factor degrees to permutations}. At good primes, the factor-orbit correspondence becomes a statement about cycle lengths.\par
\color{black}\textbf{Applies to.} \hyperlink{article:permutation-groups}{The group-theoretic step: transitivity, alternating groups, and symmetric groups}. Transitivity must be strengthened before a large symmetric or alternating group can be forced.\par
\color{black}\textbf{Contrasts with.} \hyperlink{article:multivariate-irreducibility}{Absolute irreducibility and Newton polytopes}. Absolute irreducibility in several variables survives extension of the ground field and is a different problem.\par
\endgroup
\part{Irreducibility methods}
\chapter{Irreducibility by reduction modulo primes}\label{article:modular-irreducibility}\hypertarget{article:modular-irreducibility}{}\color{black}
{\sffamily\small\color{atlasgold}method / Start here / Q / Classical theory}\par\medskip
Use one finite-field factorization, or intersect the possible degrees supplied by several primes. Understand why the method can fail for an irreducible polynomial.\par\medskip
\textbf{Questions answered}\begin{itemize}
\item When does irreducibility modulo p prove irreducibility over Q?
\item How can reducible reductions still prove rational irreducibility?
\item Can every irreducible polynomial be detected this way?
\end{itemize}
{\small\textbf{Prerequisites:} \hyperlink{article:foundations}{Fields, factors, and the action on roots}}\par
\needspace{8\baselineskip}\section{Preserve the degree}

\begin{atlasblock}[standard]
\textbf{\sffamily Theorem  /  reduction criterion}\par\smallskip
\color{black} Let \(f\in\mathbb Z[x]\) be primitive of degree \(n\ge2\). If a prime \(p\) does not divide its leading coefficient and \(\bar f\in\mathbb F_p[x]\) is irreducible, then \(f\) is irreducible over \(\mathbb Q\).\par
\end{atlasblock}\color{black}

\color{black} The leading-coefficient condition is essential. For example, \(2x^2+3x+1=(2x+1)(x+1)\) becomes the irreducible linear polynomial \(x+1\) modulo \(2\). The degree drop destroys the inference. No squarefreeness hypothesis is separately needed for the reduction criterion: an irreducible polynomial over a finite field is already separable.\par


\begin{atlasblock}[standard]
\textbf{\sffamily Why it works}\par\smallskip
\color{black} Gauss's lemma turns a putative rational factorization into \(f=gh\) in \(\mathbb Z[x]\). Since the product of the leading coefficients is not divisible by \(p\), both factors retain their positive degrees after reduction. Their reductions would factor \(\bar f\), a contradiction.\par
\end{atlasblock}\color{black}
\begin{sources}
\textbf{Sources for this section}\par\smallskip
\noindent\hyperlink{ref:ConradIrredQ}{[ConradIrredQ]} / §3 and Appendix A. Reduction modulo a prime and the Gauss-lemma proof. \href{https://kconrad.math.uconn.edu/blurbs/ringtheory/irredtestsoverQ.pdf}{Source}.\par
\end{sources}\color{black}
\needspace{8\baselineskip}\section{Keep all possible factor degrees}
\color{black} Suppose the irreducible factors of \(\bar f\), listed with multiplicity, have degrees \(d_1,\ldots,d_r\). Every rational factor degree must lie in\par


\atlasdisplay{D_p=\left\{\sum_{i\in I}d_i: I\subseteq\{1,\ldots,r\}\right\}\cap\{1,\ldots,n-1\}.}

\color{black} A rational factor reduces to a product of some of these finite-field factors. Therefore \(\bigcap_{p\in P}D_p=\varnothing\), for any finite set of degree-preserving primes \(P\), proves irreducibility. In computation one can retain only degrees at most \(n/2\), because a reducible polynomial has a factor in that range.\par

\color{black} This degree argument remains valid at a repeated reduction: multiplicities must be included. Its conclusion is weaker than a complete rational factorization, but it is exact. By contrast, reading the same factor degrees as a Frobenius cycle type requires a squarefree reduction.\par
\begin{sources}
\textbf{Sources for this section}\par\smallskip
\noindent\hyperlink{ref:Prasolov2004}{[Prasolov2004]} / Chapter 2. Reduction and coefficient methods for rational irreducibility; the subset-sum formulation follows directly from reduction of an integral factorization. \href{https://link.springer.com/book/10.1007/978-3-642-03980-5}{Source}.\par
\end{sources}\color{black}
\needspace{8\baselineskip}\section{Worked example: two incomplete tests make a proof}

\begin{atlasblock}[standard]
\textbf{\sffamily Example  /  \(f=x^4+x^3+x+2\)}\par\smallskip


\atlasdisplay{\begin{aligned}f\bmod2&=x(x^3+x^2+1),\\ f\bmod3&=(x^2+1)(x^2+x-1).\end{aligned}}

\color{black} The cubic has no root in \(\mathbb F_2\), so it is irreducible there. The two quadratics have no roots in \(\mathbb F_3\), so they are irreducible there. Modulo \(2\), a proper rational factor could have degree \(1\) or \(3\). Modulo \(3\), it could only have degree \(2\). The sets are disjoint; hence \(f\) is irreducible over \(\mathbb Q\).\par

\color{black} The proof certificate is the pair of displayed factorizations together with the root checks. Finding the primes may take a search; checking the resulting certificate is short and exact.\par
\end{atlasblock}\color{black}
\begin{sources}
\textbf{Sources for this section}\par\smallskip
\noindent\hyperlink{ref:ConradIrredQ}{[ConradIrredQ]} / §3. The reduction principle underlying the certificate; these two displayed factorizations are calculated for this entry. \href{https://kconrad.math.uconn.edu/blurbs/ringtheory/irredtestsoverQ.pdf}{Source}.\par
\end{sources}\color{black}
\needspace{8\baselineskip}\section{An obstruction that no prime search removes}
\color{black} The polynomial \(x^4+1\) is irreducible over \(\mathbb Q\): translating by \(1\) gives \(x^4+4x^3+6x^2+4x+2\), Eisenstein at \(2\). Nevertheless it is reducible modulo every prime. At \(2\) it is \((x+1)^4\). For odd \(p\), its roots are primitive eighth roots of unity; since \(p^2\equiv1\pmod8\), they lie in \(\mathbb F_{p^2}\), so all irreducible factors have degree at most two.\par

\color{black} In fact degree \(2\) remains compatible with every reduction, so even intersecting all the sets \(D_p\) does not prove irreducibility here. Group theory explains the obstruction: \(C_2\times C_2\), the Galois group, contains no \(4\)-cycle. For a separable irreducible polynomial over \(\mathbb Q\), Chebotarev and the Frobenius factorization theorem show that an irreducible good reduction exists exactly when the group contains an \(n\)-cycle.\par
\begin{sources}
\textbf{Sources for this section}\par\smallskip
\noindent\hyperlink{ref:ConradAllPrimes}{[ConradAllPrimes]}. A complementary quartic example and a proof that reduction at every prime can fail. \href{https://kconrad.math.uconn.edu/blurbs/ringtheory/reducibleallp.pdf}{Source}.\par
\noindent\hyperlink{ref:StevenhagenLenstra1996}{[StevenhagenLenstra1996]}. The relation between Frobenius conjugacy classes, factorization types, and Chebotarev density. \href{https://pub.math.leidenuniv.nl/~lenstrahw/PUBLICATIONS/1996d/art.pdf}{Source}.\par
\end{sources}\color{black}
\needspace{8\baselineskip}\section{How to choose the next method}
\color{black} Start with cheap good primes and record possible factor degrees, not just whether the reduction is irreducible. If the intersection quickly becomes empty, save the factorizations as a certificate. If the intersection stabilizes, inspect coefficient valuations and structural symmetries. Repeated reductions may be useful for Newton polygons, while reciprocal or compositional structure can predict why particular cycle types never appear.\par

\color{black} Failure of a finite prime search is an inconclusive computation. A proof that the group has no \(n\)-cycle is a mathematical explanation of failure of the single-prime test. These are different kinds of evidence.\par
\begin{sources}
\textbf{Sources for this section}\par\smallskip
\noindent\hyperlink{ref:ConradModP}{[ConradModP]}. Good-prime factorizations as Galois-group information. \href{https://kconrad.math.uconn.edu/blurbs/gradnumthy/galois-Q-factor-mod-p.pdf}{Source}.\par
\end{sources}\color{black}
\needspace{11\baselineskip}\section*{Questions and research boundaries}
\needspace{7\baselineskip}\subsection*{Does every irreducible rational polynomial have an irreducible reduction?}
{\small\sffamily\color{atlasgold}Settled in stated scope / Separable irreducible polynomials over Q}\par\smallskip
\color{black} No. The precise condition is that the Galois group contain a full-length cycle. The polynomial \(x^4+1\) is a concrete counterexample to an unconditional converse.\par
\begin{sources}
\textbf{Sources for this section}\par\smallskip
\noindent\hyperlink{ref:StevenhagenLenstra1996}{[StevenhagenLenstra1996]}. Chebotarev gives primes in each available conjugacy class. \href{https://pub.math.leidenuniv.nl/~lenstrahw/PUBLICATIONS/1996d/art.pdf}{Source}.\par
\noindent\hyperlink{ref:ConradAllPrimes}{[ConradAllPrimes]}. An explicit failure of the converse. \href{https://kconrad.math.uconn.edu/blurbs/ringtheory/reducibleallp.pdf}{Source}.\par
\end{sources}\color{black}
\needspace{6\baselineskip}\par\medskip\begingroup\footnotesize\setlength{\parskip}{1pt}\textbf{\sffamily\color{atlasgreen}Connections}\par
\color{black}\textbf{Uses.} \hyperlink{article:foundations}{Fields, factors, and the action on roots}. Gauss's lemma is the link from rational factors to integral factors.\par
\color{black}\textbf{Contrasts with.} \hyperlink{article:newton-polygons}{Newton polygons, factor degrees, and large cycles}. Valuations can extract information from primes where reduction alone has repeated factors.\par
\color{black}\textbf{Supplies examples for.} \hyperlink{article:frobenius}{Frobenius: from factor degrees to permutations}. Squarefree modular factorizations supply cycle types as well as factor-degree restrictions.\par
\color{black}\textbf{Uses.} \hyperlink{article:finite-field-irreducibility}{Irreducibility over finite fields}. Exact finite-field irreducibility tests verify the modular certificate.\par
\endgroup
\chapter{Schönemann-Eisenstein and one-edge Newton polygons}\label{article:eisenstein-dumas}\hypertarget{article:eisenstein-dumas}{}\color{black}
{\sffamily\small\color{atlasgold}method / Start here / Q / Classical theory}\par\medskip
A divisibility pattern can force every factor to have the full degree. Translation and Newton slopes extend the elementary criterion.\par\medskip
\textbf{Questions answered}\begin{itemize}
\item Why does divisibility by p, but not p\ensuremath{{}^2}, prove irreducibility?
\item What can a change of variable reveal?
\item How does the one-edge Dumas criterion extend Eisenstein?
\end{itemize}
{\small\textbf{Prerequisites:} \hyperlink{article:foundations}{Fields, factors, and the action on roots}}\par
\needspace{8\baselineskip}\section{The elementary criterion and its proof}

\begin{atlasblock}[standard]
\textbf{\sffamily Theorem  /  Schönemann-Eisenstein}\par\smallskip
\color{black} Let \(f=a_nx^n+\cdots+a_0\in\mathbb Z[x]\), with \(n\ge1\). Suppose a prime \(p\) satisfies \(p\nmid a_n\), \(p\mid a_i\) for \(0\le i<n\), and \(p^2\nmid a_0\). Then \(f\) is irreducible over \(\mathbb Q\).\par
\end{atlasblock}\color{black}


\begin{atlasblock}[standard]
\textbf{\sffamily Proof mechanism}\par\smallskip
\color{black} A factorization into positive degrees can be taken in \(\mathbb Z[x]\). Reducing it modulo \(p\) gives a factorization of the monomial \(a_nx^n\). Both reduced factors are monomials of positive degree, because their leading coefficients remain nonzero. Thus both original constant terms are divisible by \(p\), which would make \(a_0\) divisible by \(p^2\).\par
\end{atlasblock}\color{black}

\color{black} Schönemann's work of 1846 precedes Eisenstein's 1850 paper. Cox's historical account explains their different motivations and the connection with finite fields and the division of the lemniscate. The mathematical content should be separated from the later naming convention.\par
\begin{sources}
\textbf{Sources for this section}\par\smallskip
\noindent\hyperlink{ref:schoenemann1846}{[schoenemann1846]}. Original historical source. \href{https://www.degruyterbrill.com/document/doi/10.1515/crll.1846.32.93/html}{Source}.\par
\noindent\hyperlink{ref:eisenstein1850}{[eisenstein1850]}. Eisenstein's original treatment. \href{https://doi.org/10.1515/crll.1850.39.160}{Source}.\par
\noindent\hyperlink{ref:cox2011eisenstein}{[cox2011eisenstein]} / pp. 3-4. Historical attribution and the elementary form of the criterion. \href{https://people.math.ethz.ch/~halorenz/4students/Algebra/Schoenemann_Eisenstein.pdf}{Source}.\par
\end{sources}\color{black}
\needspace{8\baselineskip}\section{Translate before testing}
\color{black} For \(a\in K^\times\) and \(b\in K\), substitution \(x\mapsto ax+b\) is an automorphism of \(K[x]\) and preserves irreducibility. It can change coefficient divisibility dramatically. Reversing coefficients, \(f^*(x)=x^{\deg f}f(1/x)\), also preserves irreducibility when \(f(0)\ne0\).\par


\begin{atlasblock}[standard]
\textbf{\sffamily Example  /  a hidden Eisenstein polynomial}\par\smallskip


\atlasdisplay{(x+1)^4+1=x^4+4x^3+6x^2+4x+2.}

\color{black} The right side is Eisenstein at \(2\), proving irreducibility of \(x^4+1\). More generally, if \(p\) is prime,\par


\atlasdisplay{\Phi_p(x+1)=\frac{(x+1)^p-1}{x}=\sum_{j=1}^{p}\binom pj x^{j-1}}

\color{black} is Eisenstein at \(p\): all intermediate binomial coefficients are divisible by \(p\), while the constant coefficient is exactly \(p\).\par
\end{atlasblock}\color{black}
\begin{sources}
\textbf{Sources for this section}\par\smallskip
\noindent\hyperlink{ref:cox2011eisenstein}{[cox2011eisenstein]} / p. 4. The prime cyclotomic example after translation. \href{https://people.math.ethz.ch/~halorenz/4students/Algebra/Schoenemann_Eisenstein.pdf}{Source}.\par
\end{sources}\color{black}
\needspace{8\baselineskip}\section{Schönemann's residue-power form}
\color{black} The repeated reduction need not be a power of a linear polynomial. Let \(\phi\in\mathbb Z[x]\) be monic with irreducible reduction modulo \(p\), let \(e\ge2\), and let \(M\in\mathbb Z[x]\) satisfy \(\deg M<e\deg\phi\) and \(\gcd(\bar\phi,\bar M)=1\). Then\par

\atlasdisplay{f=\phi^e+pM}

\color{black} is irreducible over \(\mathbb Q\). Indeed, a monic integral factorization would reduce to \(\bar g=\bar\phi^a\), \(\bar h=\bar\phi^b\), where \(a,b>0\) and \(a+b=e\). Write \(g=\phi^a+pg_1\), \(h=\phi^b+ph_1\). Comparing with \(f\), then dividing by \(p\) and reducing modulo \(p\), forces \(\bar\phi\mid\bar M\), a contradiction.\par

\color{black} For example, \((x^2+x+1)^2+2=x^4+2x^3+3x^2+2x+3\) is irreducible: \(x^2+x+1\) is irreducible over \(\mathbb F_2\), and \(M=1\). This uses a higher-degree irreducible residual factor, rather than a monomial reduction.\par
\begin{sources}
\textbf{Sources for this section}\par\smallskip
\noindent\hyperlink{ref:schoenemann1846}{[schoenemann1846]}. Historical source of the residue-power criterion. \href{https://www.degruyterbrill.com/document/doi/10.1515/crll.1846.32.93/html}{Source}.\par
\noindent\hyperlink{ref:bonciocat2015multiprime}{[bonciocat2015multiprime]} / Introduction. Modern statement of Schönemann's criterion. \href{https://www.tandfonline.com/doi/abs/10.1080/00927872.2014.910800}{Source}.\par
\end{sources}\color{black}
\needspace{8\baselineskip}\section{Replace valuation one by a coprime slope}

\begin{atlasblock}[standard]
\textbf{\sffamily Theorem  /  one-edge Dumas criterion}\par\smallskip
\color{black} Let \(f=a_nx^n+\cdots+a_0\in\mathbb Q_p[x]\), with \(a_0a_n\ne0\), \(v_p(a_n)=0\), and \(v_p(a_0)=k>0\). Suppose \(\gcd(n,k)=1\) and\par


\atlasdisplay{v_p(a_i)\ge k\left(1-\frac{i}{n}\right)\qquad(0<i<n),}

\color{black} where \(v_p(0)=+\infty\). Then \(f\) is irreducible over \(\mathbb Q_p\). For a rational polynomial, this also proves irreducibility over \(\mathbb Q\).\par
\end{atlasblock}\color{black}

\color{black} The Newton polygon has one segment of slope \(-k/n\), so every root has valuation \(k/n\). A monic factor of degree \(d\) over \(\mathbb Q_p\) would have constant coefficient of valuation \(dk/n\), an integer. Coprimality forces \(n\mid d\). Hence no proper positive factor degree is possible. Eisenstein is the case \(k=1\).\par
\begin{sources}
\textbf{Sources for this section}\par\smallskip
\noindent\hyperlink{ref:dumas1906}{[dumas1906]}. Classical source for the Newton-polygon criterion. \href{https://www.numdam.org/item/JMPA_1906_6_2__191_0/}{Source}.\par
\noindent\hyperlink{ref:bonciocat2015multiprime}{[bonciocat2015multiprime]} / Introduction. Introduction states Dumas's criterion in reversed-coefficient form; the form here follows by coefficient reversal. \href{https://www.tandfonline.com/doi/abs/10.1080/00927872.2014.910800}{Source}.\par
\noindent\hyperlink{ref:Hajir2005}{[Hajir2005]} / Theorem 2.1. Roots and slopes over a discretely valued field. \href{https://jtnb.centre-mersenne.org/item/JTNB_2005__17_2_517_0/}{Source}.\par
\end{sources}\color{black}
\needspace{8\baselineskip}\section{Worked example beyond the elementary test}

\begin{atlasblock}[standard]
\textbf{\sffamily Example  /  \(f=x^5+4x+4\)}\par\smallskip
\color{black} At \(2\), the nonzero coefficient points are \((0,2),(1,2),(5,0)\). The middle point lies above the segment joining \((0,2)\) to \((5,0)\). Its slope is \(-2/5\), and \(\gcd(2,5)=1\). Thus \(f\) is irreducible over \(\mathbb Q_2\), hence over \(\mathbb Q\).\par
\color{black} The ordinary criterion at \(2\) does not apply to the displayed polynomial, since \(4\mid f(0)\). The Newton polygon uses the denominator of the root valuation, rather than insisting that the constant term have valuation one.\par
\end{atlasblock}\color{black}
\begin{sources}
\textbf{Sources for this section}\par\smallskip
\noindent\hyperlink{ref:bonciocat2015multiprime}{[bonciocat2015multiprime]} / Introduction. The one-edge criterion used here; the valuation diagram and polynomial are an explicit worked computation. \href{https://www.tandfonline.com/doi/abs/10.1080/00927872.2014.910800}{Source}.\par
\end{sources}\color{black}
\needspace{8\baselineskip}\section{What the criterion does not decide}
\color{black} A polygon with several edges, or one whose endpoint increments have a common divisor, need not imply reducibility. It may simply leave several possible degrees. Combine primes, use residual polynomials, or change the variable. Nor does local irreducibility identify the full global Galois group.\par

\color{black} Even allowing every integral translation does not make the elementary Eisenstein test complete. Conrad proves that \(x^4-10x^2+1\) is irreducible over \(\mathbb Q\), but \((x+c)^4-10(x+c)^2+1\) is never Eisenstein for any \(c\in\mathbb Z\). This is a useful reason to maintain several genuinely different proof methods.\par

\color{black} The stronger one-edge criterion succeeds after translation: at \(c=1\) the polynomial becomes \(x^4+4x^3-4x^2-16x-8\). Its \(2\)-adic coefficient valuations are \(3,4,2,2,0\), and the lower polygon is a single edge of slope \(-3/4\). Dumas therefore proves irreducibility even over \(\mathbb Q_2\). The failure of every ordinary Eisenstein shift is compatible with a successful Newton-polygon proof.\par
\begin{sources}
\textbf{Sources for this section}\par\smallskip
\noindent\hyperlink{ref:ConradAllPrimes}{[ConradAllPrimes]} / Theorems 1 and 2. Irreducibility and failure of every integral Eisenstein shift. \href{https://kconrad.math.uconn.edu/blurbs/ringtheory/reducibleallp.pdf}{Source}.\par
\noindent\hyperlink{ref:dumas1906}{[dumas1906]}. The one-edge criterion applies to the translated quartic, whose coefficients and valuations are calculated here. \href{https://www.numdam.org/item/JMPA_1906_6_2__191_0/}{Source}.\par
\end{sources}\color{black}
\needspace{11\baselineskip}\section*{Questions and research boundaries}
\needspace{7\baselineskip}\subsection*{Can every irreducible integer polynomial be made Eisenstein by an integral translation?}
{\small\sffamily\color{atlasgold}Settled in stated scope / Translations x\ensuremath{\mapsto}x+c with c\ensuremath{\in}Z}\par\smallskip
\color{black} No. The irreducible polynomial \(x^4-10x^2+1\) is a counterexample. This statement concerns integral translations; it does not classify all other changes of variable or every valuation-theoretic criterion.\par
\begin{sources}
\textbf{Sources for this section}\par\smallskip
\noindent\hyperlink{ref:ConradAllPrimes}{[ConradAllPrimes]} / Theorems 1 and 2. Explicit counterexample. \href{https://kconrad.math.uconn.edu/blurbs/ringtheory/reducibleallp.pdf}{Source}.\par
\end{sources}\color{black}
\needspace{6\baselineskip}\par\medskip\begingroup\footnotesize\setlength{\parskip}{1pt}\textbf{\sffamily\color{atlasgreen}Connections}\par
\color{black}\textbf{Uses.} \hyperlink{article:foundations}{Fields, factors, and the action on roots}. Gauss's lemma turns rational reducibility into an integral factorization.\par
\color{black}\textbf{Uses.} \hyperlink{article:newton-polygons}{Newton polygons, factor degrees, and large cycles}. The one-edge criterion is a special case of the general Newton-polygon factor-degree argument.\par
\color{black}\textbf{Contrasts with.} \hyperlink{article:modular-irreducibility}{Irreducibility by reduction modulo primes}. An Eisenstein polynomial has a highly repeated reduction at its useful prime.\par
\color{black}\textbf{Applies to.} \hyperlink{article:taylor}{Taylor truncations: two exact benchmarks}. Factorial coefficient valuations produce the polygons used for truncated exponentials.\par
\endgroup
\chapter{Newton polygons, factor degrees, and large cycles}\label{article:newton-polygons}\hypertarget{article:newton-polygons}{}\color{black}
{\sffamily\small\color{atlasgold}method / Intermediate / Q / Classical theory}\par\medskip
Read p-adic root valuations from coefficients, sharpen factor-degree restrictions with residual polynomials, and use suitable slope denominators to force a large Galois group.\par\medskip
\textbf{Questions answered}\begin{itemize}
\item What information does a Newton polygon encode?
\item How can several primes prove irreducibility?
\item What can prove irreducibility when one Newton edge still allows proper factors?
\item When does a local slope force containment of the alternating group?
\end{itemize}
{\small\textbf{Prerequisites:} \hyperlink{article:foundations}{Fields, factors, and the action on roots}; \hyperlink{article:eisenstein-dumas}{Schönemann-Eisenstein and one-edge Newton polygons}}\par
\needspace{8\baselineskip}\section{A geometric record of valuations}
\color{black} Write \(f(x)=\sum_{i=0}^n a_ix^i\in\mathbb Q_p[x]\) with \(a_0a_n\ne0\), and normalize \(v_p(p)=1\). The Newton polygon is the lower convex hull of the points \((i,v_p(a_i))\) with \(a_i\ne0\). Missing coefficients have infinite valuation and do not contribute finite points.\par


\begin{atlasblock}[standard]
\textbf{\sffamily Theorem  /  slopes and roots}\par\smallskip
\color{black} A segment of slope \(-u/v\), in lowest terms with \(v>0\), and horizontal length \(\ell\), corresponds to exactly \(\ell\) roots of valuation \(u/v\), counted with multiplicity. There is a factor over \(\mathbb Q_p\) consisting of all roots with each given valuation.\par
\end{atlasblock}\color{black}

\color{black} For a factor over \(\mathbb Q_p\), the number of its roots belonging to that slope must be a multiple of \(v\). Indeed, the constant coefficient of its corresponding monic slope factor has integer valuation. Thus the possible factor degrees are sums of multiples of the slope denominators, bounded by the lengths available on the respective segments.\par
\begin{sources}
\textbf{Sources for this section}\par\smallskip
\noindent\hyperlink{ref:Hajir2005}{[Hajir2005]} / §2.1, Theorem 2.1. Definition and root-slope theorem. \href{https://jtnb.centre-mersenne.org/item/JTNB_2005__17_2_517_0/}{Source}.\par
\end{sources}\color{black}
\needspace{8\baselineskip}\section{Local restrictions can contradict each other}
\color{black} A rational factor remains a factor over every \(\mathbb Q_p\). Consequently its degree must satisfy the restrictions from every prime used. This is the same logic as combining finite-field factor degrees, with finer local information. A particularly useful prime may be one for which all roots collapse to a repeated root modulo \(p\).\par

\color{black} The polygon-product rule explains the mechanism: the segments of the product polygon are assembled from those of the factor polygons, ordered by slope. Integral lattice lengths constrain how much of a segment a factor can use. Dumas already combined such restrictions at several primes; Bonciocat develops explicit criteria using arbitrarily many primes.\par
\begin{sources}
\textbf{Sources for this section}\par\smallskip
\noindent\hyperlink{ref:dumas1906}{[dumas1906]} / p. 238, item 4. Historical source for the polygon-product and several-prime method. \href{https://www.numdam.org/item/JMPA_1906_6_2__191_0/}{Source}.\par
\noindent\hyperlink{ref:bonciocat2015multiprime}{[bonciocat2015multiprime]} / §1, Theorem A (arXiv:1304.0874v2, p. 3). Systematic several-prime irreducibility criteria. \href{https://www.tandfonline.com/doi/abs/10.1080/00927872.2014.910800}{Source}.\par
\end{sources}\color{black}
\needspace{8\baselineskip}\section{A residual polynomial sharpens a single edge}
\color{black} A Newton edge records valuations but forgets the leading residues of the coefficients on that edge. Retaining them can eliminate factors that the slope denominator alone permits.\par


\begin{atlasblock}[standard]
\textbf{\sffamily Criterion  /  one edge with irreducible residual polynomial}\par\smallskip
\color{black} Let \(f(x)=\sum_{i=0}^{ed}a_ix^i\in\mathbb Z_p[x]\) be monic. Suppose its Newton polygon is the single segment from \((0,hd)\) to \((ed,0)\), where \(h,e,d\) are positive integers and \(\gcd(h,e)=1\). Define\par


\atlasdisplay{R_f(Y)=\sum_{j=0}^{d}\overline{\frac{a_{ej}}{p^{h(d-j)}}}\,Y^j\in\mathbb F_p[Y].}

\color{black} The polygon hypothesis makes every displayed quotient integral. If \(R_f\) is irreducible, then \(f\) is irreducible over \(\mathbb Q_p\).\par
\end{atlasblock}\color{black}


\begin{atlasblock}[standard]
\textbf{\sffamily Proof}\par\smallskip
\color{black} In a monic factorization \(f=g_1g_2\) over \(\mathbb Q_p\), the factors have integral coefficients. All roots have valuation \(h/e\), so the constant coefficients show that the factor degrees are \(ed_1,ed_2\), with \(d_1+d_2=d\). The residual-polynomial product rule gives \(R_f=R_{g_1}R_{g_2}\), with degrees \(d_1,d_2\). This rule follows directly from coefficient multiplication: terms on the edge survive after normalization and reduction; terms above it vanish. If two indices summing to a multiple of \(e\) are individually not multiples of \(e\), their integral valuations put their product strictly above the edge. Irreducibility of \(R_f\) therefore forces one factor degree to be zero.\par
\end{atlasblock}\color{black}


\begin{atlasblock}[standard]
\textbf{\sffamily Example  /  a quartic certified at a repeated reduction}\par\smallskip
\color{black} For \(f=x^4+2x^2+4\) at \(p=2\), the points \((0,2),(2,1),(4,0)\) lie on one edge of slope \(-1/2\). Its denominator leaves a possible quadratic factor. Here \(h=1,e=2,d=2\), and\par


\atlasdisplay{R_f(Y)=1+Y+Y^2.}

\color{black} This quadratic has no root in \(\mathbb F_2\). Thus \(f\) is irreducible over \(\mathbb Q_2\), and hence over \(\mathbb Q\), although its ordinary reduction is \(x^4\). The criterion does not assume tame ramification. When \(d=1\), it recovers the one-edge Dumas criterion.\par
\end{atlasblock}\color{black}
\begin{sources}
\textbf{Sources for this section}\par\smallskip
\noindent\hyperlink{ref:GuardiaMontesNart2012Newton}{[GuardiaMontesNart2012Newton]} / §1.3-1.4: Theorem 1.13 and Corollary 1.20 in arXiv:0807.2620v2. The product rule and residual-polynomial irreducibility criterion. The elementary specialization and quartic calculation are worked out here. \href{https://doi.org/10.1090/S0002-9947-2011-05442-5}{Source}.\par
\noindent\hyperlink{ref:GuardiaMontesNart2011Montes}{[GuardiaMontesNart2011Montes]} / §2.1-2.2, Theorems 2.1-2.2. The same criterion is encoded by a complete type; the source explains its ramification and residue-degree information. \href{https://www.numdam.org/articles/10.5802/jtnb.782/}{Source}.\par
\end{sources}\color{black}
\needspace{8\baselineskip}\section{Worked example: a truncated exponential of degree six}

\begin{atlasblock}[standard]
\textbf{\sffamily Example  /  an exact two-prime certificate}\par\smallskip


\atlasdisplay{F(x)=6!\sum_{j=0}^6\frac{x^j}{j!}=720+720x+360x^2+120x^3+30x^4+6x^5+x^6.}

\color{black} The coefficient valuations, from constant to leading term, are\par


\atlasdisplay{\begin{array}{c|rrrrrrr}i&0&1&2&3&4&5&6\\\hline v_2(a_i)&4&4&3&3&1&1&0\\ v_3(a_i)&2&2&2&1&1&1&0\end{array}.}

\color{black} At \(2\), the vertices are \((0,4),(4,1),(6,0)\). The slopes are \(-3/4,-1/2\), with lengths \(4,2\). The only proper factor degrees permitted are \(2,4\). At \(3\), there is one segment from \((0,2)\) to \((6,0)\), with slope \(-1/3\); the only proper factor degree permitted is \(3\). Hence \(F\) is irreducible over \(\mathbb Q\).\par
\color{black} This finite calculation illustrates the broader Newton-polygon proof for exponential Taylor polynomials. It is a proof for this polynomial, not a deduction from a numerical root plot.\par
\end{atlasblock}\color{black}
\begin{sources}
\textbf{Sources for this section}\par\smallskip
\noindent\hyperlink{ref:Coleman1987}{[Coleman1987]}. Newton-polygon treatment of exponential Taylor polynomials; the degree-six valuation table is calculated here. \href{https://doi.org/10.5169/seals-87891}{Source}.\par
\noindent\hyperlink{ref:Hajir2005}{[Hajir2005]} / Theorem 2.1. The local theorem justifying the factor-degree restrictions. \href{https://jtnb.centre-mersenne.org/item/JTNB_2005__17_2_517_0/}{Source}.\par
\end{sources}\color{black}
\needspace{8\baselineskip}\section{From a denominator to an alternating group}

\begin{atlasblock}[standard]
\textbf{\sffamily Criterion  /  a useful form of Hajir's Newton-index theorem}\par\smallskip
\color{black} Let \(f\in\mathbb Q[x]\) be irreducible of degree \(n\). Suppose a prime \(q\) divides the denominator, in lowest terms, of some Newton slope of \(f\) at a finite prime. If\par

\atlasdisplay{\frac n2<q\le n-3,}
\color{black} then \(\operatorname{Gal}(f/\mathbb Q)\) contains \(A_n\).\par
\end{atlasblock}\color{black}

\color{black} The root valuation forces \(q\) to divide a ramification index in the local splitting field, and hence the order of the global group. Cauchy's theorem gives an element of order \(q\); because \(2q>n\), this element is a single \(q\)-cycle. Transitivity comes from irreducibility. A nontrivial block system would have both its block size and its number of blocks smaller than \(q\); an element of order \(q\) could act nontrivially on neither, which is impossible. Thus the action is primitive. Jordan's theorem then gives \(A_n\).\par

\color{black} The discriminant decides between \(A_n\) and \(S_n\) only after this containment has been proved. A nonsquare discriminant alone does not imply the full symmetric group.\par
\begin{sources}
\textbf{Sources for this section}\par\smallskip
\noindent\hyperlink{ref:Hajir2005}{[Hajir2005]} / Theorem 2.2. Newton index divides the Galois-group order and a prime in the stated interval forces alternating-group containment. \href{https://jtnb.centre-mersenne.org/item/JTNB_2005__17_2_517_0/}{Source}.\par
\noindent\hyperlink{ref:Jordan1871}{[Jordan1871]}. Classical primitive-group theorem used in the final step. \href{https://www.numdam.org/item/JMPA_1871_2_16__383_0/}{Source}.\par
\noindent\hyperlink{ref:DixonMortimer1996}{[DixonMortimer1996]}. Permutation-group background for primitivity and long prime cycles. \href{https://link.springer.com/book/10.1007/978-1-4612-0731-3}{Source}.\par
\end{sources}\color{black}
\needspace{8\baselineskip}\section{Worked example: a complete symmetric-group certificate}

\begin{atlasblock}[standard]
\textbf{\sffamily Example  /  \(f=x^8+2x^5+6\) has group \(S_8\)}\par\smallskip
\color{black} Eisenstein at \(2\) proves irreducibility. At \(3\), the Newton polygon has vertices \((0,1),(5,0),(8,0)\), hence slopes \(-1/5,0\) of lengths \(5,3\). The denominator forces \(5\) to divide the Galois-group order. Since \(8<2\cdot5\), an element of order \(5\) is a single 5-cycle. Irreducibility supplies transitivity, and \(8/2<5\le8-3\) puts us exactly in the criterion above. Thus the group contains \(A_8\).\par
\color{black} The exact discriminant is\par

\atlasdisplay{\operatorname{disc}(f)=2^{12}3^9\cdot79\cdot733.}

\color{black} It is nonsquare, so the group is \(S_8\). This calculation separates the three inputs: Eisenstein for transitivity, a Newton slope for a long prime cycle, and the discriminant for parity. The useful primes \(2,3\) have highly repeated reductions: \(x^8\) and \(x^5(x-1)^3\). Reading these multiplicities as Frobenius cycle lengths would be invalid.\par
\end{atlasblock}\color{black}
\begin{sources}
\textbf{Sources for this section}\par\smallskip
\noindent\hyperlink{ref:Hajir2005}{[Hajir2005]} / Theorem 2.2. The Newton-index theorem proves alternating containment for the displayed octic. The valuations and resultant were computed exactly for this entry. \href{https://jtnb.centre-mersenne.org/item/JTNB_2005__17_2_517_0/}{Source}.\par
\noindent\hyperlink{ref:Milne2022}{[Milne2022]}. The discriminant-square criterion completes the group determination. \href{https://www.jmilne.org/math/Books/index.html}{Author's book page}.\par
\end{sources}\color{black}
\needspace{8\baselineskip}\section{What remains after drawing the polygon}
\color{black} A slope denominator is an arithmetic constraint, not automatically a cycle of that length. The prime-order argument above uses both the denominator and the interval \(n/2<q\le n-3\). If the relevant prime is smaller, an element of order \(q\) may be a product of several disjoint \(q\)-cycles. If irreducibility is unavailable, transitivity cannot be assumed.\par


\begin{atlasblock}[boundary]
\textbf{\sffamily Counterexample  /  a denominator two without a transposition}\par\smallskip
\color{black} The polynomial \(x^4+4=(x^2-2x+2)(x^2+2x+2)\) has one Newton slope \(-1/2\) at \(2\). Its roots are \(1\pm i,-1\pm i\), and its splitting field is \(\mathbb Q(i)\). The only nonidentity automorphism acts as a product of two transpositions, never a single transposition. Its discriminant \(2^{14}\) is a square, consistently with this even action. Capelli's exceptional fourth-power case explains the factorization.\par
\end{atlasblock}\color{black}

\color{black} Irreducibility does not repair this inference. The irreducible quartic \(x^4+2x^2+4\) certified above also has slope denominator two. Its roots are \((\pm\sqrt2\pm\sqrt{-6})/2\), with independent sign choices, and its splitting field is \(\mathbb Q(\sqrt2,\sqrt{-6})\). Its Galois group is \(V_4\) in the transitive degree-four action: all three nonidentity elements are double transpositions. Its discriminant \(2^{10}3^2\) is a square.\par

\color{black} For applications, separate three tasks: eliminate proper factor degrees; obtain a suitable local ramification divisor; calculate the discriminant square class. Laguerre and Taylor families are instructive because their coefficient formulas let these tasks be attacked systematically, but each has its own parameter restrictions.\par
\begin{sources}
\textbf{Sources for this section}\par\smallskip
\noindent\hyperlink{ref:Hajir2005}{[Hajir2005]} / Theorem 2.2. The irreducibility input and prime interval are explicit hypotheses. \href{https://jtnb.centre-mersenne.org/item/JTNB_2005__17_2_517_0/}{Source}.\par
\noindent\hyperlink{ref:capelli1901binomials}{[capelli1901binomials]}. The fourth-power exception explains the explicitly factored quartic. \href{https://link.springer.com/article/10.1007/BF01450724}{Source}.\par
\noindent\hyperlink{ref:GuardiaMontesNart2012Newton}{[GuardiaMontesNart2012Newton]} / Corollary 1.20 in arXiv:0807.2620v2. The earlier residual-polynomial criterion proves irreducibility of the second counterexample; its radicals, field, permutation action, and discriminant are calculated here. \href{https://doi.org/10.1090/S0002-9947-2011-05442-5}{Source}.\par
\end{sources}\color{black}
\needspace{11\baselineskip}\section*{Questions and research boundaries}
\needspace{7\baselineskip}\subsection*{How can one sharpen a Newton polygon that leaves possible factor degrees?}
{\small\sffamily\color{atlasgold}Settled in stated scope / Local factorization of a specified separable polynomial over a p-adic field}\par\smallskip
\color{black} Residual polynomials give the first refinement, illustrated above. If a residual factor is repeated, that stage may leave irreducibility undecided; higher Newton polygons refine the associated valuation and residue data. The Guàrdia-Montes-Nart theory gives a terminating local factorization procedure for separable polynomials. Choosing especially short certificates for a family is a further method-design question, rather than a general undecided local factorization problem.\par
\begin{sources}
\textbf{Sources for this section}\par\smallskip
\noindent\hyperlink{ref:GuardiaMontesNart2012Newton}{[GuardiaMontesNart2012Newton]} / Introduction; §3 and Theorem 4.18 in arXiv:0807.2620v2.  \href{https://doi.org/10.1090/S0002-9947-2011-05442-5}{Source}.\par
\noindent\hyperlink{ref:GuardiaMontesNart2011Montes}{[GuardiaMontesNart2011Montes]} / §2.2-3.  \href{https://www.numdam.org/articles/10.5802/jtnb.782/}{Source}.\par
\end{sources}\color{black}
\needspace{6\baselineskip}\par\medskip\begingroup\footnotesize\setlength{\parskip}{1pt}\textbf{\sffamily\color{atlasgreen}Connections}\par
\color{black}\textbf{Generalizes.} \hyperlink{article:eisenstein-dumas}{Schönemann-Eisenstein and one-edge Newton polygons}. The elementary criterion is the one-edge slope \ensuremath{-}1/n case.\par
\color{black}\textbf{Contrasts with.} \hyperlink{article:modular-irreducibility}{Irreducibility by reduction modulo primes}. Newton polygons can exploit repeated reductions where the squarefree Frobenius test is unavailable.\par
\color{black}\textbf{Uses.} \hyperlink{article:permutation-groups}{The group-theoretic step: transitivity, alternating groups, and symmetric groups}. Cauchy and Jordan turn suitable ramification information into a large permutation group.\par
\color{black}\textbf{Uses.} \hyperlink{article:discriminants}{Discriminants, parity, and ramification}. The discriminant square class completes the A\_\allowbreak{}n versus S\_\allowbreak{}n decision after containment.\par
\color{black}\textbf{Supplies examples for.} \hyperlink{article:capelli-composition}{Capelli: composition, binomials, and iterates}. The binomials x\ensuremath{{}^6}\ensuremath{-}72 and x\ensuremath{{}^4}+4 connect power obstructions with valuation denominators.\par
\color{black}\textbf{Applies to.} \hyperlink{article:laguerre}{Generalized Laguerre polynomials: two parameter regimes}. Factorial and binomial coefficient valuations provide controlled Newton slopes.\par
\color{black}\textbf{Develops into.} \hyperlink{article:factorization-algorithms}{Exact factorization: finite fields, Hensel lifting, and lattices}. Higher Newton polygons and residual polynomials supply the types used in the Montes local-factorization method.\par
\endgroup
\chapter{Root location, coefficient dominance, and prime values}\label{article:archimedean-criteria}\hypertarget{article:archimedean-criteria}{}\color{black}
{\sffamily\small\color{atlasgold}method / Intermediate / Q / Classical theory}\par\medskip
Use complex absolute values and integer evaluations to make any hypothetical factor impossible. These methods complement congruence and valuation arguments.\par\medskip
\textbf{Questions answered}\begin{itemize}
\item How can one large coefficient prove irreducibility?
\item When does a single prime value suffice?
\item What changes when coefficients are digits of a prime?
\end{itemize}
{\small\textbf{Prerequisites:} \hyperlink{article:foundations}{Fields, factors, and the action on roots}}\par
\needspace{8\baselineskip}\section{Perron's coefficient-dominance criterion}

\begin{atlasblock}[standard]
\textbf{\sffamily Theorem  /  Perron}\par\smallskip
\color{black} Let \(f=x^n+a_{n-1}x^{n-1}+\cdots+a_0\in\mathbb Z[x]\), where \(n\ge2\) and \(a_0\ne0\). If\par

\atlasdisplay{|a_{n-1}|>1+\sum_{i=0}^{n-2}|a_i|,}
\color{black} then \(f\) is irreducible over \(\mathbb Q\).\par
\end{atlasblock}\color{black}


\begin{atlasblock}[standard]
\textbf{\sffamily Proof mechanism}\par\smallskip
\color{black} On \(|z|=1\), the term \(a_{n-1}z^{n-1}\) strictly dominates the sum of all other terms. Rouché's theorem gives exactly \(n-1\) roots inside the unit circle and one outside. If the monic polynomial factored over \(\mathbb Z\), one positive-degree monic factor would have all roots strictly inside. Its constant coefficient would be a nonzero integer whose absolute value is the product of numbers smaller than one, an impossibility.\par
\end{atlasblock}\color{black}

\color{black} The nonzero constant term prevents a zero root; monicity makes the constant-coefficient product argument decisive. The strict inequality controls roots on the unit circle. Variants exist, but they require their own hypotheses.\par
\begin{sources}
\textbf{Sources for this section}\par\smallskip
\noindent\hyperlink{ref:perron1907}{[perron1907]} / Theorem I. Original source for coefficient-dominance irreducibility. \href{https://www.degruyterbrill.com/document/doi/10.1515/crll.1907.132.288/html}{Source}.\par
\noindent\hyperlink{ref:Prasolov2004}{[Prasolov2004]} / Theorem 2.2.5. A modern presentation of Perron's criterion. \href{https://link.springer.com/book/10.1007/978-3-642-03980-5}{Source}.\par
\end{sources}\color{black}
\needspace{8\baselineskip}\section{Worked example: reverse the polynomial}

\begin{atlasblock}[standard]
\textbf{\sffamily Example  /  \(x^5+3x+1\)}\par\smallskip
\color{black} Reverse its coefficients:\par

\atlasdisplay{f^*(x)=x^5f(1/x)=x^5+3x^4+1.}

\color{black} Now \(3>1+1\), so Perron's criterion proves \(f^*\) irreducible. Since the constant term of \(f\) is nonzero, coefficient reversal preserves factorizations and proves \(f\) irreducible. The same argument works for \(x^n+3x+1\) for every \(n\ge2\).\par
\color{black} The method uses the location of roots, but no approximations to the roots are needed. A drawing may suggest a useful circle; the coefficient inequality supplies the proof.\par
\end{atlasblock}\color{black}
\begin{sources}
\textbf{Sources for this section}\par\smallskip
\noindent\hyperlink{ref:perron1907}{[perron1907]}. Coefficient-dominance theorem; the reversal and worked family are direct applications. \href{https://www.degruyterbrill.com/document/doi/10.1515/crll.1907.132.288/html}{Source}.\par
\end{sources}\color{black}
\needspace{8\baselineskip}\section{A prime value must be sufficiently far from every root}

\begin{atlasblock}[standard]
\textbf{\sffamily Elementary prime-value criterion}\par\smallskip
\color{black} Let \(f\in\mathbb Z[x]\) be nonconstant. Suppose an integer \(m\) satisfies \(|f(m)|\) prime and \(|m-\alpha|>1\) for every complex root \(\alpha\) of \(f\). Then \(f\) is irreducible over \(\mathbb Q\).\par
\end{atlasblock}\color{black}

\color{black} In an integral factorization \(f=gh\), each positive-degree factor has an integer leading coefficient \(b\ne0\), so\par

\atlasdisplay{|g(m)|=|b|\prod_{g(\alpha)=0}|m-\alpha|>1.}

\color{black} Both factors would therefore give integer divisors of \(f(m)\) larger than one in absolute value. A practical sufficient condition is a certified root bound \(|\alpha|\le R\) and \(|m|>R+1\). This deliberately simple form is weaker than some prime-value criteria in the literature.\par

\color{black} For example, the roots of \(x^4-10x^2+1\) are \(\pm\sqrt2\pm\sqrt3\), all of absolute value less than \(4\). At \(m=6\), the value is \(937\), a prime. Every distance is greater than \(2\), so the criterion proves irreducibility. Primality can be checked by the primes at most \(\sqrt{937}<31\).\par
\begin{sources}
\textbf{Sources for this section}\par\smallskip
\noindent\hyperlink{ref:murty2002primevalues}{[murty2002primevalues]}. Prime evaluations together with root or coefficient bounds; the displayed distance criterion is proved in the text. \href{https://doi.org/10.1080/00029890.2002.11919872}{Source}.\par
\noindent\hyperlink{ref:ConradAllPrimes}{[ConradAllPrimes]}. The roots of the quartic used in the example. \href{https://kconrad.math.uconn.edu/blurbs/ringtheory/reducibleallp.pdf}{Source}.\par
\end{sources}\color{black}
\needspace{8\baselineskip}\section{A different shortcut: the digits of a prime}
\color{black} Cohn's decimal theorem says that if \(p=a_n10^n+\cdots+a_0\) is prime, with \(n\ge1\), \(0\le a_i\le9\), and \(a_n\ne0\), then \(a_nx^n+\cdots+a_0\) is irreducible over \(\mathbb Q\). Brillhart, Filaseta, and Odlyzko extend the result to arbitrary integer bases \(b\ge2\), again assuming a nonconstant digit polynomial. The digit conditions are doing mathematical work; merely knowing one evaluation is prime is insufficient.\par

\color{black} For instance, \(101\) is prime, so its decimal digits give the irreducible polynomial \(x^2+1\). The unrestricted prime-value inference fails for \((x-1)(x+1)\), which takes the prime value \(3\) at \(x=2\), yet is visibly reducible.\par
\begin{sources}
\textbf{Sources for this section}\par\smallskip
\noindent\hyperlink{ref:brillhartfilasetaodlyzko1981}{[brillhartfilasetaodlyzko1981]} / Theorem 1 (p. 1055) and Corollary 2 (p. 1058). Cohn's theorem, prime-value root separation, and extension beyond decimal digits. \href{https://www.cambridge.org/core/journals/canadian-journal-of-mathematics/article/on-an-irreducibility-theorem-of-a-cohn/F1644BDFD66CFC166841810C29CBA91A}{Source}.\par
\noindent\hyperlink{ref:murty2002primevalues}{[murty2002primevalues]}. A further exposition and proof of the arbitrary-base extension. \href{https://doi.org/10.1080/00029890.2002.11919872}{Source}.\par
\end{sources}\color{black}
\needspace{8\baselineskip}\section{Where these arguments stop}
\color{black} Root location is only useful after it is connected to an arithmetic obstruction. For example, a polynomial can have all its roots real and still have a large Galois group. The special force of Perron's criterion is that a hypothetical monic factor would have an impossible constant coefficient. The special force of a suitably large prime value is that every nonconstant factor would evaluate to an integer of absolute value greater than one.\par

\color{black} A complete factorization algorithm decides any one rational polynomial, but a uniform theorem for a sparse or parameterized family may require a new inequality or a different method. Such families provide meaningful test cases for improving these criteria.\par
\begin{sources}
\textbf{Sources for this section}\par\smallskip
\noindent\hyperlink{ref:Schinzel2000}{[Schinzel2000]}. General reducibility theory and criteria for special forms. \href{https://doi.org/10.1017/CBO9780511542916}{Source}.\par
\noindent\hyperlink{ref:GathenGerhard2013}{[GathenGerhard2013]}. Exact polynomial factorization algorithms. \href{https://doi.org/10.1017/CBO9781139856065}{Source}.\par
\end{sources}\color{black}
\needspace{11\baselineskip}\section*{Questions and research boundaries}
\needspace{7\baselineskip}\subsection*{Does one prime value prove irreducibility?}
{\small\sffamily\color{atlasgold}Settled in stated scope / Integer polynomials without additional coefficient or root restrictions}\par\smallskip
\color{black} No: \((x-1)(x+1)\) takes the prime value \(3\) at \(x=2\). A distance-to-roots bound, a sufficiently large evaluation, or digit restrictions can turn a prime value into an irreducibility criterion.\par
\begin{sources}
\textbf{Sources for this section}\par\smallskip
\noindent\hyperlink{ref:murty2002primevalues}{[murty2002primevalues]}. Precise additional hypotheses for prime-value arguments. \href{https://doi.org/10.1080/00029890.2002.11919872}{Source}.\par
\end{sources}\color{black}
\needspace{6\baselineskip}\par\medskip\begingroup\footnotesize\setlength{\parskip}{1pt}\textbf{\sffamily\color{atlasgreen}Connections}\par
\color{black}\textbf{Contrasts with.} \hyperlink{article:modular-irreducibility}{Irreducibility by reduction modulo primes}. Archimedean arguments may settle examples for which every prime reduction is reducible.\par
\color{black}\textbf{Uses.} \hyperlink{article:reciprocity}{Reciprocal symmetry and the group on root pairs}. Coefficient reversal can move a dominant coefficient into the useful position.\par
\color{black}\textbf{Applies to.} \hyperlink{article:sparse-gaps}{Sparse polynomials, factor reversal, and large gaps}. Sparse support makes root and norm estimates particularly effective.\par
\color{black}\textbf{Contrasts with.} \hyperlink{article:factorization-algorithms}{Exact factorization: finite fields, Hensel lifting, and lattices}. A short proof certificate for a family is different from a general-purpose factorization algorithm.\par
\endgroup
\chapter{Small values, unit values, and prime values}\label{article:integer-values}\hypertarget{article:integer-values}{}\color{black}
{\sffamily\small\color{atlasgold}method / Intermediate / Q / Classical theory}\par\medskip
Integer evaluations constrain every possible integral factor. Three precise criteria turn those constraints into proofs, with examples showing where the converses fail.\par\medskip
\textbf{Questions answered}\begin{itemize}
\item Why is a product of distinct linear integer factors minus one irreducible?
\item How small must many values be to force irreducibility?
\item How many prime or unit values suffice?
\end{itemize}
{\small\textbf{Prerequisites:} \hyperlink{article:foundations}{Fields, factors, and the action on roots}}\par
\needspace{8\baselineskip}\section{Unit values: a rigid sign relation}

\begin{atlasblock}[standard]
\textbf{\sffamily Theorem  /  the minus-one case of Schur's problem}\par\smallskip
\color{black} For distinct integers \(a_1,\ldots,a_m\), with \(m\ge1\), the polynomial\par

\atlasdisplay{P(x)=\prod_{i=1}^{m}(x-a_i)-1}
\color{black} is irreducible over \(\mathbb Q\).\par
\end{atlasblock}\color{black}


\begin{atlasblock}[standard]
\textbf{\sffamily Proof  /  too many roots of a smaller polynomial}\par\smallskip
\color{black} A nontrivial factorization gives monic \(g,h\in\mathbb Z[x]\) of positive degrees with \(P=gh\). At each \(a_i\), the integers \(g(a_i),h(a_i)\) multiply to \(-1\), so \(g(a_i)+h(a_i)=0\). The polynomial \(g+h\) has degree less than \(m\) and at least \(m\) distinct roots. It must vanish identically, contradicting the monicity of both factors.\par
\end{atlasblock}\color{black}

\color{black} The sign cannot simply be reversed. Two exact exceptions are\par


\atlasdisplay{x(x-2)+1=(x-1)^2,\qquad x(x-1)(x-2)(x-3)+1=(x^2-3x+1)^2.}

\color{black} Westlund and Flügel's classical solution, recalled in Győry-Hajdu-Tijdeman, shows that these are the only reducible plus-one cases, up to an integer translation of the variable. Schur posed the question in 1908; the distinction between a problem and its subsequent solution matters for attribution.\par
\begin{sources}
\textbf{Sources for this section}\par\smallskip
\noindent\hyperlink{ref:schur1908problem226}{[schur1908problem226]}. The original problem, rather than an attribution of the later solution. \href{https://link.springer.com/article/10.1007/s00605-010-0241-9}{Historical citation verification}.\par
\noindent\hyperlink{ref:gyoryhajdutijdeman2011}{[gyoryhajdutijdeman2011]} / §1, p. 416. The minus-one theorem, the complete plus-one exception list, and attribution of the 1909 solution. \href{https://link.springer.com/article/10.1007/s00605-010-0241-9}{Source}.\par
\end{sources}\color{black}
\needspace{8\baselineskip}\section{Pólya: many values below a factorial threshold}

\begin{atlasblock}[standard]
\textbf{\sffamily Theorem  /  Pólya's criterion}\par\smallskip
\color{black} Let \(P\in\mathbb Z[x]\) have degree \(n\ge2\), and put \(N=\lceil n/2\rceil\). If there are \(n\) distinct integers \(a\) such that\par

\atlasdisplay{0<|P(a)|<\frac{N!}{2^N},}
\color{black} then \(P\) is irreducible over \(\mathbb Q\).\par
\end{atlasblock}\color{black}


\begin{atlasblock}[standard]
\textbf{\sffamily Proof  /  interpolation controls a factor's leading coefficient}\par\smallskip
\color{black} If \(P=QR\) with positive-degree integral factors, choose the larger degree \(d=\deg Q\), so \(N\le d\le n-1\). At every selected integer, \(|Q(a)|\le|P(a)|\), because the other integral factor is nonzero. Choose \(d+1\) of those arguments and order them as \(b_0<\cdots<b_d\). Their integer spacing gives \(\prod_{i\ne j}|b_j-b_i|\ge j!(d-j)!\). Writing \(M=\max_j|Q(b_j)|\), Lagrange interpolation gives\par


\atlasdisplay{1\le|\operatorname{lc}(Q)|\le M\sum_{j=0}^{d}\frac1{j!(d-j)!}=M\frac{2^d}{d!}.}

\color{black} Consequently \(M\ge d!/2^d\ge N!/2^N\), contradicting the strict bound. Gauss's lemma justifies the integral factorization even when \(P\) is not monic.\par
\end{atlasblock}\color{black}


\begin{atlasblock}[standard]
\textbf{\sffamily Example  /  a degree-ten certificate}\par\smallskip
\color{black} For \(P(x)=\prod_{j=0}^{9}(x-j)+2\), the ten values at \(0,\ldots,9\) are all \(2\). Since \(2<5!/2^5=15/4\), Pólya proves irreducibility without expanding the polynomial. Nonzero values are essential: an integer zero would itself give a rational linear factor.\par
\end{atlasblock}\color{black}
\begin{sources}
\textbf{Sources for this section}\par\smallskip
\noindent\hyperlink{ref:polya1919}{[polya1919]}. The classical small-value irreducibility criterion. \href{https://eudml.org/doc/145560}{Source}.\par
\noindent\hyperlink{ref:gyoryhajdutijdeman2011}{[gyoryhajdutijdeman2011]} / §1, p. 416. Confirms the ceiling in N and the strict factorial threshold; the displayed interpolation proof is supplied here. \href{https://link.springer.com/article/10.1007/s00605-010-0241-9}{Source}.\par
\end{sources}\color{black}
\needspace{8\baselineskip}\section{Prime values: count the possible unit evaluations}

\begin{atlasblock}[standard]
\textbf{\sffamily Theorem  /  a finite prime-value certificate}\par\smallskip
\color{black} Let \(f\in\mathbb Z[x]\) be primitive of degree \(n\ge1\). If \(|f(a)|\) is either \(1\) or a prime at \(2n+1\) distinct integers \(a\), then \(f\) is irreducible over \(\mathbb Q\).\par
\end{atlasblock}\color{black}


\begin{atlasblock}[standard]
\textbf{\sffamily Proof  /  four equations}\par\smallskip
\color{black} For a putative positive-degree integral factorization \(f=gh\), every selected argument satisfies \(g(a)=1\), \(g(a)=-1\), \(h(a)=1\), or \(h(a)=-1\). Those four equations have at most \(2\deg g+2\deg h=2n\) distinct solutions in total, a contradiction.\par
\end{atlasblock}\color{black}


\begin{atlasblock}[standard]
\textbf{\sffamily Example  /  five values suffice for a quadratic}\par\smallskip
\color{black} At \(a=0,1,2,3,4\), the values of \(x^2+x+41\) are \(41,43,47,53,61\), all prime. Thus the criterion proves irreducibility. It says nothing about primality at all later arguments; at \(40\) the value is \(41^2\).\par
\end{atlasblock}\color{black}
\begin{sources}
\textbf{Sources for this section}\par\smallskip
\noindent\hyperlink{ref:ConradIrredQ}{[ConradIrredQ]} / Theorem 5.1. The finite prime-or-unit-value criterion and its elementary counting proof. \href{https://kconrad.math.uconn.edu/blurbs/ringtheory/irredtestsoverQ.pdf}{Source}.\par
\end{sources}\color{black}
\needspace{8\baselineskip}\section{What values alone do not tell you}

\begin{atlasblock}[boundary]
\textbf{\sffamily Boundary  /  sufficient conditions}\par\smallskip
\color{black} A single prime value is insufficient: \((x-1)(x+1)\) takes the prime value \(3\) at \(2\). Conversely, \(x^2+x+4\) is irreducible because its discriminant is \(-15\), yet every integer value is even and at least \(4\). It never passes the preceding prime-or-unit-value test.\par
\end{atlasblock}\color{black}

\color{black} A single prime value can suffice when extra information controls the factors: Cohn's criterion restricts the coefficients to base digits, while root bounds force nonconstant factors to have large absolute values at a chosen integer. The archimedean entry explains those mechanisms. For more general small-value criteria, Levit, Tverberg, and Győry-Hajdu-Tijdeman develop sharper bounds and extensions. Their precise hypotheses should replace any informal assertion that ``many small values'' automatically imply irreducibility.\par
\begin{sources}
\textbf{Sources for this section}\par\smallskip
\noindent\hyperlink{ref:ConradIrredQ}{[ConradIrredQ]} / §5, p. 5. The irreducible polynomial with no prime or unit evaluations. \href{https://kconrad.math.uconn.edu/blurbs/ringtheory/irredtestsoverQ.pdf}{Source}.\par
\noindent\hyperlink{ref:levit1968}{[levit1968]}. Further criteria based on low absolute values. \href{https://www.jstor.org/stable/1994841}{Source}.\par
\noindent\hyperlink{ref:tverberg1973smallvalues}{[tverberg1973smallvalues]}. Refinements of the small-value method. \href{https://journals.msp.org/mscand/article/view/2165}{Source}.\par
\noindent\hyperlink{ref:gyoryhajdutijdeman2011}{[gyoryhajdutijdeman2011]}. Extensions involving Schur-type and Pólya-type families. \href{https://link.springer.com/article/10.1007/s00605-010-0241-9}{Source}.\par
\end{sources}\color{black}
\needspace{11\baselineskip}\section*{Questions and research boundaries}
\needspace{7\baselineskip}\subsection*{Does every irreducible integer polynomial take a prime or unit value?}
{\small\sffamily\color{atlasgold}Settled in stated scope / Primitive polynomials in Z[x], with integer arguments and without changing variables}\par\smallskip
\color{black} No. The primitive irreducible polynomial \(x^2+x+4\) has only positive even values at integers, all at least \(4\). This elementary obstruction explains why the finite prime-value criterion is sufficient but cannot be necessary.\par
\begin{sources}
\textbf{Sources for this section}\par\smallskip
\noindent\hyperlink{ref:ConradIrredQ}{[ConradIrredQ]} / §5, p. 5. The same explicit obstruction is recorded after the prime-value criterion. \href{https://kconrad.math.uconn.edu/blurbs/ringtheory/irredtestsoverQ.pdf}{Source}.\par
\end{sources}\color{black}
\needspace{6\baselineskip}\par\medskip\begingroup\footnotesize\setlength{\parskip}{1pt}\textbf{\sffamily\color{atlasgreen}Connections}\par
\color{black}\textbf{Uses.} \hyperlink{article:foundations}{Fields, factors, and the action on roots}. Gauss's lemma lets integer evaluations constrain rational factors.\par
\color{black}\textbf{Contrasts with.} \hyperlink{article:archimedean-criteria}{Root location, coefficient dominance, and prime values}. The one-value criteria add coefficient or root-location information to control possible unit factors.\par
\color{black}\textbf{Applies to.} \hyperlink{article:capelli-composition}{Capelli: composition, binomials, and iterates}. Schur-type constructions lead to composed families whose irreducibility can also be studied in root fields.\par
\color{black}\textbf{Contrasts with.} \hyperlink{article:factorization-algorithms}{Exact factorization: finite fields, Hensel lifting, and lattices}. A short evaluation certificate proves irreducibility without producing factors or searching over coefficient bounds.\par
\endgroup
\chapter{Sparse polynomials, factor reversal, and large gaps}\label{article:sparse-gaps}\hypertarget{article:sparse-gaps}{}\color{black}
{\sffamily\small\color{atlasgold}method / Research / Q / Published results}\par\medskip
Separate persistent reciprocal factors from the nonreciprocal part, then use coefficient sparsity to constrain possible factorizations.\par\medskip
\textbf{Questions answered}\begin{itemize}
\item Why should reciprocal factors be separated in a sparse family?
\item What does factor reversal preserve?
\item What is actually proved for sufficiently large gaps?
\end{itemize}
{\small\textbf{Prerequisites:} \hyperlink{article:foundations}{Fields, factors, and the action on roots}; \hyperlink{article:reciprocity}{Reciprocal symmetry and the group on root pairs}}\par
\needspace{8\baselineskip}\section{The family and its unavoidable obstructions}
\color{black} A useful large-gap model is\par

\atlasdisplay{f_N(x)=x^N c(x^{-1})+d(x),\qquad N>\deg c+\deg d,}

\color{black} where \(c,d\in\mathbb Z[x]\) are fixed and \(c(0)d(0)\ne0\). The high and low blocks of coefficients are separated by a gap. Write \(h^*(x)=x^{\deg h}h(1/x)\). A nonconstant irreducible factor is reciprocal when \(h^*=\pm h\); the \emph{nonreciprocal part} is obtained by removing all such irreducible factors, including their multiplicities.\par

\color{black} This removal is different from removing only cyclotomic factors. Every cyclotomic polynomial is reciprocal up to sign, but a reciprocal polynomial need not be cyclotomic. A theorem about the nonreciprocal part should never be quoted as irreducibility of the original polynomial without checking the removed factors.\par
\begin{sources}
\textbf{Sources for this section}\par\smallskip
\noindent\hyperlink{ref:Schinzel2000}{[Schinzel2000]}. Lacunary polynomial factorization and reciprocal factors. \href{https://doi.org/10.1017/CBO9780511542916}{Source}.\par
\noindent\hyperlink{ref:sawinshustermanstoll2020}{[sawinshustermanstoll2020]} / §1. Large-gap notation and definitions. \href{https://arxiv.org/abs/1803.10811}{Source}.\par
\end{sources}\color{black}
\needspace{8\baselineskip}\section{A general theorem, with the structural exceptions stated}
\color{black} Call \((c,d)\) \emph{Capellian} if \(-d(x)/c(x^{-1})\) is a \(p\)-th power in \(\mathbb Q(x)\) for some prime \(p\), or if \(d(x)/c(x^{-1})\) is four times a fourth power. These conditions explain systematic binomial-type factorizations.\par


\begin{atlasblock}[standard]
\textbf{\sffamily Theorem  /  Schinzel, in large-gap form}\par\smallskip
\color{black} Assume \(c(0)d(0)\ne0\), \(\gcd_{\mathbb Z[x]}(c^*,d)=1\), \(c\ne\pm d\), and that \((c,d)\) is not Capellian. Then there is an effective bound \(N_0\), depending on \(c,d\), such that for \(N>N_0\) the nonreciprocal part of \(f_N\) is irreducible over \(\mathbb Q\).\par
\end{atlasblock}\color{black}

\color{black} The cited formulation is Theorem 1.2 of Sawin-Shusterman-Stoll, attributed there to Schinzel. Filaseta-Ford-Konyagin improve explicit thresholds. Under an additional robustness condition, Sawin-Shusterman-Stoll obtain substantially more practical bounds and algorithms; that extra condition cannot be dropped when quoting their improvement.\par
\begin{sources}
\textbf{Sources for this section}\par\smallskip
\noindent\hyperlink{ref:Schinzel2000}{[Schinzel2000]} / Theorem 74, as identified in Sawin-Shusterman-Stoll. The general lacunary result behind this formulation. \href{https://doi.org/10.1017/CBO9780511542916}{Source}.\par
\noindent\hyperlink{ref:filasetafordkonyagin2000}{[filasetafordkonyagin2000]}. Improved effective bounds with explicit exclusions. \href{https://www.ford126.web.illinois.edu/wwwpapers/ffk.pdf}{Source}.\par
\noindent\hyperlink{ref:sawinshustermanstoll2020}{[sawinshustermanstoll2020]} / Theorems 1.2, 1.5, and 1.6. Exact hypotheses of the large-gap theorem and the extra robustness assumption for its improved bound. \href{https://arxiv.org/abs/1803.10811}{Source}.\par
\end{sources}\color{black}
\needspace{8\baselineskip}\section{The central idea: reverse one factor}
\color{black} For \(f(0)\ne0\) and a factorization \(f=gh\), form \(G=gh^*\). The reciprocal product does not change:\par

\atlasdisplay{GG^*=ff^*.}

\color{black} Equivalently, as Laurent polynomials, one may use \(g(x)h(x^{-1})\); multiplication by a power of \(x\) only changes the presentation. On the unit circle, a real-coefficient polynomial satisfies \(|h^*(z)|=|h(z)|\). Consequently \(G\) and \(f\) have the same squared coefficient norm, by Parseval's identity:\par

\atlasdisplay{\sum_j |[x^j]G|^2=\frac1{2\pi}\int_0^{2\pi}|G(e^{it})|^2\,dt=\sum_j |[x^j]f|^2.}

\color{black} For a sparse polynomial with small integral coefficients, the right side is small. There are few ways to distribute that integer norm among nonzero coefficients. Comparing the full identity \(GG^*=ff^*\) can then force the reversal to be trivial or expose a reciprocal factor. This is the mechanism behind Ljunggren's approach and its later extensions.\par
\begin{sources}
\textbf{Sources for this section}\par\smallskip
\noindent\hyperlink{ref:Ljunggren1960}{[Ljunggren1960]}. Original factor-reversal approach to sparse trinomials and quadrinomials. \href{https://doi.org/10.7146/math.scand.a-10593}{Source}.\par
\noindent\hyperlink{ref:filaseta1999smallnorm}{[filaseta1999smallnorm]}. Small Euclidean coefficient norm as a factorization constraint. \href{https://doi.org/10.1515/9783110285581.143}{Source}.\par
\noindent\hyperlink{ref:sawinshustermanstoll2020}{[sawinshustermanstoll2020]} / §§3-4. Reversal identity and the later algorithmic development. \href{https://arxiv.org/abs/1803.10811}{Source}.\par
\end{sources}\color{black}
\needspace{8\baselineskip}\section{Worked example: a factor recurring in one residue class}

\begin{atlasblock}[standard]
\textbf{\sffamily Example  /  \(x^N+x+1\)}\par\smallskip
\color{black} Let \(\zeta\) be a primitive cube root of unity. Since \(\zeta+1=-\zeta^2\),\par

\atlasdisplay{\zeta^N+\zeta+1=0\quad\Longleftrightarrow\quad N\equiv2\pmod3.}

\color{black} Thus \(x^2+x+1\) divides \(x^N+x+1\) exactly in that residue class. In degree \(5\), for example,\par

\atlasdisplay{x^5+x+1=(x^2+x+1)(x^3-x^2+1).}

\color{black} The cubic has no rational root, hence is irreducible over \(\mathbb Q\). This illustrates why a description of the nonreciprocal part can coexist with infinitely many reducible members of a family.\par
\end{atlasblock}\color{black}

\color{black} The periodicity is general for a fixed cyclotomic factor: evaluation at a root of unity depends on \(N\) modulo its order. The large-gap literature combines this elementary observation with bounds that control other factors.\par
\begin{sources}
\textbf{Sources for this section}\par\smallskip
\noindent\hyperlink{ref:Ljunggren1960}{[Ljunggren1960]}. Sparse trinomial factorization and reciprocal obstructions; the displayed root-of-unity calculation is given here. \href{https://doi.org/10.7146/math.scand.a-10593}{Source}.\par
\noindent\hyperlink{ref:sawinshustermanstoll2020}{[sawinshustermanstoll2020]} / Corollary 1.4. Effective eventual description by finitely many exceptional exponents and arithmetic progressions. \href{https://arxiv.org/abs/1803.10811}{Source}.\par
\end{sources}\color{black}
\needspace{8\baselineskip}\section{A theorem for a concrete family and its boundary}
\color{black} Sawin-Shusterman-Stoll prove that \(x^N-7x^2+1\) is irreducible for every \(N\ge5\), and also for \(N=3\). The omitted degree \(4\) really is exceptional:\par

\atlasdisplay{x^4-7x^2+1=(x^2-3x+1)(x^2+3x+1).}

\color{black} Their general theory supplies an eventual threshold, and a finite check handles the remaining exponents in this example. This pattern is common in family problems: a uniform argument for large parameters plus a finite exact calculation, with exceptional factors recorded explicitly.\par
\begin{sources}
\textbf{Sources for this section}\par\smallskip
\noindent\hyperlink{ref:sawinshustermanstoll2020}{[sawinshustermanstoll2020]} / Example 7.1. The complete k=7 example and its exceptional quartic. \href{https://arxiv.org/abs/1803.10811}{Source}.\par
\end{sources}\color{black}
\needspace{11\baselineskip}\section*{Questions and research boundaries}
\needspace{7\baselineskip}\subsection*{How small can a uniform large-gap threshold be?}
{\small\sffamily\color{atlasgold}Open in cited source / Pairs satisfying the hypotheses of the large-gap theorem; degree sum \ensuremath{\delta} and squared coefficient weight at most w}\par\smallskip
\color{black} Sawin-Shusterman-Stoll define \(N_{\mathrm{opt}}(\delta,w)\) as the smallest uniform threshold for their Theorem 1.2 over pairs with \(\deg c+\deg d=\delta\) and \(\|c\|_2^2+\|d\|_2^2\le w\). Their Question 1.7 asks, for fixed \(\delta>1\), whether this is \(O_\delta(\log w)\); Question 1.8 asks for the uniform bound \(O(\delta\log w)\). This records the questions as posed in that paper; it is not a claim that all later work has been excluded.\par
\begin{sources}
\textbf{Sources for this section}\par\smallskip
\noindent\hyperlink{ref:sawinshustermanstoll2020}{[sawinshustermanstoll2020]} / Questions 1.7 and 1.8. Precise definition and proposed threshold questions. \href{https://arxiv.org/abs/1803.10811}{Source}.\par
\end{sources}\color{black}
\needspace{6\baselineskip}\par\medskip\begingroup\footnotesize\setlength{\parskip}{1pt}\textbf{\sffamily\color{atlasgreen}Connections}\par
\color{black}\textbf{Uses.} \hyperlink{article:reciprocity}{Reciprocal symmetry and the group on root pairs}. One must distinguish the reciprocal and nonreciprocal parts before stating a sparse-family theorem.\par
\color{black}\textbf{Applies to.} \hyperlink{article:trinomials}{Trinomials: uniform families and explicit classifications}. Trinomials are the simplest setting where factor reversal and recurring cyclotomic factors are visible.\par
\color{black}\textbf{Contrasts with.} \hyperlink{article:archimedean-criteria}{Root location, coefficient dominance, and prime values}. Coefficient norms and root size complement elementary dominance inequalities.\par
\color{black}\textbf{Uses.} \hyperlink{article:capelli-composition}{Capelli: composition, binomials, and iterates}. Capellian power obstructions are the binomial obstructions applied in a rational function field.\par
\color{black}\textbf{Contrasts with.} \hyperlink{article:composition-iteration}{Composition, iteration, and arboreal groups}. Large-gap factorization and iteration impose different parameter structures.\par
\endgroup
\chapter{Capelli: composition, binomials, and iterates}\label{article:capelli-composition}\hypertarget{article:capelli-composition}{}\color{black}
{\sffamily\small\color{atlasgold}method / Intermediate / general fields; characteristic zero for binomials / Classical theory}\par\medskip
Replace the irreducibility of a composition by a lower-degree question in a root field. Binomials and a stable quadratic illustrate how power obstructions and norms make this reduction effective.\par\medskip
\textbf{Questions answered}\begin{itemize}
\item When is a composition of polynomials irreducible?
\item Exactly when is a binomial irreducible in characteristic zero?
\item How can one prove irreducibility of every iterate?
\end{itemize}
{\small\textbf{Prerequisites:} \hyperlink{article:foundations}{Fields, factors, and the action on roots}}\par
\needspace{8\baselineskip}\section{Move the problem into a root field}

\begin{atlasblock}[standard]
\textbf{\sffamily Theorem  /  Capelli's composition lemma}\par\smallskip
\color{black} Let \(K\) be any field, let \(F\in K[x]\) be irreducible of degree \(m\ge1\), and let \(g\in K[x]\) have degree \(d\ge1\). For a root \(\alpha\) of \(F\), the composition \(F(g(x))\) is irreducible over \(K\) if and only if \(g(x)-\alpha\) is irreducible over \(K(\alpha)\).\par
\end{atlasblock}\color{black}


\begin{atlasblock}[standard]
\textbf{\sffamily Proof  /  multiplication of degrees}\par\smallskip
\color{black} Choose a root \(\beta\) of \(g(x)-\alpha\). Since \(\alpha=g(\beta)\), we have \(K(\alpha)\subseteq K(\beta)\), and\par

\atlasdisplay{[K(\beta):K]=m\,[K(\beta):K(\alpha)].}
\color{black} The polynomial \(F(g(x))\), of degree \(md\), has root \(\beta\). It is irreducible exactly when the left side is \(md\), equivalently when the minimal polynomial of \(\beta\) over \(K(\alpha)\) has degree \(d\). No separability assumption is needed for this argument.\par
\end{atlasblock}\color{black}

\color{black} The reduction replaces a large degree by arithmetic in a larger coefficient field. When \(g=x^r+c\), the new question is whether \(x^r-(\alpha-c)\) factors; power classes and norms become natural tools.\par
\begin{sources}
\textbf{Sources for this section}\par\smallskip
\noindent\hyperlink{ref:AyadMcQuillan2000}{[AyadMcQuillan2000]} / Lemma 1, p. 87. Capelli's composition lemma, stated for arbitrary fields. \href{https://doi.org/10.4064/aa-93-1-87-97}{Source}.\par
\noindent\hyperlink{ref:Schinzel2000}{[Schinzel2000]}. Composition and reducibility in extension fields. \href{https://doi.org/10.1017/CBO9780511542916}{Source}.\par
\end{sources}\color{black}
\needspace{8\baselineskip}\section{Irreducible inputs can give a reducible composition}

\begin{atlasblock}[standard]
\textbf{\sffamily Example  /  a square appears in the root field}\par\smallskip
\color{black} Both \(F=x^2-8\) and \(g=x^2-3\) are irreducible over \(\mathbb Q\), but\par

\atlasdisplay{F(g(x))=x^4-6x^2+1=(x^2-2x-1)(x^2+2x-1).}
\color{black} With \(\alpha=2\sqrt2\), the obstruction is visible before expanding: \(3+\alpha=(1+\sqrt2)^2\), so \(g(x)-\alpha\) splits in \(\mathbb Q(\sqrt2)[x]\). Checking each input over \(\mathbb Q\) misses this new square.\par
\end{atlasblock}\color{black}
\begin{sources}
\textbf{Sources for this section}\par\smallskip
\noindent\hyperlink{ref:AyadMcQuillan2000}{[AyadMcQuillan2000]} / Lemma 1. The root-field criterion explains the explicitly verified example. \href{https://doi.org/10.4064/aa-93-1-87-97}{Source}.\par
\end{sources}\color{black}
\needspace{8\baselineskip}\section{The exact binomial criterion}

\begin{atlasblock}[standard]
\textbf{\sffamily Theorem  /  Capelli's binomial criterion}\par\smallskip
\color{black} Let \(K\) have characteristic zero, \(a\in K^\times\), and \(n\ge2\). Then \(x^n-a\) is irreducible over \(K\) exactly when\par

\begin{itemize}
\item\leavevmode\color{black} \(a\notin K^p\) for every prime \(p\mid n\);
\item\leavevmode\color{black} if \(4\mid n\), then \(a\notin-4K^4\).
\end{itemize}
\color{black} Here \(K^p=\{b^p:b\in K\}\). This statement is restricted to characteristic zero; finite-field questions are treated separately.\par
\end{atlasblock}\color{black}

\color{black} The first obstruction gives a difference-of-powers factorization. The exceptional fourth-power obstruction comes from\par


\atlasdisplay{x^{4r}+4b^4=(x^{2r}-2bx^r+2b^2)(x^{2r}+2bx^r+2b^2).}

\color{black} These prove necessity. Sufficiency is the substantive part of Capelli's theorem: one builds the extension by prime-power steps and rules out a power appearing in the intermediate root fields. Odd-prime steps use norms; the power of two requires precisely the additional obstruction above.\par


\begin{atlasblock}[standard]
\textbf{\sffamily Example  /  two different obstructions}\par\smallskip
\color{black} The polynomial \(x^4+4\) factors by the displayed identity even though \(-4\) is not a rational square. In contrast, \(x^6-72\) is irreducible: \(72=2^3 3^2\) is neither a square nor a cube. In fact \(x^n-72\) is irreducible for every \(n\ge2\), because no prime divides both valuation exponents \(3,2\), and a positive rational number cannot be \(-4\) times a rational fourth power.\par
\end{atlasblock}\color{black}
\begin{sources}
\textbf{Sources for this section}\par\smallskip
\noindent\hyperlink{ref:capelli1901binomials}{[capelli1901binomials]}. The classical binomial reducibility theorem, including the exceptional fourth-power case. \href{https://link.springer.com/article/10.1007/BF01450724}{Source}.\par
\noindent\hyperlink{ref:Lang2002}{[Lang2002]}. The field-theoretic treatment of irreducible binomials and the norm argument. \href{https://link.springer.com/book/10.1007/978-1-4613-0041-0}{Source}.\par
\end{sources}\color{black}
\needspace{8\baselineskip}\section{The same binomial through valuations and fields}

\begin{atlasblock}[standard]
\textbf{\sffamily Example  /  irreducibility does not determine the group}\par\smallskip
\color{black} For \(h=x^6-72\), the Newton polygon at \(2\) has one slope \(-3/6=-1/2\), allowing proper factor degrees \(2,4\). At \(3\), its slope is \(-2/6=-1/3\), allowing only degree \(3\). The empty intersection gives a second proof of irreducibility, complementing Capelli. Neither polygon individually forces the full degree.\par
\color{black} To determine the group, choose the positive root \(\alpha=\sqrt2\sqrt[3]{3}\). The identities\par

\atlasdisplay{\sqrt2=\alpha^3/6,\qquad\sqrt[3]{3}=\alpha^4/12}

\color{black} show that the splitting field is \(L=\mathbb Q(\sqrt2,\sqrt[3]{3},\zeta_3)\). The splitting field \(M\) of \(x^3-3\) has group \(S_3\): the cubic is Eisenstein at \(3\), and its discriminant is \(-243\). Its unique quadratic subfield is \(\mathbb Q(\sqrt{-3})\). Thus \(M\cap\mathbb Q(\sqrt2)=\mathbb Q\), and restriction gives\par

\atlasdisplay{\operatorname{Gal}(h/\mathbb Q)\cong S_3\times C_2,\qquad |\operatorname{Gal}(h/\mathbb Q)|=12.}

\color{black} The discriminant \(2^{21}3^{16}\) is nonsquare, but the group is much smaller than \(S_6\). The two irreducibility proofs, the root-field calculation, and parity answer distinct questions.\par
\end{atlasblock}\color{black}
\begin{sources}
\textbf{Sources for this section}\par\smallskip
\noindent\hyperlink{ref:dumas1906}{[dumas1906]}. The slope-denominator principle applied to the explicitly computed two-prime certificate. \href{https://www.numdam.org/item/JMPA_1906_6_2__191_0/}{Source}.\par
\noindent\hyperlink{ref:Milne2022}{[Milne2022]}. Galois correspondence, composita, and the quadratic subfield determined by a discriminant; the displayed sextic calculation is supplied here. \href{https://www.jmilne.org/math/Books/index.html}{Author's book page}.\par
\noindent\hyperlink{ref:capelli1901binomials}{[capelli1901binomials]}. The complementary binomial criterion. \href{https://link.springer.com/article/10.1007/BF01450724}{Source}.\par
\end{sources}\color{black}
\needspace{8\baselineskip}\section{An all-level proof for a quadratic}

\begin{atlasblock}[standard]
\textbf{\sffamily Example  /  every iterate of \(f=x^2+1\) is irreducible over \(\mathbb Q\)}\par\smallskip
\color{black} Put \(F_r=f^{\circ r}\) and \(a_r=F_r(0)\), with \(a_0=0\). Then \(a_{r+1}=a_r^2+1\), beginning \(0,1,2,5,26,677\). The base case \(F_1=x^2+1\) is irreducible. Suppose \(F_r\) is irreducible and let \(\alpha\) be a root. Capelli reduces the next step to showing that \(\alpha-1\) is not a square in \(\mathbb Q(\alpha)\).\par
\color{black} Since \(F_r\) is monic of even degree,\par


\atlasdisplay{N_{\mathbb Q(\alpha)/\mathbb Q}(\alpha-1)=F_r(1)=a_{r+1}=a_r^2+1.}

\color{black} For \(r\ge1\), this integer lies strictly between \(a_r^2\) and \((a_r+1)^2\). It is therefore not a rational square. A square in \(\mathbb Q(\alpha)\) would have square norm, giving the contradiction and completing induction.\par
\end{atlasblock}\color{black}


\begin{atlasblock}[boundary]
\textbf{\sffamily Boundary  /  use the norm implication in its valid direction}\par\smallskip
\color{black} A nonsquare norm certifies a nonsquare element; a square norm does not certify a square element. The argument proves stability, meaning irreducibility of every iterate. Determining every Galois group in the resulting tree requires further information. The correction to Ayad-McQuillan is relevant when using broader stability criteria: an originally claimed converse does not hold.\par
\end{atlasblock}\color{black}
\begin{sources}
\textbf{Sources for this section}\par\smallskip
\noindent\hyperlink{ref:AyadMcQuillan2000}{[AyadMcQuillan2000]} / Lemma 1. Capelli's lemma underlying the supplied norm induction. \href{https://doi.org/10.4064/aa-93-1-87-97}{Source}.\par
\noindent\hyperlink{ref:AyadMcQuillan2001Correction}{[AyadMcQuillan2001Correction]} / Theorem A. The corrected sufficient stability criterion and the failure of the original converse. \href{https://geodesic.mathdoc.fr/articles/10.4064/aa99-1-9/}{Source}.\par
\noindent\hyperlink{ref:Jones2013Arboreal}{[Jones2013Arboreal]} / §5. Stability and its relation to arboreal Galois representations. \href{https://www.numdam.org/articles/10.5802/pmb.a-154/}{Source}.\par
\end{sources}\color{black}
\needspace{11\baselineskip}\section*{Questions and research boundaries}
\needspace{7\baselineskip}\subsection*{Does composition preserve irreducibility of its two inputs?}
{\small\sffamily\color{atlasgold}Settled in stated scope / Nonconstant polynomials over Q}\par\smallskip
\color{black} No. The explicit pair \(x^2-8\), \(x^2-3\) above is a counterexample. Capelli identifies the missing condition: the inner polynomial minus a root of the outer one must remain irreducible over that root field.\par
\begin{sources}
\textbf{Sources for this section}\par\smallskip
\noindent\hyperlink{ref:AyadMcQuillan2000}{[AyadMcQuillan2000]} / Lemma 1. The exact criterion, applied here to an independently expanded example. \href{https://doi.org/10.4064/aa-93-1-87-97}{Source}.\par
\end{sources}\color{black}
\needspace{6\baselineskip}\par\medskip\begingroup\footnotesize\setlength{\parskip}{1pt}\textbf{\sffamily\color{atlasgreen}Connections}\par
\color{black}\textbf{Uses.} \hyperlink{article:foundations}{Fields, factors, and the action on roots}. The tower formula turns polynomial irreducibility into an equality of extension degrees.\par
\color{black}\textbf{Applies to.} \hyperlink{article:composition-iteration}{Composition, iteration, and arboreal groups}. Stability is a prerequisite for transitivity at every level of an arboreal action.\par
\color{black}\textbf{Supplies examples for.} \hyperlink{article:newton-polygons}{Newton polygons, factor degrees, and large cycles}. The sextic x\ensuremath{{}^6}\ensuremath{-}72 has complementary Capelli and two-prime Newton-polygon proofs.\par
\color{black}\textbf{Uses.} \hyperlink{article:discriminants}{Discriminants, parity, and ramification}. The sextic example has nonsquare discriminant but a group of order only 12.\par
\color{black}\textbf{Applies to.} \hyperlink{article:integer-values}{Small values, unit values, and prime values}. Arithmetic restrictions on values can establish irreducibility of other composed families.\par
\color{black}\textbf{Contrasts with.} \hyperlink{article:finite-field-irreducibility}{Irreducibility over finite fields}. Finite fields have separate effective tests that incorporate the field size and Frobenius.\par
\endgroup
\part{Galois group detection}
\chapter{Frobenius: from factor degrees to permutations}\label{article:frobenius}\hypertarget{article:frobenius}{}\color{black}
{\sffamily\small\color{atlasgold}method / Intermediate / Q and number fields / Classical theory}\par\medskip
A factorization over a finite field reveals a cycle type in a global Galois group. Combine these local witnesses with a separate group-theoretic argument.\par\medskip
\textbf{Questions answered}\begin{itemize}
\item What does a squarefree reduction say about the Galois group?
\item When can irreducibility be witnessed by one prime?
\item How do two primes prove that a quintic has group S5?
\end{itemize}
{\small\textbf{Prerequisites:} \hyperlink{article:foundations}{Fields, factors, and the action on roots}; \hyperlink{article:modular-irreducibility}{Irreducibility by reduction modulo primes}}\par
\needspace{8\baselineskip}\section{The theorem and its hypotheses}

\begin{atlasblock}[standard]
\textbf{\sffamily Theorem  /  Dedekind-Frobenius}\par\smallskip
\color{black} Let \(f\in\mathbb Z[x]\) be monic and separable, and let \(G\) be the Galois group of its splitting field over \(\mathbb Q\), acting on its \(n\) roots. If \(p\nmid\operatorname{disc}(f)\) and\par

\atlasdisplay{\bar f=f_1\cdots f_r\quad\text{in }\mathbb F_p[x]}
\color{black} is its factorization into distinct monic irreducibles, then \(G\) contains a permutation whose cycle lengths are \(\deg f_1,\ldots,\deg f_r\).\par
\end{atlasblock}\color{black}
\color{black} Irreducibility of \(f\) is not needed for this cycle statement. It is needed when we subsequently infer that \(G\) acts transitively. For nonmonic integral polynomials, also exclude primes dividing the leading coefficient.\par
\begin{sources}
\textbf{Sources for this section}\par\smallskip
\noindent\hyperlink{ref:ConradModP}{[ConradModP]} / Theorem 1 and its proof. The note treats irreducible monic polynomials; the same reduction argument gives the separable form stated here. \href{https://kconrad.math.uconn.edu/blurbs/gradnumthy/galois-Q-factor-mod-p.pdf}{Source}.\par
\noindent\hyperlink{ref:Milne2022}{[Milne2022]} / Chapter 4, ``Computing Galois groups over Q''.  \href{https://www.jmilne.org/math/Books/index.html}{Author's book page}.\par
\end{sources}\color{black}
\needspace{8\baselineskip}\section{Why factor degrees become cycle lengths}
\color{black} Choose a prime \(\mathfrak p\) of the splitting field above \(p\). At a good prime the roots remain distinct after reduction. Frobenius acts on the residue field by \(z\mapsto z^p\). A root of an irreducible degree-\(d\) factor has an orbit of length exactly \(d\) under this map. Lifting that action to the roots gives the required permutation.\par
\color{black} A different choice of \(\mathfrak p\) conjugates Frobenius. Thus a factorization determines a cycle type, not a labeling of the roots. Permutations obtained from different primes must not be assigned compatible labels without additional information.\par
\begin{sources}
\textbf{Sources for this section}\par\smallskip
\noindent\hyperlink{ref:ConradModP}{[ConradModP]} / Proof of Theorem 1.  \href{https://kconrad.math.uconn.edu/blurbs/gradnumthy/galois-Q-factor-mod-p.pdf}{Source}.\par
\end{sources}\color{black}
\needspace{8\baselineskip}\section{A complete quintic certificate}

\begin{atlasblock}[standard]
\textbf{\sffamily Example  /  two primes give \(S_5\)}\par\smallskip
\color{black} Take \(f=x^5-x-1\). Modulo \(3\), it has no root. The only monic irreducible quadratics over \(\mathbb F_3\) are\par

\atlasdisplay{x^2+1,\quad x^2+x+2,\quad x^2+2x+2.}
\color{black} The remainders on division of \(f\) by these are \(-1,x-1,x-1\), respectively. Every reducible quintic has a factor of degree at most two, so \(\bar f\) is irreducible over \(\mathbb F_3\). Hence \(f\) is irreducible over \(\mathbb Q\), and its group contains a 5-cycle.\par
\color{black} Modulo \(2\),\par

\atlasdisplay{f=(x^2+x+1)(x^3+x^2+1),}
\color{black} with distinct irreducible factors. This gives an element of type \((2,3)\), whose cube is a transposition.\par
\color{black} A transitive action of prime degree is primitive. A primitive group containing a transposition is the full symmetric group, as proved in the permutation-group entry. Therefore \(\operatorname{Gal}(f/\mathbb Q)=S_5\). All finite-field arithmetic in this example is exact; no numerical root approximation is used.\par
\end{atlasblock}\color{black}
\begin{sources}
\textbf{Sources for this section}\par\smallskip
\noindent\hyperlink{ref:ConradSnAn}{[ConradSnAn]} / Theorem 2.1. Recognition criterion used after the explicit finite-field computations. \href{https://kconrad.math.uconn.edu/blurbs/galoistheory/galoisSnAn.pdf}{Source}.\par
\noindent\hyperlink{ref:Selmer1956}{[Selmer1956]}. Classical irreducibility of the broader trinomial family; the present degree-five proof is self-contained. \href{https://doi.org/10.7146/math.scand.a-10478}{Source}.\par
\end{sources}\color{black}
\needspace{8\baselineskip}\section{Which factorization patterns eventually appear?}
\color{black} Chebotarev gives a converse to the existence part: every conjugacy class of \(G\) occurs at infinitely many unramified primes. The density for a class \(C\) is \(|C|/|G|\); a cycle type may comprise several classes of \(G\), in which case their densities add.\par
\color{black} Consequently, an irreducible polynomial of degree \(n\) has an irreducible reduction at some good prime if and only if \(G\) contains an \(n\)-cycle. Transitivity alone does not guarantee this. For \(x^4+1\), the group is the Klein four-group and no good prime yields an irreducible quartic reduction.\par
\begin{sources}
\textbf{Sources for this section}\par\smallskip
\noindent\hyperlink{ref:StevenhagenLenstra1996}{[StevenhagenLenstra1996]}. Chebotarev density and Frobenius interpretation. \href{https://pub.math.leidenuniv.nl/~lenstrahw/PUBLICATIONS/1996d/art.pdf}{Source}.\par
\noindent\hyperlink{ref:ConradAllPrimes}{[ConradAllPrimes]}.  \href{https://kconrad.math.uconn.edu/blurbs/ringtheory/reducibleallp.pdf}{Source}.\par
\end{sources}\color{black}
\needspace{8\baselineskip}\section{What a finite prime search cannot prove}

\begin{atlasblock}[boundary]
\textbf{\sffamily Boundary}\par\smallskip
\color{black} Repeated factors invalidate this simple Frobenius reading at that prime. An observed cycle type is a valid witness; an unobserved type is not evidence of impossibility. A finite list of factorizations becomes a proof only when a theorem or a complete candidate-group argument shows that every group compatible with the witnesses is the claimed group.\par
\end{atlasblock}\color{black}
\color{black} When reduction is repeatedly bad, local ramification and Newton polygons may supply the missing permutations. When several candidate groups remain, resolvents provide a different kind of constraint.\par
\begin{sources}
\textbf{Sources for this section}\par\smallskip
\noindent\hyperlink{ref:ConradSnAn}{[ConradSnAn]} / Section 1.  \href{https://kconrad.math.uconn.edu/blurbs/galoistheory/galoisSnAn.pdf}{Source}.\par
\end{sources}\color{black}
\needspace{11\baselineskip}\section*{Questions and research boundaries}
\needspace{7\baselineskip}\subsection*{Does every irreducible rational polynomial remain irreducible modulo some prime?}
{\small\sffamily\color{atlasgold}Settled in stated scope / Monic irreducible polynomials over Q}\par\smallskip
\color{black} No. The exact criterion is that the Galois group contain a full cycle. The polynomial \(x^4+1\) supplies a counterexample.\par
\begin{sources}
\textbf{Sources for this section}\par\smallskip
\noindent\hyperlink{ref:ConradAllPrimes}{[ConradAllPrimes]}.  \href{https://kconrad.math.uconn.edu/blurbs/ringtheory/reducibleallp.pdf}{Source}.\par
\noindent\hyperlink{ref:StevenhagenLenstra1996}{[StevenhagenLenstra1996]}.  \href{https://pub.math.leidenuniv.nl/~lenstrahw/PUBLICATIONS/1996d/art.pdf}{Source}.\par
\end{sources}\color{black}
\needspace{6\baselineskip}\par\medskip\begingroup\footnotesize\setlength{\parskip}{1pt}\textbf{\sffamily\color{atlasgreen}Connections}\par
\color{black}\textbf{Uses.} \hyperlink{article:permutation-groups}{The group-theoretic step: transitivity, alternating groups, and symmetric groups}. Cycle types need a group-theoretic recognition theorem.\par
\color{black}\textbf{Uses.} \hyperlink{article:discriminants}{Discriminants, parity, and ramification}. The discriminant identifies good primes and the parity character.\par
\color{black}\textbf{Contrasts with.} \hyperlink{article:newton-polygons}{Newton polygons, factor degrees, and large cycles}. Ramified local information can remain useful when squarefree reduction fails.\par
\color{black}\textbf{Applies to.} \hyperlink{article:galois-computation}{Computing a Galois group with a proof certificate}. Modular witnesses constrain the groups examined by a certified computation.\par
\endgroup
\chapter{Discriminants, parity, and ramification}\label{article:discriminants}\hypertarget{article:discriminants}{}\color{black}
{\sffamily\small\color{atlasgold}method / Intermediate / Fields of characteristic not 2; Q / Classical theory}\par\medskip
The discriminant detects collisions of roots and the sign character of a Galois group. Its square class is powerful, but it does not measure the size of the group.\par\medskip
\textbf{Questions answered}\begin{itemize}
\item Why does a square discriminant put the group inside A\_\allowbreak{}n?
\item What distinguishes polynomial and number-field discriminants?
\item When does parity finish a Galois-group proof?
\end{itemize}
{\small\textbf{Prerequisites:} \hyperlink{article:foundations}{Fields, factors, and the action on roots}}\par
\needspace{8\baselineskip}\section{Two formulas for the same invariant}
\color{black} For \(f(x)=a_n\prod_{i=1}^n(x-\alpha_i)\) of degree \(n\), use the standard polynomial discriminant\par

\atlasdisplay{\Delta(f)=a_n^{2n-2}\prod_{i<j}(\alpha_i-\alpha_j)^2=(-1)^{n(n-1)/2}a_n^{-1}\operatorname{Res}(f,f').}
\color{black} It is nonzero exactly when the roots are distinct. The resultant formula computes it from coefficients without computing the roots. Multiplying \(f\) by a nonzero scalar changes its discriminant by an even power, so its square class is unchanged.\par
\begin{sources}
\textbf{Sources for this section}\par\smallskip
\noindent\hyperlink{ref:Milne2022}{[Milne2022]} / Chapter 4 and its appendix on resultants.  \href{https://www.jmilne.org/math/Books/index.html}{Author's book page}.\par
\end{sources}\color{black}
\needspace{8\baselineskip}\section{The sign character, proved directly}

\begin{atlasblock}[standard]
\textbf{\sffamily Theorem  /  discriminant criterion}\par\smallskip
\color{black} If \(f\in K[x]\) is separable and \(\operatorname{char}K\ne2\), then\par

\atlasdisplay{\operatorname{Gal}(f/K)\subseteq A_n\quad\Longleftrightarrow\quad\Delta(f)\in K^{\times2}.}
\end{atlasblock}\color{black}

\begin{atlasblock}[standard]
\textbf{\sffamily Proof mechanism}\par\smallskip
\color{black} Set \(D=a_n^{n-1}\prod_{i<j}(\alpha_i-\alpha_j)\), so \(D^2=\Delta(f)\). An automorphism \(\sigma\) sends \(D\) to \(\operatorname{sgn}(\sigma)D\). Thus all automorphisms are even exactly when \(D\) belongs to the ground field. Conversely, if \(D^2=b^2\) with \(b\in K\), then \(D=\pm b\), so \(D\in K\). Characteristic two is excluded because the signs coincide.\par
\end{atlasblock}\color{black}
\color{black} If the group is not contained in \(A_n\), the quadratic field \(K(\sqrt{\Delta(f)})\) is the fixed field of the kernel of the sign character.\par
\begin{sources}
\textbf{Sources for this section}\par\smallskip
\noindent\hyperlink{ref:ConradCubicQuartic}{[ConradCubicQuartic]} / Theorem 1.3.  \href{https://kconrad.math.uconn.edu/blurbs/galoistheory/cubicquartic.pdf}{Source}.\par
\noindent\hyperlink{ref:Milne2022}{[Milne2022]} / Chapter 4, ``When is G\_\allowbreak{}f contained in A\_\allowbreak{}n?''.  \href{https://www.jmilne.org/math/Books/index.html}{Author's book page}.\par
\end{sources}\color{black}
\needspace{8\baselineskip}\section{The complete cubic case}

\begin{atlasblock}[standard]
\textbf{\sffamily Example  /  parity is enough in degree three}\par\smallskip
\color{black} A separable irreducible cubic has a transitive subgroup of \(S_3\), hence group \(A_3\) or \(S_3\). For \(x^3+ax+b\), the discriminant is \(-4a^3-27b^2\). Therefore\par

\atlasdisplay{\begin{array}{c|c|c} f & \Delta(f) & \operatorname{Gal}(f/\mathbb Q)\\ \hline x^3-x-1&-23&S_3\\ x^3-3x-1&81&C_3\end{array}}
\color{black} The rational-root test proves both polynomials irreducible: neither \(1\) nor \(-1\) is a root. The table would not follow from the discriminants without this irreducibility check.\par
\end{atlasblock}\color{black}
\begin{sources}
\textbf{Sources for this section}\par\smallskip
\noindent\hyperlink{ref:ConradCubicQuartic}{[ConradCubicQuartic]} / Theorem 2.1 and Example 2.2.  \href{https://kconrad.math.uconn.edu/blurbs/galoistheory/cubicquartic.pdf}{Source}.\par
\end{sources}\color{black}
\needspace{8\baselineskip}\section{Why neither square class determines a large group}
\color{black} The polynomial \(x^4+1\) is irreducible over \(\mathbb Q\), since replacing \(x\) by \(x+1\) makes it Eisenstein at \(2\). Its discriminant is \(256\), but its splitting field is \(\mathbb Q(\zeta_8)\), whose group is \((\mathbb Z/8\mathbb Z)^\times\simeq C_2\times C_2\), strictly smaller than \(A_4\).\par
\color{black} In the other direction, \(x^4-2\) is Eisenstein at \(2\) and has nonsquare discriminant \(-2048\). Its splitting field is \(\mathbb Q(\sqrt[4]{2},i)\), of degree eight: the real field \(\mathbb Q(\sqrt[4]{2})\) has degree four and does not contain \(i\). The group is dihedral of order eight, not \(S_4\).\par

\begin{atlasblock}[boundary]
\textbf{\sffamily Correct final step}\par\smallskip
\color{black} First prove \(A_n\subseteq G\). Then a square discriminant gives \(G=A_n\), and a nonsquare discriminant gives \(G=S_n\). Without the lower containment, the discriminant only gives a parity restriction.\par
\end{atlasblock}\color{black}
\begin{sources}
\textbf{Sources for this section}\par\smallskip
\noindent\hyperlink{ref:Cox2012}{[Cox2012]}. Cyclotomic and radical extensions; discriminant interpretation. \href{https://onlinelibrary.wiley.com/doi/book/10.1002/9781118218457}{Source}.\par
\noindent\hyperlink{ref:ConradCubicQuartic}{[ConradCubicQuartic]} / Section 3.  \href{https://kconrad.math.uconn.edu/blurbs/galoistheory/cubicquartic.pdf}{Source}.\par
\end{sources}\color{black}
\needspace{8\baselineskip}\section{Polynomial discriminant versus field discriminant}
\color{black} If \(f\) is the monic minimal polynomial of an algebraic integer \(\alpha\), and \(K=\mathbb Q(\alpha)\), then\par

\atlasdisplay{\Delta(f)=\operatorname{disc}(K)\,[\mathcal O_K:\mathbb Z[\alpha]]^2.}
\color{black} The two discriminants have the same square class, but can have different prime divisors. A prime dividing \(\Delta(f)\) may come from the index rather than ramification in \(K\).\par
\color{black} For \(f=x^2-5\), the polynomial discriminant is \(20\), whereas \(\operatorname{disc}(\mathbb Q(\sqrt5))=5\). Reduction modulo \(2\) is repeated, although \(2\) is unramified in the quadratic field. The integral generator \((1+\sqrt5)/2\), with polynomial \(x^2-x-1\), removes this index obstruction.\par
\color{black} Over \(\mathbb Q\), complex conjugation contributes another parity check: with \(r\) nonreal conjugate pairs it has type \(2^r1^{n-2r}\), and \(\operatorname{sgn}\Delta(f)=(-1)^r\).\par
\begin{sources}
\textbf{Sources for this section}\par\smallskip
\noindent\hyperlink{ref:Neukirch1999}{[Neukirch1999]}. Discriminants of orders, indices, and ramification. \href{https://link.springer.com/book/10.1007/978-3-662-03983-0}{Source}.\par
\noindent\hyperlink{ref:ConradCubicQuartic}{[ConradCubicQuartic]} / Remark 3.11 and the discussion of real roots.  \href{https://kconrad.math.uconn.edu/blurbs/galoistheory/cubicquartic.pdf}{Source}.\par
\end{sources}\color{black}
\needspace{11\baselineskip}\section*{Questions and research boundaries}
\needspace{7\baselineskip}\subsection*{When can a discriminant determine the full Galois group?}
{\small\sffamily\color{atlasgold}Settled in stated scope / Separable irreducible cubics; groups already known to contain \(A_n\).}\par\smallskip
\color{black} It determines the group for irreducible cubics, and whenever an independent argument has proved \(A_n\subseteq G\). In general it determines only whether \(G\subseteq A_n\).\par
\begin{sources}
\textbf{Sources for this section}\par\smallskip
\noindent\hyperlink{ref:ConradCubicQuartic}{[ConradCubicQuartic]} / Theorems 1.3 and 2.1.  \href{https://kconrad.math.uconn.edu/blurbs/galoistheory/cubicquartic.pdf}{Source}.\par
\end{sources}\color{black}
\needspace{6\baselineskip}\par\medskip\begingroup\footnotesize\setlength{\parskip}{1pt}\textbf{\sffamily\color{atlasgreen}Connections}\par
\color{black}\textbf{Applies to.} \hyperlink{article:frobenius}{Frobenius: from factor degrees to permutations}. The polynomial discriminant specifies primes where the elementary cycle-type argument applies.\par
\color{black}\textbf{Applies to.} \hyperlink{article:resolvents}{Resolvents: turn group containment into factorization}. In degree four, the cubic resolvent adds the information parity alone lacks.\par
\color{black}\textbf{Uses.} \hyperlink{article:permutation-groups}{The group-theoretic step: transitivity, alternating groups, and symmetric groups}. Large-group containment plus one odd permutation proves equality with S\_\allowbreak{}n.\par
\color{black}\textbf{Applies to.} \hyperlink{article:random-models}{Random polynomials: models and proof strategies}. A large-group theorem and a nonsquare-discriminant theorem address distinct probabilistic obstructions.\par
\endgroup
\chapter{Ramification and inertia: cycles at bad primes}\label{article:ramification-inertia}\hypertarget{article:ramification-inertia}{}\color{black}
{\sffamily\small\color{atlasgold}method / Intermediate / Q and local fields / Classical theory}\par\medskip
Good primes reveal Frobenius cycle types. At ramified primes, inertia supplies different information: for a monic integral polynomial, an odd prime dividing its discriminant exactly once forces a transposition.\par\medskip
\textbf{Questions answered}\begin{itemize}
\item What replaces Frobenius at a ramified prime?
\item When does a discriminant valuation force a transposition?
\item Why can repeated factors modulo a prime be misleading?
\end{itemize}
{\small\textbf{Prerequisites:} \hyperlink{article:foundations}{Fields, factors, and the action on roots}; \hyperlink{article:discriminants}{Discriminants, parity, and ramification}; \hyperlink{article:frobenius}{Frobenius: from factor degrees to permutations}}\par
\needspace{8\baselineskip}\section{Two local subgroups inside the global group}
\color{black} Let \(L/\mathbb Q\) be a finite Galois extension with group \(G\), and choose a prime \(\mathfrak P\) above \(p\). Its \textbf{decomposition group} \(D_{\mathfrak P}\) consists of automorphisms preserving \(\mathfrak P\). The \textbf{inertia group} \(I_{\mathfrak P}\) consists of those acting trivially on the residue field \(k_{\mathfrak P}=\mathcal O_L/\mathfrak P\). There is an exact sequence\par

\atlasdisplay{1\longrightarrow I_{\mathfrak P}\longrightarrow D_{\mathfrak P}\longrightarrow\operatorname{Gal}(k_{\mathfrak P}/\mathbb F_p)\longrightarrow1.}
\color{black} Also \(D_{\mathfrak P}\cong\operatorname{Gal}(L_{\mathfrak P}/\mathbb Q_p)\). Thus a local calculation produces elements of the global group. When \(p\) is unramified, inertia is trivial and residue Frobenius lifts uniquely inside \(D_{\mathfrak P}\). At a ramified prime it determines a coset modulo inertia. Inertia fixes the reductions of integral roots, while possibly permuting roots that collide.\par
\begin{sources}
\textbf{Sources for this section}\par\smallskip
\noindent\hyperlink{ref:Milne2020ANT}{[Milne2020ANT]} / Theorem 7.58; Proposition 8.10. Exact sequence, maximal unramified subfield, and identification with the local Galois group. \href{https://www.jmilne.org/math/CourseNotes/ANT.pdf}{Source}.\par
\end{sources}\color{black}
\needspace{8\baselineskip}\section{Tame inertia and its cycle lengths}
\color{black} In residue characteristic \(p\), inertia has a normal wild \(p\)-subgroup and a cyclic quotient of order prime to \(p\). If the extension is \textbf{tame}, wild inertia is trivial, so inertia itself is cyclic.\par
\color{black} For irreducible \(f\), put \(K=\mathbb Q(\alpha)\). If its splitting field is tame at \(p\), write \(p\mathcal O_K=\prod_i\mathfrak p_i^{e_i}\), with residue degrees \(f_i\). A generator of inertia acts on the roots with an \(e_i\)-cycle repeated \(f_i\) times for each \(i\). Indeed, the decomposition orbit for \(\mathfrak p_i\) has size \(e_if_i\), and splits into \(f_i\) inertia orbits of size \(e_i\); a cyclic generator traverses each orbit. These are ramification indices in the ring of integers, not automatically exponents in \(f\bmod p\).\par
\color{black} For example, \(x^3-5\) is Eisenstein over \(\mathbb Q_5\). Its root has valuation \(1/3\), forcing ramification index three. Adjoining a primitive cube root of unity is unramified at \(5\), since \(x^2+x+1\) has squarefree reduction. The splitting field therefore has cyclic inertia of order three, acting as a three-cycle. Here \(\Delta(x^3-5)=-675\) has valuation two at \(5\): tame ramification yields useful cycles beyond the valuation-one criterion.\par
\begin{sources}
\textbf{Sources for this section}\par\smallskip
\noindent\hyperlink{ref:Serre1979}{[Serre1979]} / Ramification Groups, pp. 61-79. Background for tame and wild inertia; the cycle description is the orbit consequence explained here. \href{https://link.springer.com/book/10.1007/978-1-4757-5673-9}{Source}.\par
\noindent\hyperlink{ref:Milne2020ANT}{[Milne2020ANT]} / Corollary 7.59 and the following paragraph. Cyclicity of the tame quotient and the criterion for tame ramification. \href{https://www.jmilne.org/math/CourseNotes/ANT.pdf}{Source}.\par
\end{sources}\color{black}
\needspace{8\baselineskip}\section{A prime appearing once gives a transposition}

\begin{atlasblock}[standard]
\textbf{\sffamily Theorem  /  simple discriminant prime}\par\smallskip
\color{black} Let \(f\in\mathbb Z[x]\) be monic and separable, with splitting field \(L\). If \(p\) is odd and \(v_p(\Delta(f))=1\), then inertia at every prime of \(L\) above \(p\) has order two, and its nonidentity element acts as a transposition on the roots. Irreducibility is not required.\par
\end{atlasblock}\color{black}

\begin{atlasblock}[standard]
\textbf{\sffamily Proof  /  isolate one ramified quadratic factor}\par\smallskip
\color{black} Multiplication by \(f'\) on the free \(\mathbb Z_p\)-module \(\mathbb Z_p[x]/(f)\) has determinant \(\operatorname{Res}(f,f')=\pm\Delta(f)\). Valuation one means its Smith normal form has exactly one nonunit diagonal entry, so reduction modulo \(p\) has cokernel dimension one. This cokernel is \(\mathbb F_p[x]/(\bar f,\bar f')\); hence \(\deg\gcd(\bar f,\bar f')=1\).\par
\color{black} Since \(p\) is odd, necessarily \(\bar f=(x-a)^2\bar h\), where \(a\in\mathbb F_p\), \(\bar h\) is squarefree, and \(\bar h(a)\ne0\). Hensel's lemma lifts these coprime factors to \(f=qh\) over \(\mathbb Z_p\), with \(q\) quadratic. The identity\par

\atlasdisplay{\Delta(f)=\Delta(q)\Delta(h)\operatorname{Res}(q,h)^2}
\color{black} and the unit factors \(\Delta(h)\) and \(\operatorname{Res}(q,h)\) give \(v_p(\Delta(q))=1\). The splitting field \(\mathbb Q_p(\sqrt{\Delta(q)})\) is ramified quadratic; the squarefree reduction of \(h\) makes its splitting field unramified. Their compositum therefore has inertia of order two, swapping precisely the two roots of \(q\).\par
\end{atlasblock}\color{black}
\begin{sources}
\textbf{Sources for this section}\par\smallskip
\noindent\hyperlink{ref:Milne2020ANT}{[Milne2020ANT]} / Theorem 7.33; Proposition 7.50. Hensel lifting and unramified extensions used in the explicit proof; the valuation-one deduction is supplied here. \href{https://www.jmilne.org/math/CourseNotes/ANT.pdf}{Source}.\par
\noindent\hyperlink{ref:Serre1979}{[Serre1979]} / Discriminant and Different, pp. 47-60; Ramification Groups, pp. 61-79. Discriminants and local ramification framework. \href{https://link.springer.com/book/10.1007/978-1-4757-5673-9}{Source}.\par
\end{sources}\color{black}
\needspace{8\baselineskip}\section{An exact cubic certificate}

\begin{atlasblock}[standard]
\textbf{\sffamily Example  /  \(x^3-x-1\)}\par\smallskip
\color{black} The rational-root test proves \(f=x^3-x-1\) irreducible: neither \(1\) nor \(-1\) is a root. Its discriminant and bad-prime reduction are\par

\atlasdisplay{\Delta(f)=-4(-1)^3-27(-1)^2=-23,\qquad f\equiv(x-10)^2(x-3)\pmod{23}.}
\color{black} The latter identity is certified by \(f-(x-10)^2(x-3)=23(x^2-7x+13)\). The theorem supplies an inertia transposition at \(23\). Since \(G\le S_3\) is transitive, \(3\mid |G|\); the transposition gives \(2\mid |G|\), so \(G=S_3\). The valuation-one hypothesis, rather than the displayed repeated factor alone, certifies the ramification step.\par
\end{atlasblock}\color{black}
\begin{sources}
\textbf{Sources for this section}\par\smallskip
\noindent\hyperlink{ref:Serre1979}{[Serre1979]} / Ramification Groups, pp. 61-79. Local-group background; all numerical identities and the elementary group argument are verified directly here. \href{https://link.springer.com/book/10.1007/978-1-4757-5673-9}{Source}.\par
\end{sources}\color{black}
\needspace{8\baselineskip}\section{Polynomial discriminants can record a nonmaximal order}

\begin{atlasblock}[boundary]
\textbf{\sffamily Boundary  /  repeated reduction need not mean ramification}\par\smallskip
\color{black} For monic irreducible \(f\) with root \(\alpha\) and \(K=\mathbb Q(\alpha)\),\par

\atlasdisplay{\Delta(f)=\operatorname{disc}(K)\,[\mathcal O_K:\mathbb Z[\alpha]]^2.}
\color{black} Consider \(x^2-18\), with root \(3\sqrt2\). Modulo \(3\) it becomes \(x^2\), yet \(K=\mathbb Q(\sqrt2)\) has discriminant \(8\), so \(3\) is unramified and its inertia group is trivial. The polynomial discriminant is \(72=8\cdot3^2\): the repeated reduction comes from the index \(3\). Conversely, valuation one excludes any index contribution at that prime.\par
\end{atlasblock}\color{black}
\color{black} To force \(S_n\) in higher degree, a transposition still needs additional group information, such as primitivity. Transitivity alone does not suffice.\par
\begin{sources}
\textbf{Sources for this section}\par\smallskip
\noindent\hyperlink{ref:Milne2020ANT}{[Milne2020ANT]} / Remark 2.25; Theorem 3.35. Index-squared discriminant identity and the ramification-discriminant criterion. \href{https://www.jmilne.org/math/CourseNotes/ANT.pdf}{Source}.\par
\end{sources}\color{black}
\needspace{11\baselineskip}\section*{Questions and research boundaries}
\needspace{7\baselineskip}\subsection*{Does a repeated root modulo p certify nontrivial inertia?}
{\small\sffamily\color{atlasgold}Settled in stated scope / Monic integral polynomials over Q}\par\smallskip
\color{black} No: \(x^2-18\) modulo \(3\) is a counterexample, because its root field is unramified at \(3\). Control the order index or use a stronger condition such as the valuation-one theorem.\par
\begin{sources}
\textbf{Sources for this section}\par\smallskip
\noindent\hyperlink{ref:Milne2020ANT}{[Milne2020ANT]} / Remark 2.25. The index-squared discriminant identity explains the obstruction. \href{https://www.jmilne.org/math/CourseNotes/ANT.pdf}{Source}.\par
\end{sources}\color{black}
\needspace{6\baselineskip}\par\medskip\begingroup\footnotesize\setlength{\parskip}{1pt}\textbf{\sffamily\color{atlasgreen}Connections}\par
\color{black}\textbf{Contrasts with.} \hyperlink{article:frobenius}{Frobenius: from factor degrees to permutations}. Squarefree good-prime factorization records Frobenius; repeated reduction requires inertia and index analysis.\par
\color{black}\textbf{Uses.} \hyperlink{article:discriminants}{Discriminants, parity, and ramification}. The valuation of the polynomial discriminant can isolate a ramified quadratic factor.\par
\color{black}\textbf{Applies to.} \hyperlink{article:permutation-groups}{The group-theoretic step: transitivity, alternating groups, and symmetric groups}. A certified inertia transposition combines with primitivity to force the symmetric group.\par
\color{black}\textbf{Contrasts with.} \hyperlink{article:newton-polygons}{Newton polygons, factor degrees, and large cycles}. Newton polygons constrain local factor degrees; inertia records root permutations produced by ramification.\par
\endgroup
\chapter{The group-theoretic step: transitivity, alternating groups, and symmetric groups}\label{article:permutation-groups}\hypertarget{article:permutation-groups}{}\color{black}
{\sffamily\small\color{atlasgold}method / Intermediate / Permutation groups; polynomial Galois groups / Classical theory}\par\medskip
Local arithmetic supplies permutations. Block systems, transpositions, and long cycles explain when those permutations force a large group.\par\medskip
\textbf{Questions answered}\begin{itemize}
\item How is primitivity proved rather than assumed?
\item Why does a primitive group with a transposition equal S\_\allowbreak{}n?
\item Which single cycles force containment of A\_\allowbreak{}n?
\end{itemize}
{\small\textbf{Prerequisites:} \hyperlink{article:foundations}{Fields, factors, and the action on roots}; \hyperlink{article:frobenius}{Frobenius: from factor degrees to permutations}}\par
\needspace{8\baselineskip}\section{Transitivity, blocks, and primitivity}
\color{black} Let \(G\le S_n\) act on \(\Omega\). A nonempty subset \(B\subseteq\Omega\) is a block if each translate \(gB\) is either equal to \(B\) or disjoint from it. A transitive action is primitive if its only blocks have size \(1\) or \(n\).\par
\color{black} In a transitive action, translates of a block partition \(\Omega\) into equally sized pieces. Hence a transitive action of prime degree is primitive. For an irreducible separable polynomial with root \(\alpha\), the action is primitive exactly when the root field \(K(\alpha)\) has no proper intermediate field between \(K\) and \(K(\alpha)\). This follows by applying the Galois correspondence to subgroups between a root stabilizer and the whole group.\par
\color{black} This criterion concerns intermediate fields. For a polynomial map, decomposition of \(f(X)\) is instead related to the monodromy of \(f(X)-t\) over a function field. These two uses of ``primitive'' must not be interchanged.\par
\begin{sources}
\textbf{Sources for this section}\par\smallskip
\noindent\hyperlink{ref:DixonMortimer1996}{[DixonMortimer1996]}. Blocks, stabilizers, and primitive actions. \href{https://link.springer.com/book/10.1007/978-1-4612-0731-3}{Source}.\par
\noindent\hyperlink{ref:Cox2012}{[Cox2012]}. Galois correspondence and the action on conjugates. \href{https://onlinelibrary.wiley.com/doi/book/10.1002/9781118218457}{Source}.\par
\end{sources}\color{black}
\needspace{8\baselineskip}\section{A short proof using a graph}

\begin{atlasblock}[standard]
\textbf{\sffamily Theorem}\par\smallskip
\color{black} A primitive subgroup \(G\le S_n\) containing a transposition is \(S_n\).\par
\end{atlasblock}\color{black}

\begin{atlasblock}[standard]
\textbf{\sffamily Proof}\par\smallskip
\color{black} Make a graph on \(\Omega\), joining \(a\) to \(b\) when \((a\ b)\in G\). Conjugation by \(G\) preserves this graph, so its connected components form a block system. There is an edge, and primitivity forces the graph to be connected. The transpositions associated with the edges of a connected graph generate every transposition: along a path, conjugating successive edge transpositions moves an endpoint step by step. They therefore generate \(S_n\), which is contained in \(G\).\par
\end{atlasblock}\color{black}
\color{black} Primitivity cannot be replaced by transitivity. The square-symmetry group \(\langle(1234),(13)\rangle\) has order eight and contains a transposition, while preserving the block system \(\{\{1,3\},\{2,4\}\}\). It is a proper subgroup of \(S_4\).\par
\begin{sources}
\textbf{Sources for this section}\par\smallskip
\noindent\hyperlink{ref:DixonMortimer1996}{[DixonMortimer1996]}. Primitive permutation groups and generation by transpositions. \href{https://link.springer.com/book/10.1007/978-1-4612-0731-3}{Source}.\par
\noindent\hyperlink{ref:ConradSnAn}{[ConradSnAn]} / Theorem 2.1. A useful sufficient condition for the required primitivity. \href{https://kconrad.math.uconn.edu/blurbs/galoistheory/galoisSnAn.pdf}{Source}.\par
\end{sources}\color{black}
\needspace{8\baselineskip}\section{A long prime cycle proves primitivity}
\color{black} Suppose \(G\) is transitive and contains a \(q\)-cycle, with \(q\) prime and \(q>n/2\). If a nontrivial block system had \(a\) blocks of size \(b\), then \(a,b\le n/2<q\). The cycle's induced action on the blocks has order dividing \(q\); since \(a<q\), it must fix every block. Its restriction to each block is also trivial because \(b<q\). This contradicts the presence of a \(q\)-cycle. Thus \(G\) is primitive.\par
\color{black} Be careful about extracting a cycle from a permutation. If one cycle has length \(m\) coprime to every other cycle length, raising the permutation to the least common multiple of the other lengths kills those cycles and preserves an \(m\)-cycle. An element of type \((3,3)\) does not yield a single 3-cycle by this argument.\par
\begin{sources}
\textbf{Sources for this section}\par\smallskip
\noindent\hyperlink{ref:ConradSnAn}{[ConradSnAn]} / Theorem 2.1.  \href{https://kconrad.math.uconn.edu/blurbs/galoistheory/galoisSnAn.pdf}{Source}.\par
\noindent\hyperlink{ref:Jones2014Cycles}{[Jones2014Cycles]} / Remarks 1.4 and 1.6 (published version, p. 160).  \href{https://doi.org/10.1017/S000497271300049X}{Source}.\par
\end{sources}\color{black}
\needspace{8\baselineskip}\section{A cycle fixing one point gives double transitivity}

\begin{atlasblock}[standard]
\textbf{\sffamily Criterion  /  an alternative route to primitivity}\par\smallskip
\color{black} For \(n\ge3\), a transitive subgroup \(G\le S_n\) containing an \((n-1)\)-cycle is doubly transitive. If it also contains a transposition, then \(G=S_n\).\par
\end{atlasblock}\color{black}

\begin{atlasblock}[standard]
\textbf{\sffamily Proof}\par\smallskip
\color{black} If the cycle fixes \(a\), its powers act transitively on the other \(n-1\) points. Thus the stabilizer \(G_a\) is transitive there. Transitivity of \(G\) now moves any ordered pair of distinct points to any other: first move the first point, then use its stabilizer to move the second. This is double transitivity. Conjugating one transposition by these elements supplies every transposition, which generates \(S_n\).\par
\end{atlasblock}\color{black}
\color{black} No primality condition on \(n-1\) is needed. In arithmetic applications, a good reduction of factor degrees \((1,n-1)\) supplies the long cycle. Irreducibility supplies transitivity. A second squarefree pattern with one 2-cycle and all other cycles of odd lengths supplies a transposition by taking an odd power. The conditions concern cycle types, so no identification of root labels across different primes is required.\par
\begin{sources}
\textbf{Sources for this section}\par\smallskip
\noindent\hyperlink{ref:DixonMortimer1996}{[DixonMortimer1996]} / Chapter 1, The Basic Ideas. The stabilizer characterization of double transitivity; the short proof is included here. \href{https://link.springer.com/book/10.1007/978-1-4612-0731-3}{Source}.\par
\noindent\hyperlink{ref:ConradSnAn}{[ConradSnAn]} / §1 and Example 2.3. Arithmetic cycle-type witnesses and extracting a transposition by taking a power. \href{https://kconrad.math.uconn.edu/blurbs/galoistheory/galoisSnAn.pdf}{Source}.\par
\end{sources}\color{black}
\needspace{8\baselineskip}\section{Jordan and a modern strengthening}

\begin{atlasblock}[standard]
\textbf{\sffamily Jordan's prime-cycle theorem}\par\smallskip
\color{black} If a primitive subgroup \(G\le S_n\) contains a \(q\)-cycle with \(q\) prime and \(q\le n-3\), then \(A_n\subseteq G\).\par
\end{atlasblock}\color{black}
\color{black} Combining this with the preceding argument gives a frequent arithmetic criterion: transitivity plus a prime cycle with \(n/2<q\le n-3\) forces \(A_n\). Newton polygons often produce just such a prime divisor of the group order; the inequality \(q>n/2\) ensures an element of order \(q\) is one cycle.\par

\begin{atlasblock}[standard]
\textbf{\sffamily Jones's extension}\par\smallskip
\color{black} A primitive subgroup containing a nontrivial single cycle that fixes at least three points contains \(A_n\). The cycle length need not be prime.\par
\end{atlasblock}\color{black}
\color{black} This stronger result follows from the classification in Jones's paper and uses the classification of finite simple groups. The classical prime-cycle version remains useful when one wants a proof with lighter group-theoretic inputs. Neither theorem proves primitivity by itself.\par
\begin{sources}
\textbf{Sources for this section}\par\smallskip
\noindent\hyperlink{ref:Jordan1871}{[Jordan1871]}. Historical source for primitive-group criteria. \href{https://www.numdam.org/item/JMPA_1871_2_16__383_0/}{Source}.\par
\noindent\hyperlink{ref:Jones2014Cycles}{[Jones2014Cycles]} / Theorem 1.1, Theorem 1.2, and Corollary 1.3.  \href{https://doi.org/10.1017/S000497271300049X}{Source}.\par
\end{sources}\color{black}
\needspace{8\baselineskip}\section{A complete proof from a real-root count}

\begin{atlasblock}[standard]
\textbf{\sffamily Example  /  \(x^5-4x+2\)}\par\smallskip
\color{black} This polynomial is Eisenstein at \(2\), hence irreducible. Its action has prime degree five and is primitive. Put \(a=(4/5)^{1/4}\). The derivative \(5x^4-4\) vanishes only at \(-a,a\); the local maximum has value \(2+16a/5>0\), and the local minimum has value \(2-16a/5<0\), since \(a>5/8\). Monotonicity on the three intervals gives exactly three real roots.\par
\color{black} The remaining two roots form one nonreal conjugate pair. Complex conjugation is therefore a transposition. The graph theorem proves \(\operatorname{Gal}(x^5-4x+2\,/\mathbb Q)=S_5\). A plot can suggest this root count; the derivative and exact inequalities establish it.\par
\end{atlasblock}\color{black}
\begin{sources}
\textbf{Sources for this section}\par\smallskip
\noindent\hyperlink{ref:Milne2022}{[Milne2022]} / Chapter 4, ``Examples of polynomials with S\_\allowbreak{}p as Galois group over Q''.  \href{https://www.jmilne.org/math/Books/index.html}{Author's book page}.\par
\noindent\hyperlink{ref:ConradSnAn}{[ConradSnAn]} / Theorem 2.1.  \href{https://kconrad.math.uconn.edu/blurbs/galoistheory/galoisSnAn.pdf}{Source}.\par
\end{sources}\color{black}
\needspace{11\baselineskip}\section*{Questions and research boundaries}
\needspace{7\baselineskip}\subsection*{Can the prime-length assumption in Jordan's cycle criterion be removed?}
{\small\sffamily\color{atlasgold}Settled in stated scope / Primitive finite permutation groups containing a nontrivial cycle with at least three fixed points}\par\smallskip
\color{black} Yes: Jones's Corollary 1.3 gives containment of \(A_n\) without that assumption. The stronger theorem uses the classification of finite simple groups; the primitivity and fixed-point conditions remain essential.\par
\begin{sources}
\textbf{Sources for this section}\par\smallskip
\noindent\hyperlink{ref:Jones2014Cycles}{[Jones2014Cycles]} / Corollary 1.3.  \href{https://doi.org/10.1017/S000497271300049X}{Source}.\par
\end{sources}\color{black}
\needspace{6\baselineskip}\par\medskip\begingroup\footnotesize\setlength{\parskip}{1pt}\textbf{\sffamily\color{atlasgreen}Connections}\par
\color{black}\textbf{Uses.} \hyperlink{article:newton-polygons}{Newton polygons, factor degrees, and large cycles}. Slope denominators can supply the prime-cycle input.\par
\color{black}\textbf{Uses.} \hyperlink{article:discriminants}{Discriminants, parity, and ramification}. Once A\_\allowbreak{}n is contained, the discriminant decides between A\_\allowbreak{}n and S\_\allowbreak{}n.\par
\color{black}\textbf{Contrasts with.} \hyperlink{article:composition-iteration}{Composition, iteration, and arboreal groups}. Function-field monodromy has a precise connection with decomposition; fixed-polynomial root actions require a separate field test.\par
\color{black}\textbf{Contrasts with.} \hyperlink{article:reciprocity}{Reciprocal symmetry and the group on root pairs}. Inverse pairs impose a block system and change the natural ambient group.\par
\endgroup
\chapter{Resolvents: turn group containment into factorization}\label{article:resolvents}\hypertarget{article:resolvents}{}\color{black}
{\sffamily\small\color{atlasgold}method / Intermediate / Q; fields of characteristic not 2 for quartics / Classical theory}\par\medskip
A resolvent records how a group acts on selected expressions in the roots. Rational factors can detect a hidden partition or force descent to a smaller subgroup.\par\medskip
\textbf{Questions answered}\begin{itemize}
\item How does a polynomial in the roots detect a subgroup?
\item How is the cubic resolvent of a quartic constructed?
\item What distinguishes S4, A4, V4, and the C4/D4 cases?
\end{itemize}
{\small\textbf{Prerequisites:} \hyperlink{article:foundations}{Fields, factors, and the action on roots}; \hyperlink{article:discriminants}{Discriminants, parity, and ramification}}\par
\needspace{8\baselineskip}\section{A second permutation action}
\color{black} Let \(f\in K[x]\) be separable, with roots \(\alpha_1,\ldots,\alpha_n\), and the group is known to lie in \(G_0\le S_n\). Choose a polynomial \(\Phi\in K[X_1,\ldots,X_n]\) to evaluate on these roots whose formal stabilizer in \(G_0\) is a subgroup \(H\). Evaluate its orbit to obtain values \(\beta_1,\ldots,\beta_m\), where \(m=[G_0:H]\), and form\par

\atlasdisplay{R_\Phi(T)=\prod_{j=1}^m(T-\beta_j).}
\color{black} The actual Galois group permutes these values, so \(R_\Phi\) has coefficients in the ground field. If the values are distinct, its irreducible factor degrees are the orbit sizes for this new action. A root in the ground field gives a fixed coset, hence containment of the Galois group in a conjugate of \(H\).\par
\color{black} The distinctness condition prevents accidental equality between evaluated invariants from giving a false subgroup conclusion. In computations it can be checked by the discriminant of the resolvent, or addressed by changing the root expressions.\par
\begin{sources}
\textbf{Sources for this section}\par\smallskip
\noindent\hyperlink{ref:FiekerKlueners2014}{[FiekerKlueners2014]} / Section 3, Lemma 3.1 and Theorem 3.2.  \href{https://doi.org/10.1112/S1461157013000302}{Source}.\par
\noindent\hyperlink{ref:Stauduhar1973}{[Stauduhar1973]}. The subgroup-descent principle. \href{https://doi.org/10.1090/S0025-5718-1973-0327712-4}{Source}.\par
\end{sources}\color{black}
\needspace{8\baselineskip}\section{The cubic resolvent of a quartic}
\color{black} For \(f=x^4+ax^3+bx^2+cx+d\), the three partitions of its four roots into two unordered pairs give\par

\atlasdisplay{\beta_1=\alpha_1\alpha_2+\alpha_3\alpha_4,\quad\beta_2=\alpha_1\alpha_3+\alpha_2\alpha_4,\quad\beta_3=\alpha_1\alpha_4+\alpha_2\alpha_3.}
\color{black} Expanding \(\prod(T-\beta_i)\) in elementary symmetric functions yields\par

\atlasdisplay{R_f(T)=T^3-bT^2+(ac-4d)T+4bd-a^2d-c^2.}
\color{black} The action \(S_4\to S_3\) on the three pairings has kernel the normal Klein four-group. Moreover, \(\beta_1-\beta_2=(\alpha_1-\alpha_4)(\alpha_2-\alpha_3)\), with analogous formulas for the other differences. Their squared product proves \(\operatorname{disc}(R_f)=\operatorname{disc}(f)\), so this resolvent is separable whenever \(f\) is.\par
\begin{sources}
\textbf{Sources for this section}\par\smallskip
\noindent\hyperlink{ref:ConradCubicQuartic}{[ConradCubicQuartic]} / Definition 3.1, equation (3.7), and Theorem 3.4.  \href{https://kconrad.math.uconn.edu/blurbs/galoistheory/cubicquartic.pdf}{Source}.\par
\noindent\hyperlink{ref:Cox2012}{[Cox2012]}. Resolvents and their relation to permutation actions. \href{https://onlinelibrary.wiley.com/doi/book/10.1002/9781118218457}{Source}.\par
\end{sources}\color{black}
\needspace{8\baselineskip}\section{The classification and its remaining branch}
\color{black} Assume \(\operatorname{char}K\ne2\) and \(f\) is a separable irreducible quartic over \(K\). Then:\par
\begin{center}\small\begin{tabularx}{\linewidth}{@{}>{\raggedright\arraybackslash}X>{\raggedright\arraybackslash}X>{\raggedright\arraybackslash}X@{}}\toprule
\textbf{Resolvent over \(K\)} & \textbf{Discriminant of \(f\)} & \textbf{Galois group} \\
\midrule
Irreducible & Nonsquare & \(S_4\) \\
Irreducible & Square & \(A_4\) \\
Splits completely & Square & \(V_4\) \\
One linear factor and one irreducible quadratic & Nonsquare & \(C_4\) or \(D_4\) \\
\bottomrule\end{tabularx}\end{center}
\color{black} Here \(D_4\) denotes the dihedral group of order eight. The table does not distinguish its last two possibilities. If \(r\in K\) is the unique root of \(R_f\), they are separated by the following test: the group is \(C_4\) precisely when both \(T^2+aT+(b-r)\) and \(T^2-rT+d\) split over \(K(\sqrt{\Delta(f)})\); otherwise it is \(D_4\).\par
\begin{sources}
\textbf{Sources for this section}\par\smallskip
\noindent\hyperlink{ref:KappeWarren1989Quartic}{[KappeWarren1989Quartic]} / pp. 133-137. Original published quartic-group test; Conrad Theorem 4.1 gives the exact formulation used here. \href{https://www.jstor.org/stable/2323198}{Source}.\par
\noindent\hyperlink{ref:ConradCubicQuartic}{[ConradCubicQuartic]} / Theorem 3.6, Corollary 3.8, and Theorem 4.1 (Kappe-Warren).  \href{https://kconrad.math.uconn.edu/blurbs/galoistheory/cubicquartic.pdf}{Source}.\par
\end{sources}\color{black}
\needspace{8\baselineskip}\section{An exact S4 certificate}

\begin{atlasblock}[standard]
\textbf{\sffamily Example  /  \(f=x^4+2x+2\)}\par\smallskip
\color{black} Eisenstein at \(2\) proves irreducibility. The resolvent is \(R_f=T^3-8T-4\). By the rational-root test a rational root would be one of \(\pm1,\pm2,\pm4\); the six values, in the order \(-4,-2,-1,1,2,4\), are \(-36,4,3,-11,-12,28\). None is zero, so \(R_f\) is irreducible.\par
\color{black} The discriminant is\par

\atlasdisplay{256\cdot2^3-27\cdot2^4=1616=16\cdot101.}
\color{black} It is nonsquare. The first line of the classification gives \(\operatorname{Gal}(f/\mathbb Q)=S_4\). This argument verifies the candidate group using rational arithmetic and small integer tests.\par
\end{atlasblock}\color{black}
\begin{sources}
\textbf{Sources for this section}\par\smallskip
\noindent\hyperlink{ref:ConradCubicQuartic}{[ConradCubicQuartic]} / Table 5. The example also appears in the source; the arithmetic is displayed here in full. \href{https://kconrad.math.uconn.edu/blurbs/galoistheory/cubicquartic.pdf}{Source}.\par
\end{sources}\color{black}
\needspace{8\baselineskip}\section{Two examples resolving the cyclic-dihedral ambiguity}

\begin{atlasblock}[standard]
\textbf{\sffamily Example  /  the cyclic case}\par\smallskip
\color{black} For \(f=x^4+x^3+x^2+x+1\), translation gives \(f(x+1)=x^4+5x^3+10x^2+10x+5\), which is Eisenstein at \(5\). Here\par


\atlasdisplay{\Delta(f)=125,\qquad R_f(T)=(T-2)(T^2+T-1).}

\color{black} The unique rational resolvent root is \(r=2\), and the discriminant field is \(\mathbb Q(\sqrt5)\). The two test polynomials are \(T^2+T-1\) and \((T-1)^2\). Both split over that field, so the group is \(C_4\).\par
\end{atlasblock}\color{black}


\begin{atlasblock}[standard]
\textbf{\sffamily Example  /  the dihedral case}\par\smallskip
\color{black} For \(f=x^4-2\), Eisenstein at \(2\) proves irreducibility. Now\par


\atlasdisplay{\Delta(f)=-2048,\qquad R_f(T)=T(T^2+8).}

\color{black} Thus \(r=0\), the discriminant field is \(\mathbb Q(\sqrt{-2})\), and the test polynomials are \(T^2\) and \(T^2-2\). The latter does not split there: a real element of this imaginary quadratic field is rational, while \(\sqrt2\) is not. Hence the group is \(D_4\), of order eight.\par
\end{atlasblock}\color{black}

\color{black} Both resolvents have the same factor degrees \((1,2)\). The extra quadratic-field test distinguishes their groups.\par
\begin{sources}
\textbf{Sources for this section}\par\smallskip
\noindent\hyperlink{ref:KappeWarren1989Quartic}{[KappeWarren1989Quartic]} / pp. 133-137. Original group-separation method. \href{https://www.jstor.org/stable/2323198}{Source}.\par
\noindent\hyperlink{ref:ConradCubicQuartic}{[ConradCubicQuartic]} / Theorem 4.1 and Corollary 4.3. Exact test used in these independently calculated examples. \href{https://kconrad.math.uconn.edu/blurbs/galoistheory/cubicquartic.pdf}{Source}.\par
\end{sources}\color{black}
\needspace{8\baselineskip}\section{Why resolvents complement prime witnesses}
\color{black} A prime factorization proves that some cycle type occurs. A resolvent factorization proves that the group preserves an orbit or a structure in another action. These supply different constraints and often work best together. A rational linear resolvent factor can reveal a block structure invisible from the first few primes.\par
\color{black} For larger degrees, computing a single enormous absolute resolvent may be wasteful. Stauduhar's method descends through successive subgroup candidates using relative invariants. Fieker and Klüners explain how to construct invariants from group structure and how to certify the relevant factors. A numerical value close to an integer is a candidate for this process, not an exact factorization certificate.\par
\begin{sources}
\textbf{Sources for this section}\par\smallskip
\noindent\hyperlink{ref:Stauduhar1973}{[Stauduhar1973]}.  \href{https://doi.org/10.1090/S0025-5718-1973-0327712-4}{Source}.\par
\noindent\hyperlink{ref:FiekerKlueners2014}{[FiekerKlueners2014]} / Sections 3 and 7.  \href{https://doi.org/10.1112/S1461157013000302}{Source}.\par
\end{sources}\color{black}
\needspace{11\baselineskip}\section*{Questions and research boundaries}
\needspace{7\baselineskip}\subsection*{How can a quartic with a reducible cubic resolvent be distinguished as cyclic or dihedral?}
{\small\sffamily\color{atlasgold}Settled in stated scope / Separable irreducible quartics over fields of characteristic not 2}\par\smallskip
\color{black} First exclude the Klein-four case using the discriminant. In the remaining branch, the Kappe-Warren test checks the splitting of two specified quadratics over the discriminant quadratic extension.\par
\begin{sources}
\textbf{Sources for this section}\par\smallskip
\noindent\hyperlink{ref:ConradCubicQuartic}{[ConradCubicQuartic]} / Theorem 4.1.  \href{https://kconrad.math.uconn.edu/blurbs/galoistheory/cubicquartic.pdf}{Source}.\par
\noindent\hyperlink{ref:KappeWarren1989Quartic}{[KappeWarren1989Quartic]} / pp. 133-137.  \href{https://www.jstor.org/stable/2323198}{Source}.\par
\end{sources}\color{black}
\needspace{6\baselineskip}\par\medskip\begingroup\footnotesize\setlength{\parskip}{1pt}\textbf{\sffamily\color{atlasgreen}Connections}\par
\color{black}\textbf{Uses.} \hyperlink{article:discriminants}{Discriminants, parity, and ramification}. Parity finishes several branches of the quartic classification.\par
\color{black}\textbf{Applies to.} \hyperlink{article:galois-computation}{Computing a Galois group with a proof certificate}. Relative resolvents are a central tool in practical group determination.\par
\color{black}\textbf{Uses.} \hyperlink{article:permutation-groups}{The group-theoretic step: transitivity, alternating groups, and symmetric groups}. The factorization of a resolvent describes orbits in a coset action.\par
\color{black}\textbf{Supplies examples for.} \hyperlink{article:reciprocity}{Reciprocal symmetry and the group on root pairs}. The inverse-pair partition suggests natural resolvents and smaller ambient groups.\par
\endgroup
\chapter{Reciprocal symmetry and the group on root pairs}\label{article:reciprocity}\hypertarget{article:reciprocity}{}\color{black}
{\sffamily\small\color{atlasgold}method / Intermediate / Characteristic different from 2 / Published results}\par\medskip
Pass from a reciprocal polynomial to a polynomial of half the degree. The remaining quadratic extensions explain both irreducibility and the correct ambient Galois group.\par\medskip
\textbf{Questions answered}\begin{itemize}
\item How can a reciprocal polynomial be reduced to half its degree?
\item Why is irreducibility of the trace polynomial insufficient?
\item What precisely must be proved to obtain the full wreath product?
\end{itemize}
{\small\textbf{Prerequisites:} \hyperlink{article:foundations}{Fields, factors, and the action on roots}}\par
\needspace{8\baselineskip}\section{The quotient by inversion}
\color{black} Let \(K\) have characteristic different from \(2\). A monic polynomial \(f\) of degree \(2m\) is reciprocal if \(x^{2m}f(x^{-1})=f(x)\). Its constant term is then \(1\). There is a unique monic \(q\in K[y]\) of degree \(m\) such that\par

\atlasdisplay{f(x)=x^m q(x+x^{-1}).}
\color{black} To construct it, express the symmetric Laurent polynomial \(x^{-m}f(x)\) in the basis \(1,x^j+x^{-j}\). The recurrence \(s_{j+1}=ys_j-s_{j-1}\), with \(s_0=2,s_1=y\), expresses every \(s_j=x^j+x^{-j}\) as a polynomial in \(y\). The roots of \(q\) record the inverse pairs of roots of \(f\).\par
\begin{sources}
\textbf{Sources for this section}\par\smallskip
\noindent\hyperlink{ref:AndersonBertelliODorney2026}{[AndersonBertelliODorney2026]} / §2. Reciprocal polynomials and the trace substitution. \href{https://arxiv.org/abs/2406.18970v2}{Source}.\par
\end{sources}\color{black}
\needspace{8\baselineskip}\section{An exact irreducibility criterion}

\begin{atlasblock}[standard]
\textbf{\sffamily Criterion  /  quadratic lifting}\par\smallskip
\color{black} Assume \(q\) is irreducible and separable, and let \(y\) be one of its roots. Then \(x^m q(x+x^{-1})\) is irreducible over \(K\) if and only if \(y^2-4\) is not a square in \(K(y)\).\par
\end{atlasblock}\color{black}

\begin{atlasblock}[standard]
\color{black} A root \(\alpha\) above \(y\) satisfies \(\alpha^2-y\alpha+1=0\), and \(y=\alpha+\alpha^{-1}\). Thus \([K(\alpha):K]=m[K(\alpha):K(y)]\). The second factor is two exactly when the quadratic discriminant is a nonsquare. This proves the criterion by comparing the degree with \(2m\). Conversely, irreducibility of \(f\) forces irreducibility of \(q\).\par
\end{atlasblock}\color{black}
\begin{sources}
\textbf{Sources for this section}\par\smallskip
\noindent\hyperlink{ref:MinacNguyenTan2023}{[MinacNguyenTan2023]} / §4. The quadratic-lift method for reciprocal polynomials; the degree argument is included here. \href{https://arxiv.org/abs/2111.05256}{Source}.\par
\end{sources}\color{black}
\needspace{8\baselineskip}\section{The quotient group and the kernel are separate}
\color{black} Suppose \(f\) is separable and \(f(1)f(-1)\ne0\). Let \(M\) be the splitting field of \(q\), with roots \(y_1,\ldots,y_m\). The splitting field of \(f\) is\par

\atlasdisplay{L=M\bigl(\sqrt{y_1^2-4},\ldots,\sqrt{y_m^2-4}\bigr).}
\color{black} Every Galois automorphism permutes the inverse pairs and may exchange the two roots in each pair. Hence \(\operatorname{Gal}(f/K)\le C_2\wr S_m=C_2^m\rtimes S_m\). Its quotient on pairs is \(\operatorname{Gal}(q/K)\); its kernel has order \(2^r\), where \(r\) is the dimension of the span of the square classes \([y_i^2-4]\) in \(M^\times/M^{\times2}\). In particular, the full wreath product occurs exactly when the quotient is \(S_m\) and all \(m\) classes are independent.\par
\begin{sources}
\textbf{Sources for this section}\par\smallskip
\noindent\hyperlink{ref:AndersonBertelliODorney2026}{[AndersonBertelliODorney2026]} / §2. Signed permutations are the natural ambient group. \href{https://arxiv.org/abs/2406.18970v2}{Source}.\par
\noindent\hyperlink{ref:Cox2012}{[Cox2012]}. Quadratic extensions and Galois actions; the square-class formulation is standard Kummer theory. \href{https://onlinelibrary.wiley.com/doi/book/10.1002/9781118218457}{Source}.\par
\end{sources}\color{black}
\needspace{8\baselineskip}\section{Two quartics that expose the missing step}

\begin{atlasblock}[standard]
\textbf{\sffamily Example  /  quotient irreducible, lift reducible}\par\smallskip
\color{black} For \(q(y)=y^2-5\), the lift is \(f=x^4-3x^2+1=(x^2+x-1)(x^2-x-1)\). The polynomial \(q\) is irreducible, but \(y^2-4=1\) in \(\mathbb Q(y)\), so the lifting criterion fails.\par
\end{atlasblock}\color{black}

\begin{atlasblock}[standard]
\textbf{\sffamily Example  /  irreducible, smaller group}\par\smallskip
\color{black} For \(q(y)=y^2-3\), the lift is \(x^4-x^2+1=\Phi_{12}(x)\). Its splitting field is \(\mathbb Q(\sqrt3,i)\), of degree four, and its group is \(C_2\times C_2\). Both radicands equal \(-1\), so their square classes span only one dimension. Irreducibility has not forced the full order-eight wreath product.\par
\end{atlasblock}\color{black}
\begin{sources}
\textbf{Sources for this section}\par\smallskip
\noindent\hyperlink{ref:Cox2012}{[Cox2012]}. Cyclotomic splitting fields and quadratic extensions; the two factorizations are checked directly. \href{https://onlinelibrary.wiley.com/doi/book/10.1002/9781118218457}{Source}.\par
\end{sources}\color{black}
\needspace{8\baselineskip}\section{A useful discriminant identity}
\color{black} For monic \(q\), the quadratic-factor decomposition gives\par

\atlasdisplay{\operatorname{disc}(f)=q(2)q(-2)\operatorname{disc}(q)^2=(-1)^m f(1)f(-1)\operatorname{disc}(q)^2.}
\color{black} Indeed, the factors \(x^2-y_i x+1\) have discriminants \(y_i^2-4\), and the resultant of the factors indexed by \(i,j\) is \((y_i-y_j)^2\). Multiplying the discriminants and the squares of these resultants proves the identity. It supplies the parity of the action on all \(2m\) roots, but does not determine the entire kernel.\par
\color{black} The same organization applies to even polynomials \(q(x^2)\), with pairs \(\{\alpha,-\alpha\}\) and radicands \(y_i\). Affine reflection \(x\mapsto1-x\) becomes this situation after centering at \(1/2\).\par
\begin{sources}
\textbf{Sources for this section}\par\smallskip
\noindent\hyperlink{ref:AndersonBertelliODorney2026}{[AndersonBertelliODorney2026]} / §2. Discriminants in the reciprocal setting. \href{https://arxiv.org/abs/2406.18970v2}{Source}.\par
\noindent\hyperlink{ref:CullinanHajir2014}{[CullinanHajir2014]} / §1. The analogous signed action for even Legendre polynomials. \href{https://doi.org/10.1016/j.indag.2014.01.004}{Source}.\par
\end{sources}\color{black}
\needspace{11\baselineskip}\section*{Questions and research boundaries}
\needspace{7\baselineskip}\subsection*{When do all quadratic lifts remain independent?}
{\small\sffamily\color{atlasgold}Research direction / A specified reciprocal family over a field of characteristic different from 2.}\par\smallskip
\color{black} First determine the group of the trace polynomial. Then study products of the radicands in its splitting field. Relations among these square classes, rather than the degree alone, govern the possible deficit from the full wreath product.\par
\begin{sources}
\textbf{Sources for this section}\par\smallskip
\noindent\hyperlink{ref:AndersonBertelliODorney2026}{[AndersonBertelliODorney2026]} / §§2-3.  \href{https://arxiv.org/abs/2406.18970v2}{Source}.\par
\end{sources}\color{black}
\needspace{6\baselineskip}\par\medskip\begingroup\footnotesize\setlength{\parskip}{1pt}\textbf{\sffamily\color{atlasgreen}Connections}\par
\color{black}\textbf{Applies to.} \hyperlink{article:eulerian}{Eulerian polynomials: descents, shifts, and inverse pairs}. Eulerian reciprocity produces inverse pairs after the forced factor is removed.\par
\color{black}\textbf{Applies to.} \hyperlink{article:fekete}{Fekete polynomials: character sums and reciprocal arithmetic}. Character identities produce reciprocal normalizations and special-value tests.\par
\color{black}\textbf{Uses.} \hyperlink{article:discriminants}{Discriminants, parity, and ramification}. The discriminant detects one parity character of the full action.\par
\color{black}\textbf{Generalizes.} \hyperlink{article:orthogonal-families}{Hermite, Bessel, and Legendre: shared tools, different answers}. Even polynomials have an analogous action on sign pairs.\par
\endgroup
\part{Polynomial families}
\chapter{Trinomials: uniform families and explicit classifications}\label{article:trinomials}\hypertarget{article:trinomials}{}\color{black}
{\sffamily\small\color{atlasgold}family / Intermediate / Q / Published results}\par\medskip
The family x\ensuremath{{}^n}\ensuremath{-}x\ensuremath{-}1 realizes S\ensuremath{{}_n} over Q in every degree n\ensuremath{\ge}2. Power compositions lead to different root actions and recent exact classifications.\par\medskip
\textbf{Questions answered}\begin{itemize}
\item Which simple polynomial has group S\_\allowbreak{}n for every n?
\item How can two primes certify the quintic case?
\item Why does the all-degree proof use ramification?
\item How can a degree-twelve trinomial be classified using lower-degree polynomials?
\end{itemize}
{\small\textbf{Prerequisites:} \hyperlink{article:foundations}{Fields, factors, and the action on roots}}\par
\needspace{8\baselineskip}\section{A benchmark with exact attribution}

\begin{atlasblock}[standard]
\textbf{\sffamily Theorem  /  Selmer; Nart-Vila}\par\smallskip
\color{black} For every integer \(n\ge2\), \(f_n(x)=x^n-x-1\) is irreducible over \(\mathbb Q\), and its splitting field has Galois group \(S_n\).\par
\end{atlasblock}\color{black}
\color{black} Selmer proved irreducibility in 1956. Nart and Vila first proved the all-degree group statement in 1979; Osada obtained it independently as a consequence of a broader trinomial criterion. The two conclusions have different proofs and should retain their separate attribution.\par
\begin{sources}
\textbf{Sources for this section}\par\smallskip
\noindent\hyperlink{ref:Selmer1956}{[Selmer1956]}. Irreducibility of x\textasciicircum{}n\ensuremath{-}x\ensuremath{-}1. \href{https://doi.org/10.7146/math.scand.a-10478}{Source}.\par
\noindent\hyperlink{ref:NartVila1979}{[NartVila1979]}. The first all-degree symmetric-group result. \href{https://portalrecerca.uab.cat/en/publications/equations-of-the-type-xnaxb-with-absolute-galois-group-sn/}{Source}.\par
\noindent\hyperlink{ref:Osada1987}{[Osada1987]} / Corollary 3. Independent proof through a general trinomial criterion. \href{https://doi.org/10.1016/0022-314X(87)90029-1}{Source}.\par
\end{sources}\color{black}
\needspace{8\baselineskip}\section{A complete prime-degree irreducibility proof}

\begin{atlasblock}[standard]
\color{black} Let \(p\) be prime and let \(\alpha\) satisfy \(\alpha^p-\alpha-1=0\) over \(\overline{\mathbb F}_p\). Then \(\alpha^{p^k}=\alpha+k\), where \(k\) is read in \(\mathbb F_p\). The least positive \(k\) for which this equals \(\alpha\) is \(p\). Thus the Frobenius orbit has length \(p\), and \(x^p-x-1\) is irreducible over \(\mathbb F_p\). Reduction then proves irreducibility over \(\mathbb Q\).\par
\end{atlasblock}\color{black}
\color{black} This is an Artin-Schreier calculation. It is particularly short because the degree equals the residue characteristic; it does not cover arbitrary composite degrees.\par
\begin{sources}
\textbf{Sources for this section}\par\smallskip
\noindent\hyperlink{ref:LidlNiederreiter1997}{[LidlNiederreiter1997]}. Frobenius orbits and irreducible polynomials over finite fields; the orbit computation is explicit. \href{https://www.cambridge.org/core/books/finite-fields/75BDAA74ABAE713196E718392B9E5E72}{Source}.\par
\end{sources}\color{black}
\needspace{8\baselineskip}\section{A two-prime certificate for S\ensuremath{{}_5}}

\begin{atlasblock}[standard]
\textbf{\sffamily Example  /  cycle types and a graph}\par\smallskip
\color{black} For \(x^5-x-1\), reduction modulo \(5\) gives a 5-cycle. Modulo \(2\),\par

\atlasdisplay{x^5+x+1=(x^2+x+1)(x^3+x^2+1),}
\color{black} where the two factors are irreducible. The group therefore contains an element of type \((2,3)\), whose cube is a transposition. Label the roots by \(\mathbb Z/5\mathbb Z\) so that the 5-cycle adds one. Conjugating the transposition by its powers gives all edges with a fixed nonzero difference. These edges form a connected graph. Edge transpositions of a connected graph generate \(S_5\), proving the claim.\par
\end{atlasblock}\color{black}
\color{black} The conclusion uses this connectivity argument. A transitive subgroup containing one transposition need not be symmetric in general.\par
\begin{sources}
\textbf{Sources for this section}\par\smallskip
\noindent\hyperlink{ref:ConradModP}{[ConradModP]}. Factorization patterns give Frobenius cycle types. \href{https://kconrad.math.uconn.edu/blurbs/gradnumthy/galois-Q-factor-mod-p.pdf}{Source}.\par
\noindent\hyperlink{ref:ConradSnAn}{[ConradSnAn]}. Permutation-group criteria; the graph proof is included here. \href{https://kconrad.math.uconn.edu/blurbs/galoistheory/galoisSnAn.pdf}{Source}.\par
\end{sources}\color{black}
\needspace{8\baselineskip}\section{Why the uniform argument is global}
\color{black} For \(f_n\), the simultaneous equations \(f_n(a)=f_n'(a)=0\) modulo a prime give \(na^{n-1}=1\) and \((n-1)a=-n\). They force a unique possible multiple root. Such a prime divides neither \(n\) nor \(n-1\), and \(f_n''(a)=n(n-1)a^{n-2}\ne0\), so that root is double. Every nontrivial inertia element therefore exchanges just the two roots that coalesce.\par
\color{black} The normal subgroup generated by the finite inertia groups fixes an extension of \(\mathbb Q\) unramified at all finite primes. Minkowski's discriminant bound implies that this extension is trivial. Hence the whole Galois group is generated by transpositions. Selmer's irreducibility makes its action transitive, and the connected-graph argument now gives \(S_n\).\par
\begin{sources}
\textbf{Sources for this section}\par\smallskip
\noindent\hyperlink{ref:Osada1987}{[Osada1987]} / Corollary 3 and its proof mechanism. Ramification and generation by transpositions. \href{https://doi.org/10.1016/0022-314X(87)90029-1}{Source}.\par
\noindent\hyperlink{ref:Neukirch1999}{[Neukirch1999]}. Minkowski's discriminant bound and the absence of nontrivial everywhere-finitely-unramified extensions of Q. \href{https://link.springer.com/book/10.1007/978-3-662-03983-0}{Source}.\par
\end{sources}\color{black}
\needspace{8\baselineskip}\section{How to investigate a different trinomial}
\color{black} For \(x^n+ax^r+b\), first check common divisors of the exponents and the possibility of a composition. Then separate cyclotomic factors from the nonreciprocal part. The equations shared by the polynomial and its derivative tell whether a bad prime produces one double root or a more complicated collision. The latter may yield an inertia element quite different from a transposition. Selmer's and Ljunggren's results, followed by the general trinomial papers of Cohen-Movahhedi-Salinier, provide a route into this broader classification.\par
\begin{sources}
\textbf{Sources for this section}\par\smallskip
\noindent\hyperlink{ref:Ljunggren1960}{[Ljunggren1960]}. Factorization phenomena for trinomials and quadrinomials. \href{https://doi.org/10.7146/math.scand.a-10593}{Source}.\par
\noindent\hyperlink{ref:CohenMovahhediSalinier1999}{[CohenMovahhediSalinier1999]}. Broader trinomial Galois-group analysis. \href{https://doi.org/10.1006/jabr.1999.8033}{Source}.\par
\end{sources}\color{black}
\needspace{8\baselineskip}\section{A 2026 classification for x¹\ensuremath{{}^2}+ax\ensuremath{{}^6}+b}
\color{black} Suppose \(f=x^{12}+ax^6+b\in\mathbb Q[x]\) is irreducible. Chen classifies all 16 possible Galois groups. The key is the pair of decompositions\par

\atlasdisplay{f=g_4(x^3)=g_6(x^2),\qquad g_4=x^4+ax^2+b,\quad g_6=x^6+ax^3+b.}
\color{black} The groups of \(g_4,g_6\), followed by at most two additional rational-square tests, distinguish the possibilities. Theorem 3.8 gives the complete algorithm, including tests for rational cubes and reducibility of \(r=x^3-3bx+ab\). Irreducibility of \(f\) is an input hypothesis.\par

\begin{atlasblock}[standard]
\textbf{\sffamily Example}\par\smallskip
\color{black} For \(f=x^{12}+4x^6+2\), Eisenstein at 2 proves irreducibility. Here \(b(a^2-4b)=16\), \(b=2\) is not a rational cube, and \(r=x^3-6x+8\) has no rational root. The first branch of the algorithm gives \(12T39\), the transitive-group catalogue label.\par
\end{atlasblock}\color{black}
\begin{sources}
\textbf{Sources for this section}\par\smallskip
\noindent\hyperlink{ref:Chen2026Dodecic}{[Chen2026Dodecic]} / Theorem 3.8; Table 2.  \href{https://doi.org/10.1515/ms-2026-0503}{Source}.\par
\end{sources}\color{black}
\needspace{8\baselineskip}\section{A related classification for cubic power compositions}
\color{black} For \(g=x^3+ax^2+bx+c\in\mathbb Q[x]\) and irreducible \(f=g(x^3)\), Awtrey-Patane-Toone classify 11 possible transitive groups in degree nine. They use the discriminants of \(f\) and \(g\), together with two further resolvents. Their work also addresses the converse problem: whether a given number field admits a defining polynomial of this compositional shape. These polynomials usually have four terms; the connection here is the root-block structure created by power composition.\par
\begin{sources}
\textbf{Sources for this section}\par\smallskip
\noindent\hyperlink{ref:AwtreyPataneToone2025}{[AwtreyPataneToone2025]} / Abstract and Introduction.  \href{https://doi.org/10.1017/S0004972725000048}{Source}.\par
\end{sources}\color{black}
\needspace{11\baselineskip}\section*{Questions and research boundaries}
\needspace{7\baselineskip}\subsection*{Can a sparse family be handled by one uniform ramification argument?}
{\small\sffamily\color{atlasgold}Research direction / A specified family \(x^n+a_nx^{r_n}+b_n\) over \(\mathbb Q\).}\par\smallskip
\color{black} Identify which primes can create repeated roots and how many roots coalesce. A proof needs both transitivity and sufficient inertia generators; observed modular cycles alone do not give an all-degree theorem.\par
\begin{sources}
\textbf{Sources for this section}\par\smallskip
\noindent\hyperlink{ref:Osada1987}{[Osada1987]}.  \href{https://doi.org/10.1016/0022-314X(87)90029-1}{Source}.\par
\noindent\hyperlink{ref:CohenMovahhediSalinier1999}{[CohenMovahhediSalinier1999]}.  \href{https://doi.org/10.1006/jabr.1999.8033}{Source}.\par
\end{sources}\color{black}
\needspace{6\baselineskip}\par\medskip\begingroup\footnotesize\setlength{\parskip}{1pt}\textbf{\sffamily\color{atlasgreen}Connections}\par
\color{black}\textbf{Uses.} \hyperlink{article:frobenius}{Frobenius: from factor degrees to permutations}. Small good primes provide finite certificates in individual degrees.\par
\color{black}\textbf{Uses.} \hyperlink{article:permutation-groups}{The group-theoretic step: transitivity, alternating groups, and symmetric groups}. A transitive group generated by transpositions is symmetric.\par
\color{black}\textbf{Generalizes.} \hyperlink{article:sparse-gaps}{Sparse polynomials, factor reversal, and large gaps}. Sparse families may require separating reciprocal and nonreciprocal factors.\par
\color{black}\textbf{Supplies examples for.} \hyperlink{article:inverse-galois}{Inverse Galois theory and generic polynomials}. The family realizes every symmetric group over Q.\par
\color{black}\textbf{Uses.} \hyperlink{article:capelli-composition}{Capelli: composition, binomials, and iterates}. The factorization of a power composition is controlled over a root field; irreducibility must precede the classification.\par
\endgroup
\chapter{Taylor truncations: two exact benchmarks}\label{article:taylor}\hypertarget{article:taylor}{}\color{black}
{\sffamily\small\color{atlasgold}family / Intermediate / Q / Published results}\par\medskip
Exponential truncations have alternating or symmetric groups, while geometric truncations are cyclotomic. Their contrast identifies which arithmetic features a general truncation theory must track.\par\medskip
\textbf{Questions answered}\begin{itemize}
\item What exactly is known for every exponential truncation?
\item How does the discriminant computation work?
\item Why do geometric truncations behave differently?
\end{itemize}
{\small\textbf{Prerequisites:} \hyperlink{article:foundations}{Fields, factors, and the action on roots}}\par
\needspace{8\baselineskip}\section{The arithmetic question behind an analytic series}
\color{black} For \(F(x)=\sum_{j\ge0}a_jx^j\in\mathbb Q[[x]]\), let \(F_{\le n}=\sum_{j=0}^n a_jx^j\). Its degree is \(n\) only when \(a_n\ne0\). Multiplying by a nonzero rational scalar clears denominators without changing the roots or splitting field. Analytic convergence and the limiting location of complex roots provide useful structure, but do not alone prove irreducibility over \(\mathbb Q\).\par
\color{black} Katz and Rivin ask about the frequencies of Galois groups in such sequences. Rabarison, Pazuki, and Molin give a later survey and extend the experimental viewpoint to Padé approximants. In a catalogue, truncations and rational approximants should have separate definitions and normalization records.\par
\begin{sources}
\textbf{Sources for this section}\par\smallskip
\noindent\hyperlink{ref:KatzRivinSpeculations}{[KatzRivinSpeculations]} / §§1-2. Limiting-frequency questions and explicitly labeled computational evidence. \href{https://web.math.princeton.edu/~nmk/kGalois_11bis.pdf}{Source}.\par
\noindent\hyperlink{ref:RabarisonPazukiMolin2024}{[RabarisonPazukiMolin2024]}. Survey of truncations and experiments with Padé approximants. \href{https://www.numdam.org/articles/10.5802/pmb.57/}{Source}.\par
\end{sources}\color{black}
\needspace{8\baselineskip}\section{Schur's exact exponential theorem}

\begin{atlasblock}[standard]
\textbf{\sffamily Theorem  /  Schur}\par\smallskip
\color{black} For every \(n\ge2\), \(T_n(x)=\sum_{j=0}^n x^j/j!\) is irreducible over \(\mathbb Q\), and\par

\atlasdisplay{\operatorname{Gal}(T_n/\mathbb Q)=\begin{cases}A_n,&4\mid n,\\S_n,&4\nmid n.\end{cases}}
\end{atlasblock}\color{black}
\color{black} The proof separates three issues: exclude every proper factor degree, force enough permutation structure to contain \(A_n\), and finally compute the discriminant square class. Coleman gives a Newton-polygon proof, making this family a model for later factorial-coefficient arguments.\par
\begin{sources}
\textbf{Sources for this section}\par\smallskip
\noindent\hyperlink{ref:Schur1930}{[Schur1930]}. Original all-degree group theorem. \href{https://www.numdam.org/articles/10.5802/pmb.57/}{Historical citation verification}.\par
\noindent\hyperlink{ref:Coleman1987}{[Coleman1987]}. Newton polygons and the exact alternating/symmetric conclusion. \href{https://doi.org/10.5169/seals-87891}{Source}.\par
\end{sources}\color{black}
\needspace{8\baselineskip}\section{Deriving the parity test}

\begin{atlasblock}[standard]
\color{black} Set \(P_n=n!T_n\), which is monic. Since \(P_n'=P_n-x^n\), each root \(r\) satisfies \(P_n'(r)=-r^n\). The constant term is \(n!\), so multiplying these identities over the roots gives\par

\atlasdisplay{\operatorname{disc}(P_n)=(-1)^{n(n-1)/2}(n!)^n.}
\color{black} For even \(n\), this is a rational square exactly when \(4\mid n\): its magnitude is already a square and its sign decides the issue. For odd \(n\ge3\), a prime between \(n/2\) and \(n\) occurs exactly once in \(n!\), so the discriminant is a nonsquare. This step determines the group only after \(A_n\)-containment has been proved.\par
\end{atlasblock}\color{black}
\begin{sources}
\textbf{Sources for this section}\par\smallskip
\noindent\hyperlink{ref:Coleman1987}{[Coleman1987]}. The discriminant square test; the resultant computation is included explicitly. \href{https://doi.org/10.5169/seals-87891}{Source}.\par
\end{sources}\color{black}
\needspace{8\baselineskip}\section{A local proof at degree six}

\begin{atlasblock}[standard]
\textbf{\sffamily Example  /  intersecting factor degrees}\par\smallskip

\atlasdisplay{P_6=x^6+6x^5+30x^4+120x^3+360x^2+720x+720.}
\color{black} At \(2\), its lower Newton polygon has vertices \((0,4),(4,1),(6,0)\); a proper rational factor can therefore have only degree \(2\) or \(4\). At \(3\), the polygon is the segment from \((0,2)\) to \((6,0)\), so a proper factor can have only degree \(3\). These constraints have no common value. This proves irreducibility without looking for an irreducible reduction at a good prime.\par
\end{atlasblock}\color{black}
\begin{sources}
\textbf{Sources for this section}\par\smallskip
\noindent\hyperlink{ref:Coleman1987}{[Coleman1987]}. The Newton-polygon mechanism; the displayed degree-six coefficients and vertices are checked directly. \href{https://doi.org/10.5169/seals-87891}{Source}.\par
\end{sources}\color{black}
\needspace{8\baselineskip}\section{The cyclotomic contrast}
\color{black} The degree-n truncation of \((1-x)^{-1}\) is\par

\atlasdisplay{1+x+\cdots+x^n=\prod_{\substack{d\mid n+1\\d>1}}\Phi_d(x).}
\color{black} For \(n\ge1\), it is irreducible exactly when \(n+1\) is prime; its splitting field is \(\mathbb Q(\zeta_{n+1})\), with group \((\mathbb Z/(n+1)\mathbb Z)^\times\). Thus even an elementary analytic series can have arithmetic behavior quite different from the exponential.\par
\color{black} For a new series, first remove factors forced by parity, vanishing coefficients, or rational-function identities. Then seek a uniform valuation pattern, a usable specialization theorem, or a group argument tailored to the remaining family. Numerical frequency is a question to explain, not a proof of its limit.\par
\begin{sources}
\textbf{Sources for this section}\par\smallskip
\noindent\hyperlink{ref:Cox2012}{[Cox2012]}. Cyclotomic factorization and cyclotomic Galois groups. \href{https://onlinelibrary.wiley.com/doi/book/10.1002/9781118218457}{Source}.\par
\noindent\hyperlink{ref:KatzRivinSpeculations}{[KatzRivinSpeculations]} / §§1-4. Frequency questions, variants, and shifts. \href{https://web.math.princeton.edu/~nmk/kGalois_11bis.pdf}{Source}.\par
\end{sources}\color{black}
\needspace{11\baselineskip}\section*{Questions and research boundaries}
\needspace{7\baselineskip}\subsection*{When does a truncation family have a limiting Galois-group frequency?}
{\small\sffamily\color{atlasgold}Research direction / A specified rational power series, a stated normalization, and n tending to infinity.}\par\smallskip
\color{black} The exponential and geometric cases already show that no universal ``always symmetric'' conclusion is possible. Katz and Rivin propose broader frequency questions and variants. A precise new question must say which factors are removed and which group action is being counted.\par
\begin{sources}
\textbf{Sources for this section}\par\smallskip
\noindent\hyperlink{ref:KatzRivinSpeculations}{[KatzRivinSpeculations]} / §§1-4.  \href{https://web.math.princeton.edu/~nmk/kGalois_11bis.pdf}{Source}.\par
\end{sources}\color{black}
\needspace{6\baselineskip}\par\medskip\begingroup\footnotesize\setlength{\parskip}{1pt}\textbf{\sffamily\color{atlasgreen}Connections}\par
\color{black}\textbf{Uses.} \hyperlink{article:newton-polygons}{Newton polygons, factor degrees, and large cycles}. Factorial valuations provide incompatible local factor degrees.\par
\color{black}\textbf{Uses.} \hyperlink{article:discriminants}{Discriminants, parity, and ramification}. A resultant identity determines the sign character after large-group containment.\par
\color{black}\textbf{Generalizes.} \hyperlink{article:laguerre}{Generalized Laguerre polynomials: two parameter regimes}. Exponential truncations occur as a varying-parameter Laguerre specialization.\par
\color{black}\textbf{Contrasts with.} \hyperlink{article:specialization}{Specialization and Hilbert irreducibility}. A family indexed by degree is different from specialization of one fixed-degree cover.\par
\endgroup
\chapter{Generalized Laguerre polynomials: two parameter regimes}\label{article:laguerre}\hypertarget{article:laguerre}{}\color{black}
{\sffamily\small\color{atlasgold}family / Intermediate / Q / Published results}\par\medskip
Factorial products connect local slopes to large Galois groups. Fixing the parameter and increasing the degree is a different problem from fixing the degree and varying the parameter.\par\medskip
\textbf{Questions answered}\begin{itemize}
\item What changes when the parameter is fixed instead of the degree?
\item Why do primes in arithmetic progressions appear?
\item Which negative integer parameters force reducibility?
\item Which small rational parameters admit an all-degree irreducibility theorem?
\end{itemize}
{\small\textbf{Prerequisites:} \hyperlink{article:foundations}{Fields, factors, and the action on roots}}\par
\needspace{8\baselineskip}\section{The family and its arithmetic coefficients}
\color{black} For \(n\ge1\) and \(\alpha\in\mathbb Q\), set\par

\atlasdisplay{L_n^{(\alpha)}(x)=\sum_{j=0}^n\binom{n+\alpha}{n-j}\frac{(-x)^j}{j!}.}
\color{black} The generalized binomial coefficient is a polynomial in \(\alpha\). For real \(\alpha>-1\), these polynomials form a classical orthogonal family; the arithmetic question also makes sense outside that range. Their monic normalization is \(M_n^{(\alpha)}=(-1)^n n!L_n^{(\alpha)}\).\par
\color{black} If \(\alpha=a/b\) in lowest terms, \(b>0\), then an integral monic normalization with the same splitting field is\par

\atlasdisplay{b^n n!L_n^{(a/b)}(-x/b)=\sum_{j=0}^n\binom nj\Bigl(\prod_{k=j+1}^n(bk+a)\Bigr)x^j.}
\color{black} This expression exposes divisibility in an arithmetic progression.\par
\begin{sources}
\textbf{Sources for this section}\par\smallskip
\noindent\hyperlink{ref:Hajir2005}{[Hajir2005]} / §§1 and 3. Definition and the integral coefficient-product normalization. \href{https://jtnb.centre-mersenne.org/item/JTNB_2005__17_2_517_0/}{Source}.\par
\end{sources}\color{black}
\needspace{8\baselineskip}\section{Fixed rational parameter, growing degree}

\begin{atlasblock}[standard]
\textbf{\sffamily Theorem  /  Filaseta-Lam; Hajir}\par\smallskip
\color{black} Fix \(\alpha\in\mathbb Q\setminus\mathbb Z_{<0}\). For every sufficiently large \(n\), depending on \(\alpha\), the polynomial \(L_n^{(\alpha)}\) is irreducible over \(\mathbb Q\) and its group contains \(A_n\). It equals \(A_n\) or \(S_n\) according to whether its discriminant is a rational square.\par
\end{atlasblock}\color{black}
\color{black} Filaseta-Lam supply eventual irreducibility. Hajir obtains a large prime divisor of a Newton-slope denominator, then uses a permutation-group criterion. The irreducibility input is essential to that group argument.\par
\begin{sources}
\textbf{Sources for this section}\par\smallskip
\noindent\hyperlink{ref:FilasetaLam2002}{[FilasetaLam2002]}. Eventual irreducibility for a fixed rational parameter outside the negative integers. \href{https://doi.org/10.4064/aa105-2-4}{Source}.\par
\noindent\hyperlink{ref:Hajir2005}{[Hajir2005]} / Theorems 1.1 and 2.2. Eventual alternating containment and the Newton-index criterion. \href{https://jtnb.centre-mersenne.org/item/JTNB_2005__17_2_517_0/}{Source}.\par
\end{sources}\color{black}
\needspace{8\baselineskip}\section{Eight rational parameters: a uniform theorem with one exception}

\begin{atlasblock}[standard]
\textbf{\sffamily Theorem  /  Laishram-Shorey}\par\smallskip
\color{black} Let \(q\in\{\pm\tfrac13,\pm\tfrac23,\pm\tfrac14,\pm\tfrac34\}\). For every \(n\ge1\), \(L_n^{(q)}(x)\) is irreducible over \(\mathbb Q\), except when \((q,n)=(\tfrac14,2)\).\par
\end{atlasblock}\color{black}
\color{black} This is stronger than eventual irreducibility. The proof combines Newton polygons with large prime divisors of arithmetic-progression products.\par

\begin{atlasblock}[standard]
\textbf{\sffamily The exceptional pair}\par\smallskip

\atlasdisplay{L_2^{(1/4)}(x)=\frac{(4x-3)(4x-15)}{32}.}
\color{black} The exception is explicit in arXiv:1306.0736v2. The unrestricted power-substitution wording of the cited theorem needs a low-degree qualification: although \(L_1^{(-3/4)}(x)=1/4-x\) is irreducible, \(L_1^{(-3/4)}(x^4)=(1/2-x^2)(1/2+x^2)\).\par
\end{atlasblock}\color{black}
\begin{sources}
\textbf{Sources for this section}\par\smallskip
\noindent\hyperlink{ref:LaishramShorey2016III}{[LaishramShorey2016III]} / Theorem 1 and §8; arXiv:1306.0736v2. The ordinary Laguerre assertion only is used; see the explicit boundary example for power substitutions. \href{https://doi.org/10.1016/j.jnt.2016.01.003}{Source}.\par
\end{sources}\color{black}
\needspace{8\baselineskip}\section{Fixed degree, varying parameter}
\color{black} A different theorem holds over any fixed number field \(K\): for fixed \(n\ge10\), all but finitely many \(\alpha\in K\) give an irreducible polynomial \(L_n^{(\alpha)}\) with Galois group \(S_n\) over \(K\). Hajir-Wong prove this using intermediate covers and finiteness of rational points on curves of sufficiently high genus.\par
\color{black} This statement has different quantifiers from the previous one. It does not assert a common threshold in \(n\) for every rational \(\alpha\); nor does ordinary Hilbert irreducibility by itself make the exceptional set finite. A varying sequence \(\alpha_n\) requires its own argument.\par
\begin{sources}
\textbf{Sources for this section}\par\smallskip
\noindent\hyperlink{ref:HajirWong2006}{[HajirWong2006]} / Theorem 1(a). The finite exceptional set for fixed n\ensuremath{\ge}10 and fixed number field K. \href{https://www.numdam.org/item/AIF_2006__56_4_1127_0/}{Source}.\par
\end{sources}\color{black}
\needspace{8\baselineskip}\section{An exact cubic and the discriminant product}
\color{black} Schur's product formula for the monic normalization is\par

\atlasdisplay{\operatorname{disc}(M_n^{(\alpha)})=\prod_{j=2}^n j^j(\alpha+j)^{j-1}.}

\begin{atlasblock}[standard]
\textbf{\sffamily Example  /  \ensuremath{\alpha}=0, n=3}\par\smallskip
\color{black} Here \(M_3=x^3-9x^2+18x-6\). Eisenstein at \(3\) proves irreducibility. Its discriminant is \(2^3 3^5=1944\), a nonsquare. An irreducible cubic has group \(A_3\) or \(S_3\); therefore this group is \(S_3\).\par
\end{atlasblock}\color{black}
\color{black} The same product helps identify parameters where roots collide: for degree \(n\ge2\), the exceptional values are \(-2,-3,\ldots,-n\). Separability and irreducibility are separate properties; at \(\alpha=-1\) there is already a simple rational root at zero.\par
\begin{sources}
\textbf{Sources for this section}\par\smallskip
\noindent\hyperlink{ref:Schur1931}{[Schur1931]}. Classical algebraic properties and discriminants. \href{https://doi.org/10.1515/crll.1931.165.52}{Source}.\par
\noindent\hyperlink{ref:Hajir2005}{[Hajir2005]} / §1. Schur's discriminant formula for the monic normalization. \href{https://jtnb.centre-mersenne.org/item/JTNB_2005__17_2_517_0/}{Source}.\par
\end{sources}\color{black}
\needspace{8\baselineskip}\section{Why negative integers must be excluded}

\begin{atlasblock}[boundary]
\textbf{\sffamily Boundary  /  a forced power of x}\par\smallskip
\color{black} For integers \(m\ge1\) and \(n\ge m\), coefficient comparison gives\par

\atlasdisplay{L_n^{(-m)}(x)=(-x)^m\frac{(n-m)!}{n!}L_{n-m}^{(m)}(x).}
\color{black} For example, \(L_3^{(-2)}(x)=x^2(3-x)/6\). Thus a fixed negative integer parameter supplies a forced zero for every sufficiently large degree. The hypotheses of the fixed-parameter theorem are mathematically necessary.\par
\end{atlasblock}\color{black}
\color{black} On the other hand, \(\alpha=-n-1\) gives \(L_n^{(-n-1)}=(-1)^n T_n\), the exponential truncation. A parameter depending on the degree can avoid the fixed-parameter obstruction in a structured way.\par
\begin{sources}
\textbf{Sources for this section}\par\smallskip
\noindent\hyperlink{ref:Hajir2005}{[Hajir2005]} / §1, Remarks 1 and 3. The negative-integer exclusion and the exponential specialization. \href{https://jtnb.centre-mersenne.org/item/JTNB_2005__17_2_517_0/}{Source}.\par
\end{sources}\color{black}
\needspace{8\baselineskip}\section{Explicit alternating groups and the classification problem}
\color{black} There is also a uniform varying-parameter result: for every \(n>2\) with \(n\equiv2\pmod4\), the polynomial \(L_n^{(n)}\) is irreducible and has group \(A_n\). Together with Schur's other Laguerre specializations, this realizes every \(A_n\) in this family. At \(n=2\), \(L_2^{(2)}=(x-2)(x-6)/2\) is split and its group is the trivial group \(A_2\); irreducibility is not claimed there.\par
\color{black} For integer parameters, Banerjee-Filaseta-Finch-Leidy classify nonzero square discriminants. Banerjee-Bera subsequently study the actual alternating Galois groups for integer \(\alpha\ne n\), up to finitely many exceptions. The square-discriminant classification and the group classification answer different questions.\par
\begin{sources}
\textbf{Sources for this section}\par\smallskip
\noindent\hyperlink{ref:FilasetaKiddTrifonov2012}{[FilasetaKiddTrifonov2012]} / Theorem 1 and Corollary 1.  \href{https://doi.org/10.1016/j.jnt.2011.09.012}{Source}.\par
\noindent\hyperlink{ref:BanerjeeFilasetaFinchLeidy2013}{[BanerjeeFilasetaFinchLeidy2013]}. Classification of nonzero square discriminants. \href{https://www.numdam.org/articles/10.5802/jtnb.822/}{Source}.\par
\noindent\hyperlink{ref:BanerjeeBera2022}{[BanerjeeBera2022]} / Abstract and Introduction. The \ensuremath{\alpha}\ensuremath{\ne}n restriction is essential to the stated classification scope. \href{https://doi.org/10.1016/j.jnt.2021.12.009}{Source}.\par
\end{sources}\color{black}
\needspace{11\baselineskip}\section*{Questions and research boundaries}
\needspace{7\baselineskip}\subsection*{Which explicit varying-parameter sequences have irreducible Laguerre polynomials?}
{\small\sffamily\color{atlasgold}Research direction / \(L_n^{(\alpha_n)}\) over Q, with the sequence \ensuremath{\alpha}\_\allowbreak{}n specified in advance.}\par\smallskip
\color{black} Fixed-parameter eventual results cannot be applied by substituting a new parameter at every degree. A proposed sequence needs uniform control of factor degrees, or a separate theorem tailored to its coefficient products.\par
\begin{sources}
\textbf{Sources for this section}\par\smallskip
\noindent\hyperlink{ref:Hajir2009}{[Hajir2009]}.  \href{https://doi.org/10.4153/CJM-2009-031-6}{Source}.\par
\noindent\hyperlink{ref:FilasetaKiddTrifonov2012}{[FilasetaKiddTrifonov2012]}.  \href{https://doi.org/10.1016/j.jnt.2011.09.012}{Source}.\par
\end{sources}\color{black}
\needspace{7\baselineskip}\subsection*{Are there only finitely many exceptional parameters in fixed degree \(n\ge10\)?}
{\small\sffamily\color{atlasgold}Settled in stated scope / Any fixed number field \(K\) and fixed \(n\ge10\).}\par\smallskip
\color{black} Yes: the parameters in \(K\) for which the polynomial is reducible or has Galois group different from \(S_n\) form a finite set. The theorem does not automatically provide an effective list of those parameters.\par
\begin{sources}
\textbf{Sources for this section}\par\smallskip
\noindent\hyperlink{ref:HajirWong2006}{[HajirWong2006]} / Theorem 1(a) and Remark 1.  \href{https://www.numdam.org/item/AIF_2006__56_4_1127_0/}{Source}.\par
\end{sources}\color{black}
\needspace{6\baselineskip}\par\medskip\begingroup\footnotesize\setlength{\parskip}{1pt}\textbf{\sffamily\color{atlasgreen}Connections}\par
\color{black}\textbf{Uses.} \hyperlink{article:newton-polygons}{Newton polygons, factor degrees, and large cycles}. Arithmetic-progression products expose useful valuation slopes.\par
\color{black}\textbf{Uses.} \hyperlink{article:specialization}{Specialization and Hilbert irreducibility}. Fixed-degree parameter families are studied through covers and exceptional specializations.\par
\color{black}\textbf{Supplies examples for.} \hyperlink{article:taylor}{Taylor truncations: two exact benchmarks}. The parameter \ensuremath{-}n\ensuremath{-}1 gives the exponential truncation.\par
\color{black}\textbf{Generalizes.} \hyperlink{article:orthogonal-families}{Hermite, Bessel, and Legendre: shared tools, different answers}. Hermite and Bessel identities connect to Laguerre parameters and substitutions.\par
\color{black}\textbf{Connects to.} \hyperlink{article:capelli-composition}{Capelli: composition, binomials, and iterates}. Power substitutions can create factors even when the original Laguerre polynomial is irreducible.\par
\endgroup
\chapter{Hermite, Bessel, and Legendre: shared tools, different answers}\label{article:orthogonal-families}\hypertarget{article:orthogonal-families}{}\color{black}
{\sffamily\small\color{atlasgold}family / Intermediate / Q / Research overview}\par\medskip
Classical polynomial families share differential and recurrence identities, but their arithmetic factorization problems differ. Parity reduction and coefficient products expose the correct questions.\par\medskip
\textbf{Questions answered}\begin{itemize}
\item Why must odd Hermite and Legendre polynomials be normalized?
\item Which all-degree result is known for Bessel polynomials?
\item Which parts of a local proof can transfer to another classical family?
\end{itemize}
{\small\textbf{Prerequisites:} \hyperlink{article:foundations}{Fields, factors, and the action on roots}}\par
\needspace{8\baselineskip}\section{Analytic structure does not determine arithmetic factorization}
\color{black} Orthogonal polynomials occur in approximation, differential equations, and spectral problems. Their recurrence and differential identities often give exact coefficients and simple real roots. Neither property alone forces irreducibility over \(\mathbb Q\): a polynomial can have distinct real algebraic roots and still factor rationally.\par
\color{black} Schur's work on Laguerre and Hermite polynomials established local arithmetic methods for such families. Subsequent examples show both the power and the limits of transferring those methods: Bessel irreducibility has an all-degree theorem, whereas the cited Legendre paper treats the general factorization statement as a conjecture.\par
\begin{sources}
\textbf{Sources for this section}\par\smallskip
\noindent\hyperlink{ref:Schur1931}{[Schur1931]}. Foundational algebraic study of Laguerre and Hermite families. \href{https://doi.org/10.1515/crll.1931.165.52}{Source}.\par
\noindent\hyperlink{ref:FilasetaTrifonov2002}{[FilasetaTrifonov2002]}. All-degree Bessel irreducibility. \href{https://doi.org/10.1515/crll.2002.069}{Source}.\par
\noindent\hyperlink{ref:CullinanHajir2014}{[CullinanHajir2014]} / Introduction. Stieltjes' conjecture and the remaining Legendre problem in this source. \href{https://doi.org/10.1016/j.indag.2014.01.004}{Source}.\par
\end{sources}\color{black}
\needspace{8\baselineskip}\section{Hermite polynomials and Laguerre parameters}
\color{black} For the physicists' Hermite convention \(H_n(x)=(-1)^n e^{x^2}(d/dx)^n e^{-x^2}\), one has\par

\atlasdisplay{H_{2m}(x)=(-1)^m2^{2m}m!L_m^{(-1/2)}(x^2),}

\atlasdisplay{H_{2m+1}(x)=(-1)^m2^{2m+1}m!xL_m^{(1/2)}(x^2).}
\color{black} The odd polynomial has a forced factor \(x\); all remaining roots occur in opposite pairs. A theorem concerning \(L_m^{(\pm1/2)}\) alone does not automatically determine the quadratic extensions obtained when square roots of its roots are adjoined.\par

\begin{atlasblock}[standard]
\textbf{\sffamily Example  /  H\ensuremath{{}_4}}\par\smallskip
\color{black} Here \(H_4=16x^4-48x^2+12\). Its pair polynomial is \(4y^2-12y+3\), with roots \((3\pm\sqrt6)/2\). This explicitly separates the degree-two pair field from the further square-root extensions. The same procedure works in higher degree, even when the pair polynomial is harder to analyze.\par
\end{atlasblock}\color{black}

\begin{atlasblock}[standard]
\textbf{\sffamily Schur's Hermite theorem}\par\smallskip
\color{black} For every \(m\ge1\), the physicists' polynomials \(H_{2m}(x)\) and \(H_{2m+1}(x)/x\) are irreducible over \(\mathbb Q\). This uses the normalization above. Rescaling the variable by an irrational number need not preserve irreducibility over \(\mathbb Q\): for example, \(H_2=4x^2-2\) is irreducible, whereas the probabilists' polynomial \(\operatorname{He}_2=x^2-1\) is reducible.\par
\end{atlasblock}\color{black}
\begin{sources}
\textbf{Sources for this section}\par\smallskip
\noindent\hyperlink{ref:CullinanHajirSell2009}{[CullinanHajirSell2009]} / §1. The two Hermite-Laguerre identities. \href{https://www.numdam.org/articles/10.5802/jtnb.659/}{Source}.\par
\noindent\hyperlink{ref:Schur1931}{[Schur1931]}. The classical arithmetic context. \href{https://doi.org/10.1515/crll.1931.165.52}{Source}.\par
\noindent\hyperlink{ref:LaishramShorey2016III}{[LaishramShorey2016III]} / Introduction. Restates Schur's theorem for L\_\allowbreak{}n\textasciicircum{}(\ensuremath{\pm}1/2)(x\ensuremath{{}^2}) and the physicists' Hermite normalization. \href{https://doi.org/10.1016/j.jnt.2016.01.003}{Source}.\par
\end{sources}\color{black}
\needspace{8\baselineskip}\section{Wronskians: arithmetic beyond a single orthogonal polynomial}
\color{black} Define the monic probabilists' polynomials by \(e^{xt-t^2/2}=\sum_{n\ge0}\operatorname{He}_n(x)t^n/n!\). For a partition \(\lambda=(\lambda_1\ge\cdots\ge\lambda_r>0)\), put \(n_i=\lambda_{r-i+1}+i-1\) and\par

\atlasdisplay{W_\lambda(x)=\frac{\operatorname{Wr}(\operatorname{He}_{n_1},\ldots,\operatorname{He}_{n_r})}{\prod_{i<j}(n_j-n_i)}.}
\color{black} Write \(W_\lambda=x^{e_\lambda}R_\lambda\), where \(e_\lambda\) is the multiplicity at zero. Grosu-Grosu prove that \(R_{(n,1)}\) is irreducible for \(n=1\) or \(n\ge3\), \(R_{(n,2)}\) for \(n\ge2\), and \(R_{(n,n)}\) when \(n+1\) is a square.\par
\color{black} At \(n=2\), \(W_{(2,1)}=x^3\), so \(R_{(2,1)}=1\), a unit rather than an irreducible polynomial under our convention. Their methods combine Newton-slope bounds with Schur-type congruences. Irreducibility implies simple nonzero roots, connecting arithmetic to differential equations. Arbitrary partitions do not satisfy uniform irreducibility: \(R_{(6,3,2,1)}=x^2-9\). This concerns the displayed probabilists' normalization.\par
\begin{sources}
\textbf{Sources for this section}\par\smallskip
\noindent\hyperlink{ref:GrosuGrosu2021}{[GrosuGrosu2021]} / Definition 1, Theorems 1-3 and 7-8, and the discussion after Proposition 6 (arXiv:2007.00065v1).  \href{https://doi.org/10.1016/j.indag.2020.12.004}{Source}.\par
\end{sources}\color{black}
\needspace{8\baselineskip}\section{An all-degree Bessel theorem}

\begin{atlasblock}[standard]
\textbf{\sffamily Theorem  /  Filaseta-Trifonov}\par\smallskip
\color{black} Let \(n\ge1\), and let \(a_0,\ldots,a_n\in\mathbb Z\) satisfy \(|a_0|=|a_n|=1\). Then\par

\atlasdisplay{\sum_{j=0}^n a_j\frac{(n+j)!}{2^j(n-j)!j!}x^j}
\color{black} is irreducible over \(\mathbb Q\). Taking every \(a_j=1\) proves the irreducibility of the Bessel polynomial \(y_n\) in every positive degree.\par
\end{atlasblock}\color{black}
\color{black} The method bounds possible factor degrees by Newton polygons and prime divisors of factorial products. The permitted coefficient perturbations show that the valuation argument can survive changes in the intermediate coefficients.\par

\begin{atlasblock}[standard]
\textbf{\sffamily Example  /  y\ensuremath{{}_3}}\par\smallskip
\color{black} The polynomial is \(15x^3+15x^2+6x+1\). Modulo \(2\) it is \(x^3+x^2+1\), which has no root and is therefore irreducible. Its discriminant is \(-135\), a nonsquare, so its Galois group is \(S_3\).\par
\end{atlasblock}\color{black}
\begin{sources}
\textbf{Sources for this section}\par\smallskip
\noindent\hyperlink{ref:FilasetaTrifonov2002}{[FilasetaTrifonov2002]} / Main theorem and §2. All-degree coefficient-perturbation theorem and the Newton-polygon mechanism. \href{https://doi.org/10.1515/crll.2002.069}{Source}.\par
\end{sources}\color{black}
\needspace{8\baselineskip}\section{Legendre parity and the Stieltjes question}
\color{black} Define \(P_n(x)=\frac1{2^n n!}(d/dx)^n(x^2-1)^n\). Since \(P_n(-x)=(-1)^nP_n(x)\), the factor \(x\) is forced at odd indices. Stieltjes' conjecture, as presented by Cullinan-Hajir, asks whether \(P_{2m}\) and \(P_{2m+1}/x\) are irreducible over \(\mathbb Q\) for every \(m\ge1\). Their Galois-group questions concern these even polynomials and the corresponding signed-permutation action.\par

\begin{atlasblock}[standard]
\textbf{\sffamily Example  /  P\ensuremath{{}_4}}\par\smallskip
\color{black} After scaling, \(8P_4=35x^4-30x^2+3\), which is Eisenstein at \(3\). Its pair polynomial \(35y^2-30y+3\) has discriminant \(480\). Thus its pair field is \(\mathbb Q(\sqrt{30})\). The calculation is a certificate for this degree, not an argument for the full family.\par
\end{atlasblock}\color{black}
\begin{sources}
\textbf{Sources for this section}\par\smallskip
\noindent\hyperlink{ref:CullinanHajir2014}{[CullinanHajir2014]} / §1. Stieltjes' irreducibility conjecture and parity-reduced Galois groups. \href{https://doi.org/10.1016/j.indag.2014.01.004}{Source}.\par
\end{sources}\color{black}
\needspace{8\baselineskip}\section{A 2026 preprint claim changes the status record}

\begin{atlasblock}[boundary]
\textbf{\sffamily Preprint claim  /  not independently verified here}\par\smallskip
\color{black} Zikang Deng's September 2026 preprint claims a proof of Stieltjes' irreducibility conjecture in every degree. Version 2 uses local and real estimates for a resultant. The claim is recorded separately from the published 2014 account; this atlas has not independently checked the complete proof or verified a peer-reviewed publication.\par
\end{atlasblock}\color{black}
\begin{sources}
\textbf{Sources for this section}\par\smallskip
\noindent\hyperlink{ref:Deng2026Legendre}{[Deng2026Legendre]} / Theorem 1.1, arXiv:2609.22336v2 (27 September 2026). A preprint claims a complete resolution; this atlas has not independently verified the full proof. \href{https://arxiv.org/abs/2609.22336v2}{Source}.\par
\end{sources}\color{black}
\needspace{8\baselineskip}\section{What transfers between families}
\color{black} The transferable objects are explicit coefficient products, discriminant formulas, Newton slopes, and a well-defined action on pairs. A differential equation or an analogy between names is insufficient. Before applying a result, record the polynomial convention, scalar normalization, affine change of variable, forced factors, and the field over which irreducibility is claimed.\par
\color{black} For Jacobi families with a parameter, it is also essential to distinguish a fixed parameter with degree growing from a fixed degree with parameter varying. Those regimes lead to different arithmetic and geometric methods.\par
\begin{sources}
\textbf{Sources for this section}\par\smallskip
\noindent\hyperlink{ref:CullinanHajirSell2009}{[CullinanHajirSell2009]}. A fixed-degree one-parameter Jacobi family. \href{https://www.numdam.org/articles/10.5802/jtnb.659/}{Source}.\par
\noindent\hyperlink{ref:HajirWong2006}{[HajirWong2006]}. Geometric specialization methods and the distinction between parameter regimes. \href{https://www.numdam.org/item/AIF_2006__56_4_1127_0/}{Source}.\par
\end{sources}\color{black}
\needspace{11\baselineskip}\section*{Questions and research boundaries}
\needspace{7\baselineskip}\subsection*{Are \(P_{2m}\) and \(P_{2m+1}/x\) always irreducible over \(\mathbb Q\)?}
{\small\sffamily\color{atlasgold}Resolution claimed in preprint / The standard Legendre polynomials and all integers m\ensuremath{\ge}1.}\par\smallskip
\color{black} \textbf{Full proof not independently verified here.} This is Stieltjes' conjecture in the cited 2014 paper. Deng's September 2026 preprint claims a full resolution. Both records are linked; the preprint claim has not been independently verified here and is not silently promoted to an established theorem.\par
\begin{sources}
\textbf{Sources for this section}\par\smallskip
\noindent\hyperlink{ref:CullinanHajir2014}{[CullinanHajir2014]} / §1.  \href{https://doi.org/10.1016/j.indag.2014.01.004}{Source}.\par
\noindent\hyperlink{ref:Deng2026Legendre}{[Deng2026Legendre]} / Theorem 1.1, arXiv:2609.22336v2 (27 September 2026). A preprint claims a complete resolution; this atlas has not independently verified the full proof. \href{https://arxiv.org/abs/2609.22336v2}{Source}.\par
\end{sources}\color{black}
\needspace{7\baselineskip}\subsection*{Are the Bessel polynomials irreducible in every positive degree?}
{\small\sffamily\color{atlasgold}Settled in stated scope / \(y_n(x)=\sum_{j=0}^n (n+j)!x^j/(2^j(n-j)!j!)\) over Q, n\ensuremath{\ge}1.}\par\smallskip
\color{black} Yes. The Filaseta-Trifonov theorem proves this and permits arbitrary integral intermediate coefficient multipliers with endpoint multipliers \ensuremath{\pm}1.\par
\begin{sources}
\textbf{Sources for this section}\par\smallskip
\noindent\hyperlink{ref:FilasetaTrifonov2002}{[FilasetaTrifonov2002]} / Main theorem.  \href{https://doi.org/10.1515/crll.2002.069}{Source}.\par
\end{sources}\color{black}
\needspace{6\baselineskip}\par\medskip\begingroup\footnotesize\setlength{\parskip}{1pt}\textbf{\sffamily\color{atlasgreen}Connections}\par
\color{black}\textbf{Uses.} \hyperlink{article:laguerre}{Generalized Laguerre polynomials: two parameter regimes}. Hermite identities reduce the pair polynomials to Laguerre parameters \ensuremath{\pm}1/2.\par
\color{black}\textbf{Uses.} \hyperlink{article:newton-polygons}{Newton polygons, factor degrees, and large cycles}. Factorial products constrain factor degrees.\par
\color{black}\textbf{Generalizes.} \hyperlink{article:reciprocity}{Reciprocal symmetry and the group on root pairs}. Sign pairs lead to the same wreath-product architecture as inverse pairs.\par
\color{black}\textbf{Applies to.} \hyperlink{article:specialization}{Specialization and Hilbert irreducibility}. Parameterized Jacobi families connect to fixed-degree specialization theorems.\par
\endgroup
\chapter{Eulerian polynomials: descents, shifts, and inverse pairs}\label{article:eulerian}\hypertarget{article:eulerian}{}\color{black}
{\sffamily\small\color{atlasgold}family / Intermediate / Q / Published results}\par\medskip
Descent statistics lead to reciprocal polynomials. A prime-index Eisenstein argument and an exact quartic example show how combinatorics, local arithmetic, and Galois theory fit together.\par\medskip
\textbf{Questions answered}\begin{itemize}
\item Why is x+1 removed at even indices?
\item Why does an Eisenstein shift work at prime indices?
\item Why is the natural group a wreath product rather than S\_\allowbreak{}(n\ensuremath{-}1)?
\end{itemize}
{\small\textbf{Prerequisites:} \hyperlink{article:foundations}{Fields, factors, and the action on roots}}\par
\needspace{8\baselineskip}\section{Definition, motivation, and normalization}
\color{black} For \(n\ge1\), define \(A_n(x)=\sum_{\sigma\in S_n}x^{\operatorname{des}(\sigma)}\), where a descent is an index \(i\) with \(\sigma(i)>\sigma(i+1)\). These polynomials connect permutation statistics with the rational generating function for powers:\par

\atlasdisplay{\sum_{k\ge1}k^n x^{k-1}=\frac{A_n(x)}{(1-x)^{n+1}}.}
\color{black} Complementing the values of a permutation gives \(x^{n-1}A_n(x^{-1})=A_n(x)\). Set \(\widetilde A_n=A_n/(x+1)\) for even \(n\), and \(\widetilde A_n=A_n\) for odd \(n\). This notation distinguishes the normalized Eulerian polynomial from the Appell Euler polynomial \(E_n\). For \(n\ge3\), \(\widetilde A_n\) is monic, reciprocal, and has positive even degree.\par
\begin{sources}
\textbf{Sources for this section}\par\smallskip
\noindent\hyperlink{ref:FoataSchutzenberger1970}{[FoataSchutzenberger1970]} / Chapters II and IV. Permutation statistics and generating functions. \href{https://www.mat.univie.ac.at/~slc/books/foaschuetz.html}{Authorized digital retypesetting}.\par
\end{sources}\color{black}
\needspace{8\baselineskip}\section{Why there are inverse pairs and only one forced factor}
\color{black} The recurrence\par

\atlasdisplay{A_n=(1+(n-1)x)A_{n-1}+x(1-x)A'_{n-1},\qquad A_1=1,}
\color{black} gives simple negative roots by induction. At each negative root of \(A_{n-1}\), the factor \(x(1-x)\) reverses the sign of its derivative. Together with the signs at \(0\) and at negative infinity, this places one root of \(A_n\) in each intervening interval and one on each outside interval. Reciprocity then pairs roots by inversion. Since \(A_n(1)=n!\), the only possible fixed point of inversion is \(-1\). It occurs exactly when \(n\) is even, and simplicity makes its multiplicity one.\par
\begin{sources}
\textbf{Sources for this section}\par\smallskip
\noindent\hyperlink{ref:FoataSchutzenberger1970}{[FoataSchutzenberger1970]} / Chapters II, IV, and V. Eulerian identities and the value at \ensuremath{-}1; the interlacing argument is explained here. \href{https://www.mat.univie.ac.at/~slc/books/foaschuetz.html}{Authorized digital retypesetting}.\par
\end{sources}\color{black}
\needspace{8\baselineskip}\section{A recent analytic connection: the limiting zero distribution}
\color{black} Let \(\rho_{n,1},\ldots,\rho_{n,n-1}<0\) be the roots of this article's \(A_n\). Then\par

\atlasdisplay{\frac1{n-1}\sum_{j=1}^{n-1}\delta_{-\rho_{n,j}}\ \Longrightarrow\ \frac{\mathbf1_{t>0}}{t(\pi^2+\log^2t)}\,dt.}
\color{black} Melotti gives a combinatorial proof of this log-Cauchy limit, using exact stabilized moments after the transformation \(u=(1-\rho)^{-1}\). His convention uses \(xA_n(x)\), including a zero root; deleting that one atom and changing the normalization leaves the weak limit unchanged. This result describes the roots as real numbers. It does not settle their organization into Galois orbits or the irreducibility conjecture.\par
\begin{sources}
\textbf{Sources for this section}\par\smallskip
\noindent\hyperlink{ref:Melotti2025Eulerian}{[Melotti2025Eulerian]} / Theorems 1-2; arXiv:2507.15908v2. The limiting law is older; the paper supplies a new proof via stabilized moments. \href{https://arxiv.org/abs/2507.15908v2}{Source}.\par
\end{sources}\color{black}
\needspace{8\baselineskip}\section{A complete theorem at odd prime indices}

\begin{atlasblock}[standard]
\textbf{\sffamily Theorem  /  prime-index irreducibility}\par\smallskip
\color{black} For every odd prime \(p\), \(A_p\) is irreducible over \(\mathbb Q\).\par
\end{atlasblock}\color{black}

\begin{atlasblock}[standard]
\color{black} Fermat's congruence in the generating identity gives\par

\atlasdisplay{\frac{A_p(x)}{(1-x)^{p+1}}\equiv\sum_{k\ge1}kx^{k-1}=\frac1{(1-x)^2}\pmod p.}
\color{black} Therefore \(A_p(1+t)\equiv t^{p-1}\pmod p\). The shifted polynomial is monic, while its constant term is \(A_p(1)=p!\), divisible by \(p\) but not \(p^2\). Eisenstein applies. This is the \(r=1\) instance of Sharma-Straus, Theorem 6.\par
\end{atlasblock}\color{black}
\begin{sources}
\textbf{Sources for this section}\par\smallskip
\noindent\hyperlink{ref:SharmaStraus1974}{[SharmaStraus1974]} / Theorem 6, r=1. The classical prime-index case within the Euler-Frobenius family. \href{https://doi.org/10.4153/CMB-1974-052-x}{Source}.\par
\end{sources}\color{black}
\needspace{8\baselineskip}\section{A stronger factor-degree theorem for every index}

\begin{atlasblock}[standard]
\textbf{\sffamily Theorem  /  largest-prime factor bound}\par\smallskip
\color{black} For \(n\ge3\), let \(p\) be the largest prime at most \(n\). The normalized polynomial \(\widetilde A_n\) has an irreducible factor over \(\mathbb Q\) of degree at least \(p-1\). In particular, \(\widetilde A_p\) and \(\widetilde A_{p+1}\) are irreducible for every odd prime \(p\).\par
\end{atlasblock}\color{black}

\begin{atlasblock}[standard]
\color{black} Put \(r=n-p+1\). Bertrand's postulate gives \(1\le r<p\) and \(n<2p\). Fermat's congruence and the generating identity imply\par

\atlasdisplay{A_n(1+t)\equiv t^{p-1}A_r(1+t)\pmod p.}
\color{black} Since \(A_r(1)=r!\) is a p-adic unit, the coefficient of \(t^{p-1}\) is a unit and every preceding coefficient is divisible by \(p\). The constant term \(n!\) has valuation one. The first Newton edge is therefore from \((0,1)\) to \((p-1,0)\). It gives an irreducible factor of degree \(p-1\) over \(\mathbb Q_p\). Some rational irreducible factor must contain it, hence has at least this degree. At even \(n\), removing \(x+1\) removes only a local linear factor with slope zero, because \(p\) is odd.\par
\end{atlasblock}\color{black}
\color{black} This is the large-factor phenomenon associated with Heidrich's work. The argument above is given in full; the bound does not exclude an additional smaller factor at every composite index.\par
\begin{sources}
\textbf{Sources for this section}\par\smallskip
\noindent\hyperlink{ref:Heidrich1984Eulerian}{[Heidrich1984Eulerian]}. The large-factor result; original full text was not accessible in this review, so no theorem number is asserted. \href{https://doi.org/10.1016/0022-314X(84)90050-7}{Source}.\par
\noindent\hyperlink{ref:FoataSchutzenberger1970}{[FoataSchutzenberger1970]} / Chapter IV. Generating identity used in the displayed proof. \href{https://www.mat.univie.ac.at/~slc/books/foaschuetz.html}{Authorized digital retypesetting}.\par
\end{sources}\color{black}
\needspace{8\baselineskip}\section{An exact Galois group at index five}

\begin{atlasblock}[standard]
\textbf{\sffamily Example  /  a full signed group}\par\smallskip

\atlasdisplay{A_5=x^4+26x^3+66x^2+26x+1=x^2q(x+x^{-1}),\quad q(y)=y^2+26y+64.}
\color{black} The splitting field of \(q\) is \(M=\mathbb Q(\sqrt{105})\). If its roots are \(y_1,y_2\), then\par

\atlasdisplay{(y_1^2-4)(y_2^2-4)=q(2)q(-2)=1920=64\cdot30.}
\color{black} Each radicand has norm \(1920\), a nonsquare in \(\mathbb Q\), so neither is a square in \(M\). Their product is also a nonsquare in \(M\): a rational number is a square in \(\mathbb Q(\sqrt{105})\) only if it or its quotient by \(105\) is a rational square; neither condition holds for \(1920\). The two classes are independent, so the splitting field has degree \(2\cdot4=8\). Thus \(\operatorname{Gal}(A_5/\mathbb Q)=C_2\wr S_2\), the dihedral group of order eight.\par
\end{atlasblock}\color{black}
\begin{sources}
\textbf{Sources for this section}\par\smallskip
\noindent\hyperlink{ref:AndersonBertelliODorney2026}{[AndersonBertelliODorney2026]} / §2. The reciprocal group framework; this index-five calculation is worked out directly. \href{https://arxiv.org/abs/2406.18970v2}{Source}.\par
\end{sources}\color{black}
\needspace{8\baselineskip}\section{What the prime theorem does and does not settle}
\color{black} The prime and prime-plus-one cases prove transitivity on individual roots. They do not determine the group on inverse pairs or the quadratic kernel. For arbitrary indices, the largest-prime bound supplies a large factor, while the 1974 source presents broader irreducibility conjectures for its Euler-Frobenius family. Its broader conjectures should be displayed with their source and date; a finite table of irreducible examples is not an all-index proof.\par
\color{black} A practical study should record \(\widetilde A_n\), its trace polynomial, good-prime factorization patterns, discriminant square classes, and exact certificates separately. This allows an observed pattern to be replaced by a theorem without changing the underlying data.\par
\begin{sources}
\textbf{Sources for this section}\par\smallskip
\noindent\hyperlink{ref:SharmaStraus1974}{[SharmaStraus1974]} / p. 266 and Theorem 6. Conjectural scope in the original paper and one proven range. \href{https://doi.org/10.4153/CMB-1974-052-x}{Source}.\par
\end{sources}\color{black}
\needspace{11\baselineskip}\section*{Questions and research boundaries}
\needspace{7\baselineskip}\subsection*{Is the normalized Eulerian polynomial irreducible for every index n\ensuremath{\ge}3?}
{\small\sffamily\color{atlasgold}Conjecture in cited source / The classical r=1 Euler-Frobenius specialization over Q; remove x+1 at even indices.}\par\smallskip
\color{black} The source formulates irreducibility beyond its proven ranges. The prime and prime-plus-one cases and a general large-factor bound are proved above. This entry records the conjectural status of the cited 1974 source and does not assert that every subsequent advance has been excluded.\par
\begin{sources}
\textbf{Sources for this section}\par\smallskip
\noindent\hyperlink{ref:SharmaStraus1974}{[SharmaStraus1974]} / p. 266 and Theorem 6.  \href{https://doi.org/10.4153/CMB-1974-052-x}{Source}.\par
\noindent\hyperlink{ref:Heidrich1984Eulerian}{[Heidrich1984Eulerian]}. General large-factor results. \href{https://doi.org/10.1016/0022-314X(84)90050-7}{Source}.\par
\end{sources}\color{black}
\needspace{7\baselineskip}\subsection*{For which indices is the group the full signed-permutation group?}
{\small\sffamily\color{atlasgold}Research direction / The separable reciprocal polynomial \(\widetilde A_n\), of degree 2m.}\par\smallskip
\color{black} There are two tasks: determine the group of the degree-m trace polynomial and determine the rank of its quadratic radicands. The index-five example completes both tasks, but does not supply a uniform argument.\par
\begin{sources}
\textbf{Sources for this section}\par\smallskip
\noindent\hyperlink{ref:AndersonBertelliODorney2026}{[AndersonBertelliODorney2026]} / §2.  \href{https://arxiv.org/abs/2406.18970v2}{Source}.\par
\end{sources}\color{black}
\needspace{6\baselineskip}\par\medskip\begingroup\footnotesize\setlength{\parskip}{1pt}\textbf{\sffamily\color{atlasgreen}Connections}\par
\color{black}\textbf{Uses.} \hyperlink{article:eisenstein-dumas}{Schönemann-Eisenstein and one-edge Newton polygons}. The shift x\ensuremath{\mapsto}1+x exposes a prime dividing every nonleading coefficient.\par
\color{black}\textbf{Uses.} \hyperlink{article:reciprocity}{Reciprocal symmetry and the group on root pairs}. The trace polynomial separates the pair action from its quadratic kernel.\par
\color{black}\textbf{Contrasts with.} \hyperlink{article:bernoulli-euler}{Bernoulli and Euler polynomials: symmetry, irreducibility, and simple roots}. Appell Euler polynomials use a different generating function and have different forced factors.\par
\color{black}\textbf{Uses.} \hyperlink{article:frobenius}{Frobenius: from factor degrees to permutations}. Good-prime factorizations give further cycle constraints.\par
\endgroup
\chapter{Bernoulli and Euler polynomials: symmetry, irreducibility, and simple roots}\label{article:bernoulli-euler}\hypertarget{article:bernoulli-euler}{}\color{black}
{\sffamily\small\color{atlasgold}family / Intermediate / Q / Published results}\par\medskip
Reflection dictates the root pairs; denominator arithmetic gives infinite irreducible Bernoulli subfamilies. Multiple roots are completely classified, and even Bernoulli groups have an unconditional derangement theorem.\par\medskip
\textbf{Questions answered}\begin{itemize}
\item How do Euler, Bernoulli, and Eulerian polynomials differ?
\item Which rational factors are forced by the generating identities?
\item Why must repeated factors be checked before describing a root action?
\item Which infinite Bernoulli subfamilies are known to be irreducible?
\item Is the multiple-root problem settled?
\item What can one prove about even Bernoulli Galois groups without assuming irreducibility?
\end{itemize}
{\small\textbf{Prerequisites:} \hyperlink{article:foundations}{Fields, factors, and the action on roots}}\par
\needspace{8\baselineskip}\section{Definitions and the reason these families matter}
\color{black} Use the standard Appell conventions\par

\atlasdisplay{\frac{te^{xt}}{e^t-1}=\sum_{n\ge0}B_n(x)\frac{t^n}{n!},\qquad \frac{2e^{xt}}{e^t+1}=\sum_{n\ge0}E_n(x)\frac{t^n}{n!}.}
\color{black} The Bernoulli difference identity \(B_n(x+1)-B_n(x)=nx^{n-1}\) explains their role in sums of powers. The Euler identity \(E_n(x+1)+E_n(x)=2x^n\) is the corresponding alternating relation. Both families satisfy the Appell derivative rule \(P_n'=nP_{n-1}\). These identities make their roots, common factors, and multiplicities natural arithmetic questions.\par
\color{black} The notation \(E_n\) here always denotes the Euler polynomial defined above, not a normalized Eulerian descent polynomial.\par
\begin{sources}
\textbf{Sources for this section}\par\smallskip
\noindent\hyperlink{ref:DLMFBernoulliEuler}{[DLMFBernoulliEuler]} / §§24.2 and 24.4. Generating functions, difference equations, and derivative identities. \href{https://dlmf.nist.gov/24}{Source}.\par
\end{sources}\color{black}
\needspace{8\baselineskip}\section{Reflection dictates the ambient group}

\begin{atlasblock}[standard]
\textbf{\sffamily Identity  /  reflection}\par\smallskip

\atlasdisplay{B_n(1-x)=(-1)^nB_n(x),\qquad E_n(1-x)=(-1)^nE_n(x).}
\end{atlasblock}\color{black}
\color{black} Replacing \(x\) by \(z+1/2\) makes each polynomial even or odd according to \(n\). Hence for odd \(n\), \(x-1/2\) is forced. After removing that factor, the polynomial is a polynomial in \((x-1/2)^2\). Its nonfixed roots come in pairs \(\{\alpha,1-\alpha\}\). For a separable remaining degree \(2m\), the Galois action embeds in \(C_2\wr S_m\), just as for sign pairs of an even polynomial.\par
\color{black} Irreducibility and the Galois group should therefore be asked only after specifying which forced factors have been removed and whether repeated factors are retained.\par
\begin{sources}
\textbf{Sources for this section}\par\smallskip
\noindent\hyperlink{ref:DLMFBernoulliEuler}{[DLMFBernoulliEuler]} / §24.4(ii). Reflection identities; the group restriction follows from the displayed pairing. \href{https://dlmf.nist.gov/24}{Source}.\par
\noindent\hyperlink{ref:Cox2012}{[Cox2012]}. Galois actions preserving a root partition. \href{https://onlinelibrary.wiley.com/doi/book/10.1002/9781118218457}{Source}.\par
\end{sources}\color{black}
\needspace{8\baselineskip}\section{The elementary rational factors}
\color{black} For odd \(n\ge3\), the Bernoulli values at \(0\) and \(1\) vanish, so \(x(x-1)(x-1/2)\) divides \(B_n(x)\). For even \(n\ge2\), the Euler polynomial vanishes at \(0\) and \(1\), so \(x(x-1)\) divides \(E_n(x)\). These statements follow directly from the generating functions and reflection. They are the starting normalization, not a claim that all remaining factors are irreducible.\par
\color{black} Carlitz and Brillhart are classical arithmetic sources. Their results concern these Appell polynomials and cannot be transferred by name to the Euler-Frobenius or Eulerian families.\par
\begin{sources}
\textbf{Sources for this section}\par\smallskip
\noindent\hyperlink{ref:DLMFBernoulliEuler}{[DLMFBernoulliEuler]} / §§24.2 and 24.4. Special values and reflection. \href{https://dlmf.nist.gov/24}{Source}.\par
\noindent\hyperlink{ref:Carlitz1952BernoulliEuler}{[Carlitz1952BernoulliEuler]}. Classical irreducibility investigations in the Appell families. \href{https://doi.org/10.1215/S0012-7094-52-01948-0}{Source}.\par
\noindent\hyperlink{ref:Brillhart1969}{[Brillhart1969]}. Factorization and roots of Euler and Bernoulli polynomials. \href{https://doi.org/10.1515/crll.1969.234.45}{Source}.\par
\end{sources}\color{black}
\needspace{8\baselineskip}\section{An infinite irreducible subfamily: denominators and base-p digits}

\begin{atlasblock}[standard]
\textbf{\sffamily Theorem  /  Carlitz's even-index criterion}\par\smallskip
\color{black} Let \(p\) be prime and \(n=k(p-1)p^r\) be even, with \(1\le k<p\) and \(r\ge0\). Then \(B_n(x)\) is irreducible over \(\mathbb Q\). For \(p=2\), the evenness condition means \(n=2^r\) with \(r\ge1\).\par
\end{atlasblock}\color{black}

\begin{atlasblock}[standard]
\textbf{\sffamily Proof mechanism  /  an Eisenstein reciprocal}\par\smallskip
\color{black} Write \(B_j=B_j(0)\). Von Staudt-Clausen says that for even \(j\ge2\), the denominator of \(B_j\) contains \(p\) exactly when \(p-1\mid j\), and then only to the first power. The base-\(p\) digits of \(n=(p-k)p^r+(k-1)p^{r+1}\) have sum \(p-1\). If \(0<j<n\) and \(p-1\mid j\), Lucas's congruence gives \(p\mid\binom nj\): otherwise each digit of \(j\) would be bounded by the corresponding digit of \(n\), forcing its positive digit sum, a multiple of \(p-1\), to equal \(p-1\), and hence \(j=n\).\par
\color{black} Consequently every nonconstant coefficient of \(pB_n(x)\) is divisible by \(p\) in \(\mathbb Z_p\), whereas its constant coefficient \(pB_n\) is a unit. Its leading coefficient is exactly \(p\). The exceptional Bernoulli number \(B_1=-1/2\) causes no problem: for \(p=2\), its coefficient is \(-n\), which is even. Thus \(px^nB_n(1/x)\) is Eisenstein over \(\mathbb Z_p\). Taking reciprocals preserves irreducibility.\par
\end{atlasblock}\color{black}

\begin{atlasblock}[standard]
\textbf{\sffamily Example  /  degree four at p = 5}\par\smallskip

\atlasdisplay{30B_4(x)=30x^4-60x^3+30x^2-1.}
\color{black} Its reciprocal \(-x^4+30x^2-60x+30\) is Eisenstein at \(5\). This proves irreducibility over \(\mathbb Q\), despite the reflection symmetry.\par
\end{atlasblock}\color{black}
\begin{sources}
\textbf{Sources for this section}\par\smallskip
\noindent\hyperlink{ref:Carlitz1952BernoulliEuler}{[Carlitz1952BernoulliEuler]}. Original irreducibility source. The stated scope is explicitly restated in Kimura's report; no unchecked original theorem number is claimed. \href{https://doi.org/10.1215/S0012-7094-52-01948-0}{Source}.\par
\noindent\hyperlink{ref:Kimura2003Bernoulli}{[Kimura2003Bernoulli]} / §2, p. 3, theorem attributed to Carlitz (1952). The report uses the reversed coefficient convention for a p-Eisenstein polynomial; the reciprocal formulation above uses the usual convention. \href{https://www.cit.nihon-u.ac.jp/laboratorydata/kenkyu/publication/report/70.pdf}{Source}.\par
\noindent\hyperlink{ref:DLMFBernoulliEuler}{[DLMFBernoulliEuler]} / §24.10(i), equation (24.10.1). Bernoulli denominators; the displayed proof supplies the digit argument explicitly. \href{https://dlmf.nist.gov/24}{Source}.\par
\end{sources}\color{black}
\needspace{8\baselineskip}\section{Two degree-five examples prevent a false generalization}

\begin{atlasblock}[standard]
\textbf{\sffamily Example  /  Bernoulli}\par\smallskip

\atlasdisplay{B_5(x)=x(x-1)(x-\tfrac12)(x^2-x-\tfrac13).}
\color{black} After removing the three rational factors, the quadratic has discriminant \(7/3\), so its splitting field is \(\mathbb Q(\sqrt{21})\).\par
\end{atlasblock}\color{black}

\begin{atlasblock}[standard]
\textbf{\sffamily Example  /  Euler}\par\smallskip

\atlasdisplay{E_5(x)=(x-\tfrac12)(x^2-x-1)^2.}
\color{black} The residual polynomial is a square. Its distinct roots generate only \(\mathbb Q(\sqrt5)\). Thus a statement that every Euler polynomial becomes irreducible or separable after removing just its obvious rational factor would already fail in degree five. The splitting field is defined by the distinct roots, and its permutation degree must be recorded accordingly.\par
\end{atlasblock}\color{black}
\begin{sources}
\textbf{Sources for this section}\par\smallskip
\noindent\hyperlink{ref:Brillhart1969}{[Brillhart1969]}. The Euler/Bernoulli factorization literature; the displayed identities are verified by expansion. \href{https://doi.org/10.1515/crll.1969.234.45}{Source}.\par
\noindent\hyperlink{ref:DLMFBernoulliEuler}{[DLMFBernoulliEuler]} / §24.12(iv). The repeated-factor example E\_\allowbreak{}5. \href{https://dlmf.nist.gov/24}{Source}.\par
\end{sources}\color{black}
\needspace{8\baselineskip}\section{The multiple-root problem is completely settled}

\begin{atlasblock}[standard]
\textbf{\sffamily Theorem  /  Brillhart and Dilcher}\par\smallskip
\color{black} Every \(B_n(x)\), \(n\ge1\), is squarefree. Among the Euler polynomials, only \(E_5(x)\) has multiple roots; its two irrational roots each have multiplicity two.\par
\end{atlasblock}\color{black}
\color{black} Brillhart proved the odd Bernoulli case and restricted any possible even-degree multiple root to a root of \(x^2-x-b\), with \(b\) a positive odd integer. Dilcher's Theorem 2 eliminates these remaining possibilities. The method combines centered expansions, Bernoulli denominator congruences, and the requirement that a multiple root also annihilate the adjacent Appell polynomial. Brillhart's Euler classification is also recorded in Dilcher's §4.\par
\color{black} Thus multiplicity is settled for these standard families. Irreducibility of their normalized factors and identification of their Galois groups remain distinct arithmetic tasks.\par
\begin{sources}
\textbf{Sources for this section}\par\smallskip
\noindent\hyperlink{ref:Dilcher2008MultipleZeros}{[Dilcher2008MultipleZeros]} / Theorem 2, p. 150; proof in §3; Euler discussion in §4. Complete Bernoulli squarefreeness theorem and the earlier Euler classification. \href{https://doi.org/10.4064/aa134-2-6}{Source}.\par
\noindent\hyperlink{ref:Brillhart1969}{[Brillhart1969]}. Original Euler classification and earlier Bernoulli restrictions. \href{https://doi.org/10.1515/crll.1969.234.45}{Source}.\par
\noindent\hyperlink{ref:DLMFBernoulliEuler}{[DLMFBernoulliEuler]} / §24.12(iv). Authoritative concise statement of both complete multiplicity classifications. \href{https://dlmf.nist.gov/24}{Source}.\par
\end{sources}\color{black}
\needspace{8\baselineskip}\section{A group theorem that does not require global irreducibility}

\begin{atlasblock}[standard]
\textbf{\sffamily Theorem  /  Patel-Siksek}\par\smallskip
\color{black} For every even \(n\ge2\), \(\operatorname{Gal}(B_n/\mathbb Q)\) contains an element fixing none of the roots. Also \(B_n(x)\) has no root in \(\mathbb Q_2\).\par
\end{atlasblock}\color{black}
\color{black} The proof uses the \(2\)-adic Newton polygon, the pairing \(\alpha\leftrightarrow1-\alpha\), and a decomposition-group argument. Chebotarev then supplies a positive density of primes at which \(B_n\) has no root. This connects the family's Galois action to obstructions for equations involving sums of consecutive powers. A permutation without fixed points does not by itself imply that the entire Galois action is transitive.\par
\begin{sources}
\textbf{Sources for this section}\par\smallskip
\noindent\hyperlink{ref:PatelSiksek2017Powers}{[PatelSiksek2017Powers]} / Proposition 3.1, Theorem 4, and §§3.2, 5. The group and local-root conclusions are unconditional, independent of the conjectured irreducibility of every even Bernoulli polynomial. \href{https://doi.org/10.1007/s40993-016-0068-0}{Source}.\par
\end{sources}\color{black}
\needspace{8\baselineskip}\section{Formulate the remaining problem precisely}
\color{black} A family question should specify the parity and the rational factors removed. It may ask about irreducibility of the remaining polynomial in \(x\), irreducibility of its polynomial in \((x-1/2)^2\), or the Galois action on pairs. These are different questions. Taking a squarefree part changes only the exceptional Euler member \(E_5\) in the standard families discussed here.\par
\color{black} The Appell identity \(P_n'=nP_{n-1}\) turns a multiple-root question into a common-root question for adjacent members. The classification above completes that argument for Bernoulli and Euler polynomials; the same strategy is useful when investigating other Appell families.\par
\begin{sources}
\textbf{Sources for this section}\par\smallskip
\noindent\hyperlink{ref:Dilcher2008MultipleZeros}{[Dilcher2008MultipleZeros]} / Theorem 2 and §4.  \href{https://doi.org/10.4064/aa134-2-6}{Source}.\par
\noindent\hyperlink{ref:DLMFBernoulliEuler}{[DLMFBernoulliEuler]} / §§24.4(vii), 24.12(iv).  \href{https://dlmf.nist.gov/24}{Source}.\par
\end{sources}\color{black}
\needspace{11\baselineskip}\section*{Questions and research boundaries}
\needspace{7\baselineskip}\subsection*{Is every even Bernoulli polynomial irreducible over Q?}
{\small\sffamily\color{atlasgold}Conjecture in cited source / The standard Bernoulli polynomial B\_\allowbreak{}n(x), for every even n\ensuremath{\ge}2; explicitly described as a conjecture by Patel-Siksek (2017), §3.}\par\smallskip
\color{black} Carlitz's criterion settles infinite index families. The unconditional derangement theorem does not settle irreducibility. The status label records the conjecture in the cited paper; it is not a claim that every later preprint has been ruled out.\par
\begin{sources}
\textbf{Sources for this section}\par\smallskip
\noindent\hyperlink{ref:PatelSiksek2017Powers}{[PatelSiksek2017Powers]} / §3, immediately after Proposition 3.1.  \href{https://doi.org/10.1007/s40993-016-0068-0}{Source}.\par
\noindent\hyperlink{ref:Carlitz1952BernoulliEuler}{[Carlitz1952BernoulliEuler]}.  \href{https://doi.org/10.1215/S0012-7094-52-01948-0}{Source}.\par
\end{sources}\color{black}
\needspace{7\baselineskip}\subsection*{What is the Galois group after all specified forced factors are removed?}
{\small\sffamily\color{atlasgold}Research direction / A stated parity class in the Bernoulli or Appell Euler family, with an explicit squarefree normalization.}\par\smallskip
\color{black} The degree-five Euler example shows why the normalization is part of the question. Determine the pair polynomial and its quadratic kernel separately; the present entry makes no unsupported all-index irreducibility assertion.\par
\begin{sources}
\textbf{Sources for this section}\par\smallskip
\noindent\hyperlink{ref:Carlitz1952BernoulliEuler}{[Carlitz1952BernoulliEuler]}.  \href{https://doi.org/10.1215/S0012-7094-52-01948-0}{Source}.\par
\noindent\hyperlink{ref:Brillhart1969}{[Brillhart1969]}.  \href{https://doi.org/10.1515/crll.1969.234.45}{Source}.\par
\end{sources}\color{black}
\needspace{7\baselineskip}\subsection*{Which Bernoulli or Euler polynomials have repeated roots?}
{\small\sffamily\color{atlasgold}Settled in stated scope / The standard Appell families B\_\allowbreak{}n and E\_\allowbreak{}n over characteristic zero, n\ensuremath{\ge}1.}\par\smallskip
\color{black} None of the Bernoulli polynomials does. The unique Euler exception is \(E_5\), with factorization displayed above.\par
\begin{sources}
\textbf{Sources for this section}\par\smallskip
\noindent\hyperlink{ref:Dilcher2008MultipleZeros}{[Dilcher2008MultipleZeros]} / Theorem 2 and §4.  \href{https://doi.org/10.4064/aa134-2-6}{Source}.\par
\noindent\hyperlink{ref:Brillhart1969}{[Brillhart1969]}.  \href{https://doi.org/10.1515/crll.1969.234.45}{Source}.\par
\noindent\hyperlink{ref:DLMFBernoulliEuler}{[DLMFBernoulliEuler]} / §24.12(iv).  \href{https://dlmf.nist.gov/24}{Source}.\par
\end{sources}\color{black}
\needspace{6\baselineskip}\par\medskip\begingroup\footnotesize\setlength{\parskip}{1pt}\textbf{\sffamily\color{atlasgreen}Connections}\par
\color{black}\textbf{Contrasts with.} \hyperlink{article:eulerian}{Eulerian polynomials: descents, shifts, and inverse pairs}. The similarly named descent polynomials have a different definition and symmetry.\par
\color{black}\textbf{Generalizes.} \hyperlink{article:reciprocity}{Reciprocal symmetry and the group on root pairs}. Reflection becomes sign pairing after an affine change of variable.\par
\color{black}\textbf{Uses.} \hyperlink{article:discriminants}{Discriminants, parity, and ramification}. Repeated factors force a zero discriminant; squarefree parts require separate bookkeeping.\par
\color{black}\textbf{Supplies examples for.} \hyperlink{article:composition-iteration}{Composition, iteration, and arboreal groups}. Centering exhibits a quadratic composition that constrains the Galois action.\par
\color{black}\textbf{Uses.} \hyperlink{article:newton-polygons}{Newton polygons, factor degrees, and large cycles}. Denominator valuations and slopes yield local irreducible factors and Galois information.\par
\color{black}\textbf{Uses.} \hyperlink{article:frobenius}{Frobenius: from factor degrees to permutations}. A derangement produces primes with no linear factor through Chebotarev.\par
\endgroup
\chapter{Fekete polynomials: character sums and reciprocal arithmetic}\label{article:fekete}\hypertarget{article:fekete}{}\color{black}
{\sffamily\small\color{atlasgold}family / Intermediate / Q / Published results}\par\medskip
Quadratic characters impose exact rational factors and reciprocal symmetry. Trace polynomials connect their irreducibility and Galois groups to special values and square classes.\par\medskip
\textbf{Questions answered}\begin{itemize}
\item Which rational factors must be removed?
\item Why can half the degree suffice for irreducibility?
\item What extra information is needed for the full splitting group?
\end{itemize}
{\small\textbf{Prerequisites:} \hyperlink{article:foundations}{Fields, factors, and the action on roots}}\par
\needspace{8\baselineskip}\section{A character polynomial and its motivation}
\color{black} For an odd prime \(p\), let \(\chi_p(a)=(a/p)\) be the Legendre symbol and define\par

\atlasdisplay{F_p(x)=\sum_{a=1}^{p-1}\chi_p(a)x^a.}
\color{black} The coefficients encode the quadratic residues modulo \(p\). These polynomials entered the study of Dirichlet L-functions; their special values also connect to class numbers and generalized Bernoulli numbers. Their coefficients are deterministic and strongly correlated, so theorems about independent random signs cannot be applied directly.\par
\begin{sources}
\textbf{Sources for this section}\par\smallskip
\noindent\hyperlink{ref:MinacNguyenTan2023}{[MinacNguyenTan2023]} / Introduction. Historical motivation and arithmetic special values. \href{https://arxiv.org/abs/2111.05256}{Source}.\par
\end{sources}\color{black}
\needspace{8\baselineskip}\section{Remove the forced factors before asking irreducibility}
\color{black} The identity \(\chi_p(p-a)=\chi_p(-1)\chi_p(a)\) gives \(x^pF_p(x^{-1})=\chi_p(-1)F_p(x)\). Use the normalization\par

\atlasdisplay{f_p(x)=\begin{cases}\dfrac{F_p(x)}{x(1-x)},&p\equiv3\pmod4,\\[6pt]\dfrac{F_p(x)}{x(1-x)^2(x+1)},&p\equiv1\pmod4.\end{cases}}
\color{black} The factors at \(1\) and \(-1\) have precisely the indicated multiplicities. The result is monic and reciprocal. Write \(f_p=x^m g_p(x+x^{-1})\), with \(m=(p-3)/2\) or \((p-5)/2\), respectively. For \(p\ge7\), irreducibility of \(f_p\) over \(\mathbb Q\) is equivalent to irreducibility of \(g_p\). This equivalence uses arithmetic special values, not reciprocity alone.\par
\begin{sources}
\textbf{Sources for this section}\par\smallskip
\noindent\hyperlink{ref:MinacNguyenTan2023}{[MinacNguyenTan2023]} / Proposition 2.1, Definition 3.2, and Corollary 4.2. Exact forced factors and irreducibility equivalence. \href{https://arxiv.org/abs/2111.05256}{Source}.\par
\end{sources}\color{black}
\needspace{8\baselineskip}\section{Work through conductor seven}

\begin{atlasblock}[standard]
\textbf{\sffamily Example  /  an exact degree-eight splitting field}\par\smallskip
\color{black} The quadratic residues modulo \(7\) are \(1,2,4\), so\par

\atlasdisplay{F_7=x+x^2-x^3+x^4-x^5-x^6.}
\color{black} Dividing by \(x(1-x)\) gives\par

\atlasdisplay{f_7=x^4+2x^3+x^2+2x+1=x^2\bigl((x+x^{-1})^2+2(x+x^{-1})-1\bigr).}
\color{black} Thus \(g_7(y)=y^2+2y-1\), with roots \(-1\pm\sqrt2\). Their field is \(M=\mathbb Q(\sqrt2)\), and the quadratic radicands are \(-1\mp2\sqrt2\). Each has norm \(-7\). A square in \(M\) has rational norm equal to a rational square, so neither radicand is a square. Their product is \(-7\), which is not a square in the real field \(M\). The two square classes are independent, giving \([L:M]=4\) and \([L:\mathbb Q]=8\). The group of \(f_7\) is therefore \(C_2\wr S_2\).\par
\end{atlasblock}\color{black}
\color{black} This calculation supplies all parts of the certificate: the trace polynomial, its splitting field, and independence of the two quadratic lifts. Merely finding the degree-two trace polynomial would leave the final step unresolved.\par
\begin{sources}
\textbf{Sources for this section}\par\smallskip
\noindent\hyperlink{ref:Cox2012}{[Cox2012]}. Quadratic extensions and norms; this conductor-seven calculation is explicit. \href{https://onlinelibrary.wiley.com/doi/book/10.1002/9781118218457}{Source}.\par
\end{sources}\color{black}
\needspace{8\baselineskip}\section{The conjecture and the evidence must remain distinct}
\color{black} Mináč-Nguyen-Tân conjecture that the normalized prime-conductor family has full group \(C_2\wr S_m\), with \(m=\deg g_p\). Their paper separates this prediction from proven structural results and finite computations. Here it is recorded as a conjecture in that source, with the precise normalization attached. The broader generalized-character family can have additional symmetries and a smaller allowed group.\par
\begin{sources}
\textbf{Sources for this section}\par\smallskip
\noindent\hyperlink{ref:MinacNguyenTan2023}{[MinacNguyenTan2023]} / Conjecture 4.13 in arXiv:2111.05256v3. The full wreath-product prediction. \href{https://arxiv.org/abs/2111.05256}{Source}.\par
\noindent\hyperlink{ref:MinacNguyenTan2024}{[MinacNguyenTan2024]}. Generalized conductors, cyclotomic factors, and extra symmetry. \href{https://arxiv.org/abs/2206.11778}{Source}.\par
\end{sources}\color{black}
\needspace{8\baselineskip}\section{A reusable research workflow}
\color{black} Start with character identities to find reciprocal or anti-reciprocal symmetry. Next determine the exact cyclotomic factors and divide them out. Construct the trace polynomial and test irreducibility there. Finally combine good-prime cycle types with square-class information from special values to study the full splitting field. Each stage has its own possible failure: an extra cyclotomic factor, a reducible trace polynomial, a small pair group, or a relation among the quadratic radicands.\par
\color{black} This is a useful bridge between analytic character sums and algebraic Galois groups. It also explains why zero-distribution information and finite computations are valuable evidence without being substitutes for a uniform irreducibility theorem.\par
\begin{sources}
\textbf{Sources for this section}\par\smallskip
\noindent\hyperlink{ref:MinacNguyenTan2024}{[MinacNguyenTan2024]}. The organization by forced factors, trace polynomials, and modular Galois tests. \href{https://arxiv.org/abs/2206.11778}{Source}.\par
\end{sources}\color{black}
\needspace{8\baselineskip}\section{A related family: principal characters}
\color{black} For the principal character modulo \(N>1\), the corresponding polynomial is \(F_N^{\mathrm{prin}}(x)=\sum_{1\le a\le N,\,(a,N)=1}x^a\). Möbius inversion gives an exact starting point:\par

\atlasdisplay{F_N^{\mathrm{prin}}(x)=\sum_{d\mid N}\mu(d)\,\frac{x^d(1-x^N)}{1-x^d}.}
\color{black} Each quotient is a finite geometric sum, so the expression is polynomial. Chidambaram-Mináč-Nguyen-Tân study its cyclotomic factors, remaining factors, and Galois groups. This is a distinct family with coefficients 0 and 1; the prime quadratic-character conjecture above must not be transferred to it unchanged.\par
\begin{sources}
\textbf{Sources for this section}\par\smallskip
\noindent\hyperlink{ref:ChidambaramMinacNguyenTan2025}{[ChidambaramMinacNguyenTan2025]} / Introduction. Principal-character family, cyclotomic factors, and associated conjectures. \href{https://jexpmath.org/index.php/jem/article/view/Vol-1Issue-1Paper-4}{Source}.\par
\end{sources}\color{black}
\needspace{11\baselineskip}\section*{Questions and research boundaries}
\needspace{7\baselineskip}\subsection*{Does every normalized prime-conductor Fekete polynomial have the full wreath-product group?}
{\small\sffamily\color{atlasgold}Conjecture in cited source / Odd primes p\ensuremath{\ge}7, the normalization f\_\allowbreak{}p above, and m=deg(g\_\allowbreak{}p).}\par\smallskip
\color{black} The cited conjecture predicts \(\operatorname{Gal}(f_p/\mathbb Q)=C_2\wr S_m\). The prediction includes irreducibility and a maximal quadratic kernel. The site records the conjecture in the cited version, rather than treating computational verification as a proof for all primes.\par
\begin{sources}
\textbf{Sources for this section}\par\smallskip
\noindent\hyperlink{ref:MinacNguyenTan2023}{[MinacNguyenTan2023]} / Conjecture 4.13 in arXiv:2111.05256v3.  \href{https://arxiv.org/abs/2111.05256}{Source}.\par
\end{sources}\color{black}
\needspace{6\baselineskip}\par\medskip\begingroup\footnotesize\setlength{\parskip}{1pt}\textbf{\sffamily\color{atlasgreen}Connections}\par
\color{black}\textbf{Uses.} \hyperlink{article:reciprocity}{Reciprocal symmetry and the group on root pairs}. The trace substitution and quadratic kernel organize the root action.\par
\color{black}\textbf{Contrasts with.} \hyperlink{article:eulerian}{Eulerian polynomials: descents, shifts, and inverse pairs}. Both families are reciprocal after normalization, but their arithmetic inputs differ.\par
\color{black}\textbf{Uses.} \hyperlink{article:frobenius}{Frobenius: from factor degrees to permutations}. Good-prime factorizations constrain the group on the remaining roots.\par
\color{black}\textbf{Supplies examples for.} \hyperlink{article:bernoulli-euler}{Bernoulli and Euler polynomials: symmetry, irreducibility, and simple roots}. Generalized Bernoulli special values connect character sums with arithmetic invariants.\par
\endgroup
\part{Counting and probability}
\chapter{Counting exceptional Galois groups by height}\label{article:height-counting}\hypertarget{article:height-counting}{}\color{black}
{\sffamily\small\color{atlasgold}research / Research / Q / Research overview}\par\medskip
At fixed degree, almost all integer polynomials have full symmetric group. The sharp exceptional order, and its change under reciprocity, reveal which arithmetic obstructions dominate.\par\medskip
\textbf{Questions answered}\begin{itemize}
\item How frequent are exceptional groups in a coefficient box?
\item Why is the fixed-degree exceptional order sharp?
\item Why does reciprocity introduce a logarithmic factor?
\end{itemize}
{\small\textbf{Prerequisites:} \hyperlink{article:foundations}{Fields, factors, and the action on roots}; \hyperlink{article:frobenius}{Frobenius: from factor degrees to permutations}}\par
\needspace{8\baselineskip}\section{The counting problem}
\color{black} Fix \(n\ge3\), and choose a monic polynomial\par

\atlasdisplay{f=x^n+a_1x^{n-1}+\cdots+a_n,\qquad |a_i|\le H.}
\color{black} There are \((2\lfloor H\rfloor+1)^n\) choices. Let \(E_n(H)\) count choices whose Galois group on the roots is not \(S_n\), including reducible and inseparable cases. The degree is fixed throughout; the implied constants below may depend on it.\par

\begin{atlasblock}[standard]
\textbf{\sffamily Published theorem  /  Bhargava}\par\smallskip
\color{black} For fixed \(n\ge3\),\par

\atlasdisplay{E_n(H)=O_n(H^{n-1}).}
\end{atlasblock}\color{black}
\color{black} The order is best possible: the \((2\lfloor H\rfloor+1)^{n-1}\) polynomials with \(a_n=0\) have the rational factor \(x\). Dividing by the total count gives an exceptional probability of order \(H^{-1}\).\par
\begin{sources}
\textbf{Sources for this section}\par\smallskip
\noindent\hyperlink{ref:Bhargava2025}{[Bhargava2025]} / Theorem 1. Sharp upper bound in the fixed-degree box model; the safe stated range here is n\ensuremath{\ge}3. \href{https://annals.math.princeton.edu/2025/201-2/p01}{Source}.\par
\end{sources}\color{black}
\needspace{8\baselineskip}\section{Why low degrees need their own statement}

\begin{atlasblock}[standard]
\textbf{\sffamily Boundary  /  reducible quadratics have a logarithmic factor}\par\smallskip
\color{black} Consider \((x+a)(x+b)\), where \(2\le a\le\sqrt H/2\) and \(a\le b\le H/(2a)\). Both \(a+b\le H\) and \(ab\le H\) for sufficiently large \(H\). Distinct unordered pairs give distinct monic quadratics. Summing the available \(b\)'s over \(a\) yields \(\gg H\log H\) reducible examples.\par
\end{atlasblock}\color{black}
\color{black} Consequently the bound \(O(H^{n-1})\) for \emph{all} exceptional polynomials cannot be stated at \(n=2\). Counting irreducible polynomials with exceptional transitive group is a different question. This distinction matters when abbreviated introductions describe the result as holding in all degrees.\par
\needspace{8\baselineskip}\section{Allowing the leading coefficient to vary}
\color{black} Let \(R_n(H)\) count all degree-\(n\) polynomials in \(\mathbb Z[x]\), reducible over \(\mathbb Q\), whose coefficients have absolute value at most \(H\); the leading coefficient is nonzero and need not be one. Kuba proves\par

\atlasdisplay{R_n(H)\asymp_n H^n\quad(n\ge3),\qquad R_2(H)\asymp H^2\log H.}
\color{black} The ambient count is \(2\lfloor H\rfloor(2\lfloor H\rfloor+1)^n\asymp_n H^{n+1}\), so the reducible proportion has order \(H^{-1}\) for \(n\ge3\). This differs by one power of \(H\) from the monic model.\par
\color{black} The method counts factorizations using height inequalities and lattice-point estimates. More precisely, for \(n\ge3\), all but \(O_n(H^{n-1}(\log H)^2)\) reducible polynomials have an irreducible factor of degree \(n-1\); such a factor leaves a linear cofactor.\par
\begin{sources}
\textbf{Sources for this section}\par\smallskip
\noindent\hyperlink{ref:Kuba2009}{[Kuba2009]} / Theorems 1, 2 and 4; Lemma 1. The polynomial count is nonmonic and reducibility is over Q, so a nonunit integer content alone does not count as reducibility. \href{https://www.sav.sk/index.php?doc=journal-list&journal_article_no=5434&lang=en&part=article_response_page}{Source}.\par
\end{sources}\color{black}
\needspace{8\baselineskip}\section{From modular sieves to arithmetic geometry}
\color{black} At good primes, forbidden Galois groups exclude certain factorization types. A sieve counts coefficient vectors that avoid the corresponding residue classes at many primes. Gallagher's large-sieve method makes this principle quantitative.\par
\color{black} Sharper estimates must also exploit the arithmetic attached to the field generated by a root. Bhargava organizes exceptional permutation groups using their structure and the discriminants of the associated number fields, together with Fourier and geometric-sieve estimates. The counting variable is the coefficient vector; replacing coefficient height by field discriminant changes the enumeration problem because many polynomials may define the same field.\par
\begin{sources}
\textbf{Sources for this section}\par\smallskip
\noindent\hyperlink{ref:Gallagher1973}{[Gallagher1973]}. The large sieve in probabilistic Galois theory. \href{https://doi.org/10.1090/pspum/024/0332694}{Source}.\par
\noindent\hyperlink{ref:Bhargava2025}{[Bhargava2025]} / Introduction and Theorem 2. Group-sensitive bounds and the link with counting number fields. \href{https://annals.math.princeton.edu/2025/201-2/p01}{Source}.\par
\end{sources}\color{black}
\needspace{8\baselineskip}\section{Counting a prescribed permutation action}
\color{black} The same abstract group can act on root sets of different sizes. Write \(B(H;G,\Omega)\) for monic separable integer polynomials of degree \(|\Omega|\), height at most \(H\), and root action isomorphic to \((G,\Omega)\).\par

\begin{atlasblock}[standard]
\textbf{\sffamily Reported preprint theorem  /  Ben-Porath (2026)}\par\smallskip
\color{black} If \(G\le S_m\) is \(k\)-transitive and \([m]_k\) denotes ordered \(k\)-tuples of distinct elements, then, for fixed \(G,m,k\) and every \(\varepsilon>0\),\par

\atlasdisplay{|B(H;G,[m]_k)|\ll_{G,m,k,\varepsilon} H^{m+1+\varepsilon}.}
\end{atlasblock}\color{black}
\color{black} For example, \(S_5\) acting on ordered pairs has degree \(20\); the bound is \(O_\varepsilon(H^{6+\varepsilon})\). This is not the natural degree-five action.\par
\color{black} The method uses coefficients of \(\prod_{\alpha\in A}(x-\alpha)\) to generate the field fixed by the setwise stabilizer of a root subset \(A\). Suitable subset families recover the splitting field while controlling heights. The paper also treats subset actions, primitive wreath-product actions and regular actions.\par
\begin{sources}
\textbf{Sources for this section}\par\smallskip
\noindent\hyperlink{ref:BenPorath2026RootSubsets}{[BenPorath2026RootSubsets]} / Theorem 1.3; Lemma 3.1 and Proposition 3.3. The degree is the size of the specified permutation set, not the natural permutation degree m. \href{https://arxiv.org/abs/2603.14617}{Source}.\par
\end{sources}\color{black}
\needspace{8\baselineskip}\section{Reciprocity changes the dominant exception}
\color{black} For monic reciprocal polynomials of degree \(2m\), only \(m\) coefficients are free and the largest possible group is \(W_m=C_2\wr S_m\).\par

\begin{atlasblock}[standard]
\textbf{\sffamily Published theorem  /  Anderson-Bertelli-O'Dorney}\par\smallskip
\color{black} For each fixed \(m\ge3\), among separable monic reciprocal polynomials with coefficients bounded by \(H\), the number whose group is smaller than \(W_m\) is\par

\atlasdisplay{\asymp_m H^{m-1}\log H.}
\color{black} For nonmonic reciprocal polynomials, the corresponding count is \(\asymp_m H^m\log H\) for every \(m\ge1\).\par
\end{atlasblock}\color{black}
\color{black} The extra logarithm comes from a square-value condition: the leading exceptional contribution is the index-two subgroup of even permutations, rather than just rational factors. The corrected source leaves the monic quartic case \(m=2\) open.\par
\begin{sources}
\textbf{Sources for this section}\par\smallskip
\noindent\hyperlink{ref:AndersonBertelliODorney2026}{[AndersonBertelliODorney2026]} / Theorems 1.1 and 1.3; Remark 1.4, corrected v2. Degree ranges, exceptional orders and the quartic boundary. \href{https://arxiv.org/abs/2406.18970v2}{Source}.\par
\end{sources}\color{black}
\needspace{8\baselineskip}\section{A small reciprocal example}

\begin{atlasblock}[standard]
\textbf{\sffamily Example  /  monic reciprocal quadratics}\par\smallskip
\color{black} For \(f=x^2+ax+1\), the discriminant is \(a^2-4\). If this were an integer square \(d^2\), then \((a-d)(a+d)=4\), whose integer solutions force \(a=\pm2\) and \(d=0\). Thus every separable such quadratic has group \(S_2\). The generic reciprocal counting formula cannot be extended blindly to this monic degree-two case.\par
\end{atlasblock}\color{black}
\color{black} Both ordinary and reciprocal theorems allow a full generic group with density one. Their different exceptional orders record the symmetry constraints of the family.\par
\needspace{11\baselineskip}\section*{Questions and research boundaries}
\needspace{7\baselineskip}\subsection*{What is the sharp exceptional count for monic reciprocal quartics?}
{\small\sffamily\color{atlasgold}Open in cited source / Separable monic degree-four reciprocal integer polynomials, coefficient height H\ensuremath{\to}\ensuremath{\infty}.}\par\smallskip
\color{black} The corrected reciprocal counting paper proves its monic theorem only for half-degree \(m\ge3\) and explicitly leaves \(m=2\) open. Its nonmonic theorem has a different range.\par
\begin{sources}
\textbf{Sources for this section}\par\smallskip
\noindent\hyperlink{ref:AndersonBertelliODorney2026}{[AndersonBertelliODorney2026]} / Remark 1.4 and Remark 1.6, v2. Explicit low-degree exception. \href{https://arxiv.org/abs/2406.18970v2}{Source}.\par
\end{sources}\color{black}
\needspace{6\baselineskip}\par\medskip\begingroup\footnotesize\setlength{\parskip}{1pt}\textbf{\sffamily\color{atlasgreen}Connections}\par
\color{black}\textbf{Generalizes.} \hyperlink{article:specialization}{Specialization and Hilbert irreducibility}. Counting gives quantitative information beyond qualitative Hilbert irreducibility.\par
\color{black}\textbf{Uses.} \hyperlink{article:frobenius}{Frobenius: from factor degrees to permutations}. Modular factorization types provide local conditions for sieves.\par
\color{black}\textbf{Uses.} \hyperlink{article:reciprocity}{Reciprocal symmetry and the group on root pairs}. Inverse pairs change both the parameter dimension and the ambient group.\par
\color{black}\textbf{Contrasts with.} \hyperlink{article:random-models}{Random polynomials: models and proof strategies}. Fixed-degree bounds do not automatically apply when the degree grows.\par
\endgroup
\chapter{Random polynomials: models and proof strategies}\label{article:random-models}\hypertarget{article:random-models}{}\color{black}
{\sffamily\small\color{atlasgold}research / Research / Q / Research overview}\par\medskip
The coefficient law and the limiting parameter determine the problem. A guide to irreducibility, alternating-group containment, and the additional step needed for a full symmetric group.\par\medskip
\textbf{Questions answered}\begin{itemize}
\item Which asymptotic regime is being studied?
\item What separates irreducibility, A\_\allowbreak{}n-containment and S\_\allowbreak{}n?
\item Where do Fourier analysis, sieving and discriminants enter?
\end{itemize}
{\small\textbf{Prerequisites:} \hyperlink{article:foundations}{Fields, factors, and the action on roots}; \hyperlink{article:frobenius}{Frobenius: from factor degrees to permutations}}\par
\needspace{8\baselineskip}\section{Specify the probability space}
\color{black} A phrase such as ``a random polynomial has full Galois group'' is incomplete until the sampling rule and limit are fixed. Three basic models are:\par
\begin{center}\small\begin{tabularx}{\linewidth}{@{}>{\raggedright\arraybackslash}X>{\raggedright\arraybackslash}X>{\raggedright\arraybackslash}X@{}}\toprule
\textbf{Model} & \textbf{Fixed data} & \textbf{Limit} \\
\midrule
Monic coefficient box & Degree \(n\) & Height \(H\to\infty\) \\
Independent bounded coefficients & A probability measure \(\mu\) on \(\mathbb Z\) & Degree \(n\to\infty\) \\
Slowly widening box & A sequence \(L(n)\) & \(n\to\infty\), \(L(n)\to\infty\) \\
\bottomrule\end{tabularx}\end{center}
\color{black} For the second model write \(F_n=x^n+\sum_{j=0}^{n-1}\xi_jx^j\), with independent coefficients of common law \(\mu\). If \(\mu(0)>0\), then \(x\mid F_n\) with probability \(\mu(0)\). An irreducibility probability tending to one therefore requires conditioning on \(\xi_0\ne0\), fixing a nonzero constant term, or excluding zero from the law.\par
\color{black} These are genuine differences in probability spaces. A bound whose constant depends on \(n\) cannot automatically be used while \(n\) grows; a theorem requiring \(L(n)\to\infty\) does not cover a fixed two-point law.\par
\begin{sources}
\textbf{Sources for this section}\par\smallskip
\noindent\hyperlink{ref:Bhargava2025}{[Bhargava2025]} / Theorem 1. Fixed-degree height counting. \href{https://annals.math.princeton.edu/2025/201-2/p01}{Source}.\par
\noindent\hyperlink{ref:BarySorokerGoldgraber2025}{[BarySorokerGoldgraber2025]}. Full symmetric group when the coefficient interval grows, however slowly. \href{https://arxiv.org/abs/2404.13559}{Source}.\par
\end{sources}\color{black}
\needspace{8\baselineskip}\section{A general result and a stronger special case}

\begin{atlasblock}[standard]
\textbf{\sffamily Published theorem  /  Bary-Soroker-Koukoulopoulos-Kozma}\par\smallskip
\color{black} For every fixed, finitely supported, nonconstant law \(\mu\), there is \(\theta>0\) such that, conditional on \(\xi_0\ne0\), the probability of a nonconstant factor of degree less than \(\theta n\) tends to zero.\par
\color{black} If \(\mu\) is uniform on a fixed interval of at least \(35\) consecutive integers, then, under the same conditioning,\par

\atlasdisplay{\Pr(F_n\text{ irreducible and }\operatorname{Gal}(F_n/\mathbb Q)\supseteq A_n)\longrightarrow1.}
\end{atlasblock}\color{black}
\color{black} The first assertion alone does not imply irreducibility: two factors with degrees near \(n/2\) could remain. The second assertion ends with \(A_n\) or \(S_n\), leaving a parity question.\par
\begin{sources}
\textbf{Sources for this section}\par\smallskip
\noindent\hyperlink{ref:BarySorokerKoukoulopoulosKozma2023}{[BarySorokerKoukoulopoulosKozma2023]} / Theorems 1, 2 and 6. Small-factor exclusion; interval-law irreducibility; Galois-group conclusions. \href{https://arxiv.org/abs/2007.14567}{Source}.\par
\end{sources}\color{black}
\needspace{8\baselineskip}\section{Refined errors: isolate the cyclotomic factors}
\color{black} Bazin sharpens the interval-law error. For a fixed interval of \(N\ge10^8\) integers, his Theorem 1.1.2 gives\par

\atlasdisplay{\Pr(F_n\text{ reducible})\le\Pr(\xi_0=0)+O(n^{-1/2}).}
\color{black} More generally, for each fixed \(C>0\), if the interval is sufficiently long in terms of \(C\), Theorem 1.1.4 gives\par

\atlasdisplay{\Pr(F_n\text{ reducible})=\Pr(\exists\omega\in U_C:F_n(\omega)=0)+O_C(n^{-C/2}),}
\color{black} where \(U_C\) consists of zero and the roots of unity of degree less than \(C\) over \(\mathbb Q\). Separating these low-degree obstructions can reveal much smaller errors than a single irreducibility bound.\par
\begin{sources}
\textbf{Sources for this section}\par\smallskip
\noindent\hyperlink{ref:Bazin2025}{[Bazin2025]} / Theorems 1.1.2 and 1.1.4, v2. The constant-term obstruction and cyclotomic refinement are explicit. \href{https://arxiv.org/abs/2402.06643}{Source}.\par
\end{sources}\color{black}
\needspace{8\baselineskip}\section{A precise repeated-root estimate}
\color{black} To discuss separability, consider the full iid model \(P_n(x)=\sum_{j=0}^n\eta_jx^j\), where every coefficient has a fixed bounded integer law \(\mu\) and \(\max_a\mu(a)\le1/2\). Feldheim-Sen prove\par

\atlasdisplay{\Pr(P_n\text{ has a repeated root})=\Pr(P_n\text{ has a repeated root in }\{0,1,-1\})+o(n^{-2}).}
\color{black} If \(\mu(0)=0\), the probability is \(O(n^{-2})\). If \(\mu(0)>0\), the event \(\eta_0=\eta_1=0\) already has probability \(\mu(0)^2\), so a bound tending to zero is impossible without an additional restriction.\par
\color{black} This theorem randomizes the leading coefficient as well. Applying it to a different monic or conditioned ensemble requires justification. Its proof uses anti-concentration of integer evaluations and separates possible roots according to algebraic degree and size.\par
\begin{sources}
\textbf{Sources for this section}\par\smallskip
\noindent\hyperlink{ref:FeldheimSen2017}{[FeldheimSen2017]} / Theorem 1.1 and §1.1. The full iid model and the largest-atom condition are essential to the stated theorem. \href{https://arxiv.org/abs/1603.03811}{Source}.\par
\end{sources}\color{black}
\needspace{8\baselineskip}\section{Root separation and discriminant size}
\color{black} For iid mean-zero, sub-Gaussian coefficients with no atom at zero, Michelen-Yakir prove that the least distance between distinct root indices has scale \(n^{-5/4}\) and a nondegenerate limiting law. Their Corollary 1.2 implies that repeated-root probability tends to zero. It removes the older integer-law largest-atom restriction, while the earlier Littlewood theorem gives a sharper quantitative repeated-root estimate.\par
\color{black} In a separate preprint, for \(P_n(x)=\sum_{j=0}^n\eta_jx^j\) with the additional normalization \(\mathbb E|\eta_j|^2=1\), they report\par

\atlasdisplay{\frac{\log|\operatorname{disc}(P_n)|-2n\log n}{n}\xrightarrow{\Pr}-D_*,\qquad D_*>0.}
\color{black} The mechanism combines a reciprocal form of the discriminant identity, Mahler measure, and root statistics. A large discriminant can still be a perfect square: this analytic size law alone does not settle whether the Galois group is \(A_n\) or \(S_n\). The no-zero-atom hypothesis prevents the persistent repeated root at zero.\par
\begin{sources}
\textbf{Sources for this section}\par\smallskip
\noindent\hyperlink{ref:MichelenYakir2026RootSeparation}{[MichelenYakir2026RootSeparation]} / Theorem 1.1 and Corollary 1.2. The hypotheses here retain mean zero, sub-Gaussian tails and no atom at zero. \href{https://www.sciencedirect.com/science/article/pii/S0001870826002124}{Source}.\par
\noindent\hyperlink{ref:MichelenYakir2025Discriminant}{[MichelenYakir2025Discriminant]} / Theorem 1; §1.3; equation (10). The displayed convergence is a preprint theorem; the no-zero-atom and variance assumptions are explicit in the full statement. \href{https://arxiv.org/abs/2506.12206}{Source}.\par
\end{sources}\color{black}
\needspace{8\baselineskip}\section{How the proof problems separate}
\color{black} For a candidate divisor \(D\in\mathbb F_p[x]\), divisibility of a random polynomial is a random-sum problem in the finite algebra \(\mathbb F_p[x]/(D)\). Independence expresses its Fourier transform as a product. Decay of that product gives approximately uniform residues, subject to exceptional characters and algebraic structure. Several primes give incompatible rational factor degrees; combinatorial estimates control how frequently modular factorizations admit a common subset sum.\par
\color{black} To determine a large group, the modular factors also supply Frobenius cycle types. Probabilistic permutation-group arguments exclude most proper transitive groups. A different route counts roots modulo many primes and uses effective prime-ideal estimates to recover information about global Galois orbits. In the broad theorem of Breuillard-Varjú, this analytic step uses Riemann hypotheses for Dedekind zeta functions.\par
\begin{sources}
\textbf{Sources for this section}\par\smallskip
\noindent\hyperlink{ref:BarySorokerKoukoulopoulosKozma2023}{[BarySorokerKoukoulopoulosKozma2023]} / §2. Fourier estimates, factor-degree sieving and permutation groups. \href{https://arxiv.org/abs/2007.14567}{Source}.\par
\noindent\hyperlink{ref:BreuillardVarju2019}{[BreuillardVarju2019]}. Conditional prime-ideal and random-walk method. \href{https://arxiv.org/abs/1810.13360}{Source}.\par
\end{sources}\color{black}
\needspace{8\baselineskip}\section{The remaining parity obstruction}

\begin{atlasblock}[standard]
\textbf{\sffamily Reported preprint result  /  restricted model  /  Blachar-Breuillard-Kozma, v1}\par\smallskip
\color{black} Take \(F_n=x^n+\sum_{j=1}^{n-1}\xi_jx^j+1\), where \(\xi_1,\ldots,\xi_{n-1}\) are independent with a fixed finitely supported integer law \(\mu\). For\par

\atlasdisplay{H_2(\mu):=-\log\!\sum_{a\in\mathbb Z}\mu(a)^2>12,}
\color{black} the preprint reports \(\Pr(\operatorname{disc}F_n\text{ is a square})=O_\mu(1/\log n)\). This is the constant-term-one subclass of its stated Theorem 1.1.\par
\color{black} Its Corollary 1.2 also reports: when all \(\xi_0,\ldots,\xi_{n-1}\) are independent and uniform on a fixed interval of \(N>e^{12}\) integers, \(x^n+\sum_{j=0}^{n-1}\xi_jx^j\) is irreducible with group \(S_n\), conditional on \(\xi_0\ne0\), with probability \(1-O_\mu(1/\log n)\).\par
\end{atlasblock}\color{black}

\begin{atlasblock}[standard]
\textbf{\sffamily Scope correction  /  zero is a square}\par\smallskip
\color{black} The unqualified square-discriminant assertion cannot hold for every constant-term law allowed in v1. For \(R_n=x^n+\sum_{j=0}^{n-1}\xi_jx^j\), if \(\xi_0,\xi_1\) are independent with law \(\mu\) and \(\mu(0)>0\), then\par

\atlasdisplay{\Pr(\operatorname{disc}R_n=0)\ge\Pr(\xi_0=\xi_1=0)=\mu(0)^2.}
\color{black} Indeed, this event makes \(x^2\mid R_n\). It persists even for a uniform interval with \(N>e^{12}\). Section 8 invokes an \(O(1/n)\) repeated-root estimate without removing this obstruction. The statements above are recorded as preprint claims with the zero-constant obstruction excluded; this restriction does not constitute an independent verification or repair of the full proof.\par
\end{atlasblock}\color{black}
\color{black} The method studies double roots modulo primes and their lifts modulo prime squares to control discriminant valuations. Once \(A_n\)-containment is known, a nonsquare discriminant forces \(S_n\). For a uniform law on \(N\) values, \(H_2=\log N\); the fair-sign law has \(H_2=\log2\) and is outside the entropy range.\par
\begin{sources}
\textbf{Sources for this section}\par\smallskip
\noindent\hyperlink{ref:BlacharBreuillardKozma2026}{[BlacharBreuillardKozma2026]} / Theorem 1.1, Corollary 1.2, and opening of §8, v1. Reported preprint statements; the atlas excludes the zero-constant obstruction and records the scope issue in the original square-discriminant assertion. \href{https://arxiv.org/abs/2609.40123v1}{Source}.\par
\end{sources}\color{black}
\needspace{8\baselineskip}\section{Worked comparison: three short intervals}

\begin{atlasblock}[standard]
\textbf{\sffamily Example  /  read the assumptions first}\par\smallskip
\color{black} For coefficients uniform on \(\{-17,\ldots,17\}\), the 35-point theorem gives high-probability irreducibility and \(A_n\)-containment after conditioning on \(\xi_0\ne0\). Without conditioning, the factor \(x\) appears with probability \(1/35\).\par
\color{black} For coefficients uniform on \(\{-1,1\}\), there is no zero-constant obstruction, but neither the 35-consecutive-integer theorem nor the entropy bound applies. For coefficients uniform on \([-L(n),L(n)]\cap\mathbb Z\) with \(L(n)\to\infty\), Bary-Soroker-Goldgraber instead gives full \(S_n\) with probability tending to one; here the zero-constant probability itself tends to zero.\par
\end{atlasblock}\color{black}
\color{black} When reporting a numerical experiment, record the law, degree, constant-term convention, sample size, and whether the group calculation was certified. An empirical frequency of \(S_n\) is evidence about those samples, not a substitute for an asymptotic theorem.\par
\begin{sources}
\textbf{Sources for this section}\par\smallskip
\noindent\hyperlink{ref:BarySorokerGoldgraber2025}{[BarySorokerGoldgraber2025]}. The widening-interval theorem. \href{https://arxiv.org/abs/2404.13559}{Source}.\par
\end{sources}\color{black}
\needspace{8\baselineskip}\section{Dependent coefficients from a multiplicative function}
\color{black} Choose independent fair signs \(X_p\) for primes \(p\), put \(X_1=1\), and extend completely multiplicatively: \(X_n=\prod_{p^a\parallel n}X_p^a\). Define\par

\atlasdisplay{P_d(x)=\sum_{j=1}^dX_jx^j.}
\color{black} Varjú-Xu prove, assuming RH for the Dedekind zeta functions of all number fields, that for every \(\varepsilon>0\),\par

\atlasdisplay{\Pr(P_d(x)/x\text{ irreducible over }\mathbb Q)\ge1-O_\varepsilon(d^{-1/2+\varepsilon}).}
\color{black} Their theorem also covers the variant that sets \(X_j=0\) at nonsquarefree integers. The compulsory factor \(x\) must be removed; in the second model the actual degree may be less than \(d\). Dependence is substantial: for example, \(X_6=X_2X_3\). The proof combines finite-field root averages with arithmetic equidistribution. This supplies a connection to Fekete polynomials without proving their deterministic irreducibility conjectures.\par
\begin{sources}
\textbf{Sources for this section}\par\smallskip
\noindent\hyperlink{ref:VarjuXu2026}{[VarjuXu2026]} / Theorem 1.1; §1.1; Problem 1.2. Both multiplicative models are specified in the introduction; the conditional conclusion concerns P\_\allowbreak{}d/x. \href{https://arxiv.org/abs/2511.04240}{Source}.\par
\end{sources}\color{black}
\needspace{8\baselineskip}\section{The function-field analogue has a different limiting obstruction}
\color{black} Fix a prime power \(q\), and let\par

\atlasdisplay{f(x,y)=y^n+\sum_{i=0}^{n-1}a_i(x)y^i,\qquad a_i\text{ independent and uniform in }\{a\in\mathbb F_q[x]:\deg a\le d\}.}
\color{black} The splitting field here is taken over \(\mathbb F_q(x)\). Entin-Popov prove that for fixed \(n,q\), as \(d\to\infty\), the probability of irreducibility, separability and Galois group \(S_n\) tends to one. For fixed \(d\ge1,q\), as \(n\to\infty\),\par

\atlasdisplay{\Pr(f\text{ irreducible})\longrightarrow1-q^{-d},\qquad\Pr(G_f\in\{A_n,S_n\}\mid f\text{ irreducible})\longrightarrow1.}
\color{black} Thus even the limiting irreducibility probability changes with the regime. In odd characteristic, their full-\(S_n\) refinement assumes the precise uniform polynomial Chowla conjecture in Conjecture 1.5. These are function-field Galois groups; the cyclic Galois group of an irreducible polynomial over the finite field \(\mathbb F_q\) is a different setting.\par
\begin{sources}
\textbf{Sources for this section}\par\smallskip
\noindent\hyperlink{ref:EntinPopov2024}{[EntinPopov2024]} / Theorems 1.1, 1.2, 1.4 and 1.6; Conjecture 1.5. The full symmetric conclusion in the fixed-d regime is conditional, and Theorem 1.6 assumes odd q. \href{https://www.sciencedirect.com/science/article/pii/S1071579724001059}{Source}.\par
\end{sources}\color{black}
\needspace{11\baselineskip}\section*{Questions and research boundaries}
\needspace{7\baselineskip}\subsection*{Can full symmetric groups be proved for small fixed coefficient laws?}
{\small\sffamily\color{atlasgold}Research direction / Independent bounded coefficients over Q, with any constant-term obstruction removed.}\par\smallskip
\color{black} The reported high-entropy result above leaves small-support laws, especially fair signs, outside its hypotheses. Extending the parity argument requires additional input. This is a boundary of the stated results, not a claim that every small-law problem has the same current status.\par
\begin{sources}
\textbf{Sources for this section}\par\smallskip
\noindent\hyperlink{ref:BlacharBreuillardKozma2026}{[BlacharBreuillardKozma2026]} / Introduction after Theorem 1.1. The authors discuss the entropy assumption as a limitation of their method. \href{https://arxiv.org/abs/2609.40123v1}{Source}.\par
\end{sources}\color{black}
\needspace{6\baselineskip}\par\medskip\begingroup\footnotesize\setlength{\parskip}{1pt}\textbf{\sffamily\color{atlasgreen}Connections}\par
\color{black}\textbf{Applies to.} \hyperlink{article:littlewood}{Littlewood polynomials: signs, parity, and special degrees}. The fair-sign model tests the limits of general bounded-coefficient methods.\par
\color{black}\textbf{Contrasts with.} \hyperlink{article:factor-reversal}{Factor reversal for random binary polynomials}. A special identity for binary coefficients supplies a different approach.\par
\color{black}\textbf{Contrasts with.} \hyperlink{article:height-counting}{Counting exceptional Galois groups by height}. The degree stays fixed while the coefficient box grows.\par
\color{black}\textbf{Uses.} \hyperlink{article:discriminants}{Discriminants, parity, and ramification}. A nonsquare discriminant distinguishes S\_\allowbreak{}n from A\_\allowbreak{}n once containment is proved.\par
\endgroup
\chapter{Littlewood polynomials: signs, parity, and special degrees}\label{article:littlewood}\hypertarget{article:littlewood}{}\color{black}
{\sffamily\small\color{atlasgold}family / Research / Q / Research overview}\par\medskip
Independent signs create strong congruence constraints. Exact examples show both their usefulness and why the all-degree probabilistic problem needs more than reduction modulo two.\par\medskip
\textbf{Questions answered}\begin{itemize}
\item Why is reduction modulo two both useful and restrictive?
\item Which special degrees have unconditional high-probability irreducibility?
\item How can a shifted Newton polygon combine with small-factor exclusion?
\end{itemize}
{\small\textbf{Prerequisites:} \hyperlink{article:foundations}{Fields, factors, and the action on roots}; \hyperlink{article:frobenius}{Frobenius: from factor degrees to permutations}}\par
\needspace{8\baselineskip}\section{The model and its first obstruction}
\color{black} A Littlewood polynomial is \(L_n(x)=\sum_{j=0}^{n}\varepsilon_jx^j\), where each \(\varepsilon_j\in\{-1,1\}\). In the random model the signs are independent and fair. Multiplying by the leading sign makes the polynomial monic without changing its roots or irreducibility; the remaining normalized signs are again independent and fair.\par
\color{black} The only possible rational roots are \(1\) and \(-1\), by the rational-root theorem. If \(n\) is even, \(L_n(1)\) and \(L_n(-1)\) are sums of an odd number of odd integers, so neither vanishes. For \(n=2m-1\), direct counting gives\par

\atlasdisplay{\Pr(L_n(1)=0)=\frac{\binom{2m}{m}}{2^{2m}}\sim\frac1{\sqrt{\pi m}}.}
\color{black} The same marginal probability holds at \(-1\), but the two events are not independent. A reducibility estimate must also control higher-degree factors.\par
\begin{sources}
\textbf{Sources for this section}\par\smallskip
\noindent\hyperlink{ref:BarySorokerHokkenKozmaPoonen2025}{[BarySorokerHokkenKozmaPoonen2025]} / §1. Definition of the random sign model and the all-degree conjecture. \href{https://arxiv.org/abs/2308.04878}{Source}.\par
\end{sources}\color{black}
\needspace{8\baselineskip}\section{Squarefreeness is a separate question}

\begin{atlasblock}[standard]
\textbf{\sffamily Published theorem  /  Peled-Sen-Zeitouni}\par\smallskip
\color{black} For fair independent signs, write \(D_n=\Pr(L_n\text{ has a repeated complex root})\). Then\par

\atlasdisplay{D_n=\begin{cases}\dfrac{8\sqrt3}{\pi n^2}+o(n^{-2}),&4\mid n+1,\\o(n^{-2}),&4\nmid n+1.\end{cases}}
\end{atlasblock}\color{black}
\color{black} The main contribution, when present, comes from repeated roots at \(1\) and \(-1\). This estimates squarefreeness, which is weaker than irreducibility: a product of two distinct irreducible factors is squarefree. It also concerns whether the discriminant vanishes, not whether a nonzero discriminant is a square.\par
\color{black} The proof separates roots by algebraic degree and proximity to the unit circle, then uses anti-concentration and a local central limit calculation for the exceptional roots.\par
\begin{sources}
\textbf{Sources for this section}\par\smallskip
\noindent\hyperlink{ref:PeledSenZeitouni2016}{[PeledSenZeitouni2016]} / Theorems 1.1 and 1.2; §1.1. The displayed constant specializes Theorem 1.2 to mean-zero signs of variance one. \href{https://arxiv.org/abs/1409.2034}{Source}.\par
\end{sources}\color{black}
\needspace{8\baselineskip}\section{A useful deterministic reduction}

\begin{atlasblock}[standard]
\textbf{\sffamily Example  /  every degree-four Littlewood polynomial is irreducible}\par\smallskip
\color{black} Modulo \(2\), all signs become \(1\), so every degree-four Littlewood polynomial reduces to\par

\atlasdisplay{1+x+x^2+x^3+x^4=\Phi_5(x).}
\color{black} The order of \(2\) in \((\mathbb Z/5\mathbb Z)^\times\) is \(4\). Hence \(\Phi_5\) is irreducible over \(\mathbb F_2\), and all these integer polynomials are irreducible over \(\mathbb Q\).\par
\end{atlasblock}\color{black}
\color{black} Generally \(\overline L_n=(x^{n+1}-1)/(x-1)\). If \(n+1\) is an odd prime for which \(2\) is a primitive root, this gives the same deterministic conclusion. For \(n+1=2^r\), however, \(\overline L_n=(x+1)^n\); the elementary reduction test gives no irreducibility certificate.\par
\begin{sources}
\textbf{Sources for this section}\par\smallskip
\noindent\hyperlink{ref:LidlNiederreiter1997}{[LidlNiederreiter1997]}. Cyclotomic factor degrees over finite fields. \href{https://www.cambridge.org/core/books/finite-fields/75BDAA74ABAE713196E718392B9E5E72}{Source}.\par
\end{sources}\color{black}
\needspace{8\baselineskip}\section{A published subsequence theorem}

\begin{atlasblock}[standard]
\textbf{\sffamily Published theorem  /  Bary-Soroker-Hokken-Kozma-Poonen}\par\smallskip
\color{black} There is an absolute \(c>0\) such that, for \(n=2^r-1\),\par

\atlasdisplay{\Pr(L_n\text{ irreducible over }\mathbb Q)\ge1-n^{-c}.}
\color{black} Thus \(\limsup_{n\to\infty}\Pr(L_n\text{ irreducible})=1\). The paper also treats degrees \(n=p^r-1\) when \(p\ge7\) is prime and \(2\) generates \( (\mathbb Z/p^2\mathbb Z)^\times\).\par
\end{atlasblock}\color{black}
\color{black} The powers-of-two argument combines a shifted 2-adic Newton polygon with exclusion of small factors. For \(g(x)=L_n(x+1)\), every coefficient below the leading one is even. The residues of these coefficients modulo \(4\) are independent and uniform in \(\{0,2\}\). With high probability an early coefficient has valuation exactly one, forcing a local irreducible factor of degree very close to \(n\). Any other rational factor would then be small, contradicting the separate small-factor theorem.\par
\begin{sources}
\textbf{Sources for this section}\par\smallskip
\noindent\hyperlink{ref:BarySorokerHokkenKozmaPoonen2025}{[BarySorokerHokkenKozmaPoonen2025]} / Theorem 1.2 and §2. Special degrees and the Newton-polygon argument. \href{https://arxiv.org/abs/2308.04878}{Source}.\par
\end{sources}\color{black}
\needspace{8\baselineskip}\section{Worked shifted certificate in degree seven}

\begin{atlasblock}[standard]
\textbf{\sffamily Example  /  shifted Eisenstein}\par\smallskip
\color{black} Set \(L(x)=1+x+x^2+x^3+x^4+x^5+x^6-x^7\). Modulo \(2\), \(L(x+1)=x^7\). Thus every coefficient except the leading coefficient is even. Its constant coefficient is \(L(1)=6\), which is not divisible by \(4\). Eisenstein at \(2\) proves \(L(x+1)\), and hence \(L(x)\), irreducible over \(\mathbb Q\).\par
\end{atlasblock}\color{black}
\color{black} This is an exact deterministic proof for one polynomial. It illustrates the valuation mechanism, while the probabilistic theorem must handle the cases where the first coefficient of valuation one is not the constant coefficient.\par
\begin{sources}
\textbf{Sources for this section}\par\smallskip
\noindent\hyperlink{ref:ConradIrredQ}{[ConradIrredQ]}. Eisenstein and translation preserve irreducibility. \href{https://kconrad.math.uconn.edu/blurbs/ringtheory/irredtestsoverQ.pdf}{Source}.\par
\end{sources}\color{black}
\needspace{8\baselineskip}\section{What the quoted results do and do not settle}
\color{black} Breuillard-Varjú proves the all-degree high-probability irreducibility statement under Riemann hypotheses for Dedekind zeta functions. The unconditional published special-degree theorem is a limsup statement. Neither formulation alone supplies an unconditional all-degree full-\(S_n\) theorem.\par
\color{black} Kogler's 2026 preprint instead concerns monic \(0/1\) polynomials with constant term \(1\). Binary coefficients retain randomness modulo \(2\), whereas signs do not. Replacing ``binary'' by ``Littlewood'' in that theorem would change its assumptions. The all-degree question below is recorded with the status given in the cited Littlewood paper.\par
\begin{sources}
\textbf{Sources for this section}\par\smallskip
\noindent\hyperlink{ref:BreuillardVarju2019}{[BreuillardVarju2019]}. Conditional growing-degree theorem. \href{https://arxiv.org/abs/1810.13360}{Source}.\par
\noindent\hyperlink{ref:Kogler2026}{[Kogler2026]} / Theorem 1.1, v2. The binary model is different. \href{https://arxiv.org/abs/2609.26771v2}{Source}.\par
\end{sources}\color{black}
\needspace{11\baselineskip}\section*{Questions and research boundaries}
\needspace{7\baselineskip}\subsection*{Does the irreducibility probability tend to one along all degrees?}
{\small\sffamily\color{atlasgold}Open in cited source / Fair independent \ensuremath{\pm}1 coefficients, degree n tending to infinity, over Q.}\par\smallskip
\color{black} The cited paper states this as the all-degree conjecture and proves an unconditional subsequence result. The conditional theorem assumes Dedekind Riemann hypotheses. The atlas keeps these statements distinct and does not upgrade the limsup to a limit.\par
\begin{sources}
\textbf{Sources for this section}\par\smallskip
\noindent\hyperlink{ref:BarySorokerHokkenKozmaPoonen2025}{[BarySorokerHokkenKozmaPoonen2025]} / Equation (1.1), Theorems 1.1-1.2. Source-scoped conjecture and proved subsequence. \href{https://arxiv.org/abs/2308.04878}{Source}.\par
\end{sources}\color{black}
\needspace{6\baselineskip}\par\medskip\begingroup\footnotesize\setlength{\parskip}{1pt}\textbf{\sffamily\color{atlasgreen}Connections}\par
\color{black}\textbf{Uses.} \hyperlink{article:newton-polygons}{Newton polygons, factor degrees, and large cycles}. A long slope after translation forces a large local factor.\par
\color{black}\textbf{Uses.} \hyperlink{article:finite-field-irreducibility}{Irreducibility over finite fields}. The order of 2 modulo an odd prime determines cyclotomic factor degrees.\par
\color{black}\textbf{Applies to.} \hyperlink{article:random-models}{Random polynomials: models and proof strategies}. Signs are a small-support independent-coefficient law.\par
\color{black}\textbf{Contrasts with.} \hyperlink{article:factor-reversal}{Factor reversal for random binary polynomials}. Binary randomness modulo two is unavailable for sign coefficients.\par
\endgroup
\chapter{Factor reversal for random binary polynomials}\label{article:factor-reversal}\hypertarget{article:factor-reversal}{}\color{black}
{\sffamily\small\color{atlasgold}research / Research / Q / Preprint boundary}\par\medskip
A factorization can create a second binary polynomial with the same reciprocal product. This connects an old norm identity to a recent probabilistic approach.\par\medskip
\textbf{Questions answered}\begin{itemize}
\item What does reversing one factor preserve?
\item Why do the new coefficients stay binary?
\item Which probabilistic step remains after the algebraic identity?
\end{itemize}
{\small\textbf{Prerequisites:} \hyperlink{article:foundations}{Fields, factors, and the action on roots}; \hyperlink{article:frobenius}{Frobenius: from factor degrees to permutations}}\par
\needspace{8\baselineskip}\section{The binary model and the reported result}
\color{black} Consider\par

\atlasdisplay{P_n(x)=1+\sum_{j=1}^{n-1}\xi_jx^j+x^n,\qquad \Pr(\xi_j=0)=\Pr(\xi_j=1)=\tfrac12,}
\color{black} with independent \(\xi_j\).\par

\begin{atlasblock}[standard]
\textbf{\sffamily Preprint theorem  /  Kogler, v2, 14 October 2026}\par\smallskip
\color{black} Theorem 1.1 states the unconditional asymptotic\par

\atlasdisplay{\Pr(P_n\text{ reducible over }\mathbb Q)=\sqrt{\frac{2}{\pi n}}+O(n^{-1}).}
\end{atlasblock}\color{black}
\color{black} The leading contribution comes from the factor \(x+1\). This is a precise binary-model statement. It neither determines the full Galois group nor covers independent fair signs.\par
\begin{sources}
\textbf{Sources for this section}\par\smallskip
\noindent\hyperlink{ref:Kogler2026}{[Kogler2026]} / Theorem 1.1, v2. Exact model, asymptotic and preprint status. \href{https://arxiv.org/abs/2609.26771v2}{Source}.\par
\end{sources}\color{black}
\needspace{8\baselineskip}\section{The algebraic mechanism}
\color{black} For a polynomial \(P\) with nonzero constant term define its reversal by \(P^*(x)=x^{\deg P}P(1/x)\). Suppose a binary polynomial with leading and constant coefficient \(1\) factors as \(P=AB\), with \(A,B\in\mathbb Z[x]\) monic. Put \(Q=AB^*\). Then\par

\atlasdisplay{QQ^*=PP^*.}

\begin{atlasblock}[standard]
\textbf{\sffamily Why the new coefficients remain binary}\par\smallskip
\color{black} Write \(P=\sum p_jx^j\) and \(Q=\sum q_jx^j\). The central coefficient in the displayed identity gives \(\sum q_j^2=\sum p_j^2\). Reversal does not change evaluation at \(1\), so \(Q(1)=P(1)\). Since \(p_j^2=p_j\),\par

\atlasdisplay{\sum_jq_j(q_j-1)=0.}
\color{black} Each summand is a nonnegative integer. Therefore every \(q_j\) is \(0\) or \(1\). Its nonzero leading and constant coefficients are necessarily \(1\).\par
\end{atlasblock}\color{black}
\color{black} The identity is algebraic; it proves no probability estimate by itself. To turn it into an irreducibility theorem, one has to show that such alternative products are rare and handle factors for which reversal produces no new polynomial.\par
\begin{sources}
\textbf{Sources for this section}\par\smallskip
\noindent\hyperlink{ref:Ljunggren1960}{[Ljunggren1960]}. An antecedent of factor-reversal arguments. \href{https://doi.org/10.7146/math.scand.a-10593}{Source}.\par
\noindent\hyperlink{ref:filaseta1999smallnorm}{[filaseta1999smallnorm]}. The Euclidean-norm perspective on factorization. \href{https://doi.org/10.1515/9783110285581.143}{Source}.\par
\noindent\hyperlink{ref:Kogler2026}{[Kogler2026]} / Lemma 3.1. The binary twin-factorization step. \href{https://arxiv.org/abs/2609.26771v2}{Source}.\par
\end{sources}\color{black}
\needspace{8\baselineskip}\section{An exact pair of different binary products}

\begin{atlasblock}[standard]
\textbf{\sffamily Example  /  a nontrivial reciprocal-product collision}\par\smallskip
\color{black} Let \(A=1+x+x^3\) and \(B=1+x^4+x^{12}\). The exponent sets of \(P=AB\) and \(Q=AB^*\) are respectively\par

\atlasdisplay{\begin{aligned}\operatorname{supp}(P)&=\{0,1,3,4,5,7,12,13,15\},\\\operatorname{supp}(Q)&=\{0,1,3,8,9,11,12,13,15\}.\end{aligned}}
\color{black} All coefficients are binary because the shifted copies of \(A\) have disjoint supports. The products have the same reciprocal product, but \(Q\ne P\) and \(Q\ne P^*\): the support of \(P^*\) is \(\{0,2,3,8,10,11,12,14,15\}\).\par
\end{atlasblock}\color{black}
\color{black} This example explains the possible obstruction: the autocorrelation data encoded by \(PP^*\) need not determine a binary polynomial up to reversal.\par
\needspace{8\baselineskip}\section{Where the new estimate enters}
\color{black} Kogler's Theorem 1.3 bounds nontrivial binary twins even when the identity of reciprocal products is weakened to a congruence modulo \(4\). Reduction modulo \(2\) organizes potential factors; the extra congruence imposes further constraints on opposite coefficient pairs. Cyclotomic factors and reciprocal noncyclotomic factors require additional arguments.\par
\color{black} A useful way to read the proof is to keep three steps separate: factorization creates a candidate twin; congruence estimates make nontrivial twins unlikely; estimates for the exceptional factor types finish the argument. The first step is reusable in deterministic settings, including large-gap problems, without importing the later probabilistic estimates.\par
\begin{sources}
\textbf{Sources for this section}\par\smallskip
\noindent\hyperlink{ref:Kogler2026}{[Kogler2026]} / Theorem 1.3 and §3, v2. Probabilistic collision estimate and exceptional factors. \href{https://arxiv.org/abs/2609.26771v2}{Source}.\par
\noindent\hyperlink{ref:sawinshustermanstoll2020}{[sawinshustermanstoll2020]}. Related deterministic use of Ljunggren-type ideas. \href{https://arxiv.org/abs/1803.10811}{Source}.\par
\end{sources}\color{black}
\needspace{11\baselineskip}\section*{Questions and research boundaries}
\needspace{7\baselineskip}\subsection*{Can one combine binary irreducibility with an unconditional full-group argument?}
{\small\sffamily\color{atlasgold}Research direction / Monic binary polynomials with constant term 1 over Q.}\par\smallskip
\color{black} The quoted preprint addresses reducibility. A full-\(S_n\) result needs its own group-theoretic and discriminant inputs; no such conclusion is inferred here from irreducibility alone.\par
\begin{sources}
\textbf{Sources for this section}\par\smallskip
\noindent\hyperlink{ref:Kogler2026}{[Kogler2026]} / Theorem 1.1. Scope of the cited result. \href{https://arxiv.org/abs/2609.26771v2}{Source}.\par
\end{sources}\color{black}
\needspace{6\baselineskip}\par\medskip\begingroup\footnotesize\setlength{\parskip}{1pt}\textbf{\sffamily\color{atlasgreen}Connections}\par
\color{black}\textbf{Applies to.} \hyperlink{article:sparse-gaps}{Sparse polynomials, factor reversal, and large gaps}. Reversal and norm identities also appear in deterministic large-gap arguments.\par
\color{black}\textbf{Uses.} \hyperlink{article:reciprocity}{Reciprocal symmetry and the group on root pairs}. Reversal defines the reciprocal product and its exceptional factors.\par
\color{black}\textbf{Contrasts with.} \hyperlink{article:littlewood}{Littlewood polynomials: signs, parity, and special degrees}. Sign coefficients do not have the same random reduction modulo two.\par
\color{black}\textbf{Supplies examples for.} \hyperlink{article:random-models}{Random polynomials: models and proof strategies}. A model-specific mechanism can go beyond a general Fourier hypothesis.\par
\endgroup
\chapter{Characteristic polynomials of random matrices}\label{article:random-matrices}\hypertarget{article:random-matrices}{}\color{black}
{\sffamily\small\color{atlasgold}research / Research / Q / Research overview}\par\medskip
Matrix entries can be independent while characteristic-polynomial coefficients are strongly dependent. Local rank estimates replace the independent-coefficient tools.\par\medskip
\textbf{Questions answered}\begin{itemize}
\item Why do random-matrix characteristic polynomials need different tools?
\item Which entry laws admit an unconditional theorem?
\item What does irreducibility imply for rational invariant subspaces?
\end{itemize}
{\small\textbf{Prerequisites:} \hyperlink{article:foundations}{Fields, factors, and the action on roots}; \hyperlink{article:frobenius}{Frobenius: from factor degrees to permutations}}\par
\needspace{8\baselineskip}\section{The ensemble is part of the theorem}
\color{black} Let \(M_n=(m_{ij})\) have independent entries drawn from a fixed finitely supported nonconstant law \(\mu\) on \(\mathbb Z\), and set \(\chi_n(t)=\det(tI-M_n)\). The growing-dimension problem asks whether \(\chi_n\) is irreducible, and whether its Galois group is \(S_n\), with probability tending to one.\par

\begin{atlasblock}[standard]
\textbf{\sffamily Example  /  coefficients are not independent}\par\smallskip
\color{black} For \(M=\begin{pmatrix}a&b\\c&d\end{pmatrix}\),\par

\atlasdisplay{\chi_M(t)=t^2-(a+d)t+(ad-bc).}
\color{black} If \(a,b,c,d\) are signs and \(a+d=2\), the determinant is \(0\) or \(2\). If \(a+d=0\), it is \(-2\) or \(0\). The trace and determinant therefore cannot be sampled independently from their marginal laws.\par
\end{atlasblock}\color{black}
\color{black} A theorem for independent polynomial coefficients cannot simply be applied to \(\chi_n\). Symmetric, tridiagonal, and fixed-size height-bounded matrices are further distinct ensembles.\par
\needspace{8\baselineskip}\section{Two complementary results of Eberhard}

\begin{atlasblock}[standard]
\textbf{\sffamily Published theorem  /  Eberhard}\par\smallskip
\color{black} For the independent-entry model above:\par

\begin{itemize}
\item\leavevmode\color{black} If \(\mu\) is uniform modulo a product of four distinct primes, then irreducibility and \(\operatorname{Gal}(\chi_n/\mathbb Q)\supseteq A_n\) hold with probability tending to one, unconditionally.
\item\leavevmode\color{black} For arbitrary fixed nonconstant finite-support \(\mu\), assume the Riemann hypothesis for Dedekind zeta functions of all number fields. Then \(\Pr(\chi_n\text{ reducible})=O(e^{-cn})\) and \(\Pr(\operatorname{Gal}(\chi_n/\mathbb Q)\not\supseteq A_n)=O(e^{-c\sqrt n})\), with constants depending on \(\mu\).
\end{itemize}
\end{atlasblock}\color{black}
\color{black} The uniform-modulus hypothesis holds, for example, for the uniform law on \(\{0,1,\ldots,209\}\), since \(210=2\cdot3\cdot5\cdot7\). It fails for fair signs. Neither conclusion alone distinguishes \(A_n\) from \(S_n\).\par
\begin{sources}
\textbf{Sources for this section}\par\smallskip
\noindent\hyperlink{ref:Eberhard2022}{[Eberhard2022]} / Theorems 1.2 and 1.3. Unconditional four-prime theorem and conditional general theorem. \href{https://arxiv.org/abs/2008.01223}{Source}.\par
\end{sources}\color{black}
\needspace{8\baselineskip}\section{The method: eigenvalues over finite fields}
\color{black} The equation \(\chi_M(a)=0\) is equivalent to singularity of \(aI-M\). Thus local root statistics become questions about the rank of a random matrix over \(\mathbb F_p\). To control several roots at once, one needs correlations between singularity events at different values of \(a\), not only a one-point singularity estimate.\par
\color{black} The Galois interpretation is useful independently of probability. For a squarefree polynomial with root set \(\Omega\), the average number of fixed roots of elements of \(G\) equals the number of \(G\)-orbits on \(\Omega\), by the orbit-counting lemma. Frobenius identifies those fixed roots with roots modulo good primes. Effective prime-ideal estimates allow this bridge to be used quantitatively; higher root correlations correspond to actions on tuples.\par
\begin{sources}
\textbf{Sources for this section}\par\smallskip
\noindent\hyperlink{ref:Eberhard2022}{[Eberhard2022]} / §1.1 and §3. Local root statistics and the effective global-to-local step. \href{https://arxiv.org/abs/2008.01223}{Source}.\par
\noindent\hyperlink{ref:StevenhagenLenstra1996}{[StevenhagenLenstra1996]}. Frobenius and density as the general framework. \href{https://pub.math.leidenuniv.nl/~lenstrahw/PUBLICATIONS/1996d/art.pdf}{Source}.\par
\end{sources}\color{black}
\needspace{8\baselineskip}\section{Related theorems keep their own assumptions}
\begin{center}\small\begin{tabularx}{\linewidth}{@{}>{\raggedright\arraybackslash}X>{\raggedright\arraybackslash}X>{\raggedright\arraybackslash}X@{}}\toprule
\textbf{Ensemble} & \textbf{Verified conclusion and qualification} & \textbf{Source} \\
\midrule
Symmetric fair-sign matrices & High-probability irreducibility under ERH; the proof develops sharp finite-field rank estimates. & Ferber-Jain-Sah-Sawhney \\
Specified random tridiagonal matrices & Irreducibility with exponentially small failure probability and alternating-or-symmetric group, under the stated Dedekind RH assumptions. & Bary-Soroker-Garzoni-Sodin \\
Fixed-size matrices with entries bounded by height & Exceptional-group counting as height grows; this does not resolve the fixed-law growing-dimension problem. & Anderson-O'Dorney \\
\bottomrule\end{tabularx}\end{center}
\color{black} The tridiagonal paper also treats a related ensemble with a smaller Galois group. Matrix structure can impose genuine restrictions, so a ``generic full group'' expectation must be checked against the ensemble.\par
\begin{sources}
\textbf{Sources for this section}\par\smallskip
\noindent\hyperlink{ref:FerberJainSahSawhney2023}{[FerberJainSahSawhney2023]}. Conditional symmetric-matrix irreducibility and rank estimates. \href{https://arxiv.org/abs/2106.04049}{Source}.\par
\noindent\hyperlink{ref:BarySorokerGarzoniSodin2026}{[BarySorokerGarzoniSodin2026]}. Conditional tridiagonal results and contrasting ensemble. \href{https://arxiv.org/abs/2502.17218}{Source}.\par
\noindent\hyperlink{ref:AndersonODorney2026Matrices}{[AndersonODorney2026Matrices]}. Fixed-size height-counting model. \href{https://www.cambridge.org/core/journals/mathematical-proceedings-of-the-cambridge-philosophical-society/article/galois-groups-of-random-integer-matrices/B87E2B0D71138A94EF7BB1179891456A}{Source}.\par
\end{sources}\color{black}
\needspace{8\baselineskip}\section{Why irreducibility matters for linear algebra}

\begin{atlasblock}[standard]
\textbf{\sffamily Consequence  /  no proper rational invariant subspace}\par\smallskip
\color{black} If \(W\subset\mathbb Q^n\) is \(M\)-invariant and \(0<\dim W<n\), a basis adapted to \(W\) puts \(M\) in block upper-triangular form. Its characteristic polynomial factors as the product of the characteristic polynomials on \(W\) and \(\mathbb Q^n/W\). Therefore irreducibility of \(\chi_M\) excludes every such subspace.\par
\end{atlasblock}\color{black}
\color{black} Simple eigenvalues alone do not imply this property: \(\operatorname{diag}(1,2)\) has distinct eigenvalues and two rational invariant lines. Irreducibility encodes arithmetic information beyond nonsingularity or spectral simplicity.\par
\needspace{11\baselineskip}\section*{Questions and research boundaries}
\needspace{7\baselineskip}\subsection*{Can the analytic hypothesis be removed for every fixed entry law?}
{\small\sffamily\color{atlasgold}Open in cited source / Independent entries from any fixed nonconstant finitely supported law on Z; n\ensuremath{\to}\ensuremath{\infty}.}\par\smallskip
\color{black} Eberhard formulates the general irreducibility conjecture and proves the ERH theorem, together with an unconditional theorem for special laws. These stated results do not by themselves settle fair signs unconditionally or yield full \(S_n\). The question's status here is attached to that source and its stated model.\par
\begin{sources}
\textbf{Sources for this section}\par\smallskip
\noindent\hyperlink{ref:Eberhard2022}{[Eberhard2022]} / Conjecture 1.1; Theorems 1.2-1.3. Exact conjectural and proved scopes. \href{https://arxiv.org/abs/2008.01223}{Source}.\par
\end{sources}\color{black}
\needspace{6\baselineskip}\par\medskip\begingroup\footnotesize\setlength{\parskip}{1pt}\textbf{\sffamily\color{atlasgreen}Connections}\par
\color{black}\textbf{Contrasts with.} \hyperlink{article:random-models}{Random polynomials: models and proof strategies}. The characteristic coefficients have dependencies inherited from the determinant.\par
\color{black}\textbf{Uses.} \hyperlink{article:frobenius}{Frobenius: from factor degrees to permutations}. Root counts and factor degrees encode Galois action.\par
\color{black}\textbf{Contrasts with.} \hyperlink{article:height-counting}{Counting exceptional Galois groups by height}. Growing dimension and growing entry height are different limits.\par
\color{black}\textbf{Uses.} \hyperlink{article:discriminants}{Discriminants, parity, and ramification}. Separating A\_\allowbreak{}n from S\_\allowbreak{}n requires an additional parity argument.\par
\endgroup
\chapter{Random reciprocal polynomials of growing degree}\label{article:random-reciprocal}\hypertarget{article:random-reciprocal}{}\color{black}
{\sffamily\small\color{atlasgold}research / Research / Q / Preprint boundary}\par\medskip
Independent free coefficients generate paired roots and a signed-permutation group. The trace polynomial connects this structured model to ordinary factorization problems.\par\medskip
\textbf{Questions answered}\begin{itemize}
\item What is the correct largest possible Galois group?
\item Why do trace-polynomial coefficients become dependent?
\item Why can nonsquare discriminant still leave an index-two group?
\end{itemize}
{\small\textbf{Prerequisites:} \hyperlink{article:foundations}{Fields, factors, and the action on roots}; \hyperlink{article:frobenius}{Frobenius: from factor degrees to permutations}}\par
\needspace{8\baselineskip}\section{Independent free coefficients, dependent mirror pairs}
\color{black} For a positive integer \(m\), set\par

\atlasdisplay{F_m(x)=x^{2m}+1+a_0x^m+\sum_{j=1}^{m-1}a_j(x^{m-j}+x^{m+j}),}
\color{black} with \(a_0,\ldots,a_{m-1}\) independent. The two coefficients in each mirror pair are the same random draw. Consequently this is not the ordinary model with \(2m\) independent coefficients.\par
\color{black} For separable \(F_m\), its roots occur in inverse pairs. A Galois automorphism permutes the \(m\) pairs and may flip either member within each pair. The natural ambient group is \(W_m=C_2^m\rtimes S_m\), of order \(2^m m!\), acting on \(2m\) roots. It is usually a proper subgroup of \(S_{2m}\).\par
\begin{sources}
\textbf{Sources for this section}\par\smallskip
\noindent\hyperlink{ref:HokkenKoukoulopoulos2025}{[HokkenKoukoulopoulos2025]} / Equation (1.2) and §14, v2. Model and signed-permutation action. \href{https://arxiv.org/abs/2510.18857v2}{Source}.\par
\end{sources}\color{black}
\needspace{8\baselineskip}\section{The current preprint theorem in a concrete case}

\begin{atlasblock}[standard]
\textbf{\sffamily Preprint theorem  /  Hokken-Koukoulopoulos, v2}\par\smallskip
\color{black} If the free coefficients are uniform on a fixed interval of at least \(35\) consecutive integers, then \(F_m\) is irreducible with probability tending to one. With probability tending to one, its Galois group is one of\par

\atlasdisplay{W_m,\qquad C_2\wr A_m,\qquad G_2,}
\color{black} where, writing a signed permutation as \((\varepsilon,\sigma)\),\par

\atlasdisplay{G_2=\{(\varepsilon,\sigma):\operatorname{sgn}(\sigma)\prod_i\varepsilon_i=1\}.}
\end{atlasblock}\color{black}
\color{black} For \(m\ge2\), the two alternatives after \(W_m\) have index two. The theorem therefore must not be restated as full hyperoctahedral group with high probability. Its more general formulation permits varying coefficient laws under explicit Fourier hypotheses.\par
\begin{sources}
\textbf{Sources for this section}\par\smallskip
\noindent\hyperlink{ref:HokkenKoukoulopoulos2025}{[HokkenKoukoulopoulos2025]} / Theorem A; Theorem C; equation (14.3), v2. Concrete uniform-interval case, general hypothesis and exact alternatives. \href{https://arxiv.org/abs/2510.18857v2}{Source}.\par
\end{sources}\color{black}
\needspace{8\baselineskip}\section{Worked trace-polynomial calculation}
\color{black} Every monic reciprocal polynomial of degree \(2m\) can be written uniquely as \(F(x)=x^mQ(x+x^{-1})\), with \(Q\) monic of degree \(m\). For example,\par

\begin{atlasblock}[standard]
\textbf{\sffamily Example  /  the coefficient change creates dependence}\par\smallskip

\atlasdisplay{\begin{aligned}F(x)&=x^6+ax^5+bx^4+cx^3+bx^2+ax+1,\\Q(y)&=y^3+ay^2+(b-3)y+(c-2a).\end{aligned}}
\color{black} The identities \(x^2+x^{-2}=y^2-2\) and \(x^3+x^{-3}=y^3-3y\) verify the formula. Even if \(a,b,c\) are independent, the coefficients \(a\) and \(c-2a\) of \(Q\) are generally dependent. A theorem about independent coefficients of \(Q\) cannot be imported without checking this new joint law.\par
\end{atlasblock}\color{black}
\color{black} A root \(y\) of \(Q\) yields the pair of solutions of \(x^2-yx+1=0\). This separates the action on pairs from the quadratic extensions that split them.\par
\begin{sources}
\textbf{Sources for this section}\par\smallskip
\noindent\hyperlink{ref:HokkenKoukoulopoulos2025}{[HokkenKoukoulopoulos2025]} / §5. Trace polynomials and the dependence issue. \href{https://arxiv.org/abs/2510.18857v2}{Source}.\par
\end{sources}\color{black}
\needspace{8\baselineskip}\section{What must be controlled}
\color{black} Factorization of \(Q\) detects how the inverse pairs split into Galois orbits, but irreducibility of \(Q\) alone does not prove irreducibility of \(F\): the latter can split into a factor and its reversal. The proof therefore needs a separate estimate for this exceptional pattern. Fourier estimates are taken in the original reciprocal coefficient space, where independence is available.\par
\color{black} For the group, one studies the permutation of pairs and the allowed pair flips. A nonsquare discriminant of \(F\) rules out the subgroup \(W_m\cap A_{2m}\), but does not rule out every index-two subgroup of \(W_m\). This explains why a successful parity estimate can still leave the three alternatives in the theorem.\par
\begin{sources}
\textbf{Sources for this section}\par\smallskip
\noindent\hyperlink{ref:HokkenKoukoulopoulos2025}{[HokkenKoukoulopoulos2025]} / §2, §7 and §§14-16. Exceptional reciprocal factors and the group analysis. \href{https://arxiv.org/abs/2510.18857v2}{Source}.\par
\end{sources}\color{black}
\needspace{8\baselineskip}\section{Compare with fixed-degree reciprocity}
\color{black} The fixed-degree height theorem gives an exceptional proportion of order \(H^{-1}\log H\), with \(m\) fixed and \(H\to\infty\). The theorem here fixes the free-coefficient law and lets \(m\to\infty\). The same paired-root algebra appears in both settings, but the distributions and counting methods differ.\par
\begin{sources}
\textbf{Sources for this section}\par\smallskip
\noindent\hyperlink{ref:AndersonBertelliODorney2026}{[AndersonBertelliODorney2026]} / Theorem 1.1. Fixed-degree comparison. \href{https://arxiv.org/abs/2406.18970v2}{Source}.\par
\end{sources}\color{black}
\needspace{11\baselineskip}\section*{Questions and research boundaries}
\needspace{7\baselineskip}\subsection*{When can the two remaining index-two alternatives be excluded?}
{\small\sffamily\color{atlasgold}Research direction / Growing-degree reciprocal model under the hypotheses of Hokken-Koukoulopoulos.}\par\smallskip
\color{black} The quoted theorem leaves \(C_2\wr A_m\) and \(G_2\) as possible outcomes. Determining their probabilities, or proving full \(W_m\) for particular laws, is a separate question from irreducibility and from the nonsquare discriminant of \(F_m\).\par
\begin{sources}
\textbf{Sources for this section}\par\smallskip
\noindent\hyperlink{ref:HokkenKoukoulopoulos2025}{[HokkenKoukoulopoulos2025]} / Theorem A and §14, v2. Exact surviving group alternatives. \href{https://arxiv.org/abs/2510.18857v2}{Source}.\par
\end{sources}\color{black}
\needspace{6\baselineskip}\par\medskip\begingroup\footnotesize\setlength{\parskip}{1pt}\textbf{\sffamily\color{atlasgreen}Connections}\par
\color{black}\textbf{Uses.} \hyperlink{article:reciprocity}{Reciprocal symmetry and the group on root pairs}. The trace polynomial and inverse-pair structure are deterministic foundations.\par
\color{black}\textbf{Generalizes.} \hyperlink{article:random-models}{Random polynomials: models and proof strategies}. The Fourier approach is adapted to a structured coefficient space.\par
\color{black}\textbf{Contrasts with.} \hyperlink{article:height-counting}{Counting exceptional Galois groups by height}. Here the number of free coefficients grows while their law remains fixed.\par
\color{black}\textbf{Generalizes.} \hyperlink{article:twisted-reciprocity}{Twisted reciprocity and square discriminants}. The involution x\ensuremath{\mapsto}1/x can be replaced by x\ensuremath{\mapsto}b/x.\par
\endgroup
\chapter{Twisted reciprocity and square discriminants}\label{article:twisted-reciprocity}\hypertarget{article:twisted-reciprocity}{}\color{black}
{\sffamily\small\color{atlasgold}research / Research / Q / Preprint boundary}\par\medskip
The root pairing x\ensuremath{\mapsto}b/x gives a useful discriminant identity. It connects paired-root Galois groups with square values and fixed-degree counting.\par\medskip
\textbf{Questions answered}\begin{itemize}
\item How does the trace construction change under x\ensuremath{\mapsto}b/x?
\item What is the discriminant identity?
\item How can a group restriction become a square-value problem?
\end{itemize}
{\small\textbf{Prerequisites:} \hyperlink{article:foundations}{Fields, factors, and the action on roots}; \hyperlink{article:frobenius}{Frobenius: from factor degrees to permutations}}\par
\needspace{8\baselineskip}\section{The twisted involution}
\color{black} Fix an integer \(b\ne0\). A polynomial \(F\) of degree \(2m\) is \(b\)-reciprocal if\par

\atlasdisplay{F(x)=\frac{x^{2m}}{b^m}F(b/x).}
\color{black} Its nonzero roots are paired by \(\alpha\mapsto b/\alpha\). If \(F\) is separable, neither fixed point \(\alpha^2=b\) can occur: in the trace expression below, such a point would be a double root. Hence the group acts on genuine pairs and embeds in \(W_m=C_2\wr S_m\).\par
\color{black} The invariant coordinate is \(y=x+b/x\), giving \(F(x)=x^mQ(x+b/x)\) for a degree-\(m\) polynomial \(Q\). Each trace root \(y_i\) corresponds to \(x^2-y_ix+b\).\par
\begin{sources}
\textbf{Sources for this section}\par\smallskip
\noindent\hyperlink{ref:ODorney2026Twisted}{[ODorney2026Twisted]}. Twisted-reciprocal family and hyperoctahedral ambient group. \href{https://arxiv.org/abs/2609.31613v1}{Source}.\par
\end{sources}\color{black}
\needspace{8\baselineskip}\section{The exact identity}

\begin{atlasblock}[standard]
\textbf{\sffamily Algebraic identity  /  derived from the roots}\par\smallskip
\color{black} With \(F(x)=x^mQ(x+b/x)\),\par

\atlasdisplay{\operatorname{disc}(F)=b^{m(m-1)}Q(2\sqrt b)Q(-2\sqrt b)\operatorname{disc}(Q)^2.}
\end{atlasblock}\color{black}

\begin{atlasblock}[standard]
\textbf{\sffamily Derivation}\par\smallskip
\color{black} Write \(Q(y)=a\prod_{i=1}^{m}(y-y_i)\). The discriminant of \(x^2-y_ix+b\) is \(y_i^2-4b\). For \(i\ne j\), the resultant of the two quadratic factors is \(b(y_i-y_j)^2\): evaluating their difference at the two roots of the first quadratic gives the product \(b(y_i-y_j)^2\). The product formula for discriminants therefore gives\par

\atlasdisplay{a^{4m-2}\prod_i(y_i^2-4b)\,b^{m(m-1)}\prod_{i<j}(y_i-y_j)^4,}
\color{black} which is the displayed identity.\par
\end{atlasblock}\color{black}
\color{black} Since \(m(m-1)\) is even, the square class of the discriminant is governed by \(Q(2\sqrt b)Q(-2\sqrt b)\). The product is rational even when \(\sqrt b\notin\mathbb Q\).\par
\begin{sources}
\textbf{Sources for this section}\par\smallskip
\noindent\hyperlink{ref:Cox2012}{[Cox2012]}. Standard discriminant and resultant product identities. \href{https://onlinelibrary.wiley.com/doi/book/10.1002/9781118218457}{Source}.\par
\end{sources}\color{black}
\needspace{8\baselineskip}\section{Worked example with b=2}

\begin{atlasblock}[standard]
\textbf{\sffamily Example  /  a quartic and its trace quadratic}\par\smallskip
\color{black} Take \(Q(y)=y^2-y-1\). Then\par

\atlasdisplay{F(x)=x^2Q(x+2/x)=x^4-x^3+3x^2-2x+4.}
\color{black} Here \(\operatorname{disc}(Q)=5\), and\par

\atlasdisplay{Q(2\sqrt2)Q(-2\sqrt2)=(7-2\sqrt2)(7+2\sqrt2)=41.}
\color{black} Thus \(\operatorname{disc}(F)=2^2\cdot41\cdot5^2=4100\), a nonsquare. This excludes containment in \(A_4\); on its own it does not identify the full Galois group inside \(W_2\).\par
\end{atlasblock}\color{black}
\needspace{8\baselineskip}\section{A recent fixed-degree counting statement}
\color{black} O'Dorney's September 2026 preprint orders the family by an adapted Euclidean height. Its stated main result, for fixed \(b\ne0\) and \(m\ge4\), counts polynomials with group smaller than \(W_m\) as an explicit constant times \(H^m\log H\), with error \(O_{m,b}(H^m)\). The main contribution comes from an index-two subgroup. This is a preprint report; the exact counting convention and constant should be taken from the linked version before applying it.\par
\color{black} The independently derived discriminant formula above explains the mechanism: a group restriction becomes a square-value problem in a lower-dimensional parameter space. The original untwisted theorem \(b=1\) is available in the published paper of Anderson-Bertelli-O'Dorney.\par
\begin{sources}
\textbf{Sources for this section}\par\smallskip
\noindent\hyperlink{ref:ODorney2026Twisted}{[ODorney2026Twisted]} / Abstract, v1. Verified primary indexed abstract; full-text theorem details are not asserted here. \href{https://arxiv.org/abs/2609.31613v1}{Source}.\par
\noindent\hyperlink{ref:AndersonBertelliODorney2026}{[AndersonBertelliODorney2026]}. Published untwisted counting theorem. \href{https://arxiv.org/abs/2406.18970v2}{Source}.\par
\end{sources}\color{black}
\needspace{6\baselineskip}\par\medskip\begingroup\footnotesize\setlength{\parskip}{1pt}\textbf{\sffamily\color{atlasgreen}Connections}\par
\color{black}\textbf{Generalizes.} \hyperlink{article:reciprocity}{Reciprocal symmetry and the group on root pairs}. The same trace construction works with x+b/x.\par
\color{black}\textbf{Applies to.} \hyperlink{article:discriminants}{Discriminants, parity, and ramification}. The discriminant square class is expressed by two evaluations.\par
\color{black}\textbf{Applies to.} \hyperlink{article:height-counting}{Counting exceptional Galois groups by height}. Square values control exceptional-group counts.\par
\color{black}\textbf{Contrasts with.} \hyperlink{article:random-reciprocal}{Random reciprocal polynomials of growing degree}. The counting parameter here is height at a fixed degree.\par
\endgroup
\part{Specialization and inverse Galois}
\chapter{Specialization and Hilbert irreducibility}\label{article:specialization}\hypertarget{article:specialization}{}\color{black}
{\sffamily\small\color{atlasgold}method / Intermediate / number fields / Classical theory}\par\medskip
Pass from a polynomial over a rational function field to arithmetic fibers, and understand how the exceptional parameters are organized.\par\medskip
\textbf{Questions answered}\begin{itemize}
\item When does a specialization retain the generic Galois group?
\item Why can the exceptional set be infinite, and when can it be finite?
\item What changes when parameters are restricted to units?
\item Why can arithmetic and geometric Galois groups differ?
\end{itemize}
{\small\textbf{Prerequisites:} \hyperlink{article:foundations}{Fields, factors, and the action on roots}; \hyperlink{article:frobenius}{Frobenius: from factor degrees to permutations}}\par
\needspace{8\baselineskip}\section{The generic polynomial and its fibers}
\color{black} Fix a number field \(K\) and a monic, separable, irreducible polynomial \(F(T,x)\in K(T)[x]\). Its \emph{arithmetic} group is \(G=\operatorname{Gal}(F/K(T))\), acting on the roots. Remove the finitely many parameter values where a coefficient has a pole or the discriminant vanishes. At every remaining \(t\in K\), the specialized group \(G_t\) is conjugate to a subgroup of \(G\).\par
\color{black} This inclusion comes from a decomposition group above the place \(T=t\) in the generic splitting field. Because the roots stay distinct, its unramified residue action identifies with the action on the specialized roots. The generic group is an upper bound; a particular fiber can have a smaller group.\par
\begin{sources}
\textbf{Sources for this section}\par\smallskip
\noindent\hyperlink{ref:Serre2007Topics}{[Serre2007Topics]} / §3.3, polynomial interpretation. Specialization of the root permutation action. \href{https://www.routledge.com/Topics-in-Galois-Theory-2nd-Edition/Serre/p/book/9780429064401}{Source}.\par
\end{sources}\color{black}
\needspace{8\baselineskip}\section{Arithmetic and geometric groups: the role of constants}
\color{black} Let \(K\) have characteristic zero, let \(T\) be transcendental, and let \(L/K(T)\) be the splitting field of a separable polynomial. Inside a common algebraic closure, put \(C=L\cap\overline K\), the field of algebraic constants. Define\par

\atlasdisplay{G_{\mathrm{arith}}=\operatorname{Gal}(L/K(T)),\qquad G_{\mathrm{geom}}=\operatorname{Gal}(L\overline K/\overline K(T)).}

\begin{atlasblock}[standard]
\textbf{\sffamily Structure  /  the constants exact sequence}\par\smallskip
\color{black} The extension \(C/K\) is finite Galois, and restriction to constants gives\par

\atlasdisplay{1\longrightarrow G_{\mathrm{geom}}\longrightarrow G_{\mathrm{arith}}\longrightarrow\operatorname{Gal}(C/K)\longrightarrow1.}
\end{atlasblock}\color{black}
\color{black} Indeed, normality makes \(C\) stable under the arithmetic group, and automorphisms of \(C(T)/K(T)\) extend to \(L\). The kernel is \(\operatorname{Gal}(L/C(T))\). Since \(C\) is relatively algebraically closed in \(L\), the extension is regular over \(C\); extending constants to \(\overline K\) preserves this kernel. Thus the two groups agree exactly when \(C=K\), the regular case. The exact sequence need not split in general.\par

\begin{atlasblock}[standard]
\textbf{\sffamily Worked example  /  \(X^3-T\) over \(\mathbb Q(T)\)}\par\smallskip
\color{black} Choose \(\alpha^3=T\) and a primitive cube root \(\zeta\). The roots are \(\alpha,\zeta\alpha,\zeta^2\alpha\), so \(L=\mathbb Q(\zeta)(\alpha)\). Because \(\alpha\) is transcendental, this rational function field has algebraic constants exactly \(C=\mathbb Q(\zeta)=\mathbb Q(\sqrt{-3})\). Eisenstein at \(T\) over \(C[T]\) gives\par

\atlasdisplay{[L:\mathbb Q(T)]=[L:C(T)]\,[C(T):\mathbb Q(T)]=3\cdot2=6.}
\color{black} The automorphism \(r\) sends \(\alpha\mapsto\zeta\alpha\) and fixes \(\zeta\); the automorphism \(s\) fixes \(\alpha\) and sends \(\zeta\mapsto\zeta^{-1}\). They satisfy \(r^3=s^2=1\) and \(srs=r^{-1}\). Hence \(G_{\mathrm{arith}}\cong S_3\), while \(G_{\mathrm{geom}}=\langle r\rangle\cong C_3\). The quotient \(C_2\) records the action on the quadratic constants. The discriminant \(-27T^2=-3(3T)^2\) displays the same constant obstruction.\par
\end{atlasblock}\color{black}
\color{black} By comparison, \(X^3-X-T\) is irreducible over \(\overline{\mathbb Q}(T)\): apply Gauss's lemma to its irreducibility as a polynomial linear in \(T\) with unit leading coefficient in \(\overline{\mathbb Q}[X,T]\). Its discriminant \(4-27T^2\) has two simple zeros, so is not a square in that rational function field. The cubic criterion gives \(G_{\mathrm{geom}}=S_3\); consequently \(G_{\mathrm{arith}}=S_3\) and \(C=\mathbb Q\). This family is regular, whereas \(X^3-T\) is not.\par
\begin{sources}
\textbf{Sources for this section}\par\smallskip
\noindent\hyperlink{ref:Milne2022}{[Milne2022]} / Free notes v5.10, Propositions 7.7 and 7.15. Restriction and extension of Galois automorphisms and change of base field; the constants exact sequence is derived here from these facts. \href{https://www.jmilne.org/math/Books/index.html}{Author's book page}.\par
\noindent\hyperlink{ref:JensenLedetYui2002}{[JensenLedetYui2002]} / §3.3, Proposition 3.3.6 (pp. 74-75). Regular Galois extensions preserve their group after extension of constants. \href{https://library.slmath.org/books/Book45/index.html}{Source}.\par
\noindent\hyperlink{ref:ConradCubicQuartic}{[ConradCubicQuartic]} / Equation (1.1), Theorems 2.1 and 2.8. Cubic discriminant formula, cubic group criterion, and splitting field from one root and the discriminant; the two parameter examples are worked out explicitly here. \href{https://kconrad.math.uconn.edu/blurbs/galoistheory/cubicquartic.pdf}{Source}.\par
\end{sources}\color{black}
\needspace{8\baselineskip}\section{The theorem and its quantifiers}

\begin{atlasblock}[standard]
\textbf{\sffamily Theorem  /  Hilbert irreducibility}\par\smallskip
\color{black} For \(F\) as above, there is a thin set \(B\subset K\) such that every \(t\in K\setminus B\) is admissible, \(F(t,x)\) is irreducible, and \(G_t=G\) up to conjugacy. In particular, infinitely many parameters retain the generic group.\par
\end{atlasblock}\color{black}
\color{black} Here ``thin'' means contained in a finite union of finite sets and images of rational points under dominant, generically finite maps from geometrically integral curves to the parameter line of degree at least two. It does not mean finite. To detect a drop of group, take fixed fields of proper subgroups of \(G\): a smaller specialization group forces a rational point on an associated resolvent cover. Hilbert's theorem allows parameters outside all these images.\par
\begin{sources}
\textbf{Sources for this section}\par\smallskip
\noindent\hyperlink{ref:Hilbert1892}{[Hilbert1892]}. Original theorem; modern formulations use Hilbertian fields. \href{https://www.degruyterbrill.com/document/doi/10.1515/crll.1892.110.104/html}{Source}.\par
\noindent\hyperlink{ref:Serre2007Topics}{[Serre2007Topics]} / Proposition 3.3.5 and Theorem 3.4.1. Full-group specialization outside a thin set over number fields. \href{https://www.routledge.com/Topics-in-Galois-Theory-2nd-Edition/Serre/p/book/9780429064401}{Source}.\par
\end{sources}\color{black}
\needspace{8\baselineskip}\section{Example: one cubic family, infinitely many exceptions}

\begin{atlasblock}[standard]
\textbf{\sffamily Worked example}\par\smallskip
\color{black} Let \(F(T,x)=x^3-x-T\). The quotient \(\mathbb Q[T,x]/(F)\) is \(\mathbb Q[x]\), with \(T=x^3-x\). Hence \(F\) is irreducible in the polynomial ring and, by Gauss's lemma, in \(\mathbb Q(T)[x]\). Its discriminant is\par

\atlasdisplay{\Delta(T)=4-27T^2.}
\color{black} The two zeros of this rational function are simple, so it is not a square in \(\mathbb Q(T)\). A transitive cubic group with nonsquare discriminant is \(S_3\).\par
\color{black} At \(t=1\), the reduction of \(x^3-x-1\) modulo \(2\) has no root, and its discriminant is \(-23\), giving \(S_3\) again. But every \(t=u^3-u\), with \(u\in\mathbb Q\), gives the rational root \(u\). The parameters \(t=u^3-u\) form precisely the reducible-fiber locus, an infinite subset of the full-group exceptional set. The latter is larger: at \(t=1/7\), the cubic remains irreducible but its discriminant is \(169/49\), giving \(A_3\).\par
\end{atlasblock}\color{black}
\begin{sources}
\textbf{Sources for this section}\par\smallskip
\noindent\hyperlink{ref:ConradCubicQuartic}{[ConradCubicQuartic]}. The transitive cubic group and discriminant criterion; calculations here are explicit. \href{https://kconrad.math.uconn.edu/blurbs/galoistheory/cubicquartic.pdf}{Source}.\par
\end{sources}\color{black}
\needspace{8\baselineskip}\section{When geometry makes the exceptional set finite}
\color{black} Hilbert irreducibility can be strengthened in particular families. Proper-subgroup resolvent covers turn exceptional parameters into rational points on curves. If the relevant curves have genus at least two, finiteness of their rational points gives a finite exceptional set. Establishing the genus bounds requires ramification calculations and subgroup information.\par

\begin{atlasblock}[standard]
\textbf{\sffamily Published example  /  Hajir-Wong}\par\smallskip
\color{black} Fix a number field \(K\). For fixed \(n\ge5\), all but finitely many \(\alpha\in K\) give an irreducible generalized Laguerre polynomial \(L_n^{(\alpha)}\) whose group contains \(A_n\). For \(n\ge10\), the group is \(S_n\) outside a finite set.\par
\end{atlasblock}\color{black}
\color{black} The degree is fixed and the parameter varies. This differs from fixing \(\alpha\) and letting \(n\) grow. The argument does not generally compute the finite exceptional set.\par
\begin{sources}
\textbf{Sources for this section}\par\smallskip
\noindent\hyperlink{ref:HajirWong2006}{[HajirWong2006]} / Theorem 1(a), p.1129; §§2 and 6. The displayed thresholds are \ensuremath{\ge}5 and \ensuremath{\ge}10; the PDF page was checked visually. \href{https://www.numdam.org/item/AIF_2006__56_4_1127_0/}{Source}.\par
\end{sources}\color{black}
\needspace{8\baselineskip}\section{Power parameters need a geometric hypothesis}
\color{black} Hilbert irreducibility over all rational parameters does not by itself control a prescribed sequence such as \(2^n\). For \(X^2-t\), exactly the odd integers \(n\) make \(X^2-2^n\) irreducible over \(\mathbb Q\); their density is \(1/2\).\par

\begin{atlasblock}[standard]
\textbf{\sffamily Theorem  /  Bary-Soroker-Garzoni-Matei, conditional on GRH}\par\smallskip
\color{black} Let \(f(\mathbf t,X)\in\mathbb Z[t_1,\ldots,t_r,X]\), with \(\deg_Xf\ge1\), satisfy the pullback condition: for every tuple \(m_i\ge1\),\par

\atlasdisplay{f(t_1^{m_1},\ldots,t_r^{m_r},X)\text{ is irreducible in }\overline{\mathbb Q}[\mathbf t,X].}
\color{black} Fix \(a_i\in\mathbb Z\setminus\{0,1,-1\}\). Assuming GRH, the proportion of exponent vectors \(\mathbf n\in([-N,N]\cap\mathbb Z)^r\) for which \(f(a_1^{n_1},\ldots,a_r^{n_r},X)\) is irreducible over \(\mathbb Q\) tends to one as \(N\to\infty\).\par
\end{atlasblock}\color{black}
\color{black} The geometric condition rules out power-cover obstructions. For \(X^2-t\), substituting \(t=u^2\) already produces a factorization. The theorem concerns irreducibility; preservation of an entire Galois group requires applying the appropriate resolvent-cover argument with its hypotheses.\par
\color{black} \textbf{One parameter is stronger and unconditional.} Under the same pullback condition with \(r=1\), Dèbes proves that \(f(a^n,X)\) is irreducible for all but finitely many integers \(n\), without GRH. His Theorem 1.2 shows that infinitely many reducible fibers force a reducible scaled power pullback; geometric (PB) excludes this. The proof uses Siegel's theorem on integral points on curves. This finiteness assertion does not by itself give an effective list of the exceptional exponents.\par
\begin{sources}
\textbf{Sources for this section}\par\smallskip
\noindent\hyperlink{ref:BarySorokerGarzoniMatei2025Powers}{[BarySorokerGarzoniMatei2025Powers]} / §1, condition (PB), Theorem 1.1, and the following examples. Published theorem; GRH and absolute irreducibility of every power pullback are essential. \href{https://doi.org/10.1007/s00209-025-03717-7}{Source}.\par
\noindent\hyperlink{ref:Debes1992Powers}{[Debes1992Powers]} / Theorem 1.2, p. 144; see also Bary-Soroker-Garzoni-Matei, §1 after Theorem 1.1. The original theorem characterizes infinitely many reducible power specializations; geometric (PB) rules out its scaled power-pullback obstruction. \href{https://doi.org/10.1016/0022-314X(92)90018-K}{Source}.\par
\end{sources}\color{black}
\needspace{8\baselineskip}\section{A 2026 extension: exceptional unit parameters}
\color{black} \textbf{Preprint, 3 September 2026.} Let \(\pi:X\to\mathbb P^1\) be a nonconstant morphism from a smooth projective irreducible curve over a number field \(K\), let \(S\) be a finite set of non-archimedean places of \(K\), and put \(U=\mathcal O_S^\times\). Let \(\operatorname{Red}(\pi)\) be the set of rational parameters with a reducible fiber, together with all rational branch values. Siksek-Stoll prove\par

\atlasdisplay{U\cap\operatorname{Red}(\pi)=I\cup\bigcup_{j=1}^{k}\gamma_j U^{r_j},\qquad |I|<\infty,\quad \gamma_j\in U,\quad r_j\ge2.}
\color{black} The power cosets are effectively computable; the proof does not make the finite remainder \(I\) effective. This supplies structure for an exceptional set that ordinary Hilbert irreducibility only labels thin. For \(x^2-T\), the exceptional unit parameters are simply the squares \(U^2\): if a square in \(K\) is an \(S\)-unit, its square root is also an \(S\)-unit.\par
\begin{sources}
\textbf{Sources for this section}\par\smallskip
\noindent\hyperlink{ref:SiksekStoll2026}{[SiksekStoll2026]} / Theorem 3 and following remark. Version 1; this is a preprint theorem, not a complete effective algorithm. \href{https://arxiv.org/abs/2609.04551v1}{Source}.\par
\end{sources}\color{black}
\needspace{8\baselineskip}\section{What a specialization argument does not give}

\begin{atlasblock}[boundary]
\color{black} Irreducibility of one fiber proves neither a uniform theorem in degree nor a description of all exceptions. The arithmetic group over \(K(T)\) can exceed the geometric group over \(\overline K(T)\). A theorem for each finite iteration depth does not supply one parameter that works simultaneously at every depth: an intersection of infinitely many Hilbert sets needs a separate argument.\par
\end{atlasblock}\color{black}
\begin{sources}
\textbf{Sources for this section}\par\smallskip
\noindent\hyperlink{ref:FriedJarden2023}{[FriedJarden2023]}. Hilbertian fields and specialization framework. \href{https://link.springer.com/book/10.1007/978-3-031-28020-7}{Source}.\par
\noindent\hyperlink{ref:DittmannKadets2022}{[DittmannKadets2022]} / Introduction and Theorem 1.2. A concrete warning against interchanging finite-level and infinite-level quantifiers. \href{https://arxiv.org/abs/2012.03076}{Source}.\par
\end{sources}\color{black}
\needspace{11\baselineskip}\section*{Questions and research boundaries}
\needspace{7\baselineskip}\subsection*{Can the finite exceptional parameters be determined effectively?}
{\small\sffamily\color{atlasgold}Research direction / A specified one-parameter family over a fixed number field}\par\smallskip
\color{black} Finiteness and effective determination are different outputs. Hajir-Wong explain the difficulty for their curves; the 2026 unit-specialization theorem likewise leaves a finite set ineffective. A concrete project should name the family and the auxiliary curves whose rational or integral points must be found.\par
\begin{sources}
\textbf{Sources for this section}\par\smallskip
\noindent\hyperlink{ref:HajirWong2006}{[HajirWong2006]} / Remark 1.  \href{https://www.numdam.org/item/AIF_2006__56_4_1127_0/}{Source}.\par
\noindent\hyperlink{ref:SiksekStoll2026}{[SiksekStoll2026]} / Remark following Theorem 3.  \href{https://arxiv.org/abs/2609.04551v1}{Source}.\par
\end{sources}\color{black}
\needspace{6\baselineskip}\par\medskip\begingroup\footnotesize\setlength{\parskip}{1pt}\textbf{\sffamily\color{atlasgreen}Connections}\par
\color{black}\textbf{Applies to.} \hyperlink{article:inverse-galois}{Inverse Galois theory and generic polynomials}. A regular realization over K(T) specializes to realizations over K.\par
\color{black}\textbf{Uses.} \hyperlink{article:resolvents}{Resolvents: turn group containment into factorization}. Proper-subgroup fixed fields describe possible drops of Galois group.\par
\color{black}\textbf{Supplies examples for.} \hyperlink{article:laguerre}{Generalized Laguerre polynomials: two parameter regimes}. Laguerre specializations illustrate finite exceptional sets.\par
\color{black}\textbf{Contrasts with.} \hyperlink{article:composition-iteration}{Composition, iteration, and arboreal groups}. All finite depths separately do not settle one infinite tower.\par
\color{black}\textbf{Generalizes.} \hyperlink{article:height-counting}{Counting exceptional Galois groups by height}. Thin exceptional sets lead to quantitative counting questions.\par
\endgroup
\chapter{Critical points and monodromy: a geometric route to the symmetric group}\label{article:critical-points-monodromy}\hypertarget{article:critical-points-monodromy}{}\color{black}
{\sffamily\small\color{atlasgold}method / Advanced / Characteristic zero; specializations over Q / Classical theory}\par\medskip
In characteristic zero, simple critical points with distinct critical values determine the Galois group of a polynomial family. The branch permutations explain the generic group; specialization is a separate arithmetic step.\par\medskip
\textbf{Questions answered}\begin{itemize}
\item How can derivatives determine a generic Galois group?
\item Why do simple critical values produce transpositions?
\item Which parts of a generic result survive at a particular rational parameter?
\end{itemize}
{\small\textbf{Prerequisites:} \hyperlink{article:foundations}{Fields, factors, and the action on roots}; \hyperlink{article:permutation-groups}{The group-theoretic step: transitivity, alternating groups, and symmetric groups}; \hyperlink{article:ramification-inertia}{Ramification and inertia: cycles at bad primes}; \hyperlink{article:specialization}{Specialization and Hilbert irreducibility}}\par
\needspace{8\baselineskip}\section{Critical points, critical values, and the generic polynomial}
\color{black} Let \(K\) have characteristic zero and let \(F\in K[X]\) have degree \(n\ge2\). The family \(F(X)-T\) is a polynomial over \(K(T)\), with \(T\) an indeterminate. A \textbf{critical point} is a root \(c\) of \(F'\), and its \textbf{critical value} is \(F(c)\). We call \(F\) Morse when \(F'\) has \(n-1\) distinct roots over \(\overline K\) and their critical values are pairwise distinct.\par
\color{black} These are two separate tests. For \(F=X^4-2X^2\), the critical points \(0,1,-1\) are simple, but \(F(1)=F(-1)=-1\). A simple derivative alone does not satisfy the theorem. Composition \(F=(X^2)^2-2X^2\) also exposes an intermediate cover and a block system on the four generic roots.\par
\color{black} The geometric group is the splitting-field group over \(\overline K(T)\); the arithmetic group is over \(K(T)\). Both act on the same \(n\) roots. A theorem about these groups does not assert the same group for every specialization \(T=t\).\par
\begin{sources}
\textbf{Sources for this section}\par\smallskip
\noindent\hyperlink{ref:Serre2007Topics}{[Serre2007Topics]} / §4.4, Theorems 4.4.1 and 4.4.5, pp. 39-41. Morse hypotheses and the geometric construction; the quartic boundary example is computed here. \href{https://www.routledge.com/Topics-in-Galois-Theory-2nd-Edition/Serre/p/book/9780429064401}{Source}.\par
\end{sources}\color{black}
\needspace{8\baselineskip}\section{The Morse criterion}

\begin{atlasblock}[standard]
\textbf{\sffamily Theorem  /  Hilbert's Morse construction, characteristic-zero form}\par\smallskip
\color{black} If \(F\in K[X]\) is Morse and \(\operatorname{char}K=0\), then\par

\atlasdisplay{\begin{aligned}\operatorname{Gal}(F(X)-T/\overline K(T))&=S_n,\\\operatorname{Gal}(F(X)-T/K(T))&=S_n.\end{aligned}}
\color{black} The splitting field over \(K(T)\) is regular over \(K\): it introduces no new algebraic constants.\par
\end{atlasblock}\color{black}
\color{black} The geometric conclusion is the decisive one. It gives a subgroup \(S_n\) of the arithmetic group, which already lies in \(S_n\). Equality follows, and the arithmetic/geometric quotient, which measures the constant extension, is trivial.\par
\color{black} This is a sufficient criterion. Polynomials failing one of the Morse conditions can still have full symmetric monodromy; their ramification must be analyzed by another argument.\par
\begin{sources}
\textbf{Sources for this section}\par\smallskip
\noindent\hyperlink{ref:Serre2007Topics}{[Serre2007Topics]} / Theorems 4.4.1 and 4.4.5. The general theorem is applied after base change to the algebraic closure; arithmetic equality and regularity follow. \href{https://www.routledge.com/Topics-in-Galois-Theory-2nd-Edition/Serre/p/book/9780429064401}{Source}.\par
\end{sources}\color{black}
\needspace{8\baselineskip}\section{Why the branch permutations generate enough transpositions}

\begin{atlasblock}[standard]
\textbf{\sffamily Proof sketch  /  ramification plus a connected graph}\par\smallskip
\color{black} Over \(\overline K\), the degree-\(n\) cover \(X\mapsto F(X)\) is connected, so its monodromy group \(G\) is transitive. Equivalently, \(F(X)-T\) is irreducible over \(\overline K(T)\): in \(\overline K[X,T]\), a factor independent of \(T\) would have to divide the coefficient \(-1\) of \(T\).\par
\color{black} At a simple critical point \(c\), the local expansion begins \(F(c)+a(X-c)^2+\cdots\), with \(a\ne0\). The distinct-value hypothesis ensures that only these two sheets meet over \(F(c)\). Its inertia generator is therefore a single transposition.\par
\color{black} Let \(H\) be the normal subgroup generated by inertia at all finite branch values. The quotient of the Galois closure by \(H\) gives a connected degree-\(d\) cover of the projective line, unramified away from infinity. In characteristic zero all ramification is tame. If \(r\ge1\) points lie above infinity, Riemann-Hurwitz reads \(2g-2=-2d+(d-r)=-d-r\). Since \(g\ge0\), necessarily \(d=r=1\). Thus \(H=G\): the finite transpositions generate the whole group.\par
\color{black} Make a graph whose vertices are the roots and whose edges are these transpositions, including their conjugates. Its connected components are the orbits of the generated group. Transitivity makes the graph connected, and edge transpositions of a connected graph generate every permutation. Hence \(G=S_n\).\par
\end{atlasblock}\color{black}
\color{black} The generation step matters. Merely finding one transposition inside a transitive group is insufficient.\par
\begin{sources}
\textbf{Sources for this section}\par\smallskip
\noindent\hyperlink{ref:Serre2007Topics}{[Serre2007Topics]} / Theorem 4.4.5 and Proposition 4.4.6, p. 41. Inertia generation via Riemann-Hurwitz. The graph argument expands the transposition step. \href{https://www.routledge.com/Topics-in-Galois-Theory-2nd-Edition/Serre/p/book/9780429064401}{Source}.\par
\end{sources}\color{black}
\needspace{8\baselineskip}\section{A family in every degree}

\begin{atlasblock}[standard]
\textbf{\sffamily Example  /  \(F(X)=X^n-X\)}\par\smallskip
\color{black} For \(n\ge2\), the critical points satisfy \(nc^{n-1}=1\). They are nonzero and distinct, and\par

\atlasdisplay{F(c)=-\frac{n-1}{n}c.}
\color{black} The critical values are therefore distinct. It follows that \(X^n-X-T\) has geometric and arithmetic group \(S_n\) in characteristic zero. Over \(\mathbb Q\), Hilbert irreducibility then supplies infinitely many rational parameters \(t\) for which \(X^n-X-t\) has group \(S_n\).\par
\end{atlasblock}\color{black}
\color{black} This construction separates two methods: derivatives and inertia determine the generic group, while Hilbert irreducibility supplies full-group fibers. To certify a specified \(t\), use factorization, discriminants, or resolvents for that fiber.\par
\begin{sources}
\textbf{Sources for this section}\par\smallskip
\noindent\hyperlink{ref:Serre2007Topics}{[Serre2007Topics]} / §4.4 and §3.3. The Morse theorem and specialization framework; the derivative calculation is supplied here. \href{https://www.routledge.com/Topics-in-Galois-Theory-2nd-Edition/Serre/p/book/9780429064401}{Source}.\par
\end{sources}\color{black}
\needspace{8\baselineskip}\section{One cubic family, four different outcomes}
\color{black} For \(F=X^3-3X\), the critical points are \(\pm1\), with values \(\mp2\). Thus \(F(X)-T\) has generic group \(S_3\). Its discriminant is\par

\atlasdisplay{\Delta(X^3-3X-T)=27(4-T^2).}

\begin{atlasblock}[standard]
\textbf{\sffamily Exact specializations over \(\mathbb Q\)}\par\smallskip

\begin{itemize}
\item\leavevmode\color{black} \textbf{\(t=0\):} \(X(X^2-3)\) is reducible; its splitting-field group is \(C_2\), fixing the rational root.
\item\leavevmode\color{black} \textbf{\(t=1\):} \(X^3-3X-1\) has no rational root and discriminant \(81\), so its group is \(A_3=C_3\).
\item\leavevmode\color{black} \textbf{\(t=2\):} \((X-2)(X+1)^2\) has a repeated root. This is a branch value, outside the separable-fiber setting.
\item\leavevmode\color{black} \textbf{\(t=3\):} \(X^3-3X-3\) is Eisenstein at \(3\), and its discriminant is \(-135\), a nonsquare. Its group is \(S_3\).
\end{itemize}
\end{atlasblock}\color{black}
\color{black} Even away from the branch values, specialization can lower the group or destroy irreducibility. Conversely, a generic geometric argument can explain why many fibers have a large group without giving a certificate for any chosen one.\par
\begin{sources}
\textbf{Sources for this section}\par\smallskip
\noindent\hyperlink{ref:ConradCubicQuartic}{[ConradCubicQuartic]} / Equation (1.1) and Theorem 2.1. Cubic discriminant formula and the irreducible cubic group test; specialized identities are checked directly. \href{https://kconrad.math.uconn.edu/blurbs/galoistheory/cubicquartic.pdf}{Source}.\par
\noindent\hyperlink{ref:Serre2007Topics}{[Serre2007Topics]} / §4.4. The generic Morse theorem. \href{https://www.routledge.com/Topics-in-Galois-Theory-2nd-Edition/Serre/p/book/9780429064401}{Source}.\par
\end{sources}\color{black}
\needspace{8\baselineskip}\section{Characteristic and proof boundaries}
\color{black} The characteristic-zero hypothesis avoids wild ramification. In characteristic \(p\), the polynomial \(X^p-X\) has derivative \(-1\); its cover \(X^p-X=T\) is unramified at finite points, and its geometric group is the additive cyclic group of order \(p\), acting by \(X\mapsto X+a\) for \(a\in\mathbb F_p\). Ramification at infinity is wild. This shows why the tame step in the proof cannot simply be discarded.\par
\color{black} There are positive-characteristic versions of the Morse theorem, but their hypotheses must be stated separately. The criterion above is intentionally restricted to characteristic zero.\par
\begin{sources}
\textbf{Sources for this section}\par\smallskip
\noindent\hyperlink{ref:Serre2007Topics}{[Serre2007Topics]} / Theorem 4.4.5; remark after Corollary 4.4.7, pp. 41-42. Positive-characteristic scope and the Artin-Schreier boundary. \href{https://www.routledge.com/Topics-in-Galois-Theory-2nd-Edition/Serre/p/book/9780429064401}{Source}.\par
\end{sources}\color{black}
\needspace{11\baselineskip}\section*{Questions and research boundaries}
\needspace{7\baselineskip}\subsection*{Does full symmetric generic monodromy force every separable rational fiber to have the same group?}
{\small\sffamily\color{atlasgold}Settled in stated scope / The characteristic-zero Morse family X³\ensuremath{-}3X\ensuremath{-}T over Q}\par\smallskip
\color{black} No. At \(T=0\) the polynomial is reducible, and at \(T=1\) it is irreducible with group \(C_3\). Both fibers are separable. Hilbert irreducibility supplies many full-group fibers but does not remove all exceptions.\par
\begin{sources}
\textbf{Sources for this section}\par\smallskip
\noindent\hyperlink{ref:Serre2007Topics}{[Serre2007Topics]} / §3.3 and §4.4. Specialization framework; the two counterexamples are calculated in this article. \href{https://www.routledge.com/Topics-in-Galois-Theory-2nd-Edition/Serre/p/book/9780429064401}{Source}.\par
\end{sources}\color{black}
\needspace{6\baselineskip}\par\medskip\begingroup\footnotesize\setlength{\parskip}{1pt}\textbf{\sffamily\color{atlasgreen}Connections}\par
\color{black}\textbf{Precedes.} \hyperlink{article:specialization}{Specialization and Hilbert irreducibility}. Determine the generic group before asking which rational fibers retain it.\par
\color{black}\textbf{Uses.} \hyperlink{article:ramification-inertia}{Ramification and inertia: cycles at bad primes}. Inertia at a simple branch value supplies a transposition.\par
\color{black}\textbf{Uses.} \hyperlink{article:permutation-groups}{The group-theoretic step: transitivity, alternating groups, and symmetric groups}. A transitive group generated by transpositions is symmetric.\par
\color{black}\textbf{Contrasts with.} \hyperlink{article:composition-iteration}{Composition, iteration, and arboreal groups}. Composition creates intermediate covers and block systems; collisions of critical values can reveal that structure.\par
\color{black}\textbf{Constructs examples for.} \hyperlink{article:inverse-galois}{Inverse Galois theory and generic polynomials}. The family X\textasciicircum{}n\ensuremath{-}X\ensuremath{-}T realizes every symmetric group regularly over Q(T).\par
\endgroup
\chapter{Inverse Galois theory and generic polynomials}\label{article:inverse-galois}\hypertarget{article:inverse-galois}{}\color{black}
{\sffamily\small\color{atlasgold}foundation / Intermediate / Q and number fields / Research overview}\par\medskip
Construct a polynomial with a prescribed group, then distinguish one realization from regular families and universal parametrizations.\par\medskip
\textbf{Questions answered}\begin{itemize}
\item What does it mean to realize a finite group by a polynomial?
\item How do regular realizations and generic polynomials differ?
\item Which major existence statements are known?
\end{itemize}
{\small\textbf{Prerequisites:} \hyperlink{article:foundations}{Fields, factors, and the action on roots}; \hyperlink{article:specialization}{Specialization and Hilbert irreducibility}}\par
\needspace{8\baselineskip}\section{Four construction problems}
\color{black} The inverse Galois problem over \(K\) asks whether a given finite group \(G\) is \(\operatorname{Gal}(L/K)\) for some finite Galois extension. An explicit polynomial whose splitting field is \(L\) is a stronger practical output. A \emph{regular realization} over \(K(T)\) requires \(L\cap\overline K=K\). It is a family without an enlarged constant field.\par
\color{black} A polynomial \(P(\mathbf T,x)\) is \emph{parametric} for \(G\) over \(K\) if its generic splitting group is \(G\) and every \(G\)-extension of \(K\) is obtained by specialization. It is \emph{generic} if this remains true over every extension field of \(K\). A family with group \(G\) at infinitely many parameters need not parametrize all \(G\)-extensions.\par
\begin{sources}
\textbf{Sources for this section}\par\smallskip
\noindent\hyperlink{ref:JensenLedetYui2002}{[JensenLedetYui2002]} / §0.1, Definition 0.1.1. Parametric versus generic; existence versus construction. \href{https://library.slmath.org/books/Book45/index.html}{Source}.\par
\end{sources}\color{black}
\needspace{8\baselineskip}\section{From a cover to arithmetic realizations}

\begin{atlasblock}[standard]
\textbf{\sffamily Method  /  regular realization and specialization}\par\smallskip
\color{black} A regular finite Galois extension of \(K(T)\) with nontrivial group \(G\), for \(K\) a number field, gives infinitely many linearly disjoint \(G\)-extensions of \(K\) by Hilbert irreducibility.\par
\end{atlasblock}\color{black}
\color{black} One first builds a cover with a controlled geometric group. Possible tools include invariant fields, branched covers and rigidity, embedding problems, and arithmetic families such as torsion on elliptic curves. One then proves descent to the desired base field and regularity, and finally chooses arithmetic fibers. The descent step is substantial: constructing a complex cover with group \(G\) does not by itself construct a \(G\)-extension of \(\mathbb Q\).\par
\begin{sources}
\textbf{Sources for this section}\par\smallskip
\noindent\hyperlink{ref:Serre2007Topics}{[Serre2007Topics]} / Proposition 3.3.3, Corollary 3.3.4, and §4.1. Linear disjointness from a regular family. \href{https://www.routledge.com/Topics-in-Galois-Theory-2nd-Edition/Serre/p/book/9780429064401}{Source}.\par
\noindent\hyperlink{ref:MalleMatzat2018}{[MalleMatzat2018]}. Rigidity, embedding problems, and geometric constructions. \href{https://link.springer.com/book/10.1007/978-3-662-55420-3}{Source}.\par
\end{sources}\color{black}
\needspace{8\baselineskip}\section{Example: a generic quadratic and a cyclic cubic}

\begin{atlasblock}[standard]
\textbf{\sffamily Worked examples}\par\smallskip
\color{black} If \(\operatorname{char}K\ne2\), the polynomial \(x^2-T\) is generic for \(C_2\): every quadratic extension of every field \(E\supseteq K\) has the form \(E(\sqrt a)\) for a nonsquare \(a\in E\). The exceptional specializations \(T=a\in E^2\) split. Thus ``generic for a group'' never means that every parameter has that group.\par
\color{black} For a single cyclic cubic realization over \(\mathbb Q\), take\par

\atlasdisplay{f=x^3+x^2-2x-1.}
\color{black} The rational-root candidates \(\pm1\) are not roots, so \(f\) is irreducible. Its discriminant is \(49\), hence the transitive group lies in \(A_3=C_3\) and is exactly \(C_3\). In fact a root is \(\zeta_7+\zeta_7^{-1}\), displaying the real cubic subfield of \(\mathbb Q(\zeta_7)\).\par
\end{atlasblock}\color{black}
\begin{sources}
\textbf{Sources for this section}\par\smallskip
\noindent\hyperlink{ref:ConradCubicQuartic}{[ConradCubicQuartic]}. Cubic discriminant test; the explicit calculations are included here. \href{https://kconrad.math.uconn.edu/blurbs/galoistheory/cubicquartic.pdf}{Source}.\par
\noindent\hyperlink{ref:Milne2022}{[Milne2022]}. Cyclotomic extensions and quadratic extensions. \href{https://www.jmilne.org/math/Books/index.html}{Author's book page}.\par
\end{sources}\color{black}
\needspace{8\baselineskip}\section{Established existence results and remaining scale}
\color{black} Every finite solvable group occurs over \(\mathbb Q\) by Shafarevich's theorem. Symmetric and alternating groups also occur over \(\mathbb Q\), as Hilbert proved. These cover many groups but do not settle the assertion for every finite group. The regular inverse problem is stronger still; the literature treats it separately from ordinary existence.\par
\color{black} For a concrete target, record the faithful permutation action as well as the abstract group. A transitive action of degree \(n\) corresponds to a degree-\(n\) irreducible polynomial whose splitting group acts in that way. The regular action has degree \(|G|\), often much larger than the smallest useful polynomial degree. Databases of number fields therefore organize groups by degree and transitive action, not only abstract isomorphism type.\par
\begin{sources}
\textbf{Sources for this section}\par\smallskip
\noindent\hyperlink{ref:JensenLedetYui2002}{[JensenLedetYui2002]} / Theorems 0.2.1 and 0.2.4. Hilbert and Shafarevich existence theorems. \href{https://library.slmath.org/books/Book45/index.html}{Source}.\par
\noindent\hyperlink{ref:KluenersMalle2000}{[KluenersMalle2000]}. Explicit realizations of transitive permutation groups. \href{https://doi.org/10.1006/jsco.2000.0378}{Source}.\par
\noindent\hyperlink{ref:MalleMatzat2018}{[MalleMatzat2018]}. The ordinary and regular inverse problems. \href{https://link.springer.com/book/10.1007/978-3-662-55420-3}{Source}.\par
\end{sources}\color{black}
\needspace{8\baselineskip}\section{What a list of examples proves}

\begin{atlasblock}[boundary]
\color{black} An example realizes its own group and action. Computations for every group up to a degree bound do not prove a theorem for all finite groups. A failure to find a polynomial is not a nonexistence proof. A generic polynomial demands a universality argument, not a successful search over many parameter values.\par
\end{atlasblock}\color{black}
\color{black} For a new family, an effective route is to state the intended group and action; prove an upper bound from symmetry; prove irreducibility; find sufficient local cycle types or resolvent information for the lower bound; then explain which extensions the family is known to parametrize.\par
\begin{sources}
\textbf{Sources for this section}\par\smallskip
\noindent\hyperlink{ref:JensenLedetYui2002}{[JensenLedetYui2002]} / §0.1. The separate universality requirement. \href{https://library.slmath.org/books/Book45/index.html}{Source}.\par
\noindent\hyperlink{ref:FiekerKlueners2014}{[FiekerKlueners2014]}. Group determination for explicit rational polynomials. \href{https://doi.org/10.1112/S1461157013000302}{Source}.\par
\end{sources}\color{black}
\needspace{11\baselineskip}\section*{Questions and research boundaries}
\needspace{7\baselineskip}\subsection*{Does every finite group occur as a Galois group over Q?}
{\small\sffamily\color{atlasgold}Open in cited source / Finite Galois extensions of Q; no fixed degree bound}\par\smallskip
\color{black} This is the classical inverse Galois problem. The cited monographs give constructions and known families, rather than a general solution. A group-specific project must also distinguish an abstract realization from a prescribed transitive action.\par
\begin{sources}
\textbf{Sources for this section}\par\smallskip
\noindent\hyperlink{ref:MalleMatzat2018}{[MalleMatzat2018]}. Inverse Galois problem and known realizations. \href{https://link.springer.com/book/10.1007/978-3-662-55420-3}{Source}.\par
\end{sources}\color{black}
\needspace{7\baselineskip}\subsection*{Does every finite group admit a regular realization over Q(T)?}
{\small\sffamily\color{atlasgold}Conjecture in cited source / Regular finite Galois extensions of Q(T)}\par\smallskip
\color{black} This is Serre's Conjecture 4.1.1. A positive answer would imply the ordinary inverse Galois problem by specialization. The conjecture concerns finite groups, not arbitrary profinite groups.\par
\begin{sources}
\textbf{Sources for this section}\par\smallskip
\noindent\hyperlink{ref:Serre2007Topics}{[Serre2007Topics]} / Conjecture 4.1.1.  \href{https://www.routledge.com/Topics-in-Galois-Theory-2nd-Edition/Serre/p/book/9780429064401}{Source}.\par
\end{sources}\color{black}
\needspace{6\baselineskip}\par\medskip\begingroup\footnotesize\setlength{\parskip}{1pt}\textbf{\sffamily\color{atlasgreen}Connections}\par
\color{black}\textbf{Uses.} \hyperlink{article:specialization}{Specialization and Hilbert irreducibility}. Hilbert irreducibility transfers a regular function-field construction to number fields.\par
\color{black}\textbf{Uses.} \hyperlink{article:galois-computation}{Computing a Galois group with a proof certificate}. A proposed realization needs a certified group calculation.\par
\color{black}\textbf{Uses.} \hyperlink{article:resolvents}{Resolvents: turn group containment into factorization}. Resolvents distinguish candidate permutation groups.\par
\color{black}\textbf{Contrasts with.} \hyperlink{article:finite-field-irreducibility}{Irreducibility over finite fields}. Over a finite base field, every finite Galois group is cyclic.\par
\color{black}\textbf{Applies to.} \hyperlink{article:composition-iteration}{Composition, iteration, and arboreal groups}. Arboreal existence problems prescribe an entire compatible tower of groups.\par
\endgroup
\part{Composition and iteration}
\chapter{Composition, iteration, and arboreal groups}\label{article:composition-iteration}\hypertarget{article:composition-iteration}{}\color{black}
{\sffamily\small\color{atlasgold}method / Intermediate / general fields and number fields / Research overview}\par\medskip
Use extension towers and rooted trees to study repeated polynomial composition. Separate stability from maximality of the Galois action.\par\medskip
\textbf{Questions answered}\begin{itemize}
\item When is a composition irreducible?
\item Why do iterated Galois groups lie in wreath products?
\item What extra arithmetic is needed for all iteration levels?
\end{itemize}
{\small\textbf{Prerequisites:} \hyperlink{article:foundations}{Fields, factors, and the action on roots}; \hyperlink{article:permutation-groups}{The group-theoretic step: transitivity, alternating groups, and symmetric groups}}\par
\needspace{8\baselineskip}\section{Multiplication and composition are different}
\color{black} A factorization \(f=gh\) and a decomposition \(f=g\circ h\) answer different questions. In composition theory the factors are nonconstant maps, and the units are linear changes of variable. For example, \(x^4=(x^2)\circ(x^2)\), while over \(\mathbb C\) every polynomial already factors multiplicatively into linear factors. Ritt's ``prime'' polynomials are indecomposable under composition.\par
\color{black} For \(f\in\mathbb C[x]\), composition is connected to intermediate fields of \(\mathbb C(x)/\mathbb C(f(x))\) and to blocks in the monodromy action of \(f(x)-T\). This is a function-field group, not the splitting group of the particular polynomial \(f(x)\) over \(\mathbb Q\).\par
\begin{sources}
\textbf{Sources for this section}\par\smallskip
\noindent\hyperlink{ref:Ritt1922}{[Ritt1922]}. Functional decomposition; consult the recorded erratum. \href{https://www.jstor.org/stable/1988911}{Source}.\par
\noindent\hyperlink{ref:Mueller1995Monodromy}{[Mueller1995Monodromy]}. Primitive monodromy of complex polynomial maps. \href{https://ypfmde.github.io/Papers/mon.pdf}{Source}.\par
\end{sources}\color{black}
\needspace{8\baselineskip}\section{The composition test}

\begin{atlasblock}[standard]
\textbf{\sffamily Theorem  /  Capelli's composition lemma}\par\smallskip
\color{black} Let \(F\in K[x]\) be irreducible of degree \(m\), let \(g\in K[x]\) have degree \(d\ge1\), and let \(\alpha\) be a root of \(F\). Then \(F(g(x))\) is irreducible over \(K\) if and only if \(g(x)-\alpha\) is irreducible over \(K(\alpha)\).\par
\end{atlasblock}\color{black}

\begin{atlasblock}[standard]
\color{black} For a root \(\beta\) of \(g(x)-\alpha\), we have \(\alpha=g(\beta)\), so \(K(\alpha)\subseteq K(\beta)\). The tower law gives \([K(\beta):K]=m[K(\beta):K(\alpha)]\). This equals \(md=\deg(F\circ g)\) exactly when \(g(x)-\alpha\) has degree \(d\) as a minimal polynomial. This proves both implications.\par
\end{atlasblock}\color{black}
\color{black} The challenge has moved to an extension field; it has not vanished. In quadratic iteration, deciding whether certain elements become squares is the essential new step.\par
\begin{sources}
\textbf{Sources for this section}\par\smallskip
\noindent\hyperlink{ref:Schinzel2000}{[Schinzel2000]}. Composition and reducibility; the field-degree proof is given in full here. \href{https://doi.org/10.1017/CBO9780511542916}{Source}.\par
\noindent\hyperlink{ref:Odoni1985}{[Odoni1985]}. Composition as the starting point for iterated Galois theory. \href{https://londmathsoc.onlinelibrary.wiley.com/doi/10.1112/plms/s3-51.3.385}{Source}.\par
\end{sources}\color{black}
\needspace{8\baselineskip}\section{The natural ambient group is a rooted-tree group}
\color{black} Let \(f\in K[x]\) have degree \(d\ge2\), and suppose each iterate \(f^{\circ n}\) is separable. Put a separate copy of its roots at level \(n\), with an edge from \(\beta\) to \(f(\beta)\). The resulting regular rooted \(d\)-ary tree is preserved by Galois, because every automorphism commutes with evaluation of \(f\).\par

\atlasdisplay{G_n\ \le\ \operatorname{Aut}(T_{d,n})\cong [S_d]^n,\qquad |\operatorname{Aut}(T_{d,n})|=(d!)^{1+d+\cdots+d^{n-1}}.}
\color{black} The right side is an iterated wreath product, generally far smaller than \(S_{d^n}\). Irreducibility of every iterate, called \emph{stability}, says the action is transitive at every level. Surjectivity onto the tree automorphism group is stronger.\par
\begin{sources}
\textbf{Sources for this section}\par\smallskip
\noindent\hyperlink{ref:Jones2013Arboreal}{[Jones2013Arboreal]} / §1. The preimage tree, arboreal representation, and natural maximal group. \href{https://www.numdam.org/articles/10.5802/pmb.a-154/}{Source}.\par
\end{sources}\color{black}
\needspace{8\baselineskip}\section{Example: the first two levels of x\ensuremath{{}^2}+1}

\begin{atlasblock}[standard]
\textbf{\sffamily Worked example}\par\smallskip
\color{black} For \(f=x^2+1\), the first iterate is irreducible over \(\mathbb Q\), and\par

\atlasdisplay{f^{\circ2}=x^4+2x^2+2}
\color{black} is Eisenstein at \(2\). Its four roots fall into two pairs according to whether \(f(\beta)=i\) or \(-i\). Thus its group is contained in \(C_2\wr C_2\), of order eight.\par
\color{black} The quadratic critical-orbit criterion uses \(c_1=-a\) and \(c_n=f_a^{\circ n}(0)\) for \(n\ge2\), where \(f_a=x^2+a\) and \(a\in\mathbb Q\setminus\{0,-1\}\). It says that the level-\(n\) group is full precisely when \(c_1,\ldots,c_n\) are independent in \(\mathbb Q^\times/\mathbb Q^{\times2}\). Here \((c_1,c_2)=(-1,2)\); none of \(-1,2,-2\) is a rational square. Consequently \(G_2=C_2\wr C_2\), the dihedral group of order eight. Checking two levels alone does not prove maximality at every level.\par
\end{atlasblock}\color{black}
\begin{sources}
\textbf{Sources for this section}\par\smallskip
\noindent\hyperlink{ref:Stoll1992}{[Stoll1992]} / Theorem in §1, part 1, equivalence (a)-(b). Square-class criterion for quadratic iterates; restated in Stoll2026 Proposition1.7. \href{https://link.springer.com/article/10.1007/BF01197321}{Source}.\par
\end{sources}\color{black}
\needspace{8\baselineskip}\section{Stable does not mean maximal: x\ensuremath{{}^2}\ensuremath{-}2 at every level}
\color{black} A polynomial is \emph{postcritically finite} (PCF) when every critical point has finite forward orbit. For \(f=x^2-2\), the finite critical orbit is \(0\mapsto-2\mapsto2\mapsto2\). Nevertheless every iterate is irreducible over \(\mathbb Q\).\par

\begin{atlasblock}[standard]
\textbf{\sffamily Exact cyclotomic description}\par\smallskip
\color{black} The identity \(f(z+z^{-1})=z^2+z^{-2}\) gives\par

\atlasdisplay{f^{\circ n}(z+z^{-1})=z^{2^n}+z^{-2^n}.}
\color{black} For a primitive \(2^{n+2}\)-th root of unity \(\zeta\), the element \(\zeta+\zeta^{-1}\) is therefore a root of \(f^{\circ n}\). It has degree \(\varphi(2^{n+2})/2=2^n\), the degree of the iterate. All its conjugates lie in the same real cyclotomic field, so\par

\atlasdisplay{K_n=\mathbb Q(\zeta+\zeta^{-1}),\qquad G_n\cong C_{2^n}.}
\color{black} Thus the action is transitive at every level, whereas \(|G_n|=2^n\) is much smaller than \(|\operatorname{Aut}(T_{2,n})|=2^{2^n-1}\) for \(n\ge2\).\par
\end{atlasblock}\color{black}
\color{black} The recent PCF work of Benedetto-Ghioca-Juul-Tucker describes the constrained ambient groups and gives criteria for attaining them: the 2025 paper treats periodic critical points; its 2026 companion treats strictly preperiodic ones. These groups depend on critical-orbit relations, the base field, and the target point. The classical Chebyshev case above is separated from the general constructions in the latter paper.\par
\begin{sources}
\textbf{Sources for this section}\par\smallskip
\noindent\hyperlink{ref:BenedettoGhiocaJuulTucker2025Periodic}{[BenedettoGhiocaJuulTucker2025Periodic]} / §§1.2-1.6, Theorem 1.5. Periodic PCF case; maximality is inside the constrained group M, not the full binary-tree group. \href{https://doi.org/10.1112/jlms.70257}{Source}.\par
\noindent\hyperlink{ref:BenedettoGhiocaJuulTucker2026Preperiodic}{[BenedettoGhiocaJuulTucker2026Preperiodic]} / §1.4; Theorems 1.2 and 7.6. §1.4 explicitly separates x\ensuremath{{}^2}\ensuremath{-}2 and records the Chebyshev identity; the cyclotomic calculation at target zero is supplied here. \href{https://doi.org/10.4153/S0008414X2610217X}{Source}.\par
\end{sources}\color{black}
\needspace{8\baselineskip}\section{Generic maximality versus an arithmetic tower}
\color{black} For the generic monic degree-\(d\) polynomial \(F=x^d+s_{d-1}x^{d-1}+\cdots+s_0\) over a characteristic-zero field \(k\), Odoni proves \(\operatorname{Gal}(F^{\circ n}/k(\mathbf s))=[S_d]^n\) for every fixed \(n\). Juul extends this to positive characteristic, excluding \((d,\operatorname{char}k)=(2,2)\) for polynomials. Disjoint ramification in the new branches supplies the independent permutations needed at each induction step.\par
\color{black} Over number fields, Benedetto-Juul prove all-level existence in every even degree and in odd degree when \([K:\mathbb Q]\) is odd. Specter's 2018 preprint states all degrees over every number field. In contrast, Dittmann-Kadets construct characteristic-zero Hilbertian fields admitting no finite-index arboreal image. The base-field quantifier matters.\par
\begin{sources}
\textbf{Sources for this section}\par\smallskip
\noindent\hyperlink{ref:Odoni1985}{[Odoni1985]}. Characteristic-zero generic theorem. \href{https://londmathsoc.onlinelibrary.wiley.com/doi/10.1112/plms/s3-51.3.385}{Source}.\par
\noindent\hyperlink{ref:Juul2019}{[Juul2019]} / Theorem 1.1. Positive-characteristic exception and generic full wreath products. \href{https://arxiv.org/abs/1410.3814}{Source}.\par
\noindent\hyperlink{ref:BenedettoJuul2019}{[BenedettoJuul2019]} / Abstract and main theorem. Published number-field cases. \href{https://londmathsoc.onlinelibrary.wiley.com/doi/10.1112/blms.12225}{Source}.\par
\noindent\hyperlink{ref:Specter2018}{[Specter2018]} / Version 2. All-degree number-field statement, labeled as a preprint. \href{https://arxiv.org/abs/1803.00434v2}{Source}.\par
\noindent\hyperlink{ref:DittmannKadets2022}{[DittmannKadets2022]} / Theorem 1.2. General Hilbertian-field counterexample. \href{https://arxiv.org/abs/2012.03076}{Source}.\par
\end{sources}\color{black}
\needspace{8\baselineskip}\section{A recent boundary: how many levels suffice?}
\color{black} \textbf{Preprint submitted 10 November 2026.} Stoll gives revised proofs and rational-parameter extensions of maximality criteria for \(x^2+a\). Proposition 1.7 restates the square-class test. Theorem 6.1 gives explicit regions and congruence classes where level two controls the full tree; it is not a theorem for every rational \(a\).\par
\color{black} The same paper formulates a uniform finite-level detection conjecture and reports supporting computations. Computational verification to a finite height or depth is recorded as evidence, not as an all-parameter theorem. Critical-orbit square relations explain why a tower may first fail at a later level.\par

\begin{atlasblock}[standard]
\textbf{\sffamily A later failure with an exact certificate}\par\smallskip
\color{black} For \(a=2/3\), the first four adjusted critical values are\par

\atlasdisplay{(c_1,c_2,c_3,c_4)=\left(-\frac23,\frac{10}{9},\frac{154}{81},\frac{28090}{6561}\right).}
\color{black} The first three square classes are independent: the sign detects \(c_1\), the prime \(5\) detects \(c_2\), and the prime \(7\) detects \(c_3\). But \(c_2c_4=(530/243)^2\). Hence the level-three group is full and the level-four group is not. A universal detection level, if one exists, must therefore be at least four.\par
\end{atlasblock}\color{black}
\begin{sources}
\textbf{Sources for this section}\par\smallskip
\noindent\hyperlink{ref:Stoll2026Arboreal}{[Stoll2026Arboreal]} / Proposition 1.7, Theorem 6.1, Proposition 8.2, and Remark 8.5. arXiv:2611.14014v1, 10 November 2026. The a=2/3 relation is checked directly; it is an exact obstruction, not evidence from a finite numerical search. \href{https://arxiv.org/abs/2611.14014v1}{Source}.\par
\end{sources}\color{black}
\needspace{11\baselineskip}\section*{Questions and research boundaries}
\needspace{7\baselineskip}\subsection*{Does a universal finite level detect maximality for x\ensuremath{{}^2}+a over Q?}
{\small\sffamily\color{atlasgold}Conjecture in cited source / a\ensuremath{\in}Q, base point 0, full binary preimage tree}\par\smallskip
\color{black} Stoll's Conjecture 8.1 asks whether one integer \(N\) suffices: if the level-\(N\) representation is surjective, must the full arboreal representation be surjective? The 2026 preprint discusses evidence for \(N=4\), not a proof for all rational parameters. The exact example \(a=2/3\) above rules out every \(N<4\). Changing the base point gives a different question.\par
\begin{sources}
\textbf{Sources for this section}\par\smallskip
\noindent\hyperlink{ref:Stoll2026Arboreal}{[Stoll2026Arboreal]} / Conjecture 8.1 and Remark 8.5.  \href{https://arxiv.org/abs/2611.14014v1}{Source}.\par
\end{sources}\color{black}
\needspace{7\baselineskip}\subsection*{Does Hilbertianity alone force an all-level full arboreal realization?}
{\small\sffamily\color{atlasgold}Settled in stated scope / General Hilbertian fields of characteristic zero}\par\smallskip
\color{black} No. Dittmann-Kadets give counterexample fields. This does not contradict the positive existence results over number fields, and does not negate generic maximality at every fixed finite level.\par
\begin{sources}
\textbf{Sources for this section}\par\smallskip
\noindent\hyperlink{ref:DittmannKadets2022}{[DittmannKadets2022]} / Theorem 1.2.  \href{https://arxiv.org/abs/2012.03076}{Source}.\par
\end{sources}\color{black}
\needspace{6\baselineskip}\par\medskip\begingroup\footnotesize\setlength{\parskip}{1pt}\textbf{\sffamily\color{atlasgreen}Connections}\par
\color{black}\textbf{Uses.} \hyperlink{article:specialization}{Specialization and Hilbert irreducibility}. Generic finite-level groups specialize, while all-level arithmetic needs an additional argument.\par
\color{black}\textbf{Uses.} \hyperlink{article:permutation-groups}{The group-theoretic step: transitivity, alternating groups, and symmetric groups}. Blocks and wreath products identify the ambient action.\par
\color{black}\textbf{Uses.} \hyperlink{article:discriminants}{Discriminants, parity, and ramification}. Square classes and ramification control new quadratic branches.\par
\color{black}\textbf{Generalizes.} \hyperlink{article:inverse-galois}{Inverse Galois theory and generic polynomials}. The target is a compatible tower of prescribed finite groups.\par
\color{black}\textbf{Uses.} \hyperlink{article:newton-polygons}{Newton polygons, factor degrees, and large cycles}. Valuations can prove irreducibility or produce new ramification at each level.\par
\endgroup
\part{Finite fields}
\chapter{Irreducibility over finite fields}\label{article:finite-field-irreducibility}\hypertarget{article:finite-field-irreducibility}{}\color{black}
{\sffamily\small\color{atlasgold}method / Start here / finite fields / Classical theory}\par\medskip
Turn Frobenius powers into exact irreducibility certificates, count irreducibles, and distinguish irreducible from primitive polynomials.\par\medskip
\textbf{Questions answered}\begin{itemize}
\item How can irreducibility be certified without a full factorization?
\item How many monic irreducibles of a given degree exist?
\item Why is the Galois group cyclic, and what does primitive mean?
\end{itemize}
{\small\textbf{Prerequisites:} \hyperlink{article:foundations}{Fields, factors, and the action on roots}}\par
\needspace{8\baselineskip}\section{The identity behind the algorithms}
\color{black} Let \(q\) be a prime power. In \(\mathbb F_q[x]\),\par

\atlasdisplay{x^{q^m}-x=\prod_{\substack{P\text{ monic irreducible}\\\deg P\mid m}}P(x).}
\color{black} Every element of \(\mathbb F_{q^m}\) is a root, the derivative is \(-1\), and an element of degree \(d\) over \(\mathbb F_q\) belongs to \(\mathbb F_{q^m}\) exactly when \(d\mid m\). Thus each indicated factor occurs once.\par
\color{black} If \(f\) is irreducible of degree \(n\), adjoining one root gives \(\mathbb F_{q^n}\), which already contains every conjugate. Its Galois group is \(C_n\), generated by \(z\mapsto z^q\). Noncyclic monodromy groups in finite-field research concern function fields such as \(\mathbb F_q(T)\), not this constant-field extension.\par
\begin{sources}
\textbf{Sources for this section}\par\smallskip
\noindent\hyperlink{ref:LidlNiederreiter1997}{[LidlNiederreiter1997]}. Finite-field structure and irreducible polynomials. \href{https://www.cambridge.org/core/books/finite-fields/75BDAA74ABAE713196E718392B9E5E72}{Source}.\par
\noindent\hyperlink{ref:GuralnickZieve2010}{[GuralnickZieve2010]}. A contrasting function-field monodromy setting. \href{https://annals.math.princeton.edu/2010/172-2/p11}{Source}.\par
\end{sources}\color{black}
\needspace{8\baselineskip}\section{A necessary and sufficient test}

\begin{atlasblock}[standard]
\textbf{\sffamily Theorem  /  Frobenius-power criterion}\par\smallskip
\color{black} A monic polynomial \(f\in\mathbb F_q[x]\) of degree \(n\ge1\) is irreducible if and only if\par

\atlasdisplay{x^{q^n}\equiv x\pmod f,\qquad \gcd\bigl(f,x^{q^{n/\ell}}-x\bigr)=1\quad\text{for every prime }\ell\mid n.}
\end{atlasblock}\color{black}
\color{black} The first condition makes \(f\) squarefree and forces every factor degree to divide \(n\). If a proper divisor \(d\mid n\) occurred as a factor degree, some prime \(\ell\mid n/d\) would give \(d\mid n/\ell\), contradicting the corresponding gcd. The reverse implication follows from the same Frobenius identity. Modular powering and Euclidean gcds suffice; there is no need to expand the huge power.\par
\begin{sources}
\textbf{Sources for this section}\par\smallskip
\noindent\hyperlink{ref:Rabin1980}{[Rabin1980]}. Irreducibility testing in finite-field algorithms. \href{https://doi.org/10.1137/0209024}{Source}.\par
\noindent\hyperlink{ref:FLINTFactorizationDocumentation}{[FLINTFactorizationDocumentation]} / Irreducibility testing over finite fields, Theorem 1. Official documentation states the criterion and its exact algorithmic use. \href{https://flintlib.org/doc/algorithms/polynomials/factoring-finite-fields.html}{Source}.\par
\end{sources}\color{black}
\needspace{8\baselineskip}\section{Example: a four-degree certificate over F\ensuremath{{}_2}}

\begin{atlasblock}[standard]
\textbf{\sffamily Worked example}\par\smallskip
\color{black} Take \(f=x^4+x+1\in\mathbb F_2[x]\). Repeated squaring modulo \(f\) gives\par

\atlasdisplay{x^4\equiv x+1,\qquad x^8\equiv x^2+1,\qquad x^{16}\equiv x.}
\color{black} The only prime divisor of \(n=4\) is \(2\), and \(x^4-x\equiv1\pmod f\). Thus the required gcd is one, and \(f\) is irreducible. Its splitting field has 16 elements, and the Galois group over \(\mathbb F_2\) has order four.\par
\color{black} By contrast, the congruence \(x^{16}\equiv x\pmod f\) alone would not prove irreducibility: the reducible quartic \(x^4+x=x(x+1)(x^2+x+1)\), for instance, also divides \(x^{16}-x\). The gcd conditions eliminate smaller-degree factors.\par
\end{atlasblock}\color{black}
\begin{sources}
\textbf{Sources for this section}\par\smallskip
\noindent\hyperlink{ref:Rabin1980}{[Rabin1980]}. Application of the exact criterion; the displayed remainders are direct calculations. \href{https://doi.org/10.1137/0209024}{Source}.\par
\end{sources}\color{black}
\needspace{8\baselineskip}\section{Counting and constructing irreducibles}
\color{black} If \(N_q(n)\) counts monic irreducibles of degree \(n\), taking degrees in the Frobenius identity and applying Möbius inversion gives\par

\atlasdisplay{N_q(n)=\frac1n\sum_{d\mid n}\mu(d)q^{n/d}.}
\color{black} For example, \(N_2(4)=(16-4)/4=3\). For fixed \(q\) and large \(n\), the proportion of monic degree-\(n\) polynomials that are irreducible is asymptotic to \(1/n\). Random sampling followed by exact testing therefore supplies a simple construction method. Deterministic constructions and fast implementations are separate algorithmic questions, developed by Shoup and others.\par
\begin{sources}
\textbf{Sources for this section}\par\smallskip
\noindent\hyperlink{ref:LidlNiederreiter1997}{[LidlNiederreiter1997]}. Counting formula from Möbius inversion. \href{https://www.cambridge.org/core/books/finite-fields/75BDAA74ABAE713196E718392B9E5E72}{Source}.\par
\noindent\hyperlink{ref:Shoup1990Irreducible}{[Shoup1990Irreducible]}. Algorithms for constructing irreducible polynomials. \href{https://www.shoup.net/papers/detirred.pdf}{Source}.\par
\noindent\hyperlink{ref:Shoup1994}{[Shoup1994]}. Fast constructions. \href{https://doi.org/10.1006/jsco.1994.1025}{Source}.\par
\end{sources}\color{black}
\needspace{8\baselineskip}\section{Primitive polynomials and prescribed coefficients}
\color{black} An irreducible polynomial of degree \(n\) is \emph{primitive} if a root generates \(\mathbb F_{q^n}^{\times}\), a group of order \(q^n-1\). This is a statement about the multiplicative order of a root, not about the order of the Galois group. For the quartic above, the root \(\alpha\) has order 15: \(\alpha^3\ne1\) and \(\alpha^5=\alpha^2+\alpha\ne1\), excluding the proper maximal divisors 5 and 3 of 15.\par
\color{black} Prescribing coefficients, insisting on a sparse polynomial, or asking for a primitive root changes the counting problem. The Hansen-Mullen irreducibility problem and the stronger primitivity problem are different. Wan's theorem and the finite checks of Ham-Mullen concern irreducibility; Tuxanidy-Wang give a later proof. Results prescribing several coefficients have parameter restrictions; in particular a prescribed constant coefficient zero rules out irreducibility in degree greater than one.\par
\begin{sources}
\textbf{Sources for this section}\par\smallskip
\noindent\hyperlink{ref:HansenMullen1992}{[HansenMullen1992]}. Primitive polynomials and prescribed coefficients. \href{https://doi.org/10.1090/S0025-5718-1992-1134730-7}{Source}.\par
\noindent\hyperlink{ref:Wan1997}{[Wan1997]}. Hansen-Mullen irreducibility theorem outside finitely many cases. \href{https://doi.org/10.1090/S0025-5718-97-00835-1}{Source}.\par
\noindent\hyperlink{ref:HamMullen1998}{[HamMullen1998]}. Completion of the remaining finite cases. \href{https://doi.org/10.1090/S0025-5718-98-00904-1}{Source}.\par
\noindent\hyperlink{ref:TuxanidyWang2018}{[TuxanidyWang2018]}. A later proof of the irreducibility theorem. \href{https://arxiv.org/abs/1604.04023}{Source}.\par
\noindent\hyperlink{ref:Pollack2013Prescribed}{[Pollack2013Prescribed]}. Several prescribed coefficients; the range is approximately a square root of the degree, not the full degree. \href{https://www.pollack-math.net/prescribed.pdf}{Source}.\par
\end{sources}\color{black}
\needspace{8\baselineskip}\section{Parity gives an obstruction, not a certificate}
\color{black} For a squarefree polynomial of degree \(n\) with \(r\) irreducible factors over a finite field of odd order, Frobenius has sign \((-1)^{n-r}\). Its discriminant is a square exactly when \(n-r\) is even. This is the Stickelberger discriminant-parity theorem. Swan develops its characteristic-two lifting form and applies it to trinomials. An admissible parity of the number of factors still allows several factors; it does not prove irreducibility.\par
\begin{sources}
\textbf{Sources for this section}\par\smallskip
\noindent\hyperlink{ref:Swan1962}{[Swan1962]} / Theorem 1 and added-in-proof attribution, pp.1099-1100. The basic theorem is credited by Swan to Stickelberger, with Dalen's prior proof noted. \href{https://msp.org/pjm/1962/12-3/pjm-v12-n3-p27-p.pdf}{Source}.\par
\end{sources}\color{black}
\needspace{11\baselineskip}\section*{Questions and research boundaries}
\needspace{7\baselineskip}\subsection*{Which prescribed supports admit irreducibles over a fixed finite field?}
{\small\sffamily\color{atlasgold}Research direction / A specified finite field, degree, and support}\par\smallskip
\color{black} Existence without coefficient restrictions is settled by the counting formula. Prescribing a support or coefficient pattern can introduce obstructions and is a separate question. A useful project should identify the exact support, exclude forced factors, and publish reproducible certificates together with any general theorem.\par
\begin{sources}
\textbf{Sources for this section}\par\smallskip
\noindent\hyperlink{ref:HansenMullen1992}{[HansenMullen1992]}.  \href{https://doi.org/10.1090/S0025-5718-1992-1134730-7}{Source}.\par
\noindent\hyperlink{ref:Swan1962}{[Swan1962]}.  \href{https://msp.org/pjm/1962/12-3/pjm-v12-n3-p27-p.pdf}{Source}.\par
\noindent\hyperlink{ref:MullenPanario2013}{[MullenPanario2013]}.  \href{https://doi.org/10.1201/b15006}{Source}.\par
\end{sources}\color{black}
\needspace{6\baselineskip}\par\medskip\begingroup\footnotesize\setlength{\parskip}{1pt}\textbf{\sffamily\color{atlasgreen}Connections}\par
\color{black}\textbf{Supplies examples for.} \hyperlink{article:modular-irreducibility}{Irreducibility by reduction modulo primes}. Finite-field certificates can prove rational irreducibility at good primes.\par
\color{black}\textbf{Supplies examples for.} \hyperlink{article:frobenius}{Frobenius: from factor degrees to permutations}. Finite-field factor degrees record global Frobenius cycle types.\par
\color{black}\textbf{Uses.} \hyperlink{article:factorization-algorithms}{Exact factorization: finite fields, Hensel lifting, and lattices}. Modular exponentiation and gcds are shared algorithmic components.\par
\color{black}\textbf{Contrasts with.} \hyperlink{article:inverse-galois}{Inverse Galois theory and generic polynomials}. Only cyclic groups occur over a finite constant field.\par
\color{black}\textbf{Uses.} \hyperlink{article:discriminants}{Discriminants, parity, and ramification}. Factor-count parity is encoded by the discriminant in odd characteristic.\par
\endgroup
\chapter{Finite fields: construction, prescribed coefficients, and distribution}\label{article:finite-field-constructions}\hypertarget{article:finite-field-constructions}{}\color{black}
{\sffamily\small\color{atlasgold}method / Intermediate / finite fields / Published results}\par\medskip
Count and construct finite-field irreducibles, prescribe selected coefficients, or restrict every coefficient to an allowed set. Keep the coefficient distribution and the limiting regime explicit.\par\medskip
\textbf{Questions answered}\begin{itemize}
\item How many monic irreducibles of degree n exist over a finite field?
\item How does counting produce a certified construction algorithm?
\item Which arbitrary coefficients can be prescribed simultaneously?
\item Can every coefficient be restricted to an allowed set, and how does this differ from prescribing a few coefficients?
\item Why do factor degrees resemble permutation cycles, and in which limit?
\end{itemize}
{\small\textbf{Prerequisites:} \hyperlink{article:foundations}{Fields, factors, and the action on roots}; \hyperlink{article:finite-field-irreducibility}{Irreducibility over finite fields}}\par
\needspace{8\baselineskip}\section{Counting by Frobenius orbits}
\color{black} Let \(N_q(n)\) count monic irreducibles of degree \(n\) over \(\mathbb F_q\), where \(q\) is a prime power. Every element of \(\mathbb F_{q^m}\) is a root of \(x^{q^m}-x\), whose derivative is \(-1\). A monic irreducible divides this polynomial exactly when its degree divides \(m\).\par

\begin{atlasblock}[standard]
\textbf{\sffamily Proof mechanism  /  partition by minimal polynomial}\par\smallskip

\atlasdisplay{q^m=\sum_{d\mid m}dN_q(d),\qquad N_q(n)=\frac1n\sum_{d\mid n}\mu(d)q^{n/d}.}
\color{black} The second identity is Möbius inversion; \(\mu\) is the Möbius function. Every proper divisor of \(n\) is at most \(n/2\), giving \(N_q(n)=q^n/n+O(q^{n/2}/n)\), uniformly for \(q\ge2\). Thus a uniform monic polynomial is irreducible with probability asymptotic to \(1/n\) as \(n\to\infty\), even for fixed \(q\).\par
\end{atlasblock}\color{black}

\begin{atlasblock}[standard]
\textbf{\sffamily Exact example}\par\smallskip

\atlasdisplay{N_2(6)=\frac{64-8-4+2}{6}=9.}
\color{black} There are nine successes among all 64 monic binary sextics, or among the 32 with constant coefficient one.\par
\end{atlasblock}\color{black}
\begin{sources}
\textbf{Sources for this section}\par\smallskip
\noindent\hyperlink{ref:LidlNiederreiter1997}{[LidlNiederreiter1997]}. Finite-field structure and irreducible-polynomial counting; the derivation is included explicitly. \href{https://www.cambridge.org/core/books/finite-fields/75BDAA74ABAE713196E718392B9E5E72}{Source}.\par
\noindent\hyperlink{ref:Pollack2013Prescribed}{[Pollack2013Prescribed]} / §1 and Lemma 4. Counting normalization and the standard estimate for degree-n irreducibles. \href{https://www.pollack-math.net/prescribed.pdf}{Source}.\par
\end{sources}\color{black}
\needspace{8\baselineskip}\section{From abundance to construction}
\color{black} Sample monic polynomials independently and apply the exact Frobenius-power test in the companion article. Stop only after the test proves irreducibility. This is a Las Vegas construction: runtime is random, but every returned answer is certified. For binary sextics, sampling only constant-one candidates takes exactly \(32/9\) trials on average.\par
\color{black} If an extension \(\mathbb F_{q^n}\) is already represented, choose \(\alpha\) outside its maximal proper subfields. Equivalently, require \(\alpha^{q^{n/\ell}}\ne\alpha\) for every prime \(\ell\mid n\). Its minimal polynomial is then \(\prod_{i=0}^{n-1}(x-\alpha^{q^i})\). This is a construction inside an available extension, not a way to obtain that initial representation for free.\par
\color{black} Shoup develops deterministic constructions, including bounds polynomial for small characteristic, and faster construction methods. A complexity bound polynomial in the characteristic \(p\) must not be read as polynomial in its input length \(\log p\).\par
\begin{sources}
\textbf{Sources for this section}\par\smallskip
\noindent\hyperlink{ref:Rabin1980}{[Rabin1980]}. Probabilistic finite-field algorithms, with exact irreducibility verification. \href{https://doi.org/10.1137/0209024}{Source}.\par
\noindent\hyperlink{ref:Shoup1990Irreducible}{[Shoup1990Irreducible]} / Introduction. Deterministic constructions and dependence on the characteristic. \href{https://www.shoup.net/papers/detirred.pdf}{Source}.\par
\noindent\hyperlink{ref:Shoup1994}{[Shoup1994]}. Faster construction of irreducible polynomials of prescribed degree. \href{https://doi.org/10.1006/jsco.1994.1025}{Source}.\par
\end{sources}\color{black}
\needspace{8\baselineskip}\section{One coefficient: a complete existence theorem}

\begin{atlasblock}[standard]
\textbf{\sffamily Theorem  /  Hansen-Mullen, including the genuine exceptions}\par\smallskip
\color{black} Given \(n\ge2\), \(0\le j<n\), and \(c\in\mathbb F_q\), a monic irreducible degree-\(n\) polynomial with coefficient \([x^j]f=c\) exists, except precisely when \(j=0,c=0\), or when \(q\) is even and \((n,j,c)=(2,1,0)\).\par
\end{atlasblock}\color{black}
\color{black} The exceptions have elementary explanations: a zero constant term supplies the factor \(x\); in characteristic two, \(x^2+c\) is a square because finite fields are perfect. For \(n\ge3\), the constant-term obstruction is the only one.\par
\color{black} Wan proved a large-parameter range and Ham-Mullen finished the remaining cases. Tuxanidy-Wang give another proof: sufficiently long periods of finite-field discrete Fourier transforms force elements of full degree in supports of functions encoding elementary symmetric coefficients. Irreducibility here does not require a root to generate the multiplicative group.\par
\begin{sources}
\textbf{Sources for this section}\par\smallskip
\noindent\hyperlink{ref:Wan1997}{[Wan1997]}. The large-parameter part of the Hansen-Mullen irreducibility theorem. \href{https://doi.org/10.1090/S0025-5718-97-00835-1}{Source}.\par
\noindent\hyperlink{ref:HamMullen1998}{[HamMullen1998]}. Completion of the remaining cases. \href{https://doi.org/10.1090/S0025-5718-98-00904-1}{Source}.\par
\noindent\hyperlink{ref:TuxanidyWang2018}{[TuxanidyWang2018]} / Theorem 1.1 and Introduction. Exact single-coefficient statement, exceptions, and the Fourier-period proof mechanism. \href{https://arxiv.org/abs/1604.04023}{Source}.\par
\end{sources}\color{black}
\needspace{8\baselineskip}\section{Several coefficients: uniformity and explicit bounds}
\color{black} Let \(I\subset\{0,\ldots,n-1\}\) have size \(r\); prescribe arbitrary \(a_i\in\mathbb F_q\), requiring \(a_0\ne0\) if \(0\in I\). Pollack proves, for fixed \(\varepsilon>0\) and \(r\le(1-\varepsilon)\sqrt n\),\par

\atlasdisplay{\#\{f\text{ irreducible with these coefficients}\}\sim\sigma N_q(n),\quad \sigma=\begin{cases}q^{-r},&0\notin I,\\q^{1-r}/(q-1),&0\in I,\end{cases}}
\color{black} as \(n\to\infty\), uniformly in the field and prescriptions.\par

\begin{atlasblock}[standard]
\textbf{\sffamily Theorem  /  Ha, an explicit linear range}\par\smallskip
\color{black} For every prime power \(q\), if \(n\ge52\) and \(r\le n/10\), every such prescription is realized by a monic irreducible. Thus any six chosen coefficients of a degree-60 polynomial over \(\mathbb F_2\) can be prescribed, subject to the constant-term condition.\par
\end{atlasblock}\color{black}
\color{black} Ha also obtains \(\lfloor(1/4-\varepsilon)n\rfloor\) prescriptions for \(n\ge8\), \(0<\varepsilon<1/4\), and \(q\ge q_0(\varepsilon)\). These methods use exponential-sum estimates and the function-field circle method to compare a positive main term with an error. They provide existence beyond unrestricted random sampling; they do not assert that every sparse pattern is possible.\par
\begin{sources}
\textbf{Sources for this section}\par\smallskip
\noindent\hyperlink{ref:Pollack2013Prescribed}{[Pollack2013Prescribed]} / Theorem 1 and Corollary 2. Uniform growing-degree asymptotic, including the constant-coefficient normalization. \href{https://www.pollack-math.net/prescribed.pdf}{Source}.\par
\noindent\hyperlink{ref:Ha2016}{[Ha2016]} / Theorem 1.3, Corollary 1.2(2), and Introduction. Explicit all-fields linear range, large-field quarter-degree range, and proof strategy. \href{https://arxiv.org/abs/1601.06867}{Source}.\par
\end{sources}\color{black}
\needspace{8\baselineskip}\section{Restricting every coefficient: forbidden values}
\color{black} Prescribing a few coefficients and restricting \emph{all} coefficients to large sets are different problems. Porritt treats a common forbidden set using the function-field circle method. Cheng allows the forbidden set to depend on the coefficient position.\par

\begin{atlasblock}[standard]
\textbf{\sffamily Theorem  /  Cheng, fixed field and growing degree}\par\smallskip
\color{black} Fix a prime power \(q\ge300\). For each \(0\le i<n\), let \(R_i\subseteq\mathbb F_q\) have \(|R_i|\le\sqrt q/2\), and put \(A_i=\mathbb F_q\setminus R_i\). The number of monic irreducibles \(x^n+\sum_{i<n}a_ix^i\) with \(a_i\in A_i\) satisfies\par

\atlasdisplay{\#\mathcal I=\frac{q}{n(q-1)}\prod_{i=0}^{n-1}|A_i|\left(\Delta+O\bigl(q^{-\sqrt{n/45}}\bigr)\right),\qquad \Delta=\frac{|A_0\setminus\{0\}|}{|A_0|},}
\color{black} for sufficiently large \(n\). The cited theorem gives an explicit threshold. These are large allowed sets; the statement does not permit arbitrary tiny alphabets.\par
\end{atlasblock}\color{black}
\color{black} The factor \(\Delta\) removes the constant-zero obstruction. For example, if every coefficient must be nonzero, the count is asymptotic to \(q(q-1)^{n-1}/n\). If each position must avoid a specified value \(c_i\), multiply this main term by \(1\) when \(c_0=0\), and by \((q-2)/(q-1)\) otherwise. These are asymptotic counts, not assertions for every small degree.\par
\begin{sources}
\textbf{Sources for this section}\par\smallskip
\noindent\hyperlink{ref:Porritt2019Restricted}{[Porritt2019Restricted]} / Theorem 1 and §4. Common forbidden set; q\ensuremath{\ge}500, |R|<\(\sqrt{q}\)/2, n\ensuremath{\ge}100(log q)\ensuremath{{}^2} in the displayed theorem. \href{https://doi.org/10.4153/CMB-2018-027-x}{Source}.\par
\noindent\hyperlink{ref:Cheng2025Restricted}{[Cheng2025Restricted]} / Theorem 1.1 and Corollary 1.3. Position-dependent forbidden sets, exact constant-term factor, and explicit sufficient threshold. \href{https://doi.org/10.1016/j.ffa.2025.102691}{Source}.\par
\end{sources}\color{black}
\needspace{8\baselineskip}\section{Square coefficients: a different limiting regime}

\begin{atlasblock}[standard]
\textbf{\sffamily Theorem  /  Bary-Soroker-Shmueli, published 2026}\par\smallskip
\color{black} Fix \(n\ge1\). Choose the \(n\) nonleading coefficients of a monic degree-\(n\) polynomial independently and uniformly from the \emph{set of square values} \(S_q=\{u^2:u\in\mathbb F_q\}\), including zero. Then\par

\atlasdisplay{\Pr(f\text{ is irreducible over }\mathbb F_q)=\frac1n+O_n(q^{-1/2}),\qquad q\longrightarrow\infty.}
\end{atlasblock}\color{black}
\color{black} For odd \(q\), uniform sampling from \(S_q\) differs from squaring a uniform element of \(\mathbb F_q\): the latter gives zero half the weight of each nonzero square. For even \(q\), squaring is a bijection and this is the unrestricted model.\par
\color{black} The proof uses finite-field Chebotarev and disjointness of suitable function-field extensions. The paper also handles coefficient sets of bounded algebraic complexity. Its fixed-degree error term does not establish a fixed-field, growing-degree analogue.\par

\begin{atlasblock}[standard]
\textbf{\sffamily Exact small example}\par\smallskip
\color{black} Over \(\mathbb F_5\), the square values are \(0,1,4\). The polynomial \(x^2+x+1\) has square coefficients and is irreducible because its discriminant is \(2\), a nonsquare. By contrast, \(x^2+1=(x-2)(x+2)\). Membership in the coefficient set alone is no certificate of irreducibility.\par
\end{atlasblock}\color{black}
\begin{sources}
\textbf{Sources for this section}\par\smallskip
\noindent\hyperlink{ref:BarySorokerShmueli2026Squares}{[BarySorokerShmueli2026Squares]} / Theorem 1; Theorem 2 and §§3-5 for the broader method. Monic polynomial, independent uniform square values, fixed n, q\ensuremath{\to}\ensuremath{\infty}; the O-constant depends on n. \href{https://doi.org/10.1016/j.ffa.2025.102714}{Source}.\par
\end{sources}\color{black}
\needspace{8\baselineskip}\section{Factor degrees and permutation cycles}
\color{black} For a partition \(\lambda=1^{m_1}2^{m_2}\cdots\) of \(n\), unique factorization gives the exact number of squarefree monic polynomials of type \(\lambda\):\par

\atlasdisplay{\prod_{j=1}^n\binom{N_q(j)}{m_j}\sim\frac{q^n}{z_\lambda},\qquad z_\lambda=\prod_j j^{m_j}m_j!,\qquad q\to\infty\quad(n\text{ fixed}).}
\color{black} Indeed, choose the distinct factors of each degree independently. The factor \(1/z_\lambda\) is also the probability of cycle type \(\lambda\) in \(S_n\).\par

\begin{atlasblock}[standard]
\textbf{\sffamily Theorem  /  distribution in short polynomial intervals}\par\smallskip
\color{black} Let \(f\) be monic of degree \(n\) and \(3\le m<n\). Bank-Bary-Soroker-Rosenzweig prove that among \(f+\{h:\deg h\le m\}\), the number with factorization type \(\lambda\), counting factor degrees with multiplicity, equals\par

\atlasdisplay{\frac{q^{m+1}}{z_\lambda}+O_n(q^{m+1/2}).}
\end{atlasblock}\color{black}
\color{black} The proof identifies \(S_n\) as geometric monodromy of the parameter family and translates Frobenius cycle types into point counts. This group belongs to a function-field cover; an individual finite-field splitting extension still has cyclic Galois group.\par

\begin{atlasblock}[boundary]
\textbf{\sffamily Keep the limiting regime}\par\smallskip
\color{black} The interval estimate has a constant depending on \(n\); its stated asymptotic fixes \(n\) and lets \(q\to\infty\). It does not establish the corresponding fixed-\(q\), growing-degree assertion. The unrestricted \(1/n\) probability in that regime follows separately from Möbius inversion.\par
\end{atlasblock}\color{black}
\begin{sources}
\textbf{Sources for this section}\par\smallskip
\noindent\hyperlink{ref:Cohen1970}{[Cohen1970]}. Foundational use of function-field Galois groups to study factorization distributions. \href{https://doi.org/10.4064/aa-17-3-255-271}{Source}.\par
\noindent\hyperlink{ref:BankBarySorokerRosenzweig2015}{[BankBarySorokerRosenzweig2015]} / Theorem 2.3 and §4. Uniform interval estimate for 3\ensuremath{\le}m<n and its geometric-monodromy proof; smaller m requires extra hypotheses. \href{https://arxiv.org/abs/1302.0625}{Source}.\par
\end{sources}\color{black}
\needspace{11\baselineskip}\section*{Questions and research boundaries}
\needspace{7\baselineskip}\subsection*{Can any one coefficient be prescribed?}
{\small\sffamily\color{atlasgold}Settled in stated scope / Monic degree n\ensuremath{\ge}2 polynomials over every finite field}\par\smallskip
\color{black} Yes, with exactly the constant-zero and even-characteristic quadratic exceptions stated above. Primitive-polynomial prescriptions are a different problem with additional multiplicative conditions.\par
\begin{sources}
\textbf{Sources for this section}\par\smallskip
\noindent\hyperlink{ref:TuxanidyWang2018}{[TuxanidyWang2018]} / Theorem 1.1.  \href{https://arxiv.org/abs/1604.04023}{Source}.\par
\end{sources}\color{black}
\needspace{7\baselineskip}\subsection*{Which coefficient positions allow stronger bounds?}
{\small\sffamily\color{atlasgold}Research direction / Fixed finite field, growing degree, and a specified family of prescription patterns}\par\smallskip
\color{black} Fix the positions and values before comparing results. Ask whether their structure improves the guaranteed linear proportion, yields an asymptotic count, or enables an efficient construction. The bounds above are proved ranges, not a claim that they are optimal or the latest possible for every pattern.\par
\begin{sources}
\textbf{Sources for this section}\par\smallskip
\noindent\hyperlink{ref:Ha2016}{[Ha2016]} / Theorem 1.1 and Corollary 1.2.  \href{https://arxiv.org/abs/1601.06867}{Source}.\par
\noindent\hyperlink{ref:Pollack2013Prescribed}{[Pollack2013Prescribed]} / Corollary 2.  \href{https://www.pollack-math.net/prescribed.pdf}{Source}.\par
\end{sources}\color{black}
\needspace{7\baselineskip}\subsection*{What happens when both degree and coefficient restrictions vary?}
{\small\sffamily\color{atlasgold}Research direction / A specified sequence of fields, degrees, and allowed coefficient sets}\par\smallskip
\color{black} Specify which of \(q,n\), and the sets \(A_i\) vary. Cheng's large-degree theorem assumes large allowed sets; the square-coefficient theorem fixes the degree. A useful project is to determine quantitative overlap or stronger uniformity for a chosen model. This is a comparison of proved regimes, not an assertion that every intermediate regime is open.\par
\begin{sources}
\textbf{Sources for this section}\par\smallskip
\noindent\hyperlink{ref:Cheng2025Restricted}{[Cheng2025Restricted]} / Theorem 1.1.  \href{https://doi.org/10.1016/j.ffa.2025.102691}{Source}.\par
\noindent\hyperlink{ref:BarySorokerShmueli2026Squares}{[BarySorokerShmueli2026Squares]} / Theorems 1-2.  \href{https://doi.org/10.1016/j.ffa.2025.102714}{Source}.\par
\end{sources}\color{black}
\needspace{6\baselineskip}\par\medskip\begingroup\footnotesize\setlength{\parskip}{1pt}\textbf{\sffamily\color{atlasgreen}Connections}\par
\color{black}\textbf{Uses.} \hyperlink{article:finite-field-irreducibility}{Irreducibility over finite fields}. The Frobenius-power test certifies each candidate produced by a construction.\par
\color{black}\textbf{Generalizes.} \hyperlink{article:frobenius}{Frobenius: from factor degrees to permutations}. Cycle types record factor degrees both for arithmetic reduction and for finite-field parameter families.\par
\color{black}\textbf{Applies to.} \hyperlink{article:factorization-algorithms}{Exact factorization: finite fields, Hensel lifting, and lattices}. Extension-field representations and irreducibility testing are basic inputs to factorization algorithms.\par
\color{black}\textbf{Contrasts with.} \hyperlink{article:random-models}{Random polynomials: models and proof strategies}. The finite-field coefficient model and its q-versus-degree limits differ from random integer-polynomial models.\par
\color{black}\textbf{Supplies examples for.} \hyperlink{article:specialization}{Specialization and Hilbert irreducibility}. Finite-field parameter families make the distinction between generic monodromy and specialized groups explicit.\par
\endgroup
\part{Algorithms and certificates}
\chapter{Exact factorization: finite fields, Hensel lifting, and lattices}\label{article:factorization-algorithms}\hypertarget{article:factorization-algorithms}{}\color{black}
{\sffamily\small\color{atlasgold}algorithm / Intermediate / Finite fields, Q, and p-adic fields / Published results}\par\medskip
Factorization algorithms turn structural ideas into certificates: Frobenius and linear algebra over finite fields, then p-adic lifting and lattice reduction over the rationals.\par\medskip
\textbf{Questions answered}\begin{itemize}
\item What is the mathematics behind a factorization command?
\item Why is modular factorization followed by recombination?
\item How can a randomized computation still produce exact factors?
\item How do recent modular-composition advances relate to factorization?
\item How are local factors found when reduction has repeated factors?
\end{itemize}
{\small\textbf{Prerequisites:} \hyperlink{article:modular-irreducibility}{Irreducibility by reduction modulo primes}; \hyperlink{article:finite-field-irreducibility}{Irreducibility over finite fields}}\par
\needspace{8\baselineskip}\section{The finite-field pipeline}
\color{black} Start with a polynomial over \(\mathbb F_q\). A squarefree decomposition separates repeated factors, using gcds with the derivative. If the derivative vanishes, the polynomial is a \(p\)-th power over the finite field and one first takes its \(p\)-th root.\par
\color{black} For a squarefree polynomial, \(\gcd(f,x^{q^d}-x)\) selects its irreducible factors whose degrees divide \(d\). Repeating this with the factors already found removed gives a distinct-degree decomposition. Equal-degree splitting then separates factors of the same degree. Cantor-Zassenhaus uses random choices for this splitting step; the output consists of exact polynomial factors whose product and irreducibility can be checked.\par
\begin{sources}
\textbf{Sources for this section}\par\smallskip
\noindent\hyperlink{ref:CantorZassenhaus1981}{[CantorZassenhaus1981]}.  \href{https://doi.org/10.1090/S0025-5718-1981-0606517-5}{Source}.\par
\noindent\hyperlink{ref:GathenGerhard2013}{[GathenGerhard2013]}. Squarefree, distinct-degree, and equal-degree factorization. \href{https://doi.org/10.1017/CBO9781139856065}{Source}.\par
\end{sources}\color{black}
\needspace{8\baselineskip}\section{Berlekamp: factors as Frobenius-fixed functions}

\begin{atlasblock}[standard]
\textbf{\sffamily The structural observation}\par\smallskip
\color{black} If \(f\) is squarefree with \(r\) distinct irreducible factors, then\par

\atlasdisplay{\dim_{\mathbb F_q}\{h\in\mathbb F_q[x]/(f):h^q=h\}=r.}
\end{atlasblock}\color{black}
\color{black} Indeed, the Chinese remainder theorem identifies the quotient with a product \(\prod_{i=1}^r\mathbb F_{q^{d_i}}\). The fixed elements of Frobenius on each component are exactly \(\mathbb F_q\). Finding the fixed subspace is linear algebra. A fixed element assumes constant values on the components; gcds \(\gcd(f,h-c)\), for \(c\in\mathbb F_q\), separate components with different values.\par

\begin{atlasblock}[standard]
\textbf{\sffamily Example  /  a three-dimensional fixed subspace}\par\smallskip
\color{black} Over \(\mathbb F_2\), take \(f=x^4+x\). Since \(f'=1\), it is squarefree. In the quotient, \(x^4=x\), and solving \(h^2=h\) gives the basis \(1,x^2+x,x^3\). Thus there are three irreducible factors. The gcds from \(h=x^3\) first split off \(x\); those from \(h=x^2+x\) split the remaining cubic. The result is\par

\atlasdisplay{x^4+x=x(x+1)(x^2+x+1).}
\color{black} The quadratic has no root in \(\mathbb F_2\), so this is an independently checkable complete factorization.\par
\end{atlasblock}\color{black}
\begin{sources}
\textbf{Sources for this section}\par\smallskip
\noindent\hyperlink{ref:Berlekamp1967}{[Berlekamp1967]}. Frobenius-fixed algebra and the linear-algebra factorization method. \href{https://doi.org/10.1002/j.1538-7305.1967.tb03174.x}{Source}.\par
\end{sources}\color{black}
\needspace{8\baselineskip}\section{Modular composition: applying Frobenius efficiently}
\color{black} The operation \emph{modular composition} takes \(a,g,f\in K[x]\), with \(f\) monic of degree \(n\) and \(\deg a,\deg g<n\), and computes \(g(a(x))\bmod f(x)\). It avoids storing a composition whose unreduced degree can be nearly \(n^2\).\par

\begin{atlasblock}[standard]
\textbf{\sffamily Why this operation belongs in factorization}\par\smallskip
\color{black} In \(A=\mathbb F_q[x]/(f)\), put \(a=x^q\bmod f\). For every \(h\in\mathbb F_q[x]\),\par

\atlasdisplay{h(x)^q\equiv h(a(x))\pmod f.}
\color{black} The coefficients are fixed by the \(q\)-power map. Thus one application of Frobenius is a modular composition. Iterates of this map supply the powers used in distinct-degree factorization and irreducibility tests.\par
\end{atlasblock}\color{black}

\begin{atlasblock}[standard]
\textbf{\sffamily Example  /  Frobenius without a large unreduced power}\par\smallskip
\color{black} Over \(\mathbb F_2\), take \(f=x^3+x+1\). Since \(x^3=x+1\) in the quotient, applying Frobenius to \(h=x^2+x\) gives\par

\atlasdisplay{h^2=x^4+x^2=(x^2+x)+x^2=x.}
\color{black} Equivalently, compose \(h\) with \(a=x^2\) and reduce modulo \(f\). This identity is exact; a complexity theorem concerns how efficiently such identities can be computed in large degree.\par
\end{atlasblock}\color{black}
\color{black} Brent-Kung's baby-step/giant-step method groups powers of \(a\) into blocks and combines them by matrix multiplication. Later methods encode small algebraic relations between \(x\) and \(a\) in polynomial matrices, or reduce the task to multivariate evaluation. These are different ways of avoiding the full intermediate composition.\par
\begin{sources}
\textbf{Sources for this section}\par\smallskip
\noindent\hyperlink{ref:BrentKung1978}{[BrentKung1978]} / Composition algorithms; see also Neiger-Salvy-Schost-Villard (2024), Section 1.1, for modular composition.  \href{https://doi.org/10.1145/322092.322099}{Source}.\par
\noindent\hyperlink{ref:KedlayaUmans2011}{[KedlayaUmans2011]} / Sections 3 and 7.  \href{https://doi.org/10.1137/08073408X}{Source}.\par
\noindent\hyperlink{ref:NeigerSalvySchostVillard2024}{[NeigerSalvySchostVillard2024]} / Section 1.1 and Theorem 1.1.  \href{https://doi.org/10.1145/3638349}{Source}.\par
\end{sources}\color{black}
\needspace{8\baselineskip}\section{Recent progress: compare the same computational problem and model}
\color{black} A \emph{field-operation count} charges unit cost for arithmetic in \(K\). A \emph{bit-operation count} includes the representation and size of coefficients. Over \(\mathbb F_q\), these are not interchangeable uniformly in \(q\). Also distinguish modular composition from the entire factorization pipeline.\par
\begin{center}\small\begin{tabularx}{\linewidth}{@{}>{\raggedright\arraybackslash}X>{\raggedright\arraybackslash}X>{\raggedright\arraybackslash}X@{}}\toprule
\textbf{Source} & \textbf{Precise scope} & \textbf{Contribution} \\
\midrule
Kedlaya-Umans (2011) & Modular composition over \(\mathbb F_q\); bit complexity & For each fixed \(\varepsilon>0\), \(n^{1+\varepsilon}\log^{1+o(1)}q\) bit operations in sufficiently large degree. This leads to a randomized finite-field factorization bound, displayed below. \\
Neiger-Salvy-Schost-Villard (2024) & Arbitrary fields of known characteristic; arithmetic complexity & A Las Vegas method with expected \(O(n^{1.43})\) field operations. Small finite fields are handled through a suitably enlarged extension; a failed randomized trial is detected. \\
Neiger-Salvy-Schost-Villard (2026) & Fixed \(f\), generic inner polynomial \(a\); arithmetic complexity & Two relation matrices give \(\widetilde O(n^{(\omega+3)/4})\) field operations. Here \(\omega\) is a feasible square-matrix multiplication exponent. The genericity requirement is essential. \\
van der Hoeven-Lecerf (2026) & Modular composition over \(\mathbb Z/r\mathbb Z\); bit complexity & A sharper nearly linear bound via spiroids and Fourier transforms with proved precision control. Prime finite fields are a special case. \\
\bottomrule\end{tabularx}\end{center}
\color{black} For orientation, the Kedlaya-Umans randomized \emph{factorization} bound is\par

\atlasdisplay{O\!\left(n^{3/2+o(1)}\log^{1+o(1)}q+n^{1+o(1)}\log^{2+o(1)}q\right).}
\color{black} A faster modular-composition subroutine does not by itself turn this exponent into \(1\). Guo-Narayanan-Umans (2016) organize related algebraic problems through reductions that would transfer an improvement below the factorization exponent \(3/2\). These reductions describe a research target, not a lower-bound theorem forbidding faster algorithms.\par
\color{black} \textbf{Conditional direction.} Umans-Wang (2025) propose structured sets of integers with a divisibility-covering property as a route to faster factorization. Their proposed improvement depends on a number-theoretic conjecture; it is not an unconditional algorithm beating \(3/2\). The paper was listed for ISAAC in December 2026 when checked on 12 November; the reference below identifies the available preprint.\par
\begin{sources}
\textbf{Sources for this section}\par\smallskip
\noindent\hyperlink{ref:KedlayaUmans2011}{[KedlayaUmans2011]} / Corollary 7.2 and finite-field factorization results.  \href{https://doi.org/10.1137/08073408X}{Source}.\par
\noindent\hyperlink{ref:NeigerSalvySchostVillard2024}{[NeigerSalvySchostVillard2024]} / Theorem 1.1 and the following discussion of sampling and small finite fields.  \href{https://doi.org/10.1145/3638349}{Source}.\par
\noindent\hyperlink{ref:NeigerSalvySchostVillard2026}{[NeigerSalvySchostVillard2026]} / Theorem 1.3 and the paragraph immediately following it.  \href{https://doi.org/10.1145/3815436.3815451}{Source}.\par
\noindent\hyperlink{ref:HoevenLecerf2026}{[HoevenLecerf2026]} / Theorem 1 and Sections 1.3-1.4; author version dated 12 August 2026.  \href{https://doi.org/10.1016/j.jco.2026.102073}{Source}.\par
\noindent\hyperlink{ref:GuoNarayananUmans2016}{[GuoNarayananUmans2016]} / Introduction and reduction theorems.  \href{https://doi.org/10.4230/LIPIcs.MFCS.2016.47}{Source}.\par
\noindent\hyperlink{ref:UmansWang2025}{[UmansWang2025]} / Abstract and introduction; conditional strategy.  \href{https://arxiv.org/abs/2511.10851}{Source}.\par
\end{sources}\color{black}
\needspace{8\baselineskip}\section{Lifting a modular factorization}
\color{black} For clarity, let \(f\in\mathbb Z[x]\) be monic and squarefree. Choose a prime where its reduction remains squarefree, factor there, and lift the pairwise coprime factors modulo \(p,p^2,p^4,\ldots\) by Hensel's lemma. Coprimality permits the correction equations to be solved at each stage.\par
\color{black} A rational factor can combine several of the lifted factors. The task is therefore not merely to lift but to discover which combinations descend to \(\mathbb Z[x]\). Bounds for the coefficients of a genuine factor turn sufficiently accurate modular coefficients into unique candidate integers. Exact division then accepts or rejects a candidate.\par
\color{black} This explains why ``reducible modulo \(p\)'' gives no rational factor by itself: the finite-field pieces need not descend separately. Conversely, when one good reduction is irreducible, the rational computation can stop immediately.\par
\begin{sources}
\textbf{Sources for this section}\par\smallskip
\noindent\hyperlink{ref:Zassenhaus1969}{[Zassenhaus1969]}.  \href{https://doi.org/10.1016/0022-314X(69)90047-X}{Source}.\par
\noindent\hyperlink{ref:GathenGerhard2013}{[GathenGerhard2013]}. Hensel lifting, coefficient bounds, and modular factorization. \href{https://doi.org/10.1017/CBO9781139856065}{Source}.\par
\end{sources}\color{black}
\needspace{8\baselineskip}\section{Local factorization beyond coprime reduction}
\color{black} Factoring over \(\mathbb Q_p\) is itself a complete task, even when every factor of \(\bar f\) is repeated. A residual polynomial separates roots that share a Newton slope; higher Newton polygons continue when that residual information is insufficient. The Montes method organizes successive slopes, residue-field extensions, and polynomial representatives into \emph{types}. Complete types identify irreducible local factors and their ramification indices and residue degrees.\par
\color{black} Guàrdia-Montes-Nart give the underlying theory and an algorithm for prime ideal decomposition and discriminants. Poteaux-Weimann's ISSAC 2022 paper improves local factorization through Hensel lifting adapted to these valuations and a divide-and-conquer strategy. Its sharper complexity bound assumes residue characteristic zero or greater than the polynomial degree; residual-field factorization costs are accounted for separately.\par
\color{black} A local factor is represented to a stated \(p\)-adic precision together with exact structural data; its coefficients need not be rational. Recovering rational factors still requires global recombination and certification.\par
\begin{sources}
\textbf{Sources for this section}\par\smallskip
\noindent\hyperlink{ref:GuardiaMontesNart2012Newton}{[GuardiaMontesNart2012Newton]} / §§2-4. Higher-order polygon and residual-polynomial theory. \href{https://doi.org/10.1090/S0002-9947-2011-05442-5}{Source}.\par
\noindent\hyperlink{ref:GuardiaMontesNart2011Montes}{[GuardiaMontesNart2011Montes]} / §2.2, Theorem 2.2; §3. Complete types and the optimized Montes method. \href{https://www.numdam.org/articles/10.5802/jtnb.782/}{Source}.\par
\noindent\hyperlink{ref:PoteauxWeimann2022Local}{[PoteauxWeimann2022Local]} / Assumption 1; §§4-5, Theorems 3-4. Published algorithmic improvements; assumptions and residual-field costs remain explicit. \href{https://doi.org/10.1145/3476446.3535487}{Source}.\par
\end{sources}\color{black}
\needspace{8\baselineskip}\section{Why lattice reduction changes the recombination step}
\color{black} Trying every subset of lifted factors can be exponentially expensive in their number. The Lenstra-Lenstra-Lovász theorem gives a deterministic algorithm for factoring rational polynomials in time polynomial in their degree and coefficient bit size, using lattice reduction. The lattice encodes congruence information; short vectors reveal relations compatible with a genuine factor.\par
\color{black} Van Hoeij recasts practical recombination through a structured knapsack problem. Modern variants, including FLINT's implementation, use logarithmic derivatives to turn products into sums, so a prospective factor corresponds to a vector of choices among local factors. Lattice reduction narrows the possible choice vectors before expensive exact tests. Belabas-van Hoeij-Klüners-Steel (2009) proved polynomial-time complexity for van Hoeij's original rational-polynomial algorithm and natural variants. This recombination method uses local information to find exact global factors.\par
\begin{sources}
\textbf{Sources for this section}\par\smallskip
\noindent\hyperlink{ref:LLL1982}{[LLL1982]}. Polynomial-time factorization over Q. \href{https://doi.org/10.1007/BF01457454}{Source}.\par
\noindent\hyperlink{ref:VanHoeij2002}{[VanHoeij2002]}. Recombination by knapsack and lattice methods. \href{https://doi.org/10.1006/jnth.2001.2763}{Source}.\par
\noindent\hyperlink{ref:BHKS2009}{[BHKS2009]}. Polynomial-time analysis of van Hoeij's original rational-polynomial factorization algorithm and natural variants. \href{https://www.numdam.org/articles/10.5802/jtnb.655/}{Source}.\par
\noindent\hyperlink{ref:FLINTFactorizationDocumentation}{[FLINTFactorizationDocumentation]} / Factoring polynomials over the integers: Van Hoeij's algorithm and Complexity. Implementation-oriented explanation of modular factorization. \href{https://flintlib.org/doc/algorithms/polynomials/factoring-finite-fields.html}{Source}.\par
\end{sources}\color{black}
\needspace{8\baselineskip}\section{What a complete factorization certificate contains}
\color{black} The identity \(f=c\prod_i g_i^{e_i}\) proves a factorization, but not its completeness. Completeness additionally requires proof that each \(g_i\) is irreducible over the stated field. A modular irreducibility witness, a proved algorithm, or a suitable exact criterion can supply that step. Record the coefficient field, normalization, multiplicities, and the irreducibility method along with the factors.\par
\color{black} Randomized search and uncertain conclusions are different issues. Random choices may affect runtime while every returned factor and certificate is checked exactly. Conversely, a deterministic-looking numerical procedure can still require an explicit certification step. This distinction is essential when computational output becomes evidence in a theorem about a family.\par
\begin{sources}
\textbf{Sources for this section}\par\smallskip
\noindent\hyperlink{ref:GathenGerhard2013}{[GathenGerhard2013]}.  \href{https://doi.org/10.1017/CBO9781139856065}{Source}.\par
\noindent\hyperlink{ref:Rabin1980}{[Rabin1980]}. Finite-field algorithms and exact irreducibility tests. \href{https://doi.org/10.1137/0209024}{Source}.\par
\end{sources}\color{black}
\needspace{11\baselineskip}\section*{Questions and research boundaries}
\needspace{7\baselineskip}\subsection*{Is irreducibility over Q algorithmically decidable in polynomial time?}
{\small\sffamily\color{atlasgold}Settled in stated scope / Dense univariate polynomials given by rational coefficients in binary}\par\smallskip
\color{black} Yes: LLL gives polynomial-time factorization in degree and coefficient bit size. This is a complexity theorem for this input model; it does not imply equally efficient factorization for every sparse encoding or every extension field representation.\par
\begin{sources}
\textbf{Sources for this section}\par\smallskip
\noindent\hyperlink{ref:LLL1982}{[LLL1982]}.  \href{https://doi.org/10.1007/BF01457454}{Source}.\par
\end{sources}\color{black}
\needspace{7\baselineskip}\subsection*{Can the 2026 generic composition bound be extended to arbitrary input?}
{\small\sffamily\color{atlasgold}Research direction / Algebraic modular composition over a field; the two-relation-matrix method}\par\smallskip
\color{black} For each degree-\(n\) modulus \(f\), Theorem 1.3 supplies a nonempty Zariski-open set of coefficient vectors for \(a\) on which the faster bound holds. The algorithm can detect failure of its genericity conditions and fall back on another method. The authors explicitly leave a Las Vegas reduction for arbitrary input to future work. A nonempty Zariski-open set over an algebraic closure need not contain a point over a small finite base field; ``generic'' must not be silently replaced by ``every.''\par
\begin{sources}
\textbf{Sources for this section}\par\smallskip
\noindent\hyperlink{ref:NeigerSalvySchostVillard2026}{[NeigerSalvySchostVillard2026]} / Theorem 1.3 and following paragraph.  \href{https://doi.org/10.1145/3815436.3815451}{Source}.\par
\end{sources}\color{black}
\needspace{6\baselineskip}\par\medskip\begingroup\footnotesize\setlength{\parskip}{1pt}\textbf{\sffamily\color{atlasgreen}Connections}\par
\color{black}\textbf{Uses.} \hyperlink{article:finite-field-irreducibility}{Irreducibility over finite fields}. Frobenius tests certify the finite-field factors.\par
\color{black}\textbf{Generalizes.} \hyperlink{article:modular-irreducibility}{Irreducibility by reduction modulo primes}. One-prime criteria are early exits within a broader exact algorithm.\par
\color{black}\textbf{Applies to.} \hyperlink{article:galois-computation}{Computing a Galois group with a proof certificate}. A group computation begins with exact factorization and separability checks.\par
\color{black}\textbf{Uses.} \hyperlink{article:newton-polygons}{Newton polygons, factor degrees, and large cycles}. Residual polynomials and higher Newton polygons refine valuations into local irreducibility and factorization data.\par
\endgroup
\chapter{Computing a Galois group with a proof certificate}\label{article:galois-computation}\hypertarget{article:galois-computation}{}\color{black}
{\sffamily\small\color{atlasgold}algorithm / Intermediate / Q, number fields, and p-adic fields / Published results}\par\medskip
A reliable computation combines exact arithmetic, subgroup constraints, and certification. Save enough evidence that the conclusion survives a change of software.\par\medskip
\textbf{Questions answered}\begin{itemize}
\item What should be stored with a computed Galois group?
\item How do resolvents and modular witnesses work together?
\item Which official software limitations and proof settings matter?
\end{itemize}
{\small\textbf{Prerequisites:} \hyperlink{article:frobenius}{Frobenius: from factor degrees to permutations}; \hyperlink{article:resolvents}{Resolvents: turn group containment into factorization}; \hyperlink{article:factorization-algorithms}{Exact factorization: finite fields, Hensel lifting, and lattices}}\par
\needspace{8\baselineskip}\section{Specify the mathematical object first}
\color{black} The input is a polynomial and a ground field. For \(f\in\mathbb Q[x]\), clear denominators, remove content, and compute a squarefree factorization. A transitive-group label describes the action for an irreducible factor. If \(f\) is reducible, its whole splitting field is a compositum; the group need not be the direct product of the groups of the factors because their splitting fields can intersect.\par
\color{black} For an irreducible degree-\(n\) input, the useful output is the subgroup of \(S_n\) up to conjugacy, together with its action on the roots. An abstract isomorphism type alone may lose information needed for cycle types, resolvents, subfields, or iteration.\par
\begin{sources}
\textbf{Sources for this section}\par\smallskip
\noindent\hyperlink{ref:Cox2012}{[Cox2012]}.  \href{https://onlinelibrary.wiley.com/doi/book/10.1002/9781118218457}{Source}.\par
\noindent\hyperlink{ref:FiekerKlueners2014}{[FiekerKlueners2014]} / Introduction and Section 2.  \href{https://doi.org/10.1112/S1461157013000302}{Source}.\par
\end{sources}\color{black}
\needspace{8\baselineskip}\section{A practical proof workflow}

\begin{enumerate}
\item\leavevmode\color{black} \textbf{Establish irreducibility and separability.} Keep the factorization or irreducibility certificate, rather than inferring it from an expected group.
\item\leavevmode\color{black} \textbf{Compute the discriminant and structural restrictions.} Reciprocity, composition, and known intermediate fields can reduce the ambient permutation group.
\item\leavevmode\color{black} \textbf{Gather exact local witnesses.} Record good primes, the complete finite-field factorization, and the resulting cycle type. Verify the factors and their irreducibility.
\item\leavevmode\color{black} \textbf{Eliminate candidates.} Use primitive-group theorems when applicable; otherwise use relative resolvents and subgroup descent.
\item\leavevmode\color{black} \textbf{Certify the final descent.} Save the successful proof status and the exact data needed to check the decisive tests.
\end{enumerate}
\color{black} Stauduhar's method makes the fourth step systematic. Starting from a known overgroup, it tests relative invariants for maximal subgroups and descends whenever a certified resolvent condition permits it.\par
\begin{sources}
\textbf{Sources for this section}\par\smallskip
\noindent\hyperlink{ref:Stauduhar1973}{[Stauduhar1973]}.  \href{https://doi.org/10.1090/S0025-5718-1973-0327712-4}{Source}.\par
\noindent\hyperlink{ref:FiekerKlueners2014}{[FiekerKlueners2014]} / Sections 3 and 7.  \href{https://doi.org/10.1112/S1461157013000302}{Source}.\par
\end{sources}\color{black}
\needspace{8\baselineskip}\section{A decision problem can be easier than constructing the splitting field}
\color{black} Landau-Miller (1985) give a polynomial-time decision procedure for solvability by radicals over \(\mathbb Q\), in the dense binary input model. Their paper also constructs radical expressions in a compact straight-line-program representation. The encoding is part of the theorem: expanding every repeated subexpression can destroy the size bound.\par
\color{black} The method exploits blocks of imprimitivity and the restricted size of primitive solvable permutation groups. Thus one can answer a structural question about a Galois group without first listing all its elements. This is different from constructing the splitting field in an explicit basis: its degree may be \(n!\) for a degree-\(n\) polynomial with group \(S_n\).\par
\begin{sources}
\textbf{Sources for this section}\par\smallskip
\noindent\hyperlink{ref:LandauMiller1985}{[LandauMiller1985]} / Introduction and Sections 2-4.  \href{https://doi.org/10.1016/0022-0000(85)90013-3}{Source}.\par
\end{sources}\color{black}
\needspace{8\baselineskip}\section{Read the proof contract of the software}
\begin{center}\small\begin{tabularx}{\linewidth}{@{}>{\raggedright\arraybackslash}X>{\raggedright\arraybackslash}X>{\raggedright\arraybackslash}X@{}}\toprule
\textbf{Tool} & \textbf{Relevant contract} & \textbf{What to retain} \\
\midrule
Magma & Native \texttt{GaloisGroup} accepts integer, rational, and number-field polynomials without a fixed degree cap, but does not automatically prove its conditional steps. The documented \texttt{GaloisProof(f,S)} contract covers integer polynomials; an absolute number-field overload is also available. & Input, options, \texttt{GaloisData}, successful proof result, and version. \\
PARI/GP 2.19 & \texttt{polgalois} requires irreducible input over \(\mathbb Q\), of degree at most seven, or eleven with \texttt{galdata}. Its resolvent method can rarely return a wrong answer when initial numerical precision is insufficient. & Polynomial, precision, complete output, numbering convention, and version. \\
SageMath 10.10 & The rational-polynomial \texttt{galois\_\allowbreak{}group} interface requires irreducible input. It uses PARI through degree eleven and attempts GAP above that; GAP's documented algorithm range is through degree fifteen, with some higher cases possible. The Magma wrapper identifies groups only through degree thirty. & Routine, algorithm, backend version, permutation action, and separate certification outcome. \\
\bottomrule\end{tabularx}\end{center}
\color{black} Native Magma's degree-independent algorithm and Sage's wrapper have different limits. A database containing transitive groups of a given degree does not guarantee that the group-computation algorithm handles every polynomial of that degree. Runtime and memory can be decisive even without a formal degree cap.\par
\color{black} For Sage number fields, \texttt{gc\_\allowbreak{}numbering} controls whether permutations act on roots of the original defining polynomial or of the Galois closure's defining polynomial. The documentation also warns that calling \texttt{galois\_\allowbreak{}group} on a relative extension currently computes its absolute Galois closure over \(\mathbb Q\); specify the intended ground field explicitly.\par
\color{black} Sage also documents \texttt{galois\_\allowbreak{}group\_\allowbreak{}davenport\_\allowbreak{}smith\_\allowbreak{}test}: return values \(1\) and \(2\) certify \(S_n\) and \(A_n\), while \(0\) means inconclusive. Its default checks irreducibility; setting \texttt{assume\_\allowbreak{}irreducible=True} requires an independently established hypothesis. This specialized test is different from unrestricted group determination.\par
\begin{sources}
\textbf{Sources for this section}\par\smallskip
\noindent\hyperlink{ref:MagmaGaloisDocumentation}{[MagmaGaloisDocumentation]} / GaloisGroup and GaloisProof. Input fields, degree independence, conditional output, and documented certification scope. \href{https://docs.magma-maths.org/GlobalFields/GaloisGroupsAndAutomorphisms/galois.html}{Source}.\par
\noindent\hyperlink{ref:PARIDocumentation}{[PARIDocumentation]} / polgalois, PARI/GP 2.19.0. Degree limits, output conventions, and explicit precision warning. \href{https://pari.math.u-bordeaux.fr/dochtml/ref/General_number_fields.html}{Source}.\par
\noindent\hyperlink{ref:SageGaloisDocumentation}{[SageGaloisDocumentation]} / Polynomial\_\allowbreak{}rational\_\allowbreak{}flint.galois\_\allowbreak{}group; galois\_\allowbreak{}group\_\allowbreak{}davenport\_\allowbreak{}smith\_\allowbreak{}test; NumberField.galois\_\allowbreak{}group, SageMath 10.10. Backend limits, certified specialized test, root-action conventions, and relative-field warning. \href{https://doc.sagemath.org/html/en/reference/polynomial_rings/sage/rings/polynomial/polynomial_rational_flint.html}{Source}.\par
\end{sources}\color{black}
\needspace{8\baselineskip}\section{A portable quartic certificate}

\begin{atlasblock}[standard]
\textbf{\sffamily Example  /  exact arithmetic first}\par\smallskip
\color{black} For \(f=x^4+2x+2\), retain three pieces of data:\par

\begin{enumerate}
\item\leavevmode\color{black} Eisenstein at \(2\), proving irreducibility.
\item\leavevmode\color{black} The resolvent \(R=T^3-8T-4\), with no rational root among \(\pm1,\pm2,\pm4\).
\item\leavevmode\color{black} The discriminant \(\Delta(f)=1616=16\cdot101\), a nonsquare.
\end{enumerate}
\color{black} The quartic classification proves the group is \(S_4\). The proof can be checked independently of the program that first suggested the group. The companion exact-arithmetic script verifies these identities, the root tests, and additional small examples.\par
\end{atlasblock}\color{black}
\color{black} A corresponding Magma command sequence is:\par

\begin{quote}\small\ttfamily\raggedright Zx<x> := PolynomialRing(Integers());\par 
f := x\textasciicircum{}4 + 2*x + 2;\par 
G, roots, S := GaloisGroup(f);\par 
assert GaloisProof(f, S);\par 
Order(G);\end{quote}
\color{black} This is a reproducible command template, not a claimed Magma execution from this atlas. The algebraic proof above establishes \(|G|=24\) without relying on that execution.\par
\begin{sources}
\textbf{Sources for this section}\par\smallskip
\noindent\hyperlink{ref:ConradCubicQuartic}{[ConradCubicQuartic]} / Table 5 and Theorem 3.6.  \href{https://kconrad.math.uconn.edu/blurbs/galoistheory/cubicquartic.pdf}{Source}.\par
\noindent\hyperlink{ref:MagmaGaloisDocumentation}{[MagmaGaloisDocumentation]} / GaloisGroup and GaloisProof.  \href{https://docs.magma-maths.org/GlobalFields/GaloisGroupsAndAutomorphisms/galois.html}{Source}.\par
\end{sources}\color{black}
\needspace{8\baselineskip}\section{Local Galois groups require a local computation}
\color{black} For a rational polynomial, computing its group over \(\mathbb Q_p\) asks a different question from computing its group over \(\mathbb Q\). If \(L/\mathbb Q\) is the global splitting field and \(\mathfrak P\mid p\), the local splitting-field group identifies with the decomposition subgroup \(D_{\mathfrak P}\subseteq\operatorname{Gal}(L/\mathbb Q)\). The embedding depends on the choice of \(\mathfrak P\), up to conjugacy.\par
\color{black} Doris (2020) develops local algorithms using ramification towers, global models, and resolvents. A global model supplies a tractable overgroup; local factorization of resolvents distinguishes the desired subgroup. The base field therefore remains part of the result even when global auxiliary calculations are used.\par
\color{black} Certification must also follow the actual implementation. Remark 3.5 of the paper identifies a heuristic complex-root bound in the then-current implementation. The author's package documentation later describes an exact \(p\)-adic input mode that produces proved results. Record the mode and verification outcome; the existence of a mathematical algorithm alone does not certify a particular finite-precision run.\par
\begin{sources}
\textbf{Sources for this section}\par\smallskip
\noindent\hyperlink{ref:Neukirch1999}{[Neukirch1999]} / Decomposition groups and completions.  \href{https://link.springer.com/book/10.1007/978-3-662-03983-0}{Source}.\par
\noindent\hyperlink{ref:Doris2020}{[Doris2020]} / Sections 1.1-4 and Remark 3.5; linked author package documentation for the exact-input mode.  \href{https://arxiv.org/abs/2003.05834}{Source}.\par
\end{sources}\color{black}
\needspace{8\baselineskip}\section{From one calculation to a family}
\color{black} For a parameterized family, save the parameter range, normalization, exceptions, and every input polynomial. A table of successful cases is evidence for a conjecture; it does not establish the same conclusion for all parameters. Separate a proved general result from a finite verified range and from a numerical experiment.\par
\color{black} Factor degrees usually identify the Frobenius cycle type rather than its precise conjugacy class in the group. Dokchitser and Dokchitser give finer invariants for distinguishing classes when cycle type is insufficient. This matters when arithmetic applications require more than the abstract group or its order.\par
\color{black} A useful downloadable record includes exact coefficients, ground field, source theorem or algorithm, witness primes, discriminant, group action, certification outcome, and the date of the run. Databases supply comparison examples; the original computation or theorem supplies the proof.\par
\begin{sources}
\textbf{Sources for this section}\par\smallskip
\noindent\hyperlink{ref:DokchitserDokchitser2013}{[DokchitserDokchitser2013]}. Distinguishing Frobenius classes beyond their cycle types. \href{https://doi.org/10.2140/ant.2013.7.1325}{Source}.\par
\noindent\hyperlink{ref:LMFDBFieldsGroups}{[LMFDBFieldsGroups]}. Comparison data for number fields and transitive groups. \href{https://www.lmfdb.org/NumberField/}{Source}.\par
\end{sources}\color{black}
\needspace{11\baselineskip}\section*{Questions and research boundaries}
\needspace{7\baselineskip}\subsection*{What turns a plausible computed Galois group into a proved result?}
{\small\sffamily\color{atlasgold}Research direction / Reproducible computations for explicit polynomial families}\par\smallskip
\color{black} Design short certificates for the decisive containment and noncontainment tests, and retain their verifiers. The quartic example already has such a certificate. Extending this style to a chosen family is a concrete project; it is not a claim that certification of Galois groups is itself unsolved.\par
\begin{sources}
\textbf{Sources for this section}\par\smallskip
\noindent\hyperlink{ref:FiekerKlueners2014}{[FiekerKlueners2014]} / Section 7.  \href{https://doi.org/10.1112/S1461157013000302}{Source}.\par
\noindent\hyperlink{ref:MagmaGaloisDocumentation}{[MagmaGaloisDocumentation]} / GaloisProof.  \href{https://docs.magma-maths.org/GlobalFields/GaloisGroupsAndAutomorphisms/galois.html}{Source}.\par
\end{sources}\color{black}
\needspace{6\baselineskip}\par\medskip\begingroup\footnotesize\setlength{\parskip}{1pt}\textbf{\sffamily\color{atlasgreen}Connections}\par
\color{black}\textbf{Uses.} \hyperlink{article:resolvents}{Resolvents: turn group containment into factorization}. Relative invariants certify containment in smaller subgroups.\par
\color{black}\textbf{Uses.} \hyperlink{article:frobenius}{Frobenius: from factor degrees to permutations}. Small-prime witnesses eliminate candidate groups inexpensively.\par
\color{black}\textbf{Uses.} \hyperlink{article:factorization-algorithms}{Exact factorization: finite fields, Hensel lifting, and lattices}. Exact factors and irreducibility checks establish the input assumptions.\par
\color{black}\textbf{Applies to.} \hyperlink{article:eulerian}{Eulerian polynomials: descents, shifts, and inverse pairs}. Family tables must separate finite certified calculations from a theorem in every degree.\par
\endgroup
\part{Several variables}
\chapter{Absolute irreducibility and Newton polytopes}\label{article:multivariate-irreducibility}\hypertarget{article:multivariate-irreducibility}{}\color{black}
{\sffamily\small\color{atlasgold}method / Intermediate / arbitrary fields / Published results}\par\medskip
Use the shape of a multivariate support to forbid factors over every algebraic extension, then understand the limits of a coefficient-free test.\par\medskip
\textbf{Questions answered}\begin{itemize}
\item How does absolute irreducibility differ from irreducibility over the ground field?
\item When does the Newton polytope certify absolute irreducibility?
\item Why is a decomposable polytope only an inconclusive result?
\end{itemize}
{\small\textbf{Prerequisites:} \hyperlink{article:foundations}{Fields, factors, and the action on roots}}\par
\needspace{8\baselineskip}\section{Why the multivariate question is different}
\color{black} A nonconstant \(f\in K[x_1,\ldots,x_r]\) is \emph{absolutely irreducible} if it remains irreducible over \(\overline K\). For a univariate polynomial this forces degree one. With several variables it is a rich condition: the coordinate ring of the hypersurface defined by \(f\) stays a domain after extending to \(\overline K\); the hypersurface is geometrically integral.\par
\color{black} For example, \(x^2+y^2\) is irreducible over \(\mathbb Q\) but factors over \(\mathbb Q(i)\) as \((x+iy)(x-iy)\). A certificate of rational irreducibility therefore does not automatically certify absolute irreducibility. One must also keep irreducibility of a polynomial separate from smoothness of its zero set.\par
\begin{sources}
\textbf{Sources for this section}\par\smallskip
\noindent\hyperlink{ref:Gao2001Polytopes}{[Gao2001Polytopes]}. Absolute irreducibility of multivariate polynomials. \href{https://www.sciencedirect.com/science/article/pii/S0021869300985867}{Source}.\par
\noindent\hyperlink{ref:Lecerf2013}{[Lecerf2013]}. Rational and absolute factorization in several variables. \href{https://www.numdam.org/articles/10.5802/ccirm.18/}{Source}.\par
\end{sources}\color{black}
\needspace{8\baselineskip}\section{Support, Minkowski sums, and factors}
\color{black} For \(f=\sum_{\mathbf a}c_{\mathbf a}x^{\mathbf a}\), its Newton polytope is\par

\atlasdisplay{\operatorname{Newt}(f)=\operatorname{conv}\{\mathbf a\in\mathbb Z_{\ge0}^{r}:c_{\mathbf a}\ne0\}.}
\color{black} A product obeys \(\operatorname{Newt}(gh)=\operatorname{Newt}(g)+\operatorname{Newt}(h)\), where the right side is the Minkowski sum. To see why cancellation cannot destroy a vertex, choose a linear weight uniquely exposing the corresponding vertices of the two factor polytopes. The product of the two highest-weight terms is then nonzero and uniquely highest. Varying the weight determines the entire polytope.\par
\color{black} A lattice polytope is integrally indecomposable if every expression as a sum of two lattice polytopes has a point as one summand. A point polytope corresponds to a monomial, not necessarily to a constant.\par
\begin{sources}
\textbf{Sources for this section}\par\smallskip
\noindent\hyperlink{ref:GaoLauder2001}{[GaoLauder2001]} / §2, Lemma 1 and Corollary 2. Product formula, with its origin attributed to Ostrowski, and the monomial caveat. \href{https://link.springer.com/article/10.1007/s00454-001-0024-0}{Source}.\par
\end{sources}\color{black}
\needspace{8\baselineskip}\section{A certificate that survives all coefficient fields}

\begin{atlasblock}[standard]
\textbf{\sffamily Theorem  /  Newton-polytope criterion}\par\smallskip
\color{black} Let \(f\in K[x_1,\ldots,x_r]\) be nonconstant and divisible by no variable \(x_i\). If \(\operatorname{Newt}(f)\) is integrally indecomposable, then \(f\) is absolutely irreducible.\par
\end{atlasblock}\color{black}

\begin{atlasblock}[standard]
\color{black} Any factorization over \(\overline K\) decomposes the same lattice polytope. One factor must therefore have point support and be a monomial. Since no variable divides \(f\), that factor is a nonzero constant. This excludes every nontrivial factorization.\par
\end{atlasblock}\color{black}
\color{black} The coefficients matter only through which terms are nonzero. The certificate remains valid if coefficients change in any field while the relevant support and its indecomposability remain intact.\par
\begin{sources}
\textbf{Sources for this section}\par\smallskip
\noindent\hyperlink{ref:Gao2001Polytopes}{[Gao2001Polytopes]}. Integral indecomposability criterion; the elementary proof is included. \href{https://www.sciencedirect.com/science/article/pii/S0021869300985867}{Source}.\par
\noindent\hyperlink{ref:GaoLauder2001}{[GaoLauder2001]} / Corollary 2. Explicit statement excluding divisibility by any variable. \href{https://link.springer.com/article/10.1007/s00454-001-0024-0}{Source}.\par
\end{sources}\color{black}
\needspace{8\baselineskip}\section{Example: a sparse triangle in every characteristic}

\begin{atlasblock}[standard]
\textbf{\sffamily Worked example}\par\smallskip
\color{black} Consider \(f=1+x^2y+xy^3\) over any field. Its Newton polygon has vertices \((0,0),(2,1),(1,3)\). The oriented edge vectors are \((2,1),(-1,2),(-1,-3)\), all primitive integer vectors.\par
\color{black} A nontrivial Minkowski summand of a triangle must itself be a triangle with parallel corresponding edges; a segment summand would introduce another edge direction. The two triangular summands would therefore divide each edge into two positive integral lengths. A primitive edge has lattice length one, so this is impossible. The polygon is integrally indecomposable. The constant term excludes monomial factors, and the criterion proves absolute irreducibility in every characteristic.\par
\end{atlasblock}\color{black}
\begin{sources}
\textbf{Sources for this section}\par\smallskip
\noindent\hyperlink{ref:Gao2001Polytopes}{[Gao2001Polytopes]}. Triangle indecomposability criteria; this particular support is checked explicitly here. \href{https://www.sciencedirect.com/science/article/pii/S0021869300985867}{Source}.\par
\end{sources}\color{black}
\needspace{8\baselineskip}\section{Two indispensable limitations}

\begin{atlasblock}[boundary]
\color{black} \textbf{Do not forget monomials.} The polynomial \(x(1+x+y)\) has a translated indecomposable triangle as its Newton polygon, yet visibly factors. The hypothesis excluding variable divisors is essential.\par
\color{black} \textbf{The converse is false.} Over a field of characteristic different from two, \(x^2+y^2+1\) is absolutely irreducible, although its polygon is twice the standard simplex and hence decomposable. Its projective homogenization \(X^2+Y^2+Z^2\) is nonsingular; a reducible projective conic over an algebraically closed field is a pair of lines, possibly repeated, and is singular. In characteristic two this example becomes \((x+y+1)^2\), illustrating the importance of the field.\par
\end{atlasblock}\color{black}
\begin{sources}
\textbf{Sources for this section}\par\smallskip
\noindent\hyperlink{ref:GaoLauder2001}{[GaoLauder2001]} / §5. A decomposable polytope makes the shape test inconclusive; the examples are direct checks. \href{https://link.springer.com/article/10.1007/s00454-001-0024-0}{Source}.\par
\end{sources}\color{black}
\needspace{8\baselineskip}\section{When geometry alone is not enough}
\color{black} If the polytope decomposes, actual coefficients decide whether the potential factors exist. Effective Noether forms describe the reducibility locus precisely. Fix a total-degree bound \(d\). There are finitely many integer polynomials in the coefficient variables such that all of them vanish exactly when the specialized polynomial is reducible over the algebraic closure or its total degree drops below \(d\). Thus, for a polynomial of \emph{exact} degree \(d\), absolute irreducibility is certified by at least one nonzero Noether form. In positive characteristic, reduce the integer coefficients of the forms modulo that characteristic. Differential methods instead recover factor information through linear systems. For a squarefree bivariate \(f=\prod_i f_i\) in characteristic zero, the identity\par

\atlasdisplay{\frac{df}{f}=\sum_i\frac{df_i}{f_i}}
\color{black} suggests looking for closed rational one-forms with denominator \(f\). Writing such a form as \((g/f)\,dx+(h/f)\,dy\) gives the polynomial linear equation\par

\atlasdisplay{f\,(g_y-h_x)=g f_y-h f_x.}
\color{black} Degree bounds and characteristic hypotheses are essential when turning this observation into an algorithm. Gao's paper develops the method in characteristic zero or sufficiently large characteristic; arbitrary small characteristic requires separate care.\par
\color{black} The practical strategy is to try a cheap polytope certificate first, then use coefficient-sensitive factorization when the shape test is inconclusive. A single reducible specialization \(f(x,t)\) does not establish that \(f(x,y)\) itself is reducible.\par
\begin{sources}
\textbf{Sources for this section}\par\smallskip
\noindent\hyperlink{ref:Kaltofen1995}{[Kaltofen1995]} / §5, Theorem 7. Vanishing describes absolute reducibility together with degree drop; the fixed-degree qualification is essential. \href{https://www.sciencedirect.com/science/article/pii/S0022000085710239}{Source}.\par
\noindent\hyperlink{ref:Gao2003PDE}{[Gao2003PDE]}. Differential-equation factorization with characteristic restrictions. \href{https://doi.org/10.1090/S0025-5718-02-01428-X}{Source}.\par
\noindent\hyperlink{ref:Lecerf2013}{[Lecerf2013]}. Multivariate factorization methods and their hypotheses. \href{https://www.numdam.org/articles/10.5802/ccirm.18/}{Source}.\par
\noindent\hyperlink{ref:FLINTFactorizationDocumentation}{[FLINTFactorizationDocumentation]} / Multivariate factorisation: irreducibility from the Newton polytope. The polytope test is used as a certificate/pretest, not as an absolute factorization routine. \href{https://flintlib.org/doc/algorithms/polynomials/factoring-finite-fields.html}{Source}.\par
\end{sources}\color{black}
\needspace{11\baselineskip}\section*{Questions and research boundaries}
\needspace{7\baselineskip}\subsection*{Which coefficient conditions resolve a decomposable support?}
{\small\sffamily\color{atlasgold}Research direction / A fixed multivariate support and a specified coefficient field}\par\smallskip
\color{black} The polytope gives possible factor shapes but generally does not decide whether those shapes occur. One can study the corresponding reducibility locus in coefficient space or combine sparse geometry with differential or lifting methods. This is a program for a specified family, not a claim that absolute factorization is algorithmically undecidable.\par
\begin{sources}
\textbf{Sources for this section}\par\smallskip
\noindent\hyperlink{ref:GaoLauder2001}{[GaoLauder2001]} / §5.  \href{https://link.springer.com/article/10.1007/s00454-001-0024-0}{Source}.\par
\noindent\hyperlink{ref:Kaltofen1995}{[Kaltofen1995]}.  \href{https://www.sciencedirect.com/science/article/pii/S0022000085710239}{Source}.\par
\noindent\hyperlink{ref:Gao2003PDE}{[Gao2003PDE]}.  \href{https://doi.org/10.1090/S0025-5718-02-01428-X}{Source}.\par
\end{sources}\color{black}
\needspace{6\baselineskip}\par\medskip\begingroup\footnotesize\setlength{\parskip}{1pt}\textbf{\sffamily\color{atlasgreen}Connections}\par
\color{black}\textbf{Contrasts with.} \hyperlink{article:newton-polygons}{Newton polygons, factor degrees, and large cycles}. Valuation polygons use coefficient valuations; Newton polytopes use exponent support.\par
\color{black}\textbf{Generalizes.} \hyperlink{article:sparse-gaps}{Sparse polynomials, factor reversal, and large gaps}. Sparse support can constrain factors through convex geometry.\par
\color{black}\textbf{Applies to.} \hyperlink{article:factorization-algorithms}{Exact factorization: finite fields, Hensel lifting, and lattices}. Shape certificates can precede more expensive coefficient-sensitive factoring.\par
\color{black}\textbf{Contrasts with.} \hyperlink{article:specialization}{Specialization and Hilbert irreducibility}. Specialized factorization and generic multivariate factorization have different quantifiers.\par
\endgroup
\backmatter\setcounter{secnumdepth}{-1}\part{Annotated reference catalogue}\chapter{232 sources, with links and annotations}
The stable keys below are used in the section citations throughout the handbook. The reference list combines the original 188-source collection with subsequent verified additions. Publication and preprint status are preserved record by record.\par\medskip
\needspace{12\baselineskip}\section{Foundations and reference books}
\needspace{7\baselineskip}\hypertarget{ref:Serre1979}{}
{\sffamily\bfseries\color{atlasgreen}[Serre1979] Local Fields}\par
\color{black} Serre, Jean-Pierre. 1979. Springer. 67. Graduate Texts in Mathematics.\par
{\small\sffamily\color{atlasgold}Book / Published book}\par
\color{black} Reference for valuations, ramification, the different, and discriminants. Supplies the local field framework behind Newton polygons and arguments that produce inertia elements in polynomial Galois groups.\par
{\small English translation of Corps locaux by Marvin Jay Greenberg.}\par
{\footnotesize\url{https://link.springer.com/book/10.1007/978-1-4757-5673-9}}\par
{\footnotesize\url{https://doi.org/10.1007/978-1-4757-5673-9}}\par
\vspace{9pt}
\needspace{7\baselineskip}\hypertarget{ref:LidlNiederreiter1997}{}
{\sffamily\bfseries\color{atlasgreen}[LidlNiederreiter1997] Finite Fields}\par
\color{black} Lidl, Rudolf and Niederreiter, Harald. 1997. Cambridge University Press. 20. Edition 2. Encyclopedia of Mathematics and its Applications.\par
{\small\sffamily\color{atlasgold}Book / Published book}\par
\color{black} Standard reference for irreducible and primitive polynomials over finite fields, factorization, and constructions. Keep finite-field questions separate from Galois groups over the rationals: a degree-n irreducible polynomial over a finite field has cyclic splitting-field group of order n.\par
{\small The standard bibliographic year is 1997 (BNF and AMS review); the Cambridge product page lists a 24 October 1996 release date.}\par
{\footnotesize\url{https://www.cambridge.org/core/books/finite-fields/75BDAA74ABAE713196E718392B9E5E72}}\par
{\footnotesize\url{https://doi.org/10.1017/CBO9780511525926}}\par
\vspace{9pt}
\needspace{7\baselineskip}\hypertarget{ref:Neukirch1999}{}
{\sffamily\bfseries\color{atlasgreen}[Neukirch1999] Algebraic Number Theory}\par
\color{black} Neukirch, Jürgen. 1999. Springer. 322. Grundlehren der mathematischen Wissenschaften.\par
{\small\sffamily\color{atlasgold}Book / Published book}\par
\color{black} A source for prime decomposition, ramification, and global arithmetic of splitting fields. Useful when a polynomial argument passes from reductions modulo primes to arithmetic properties of number fields.\par
{\small Translated from German by Norbert Schappacher.}\par
{\footnotesize\url{https://link.springer.com/book/10.1007/978-3-662-03983-0}}\par
{\footnotesize\url{https://doi.org/10.1007/978-3-662-03983-0}}\par
\vspace{9pt}
\needspace{7\baselineskip}\hypertarget{ref:Schinzel2000}{}
{\sffamily\bfseries\color{atlasgreen}[Schinzel2000] Polynomials with Special Regard to Reducibility}\par
\color{black} Schinzel, Andrzej. 2000. Cambridge University Press. 77. Encyclopedia of Mathematics and its Applications.\par
{\small\sffamily\color{atlasgold}Book / Published book}\par
\color{black} The central research monograph for this bibliography. It treats arbitrary and lacunary polynomials over several classes of fields; its bibliography is an essential starting point for backward searches into the classical literature.\par
{\small The monograph treats several broad field classes. Its publisher explicitly excludes results valid only over finite fields, local fields, or the rational field; use it together with field-specific references.}\par
{\footnotesize\url{https://doi.org/10.1017/CBO9780511542916}}\par
{\footnotesize\url{https://www.cambridge.org/core/books/polynomials-with-special-regard-to-reducibility/FBC37D46462D46A6E121FFC9A09391C5}}\par
\vspace{9pt}
\needspace{7\baselineskip}\hypertarget{ref:Lang2002}{}
{\sffamily\bfseries\color{atlasgreen}[Lang2002] Algebra}\par
\color{black} Lang, Serge. 2002. Springer. 211. Edition 3. Graduate Texts in Mathematics.\par
{\small\sffamily\color{atlasgold}Book / Published book}\par
\color{black} A reference for the algebra underlying irreducibility, field extensions, and Galois theory. Use it to check hypotheses over general fields and the distinction between separability and irreducibility.\par
{\footnotesize\url{https://link.springer.com/book/10.1007/978-1-4613-0041-0}}\par
{\footnotesize\url{https://doi.org/10.1007/978-1-4613-0041-0}}\par
\vspace{9pt}
\needspace{7\baselineskip}\hypertarget{ref:Prasolov2004}{}
{\sffamily\bfseries\color{atlasgreen}[Prasolov2004] Polynomials}\par
\color{black} Prasolov, Victor V. 2004. Springer. 11. Algorithms and Computation in Mathematics.\par
{\small\sffamily\color{atlasgold}Book / Published book}\par
\color{black} A broad treatment of polynomial techniques: irreducibility appears on pp. 47-76, special polynomial forms on pp. 77-131, and Galois theory on pp. 181-218. Useful for connecting algebraic arguments with root location and constructive examples.\par
{\small The English first edition is ©2004 (hardcover published 9 July 2004); the linked eBook was released 23 September 2009. These are editions/formats of the same book.}\par
{\footnotesize\url{https://link.springer.com/book/10.1007/978-3-642-03980-5}}\par
{\footnotesize\url{https://doi.org/10.1007/978-3-642-03980-5}}\par
\vspace{9pt}
\needspace{7\baselineskip}\hypertarget{ref:Cox2012}{}
{\sffamily\bfseries\color{atlasgreen}[Cox2012] Galois Theory}\par
\color{black} Cox, David A. 2012. John Wiley \& Sons. Edition 2.\par
{\small\sffamily\color{atlasgold}Book / Published book}\par
\color{black} A detailed route from field extensions to explicit Galois groups. Chapters 13 and 14 treat resolvents, computation, and solvable permutation groups; especially useful for building worked examples.\par
{\footnotesize\url{https://onlinelibrary.wiley.com/doi/book/10.1002/9781118218457}}\par
{\footnotesize\url{https://doi.org/10.1002/9781118218457}}\par
\vspace{9pt}
\needspace{7\baselineskip}\hypertarget{ref:Milne2022}{}
{\sffamily\bfseries\color{atlasgreen}[Milne2022] Fields and Galois Theory}\par
\color{black} Milne, James S. 2022. Kea Books.\par
{\small\sffamily\color{atlasgold}Book / Published book; author errata available}\par
\color{black} A concise graduate reference with proofs and examples; Chapter 4 explains computation of Galois groups. The linked file is the 2022 book. Consult the author's errata, including the distinction between irreducibility in Q[X] and Z[X] and the degree-two exception in the polynomial-counting discussion. Frozen course notes v5.10 have different pagination.\par
{\small Author errata: https://www.jmilne.org/math/Books/fterrata.html. Course-note version 5.10 is frozen; the author directs readers to the book for corrections and updates.}\par
{\footnotesize\url{https://www.jmilne.org/math/Books/index.html}}\par
{\footnotesize\url{https://www.jmilne.org/math/Books/FT0.pdf}}\par
{\footnotesize\url{https://www.jmilne.org/math/CourseNotes/ft.html}}\par
{\footnotesize\url{https://www.jmilne.org/math/Books/fterrata.html}}\par
\vspace{9pt}
\needspace{7\baselineskip}\hypertarget{ref:Milne2020ANT}{}
{\sffamily\bfseries\color{atlasgreen}[Milne2020ANT] Algebraic Number Theory (v3.08)}\par
\color{black} Milne, James S. 2020. Author course notes, 166 pages.\par
{\small\sffamily\color{atlasgold}Exposition / Author course notes; version 3.08, 19 July 2020}\par
\color{black} Free notes supplying precise proofs of the discriminant-index identity, prime decomposition, Hensel lifting, unramified extensions, decomposition groups, and tame/wild inertia. Version 3.08 is dated 19 July 2020; this is a different work from Fields and Galois Theory.\par
{\footnotesize\url{https://www.jmilne.org/math/CourseNotes/ANT.pdf}}\par
\vspace{9pt}
\needspace{12\baselineskip}\section{Classical criteria and sparse-polynomial methods}
\needspace{7\baselineskip}\hypertarget{ref:schoenemann1846}{}
{\sffamily\bfseries\color{atlasgreen}[schoenemann1846] Von denjenigen Moduln, welche Potenzen von Primzahlen sind. (Fortsetzung)}\par
\color{black} Schönemann, Theodor. 1846. Journal für die reine und angewandte Mathematik. 32. pp. 93-105.\par
{\small\sffamily\color{atlasgold}Published work / published}\par
\color{black} Historical source preceding Eisenstein for the criterion now often called Schönemann-Eisenstein. Read alongside Cox's historical account for the role of higher congruences and the precise translation into modern polynomial language.\par
{\footnotesize\url{https://www.degruyterbrill.com/document/doi/10.1515/crll.1846.32.93/html}}\par
{\footnotesize\url{https://doi.org/10.1515/crll.1846.32.93}}\par
{\footnotesize\url{https://people.math.ethz.ch/~halorenz/4students/Algebra/Schoenemann_Eisenstein.pdf}}\par
\vspace{9pt}
\needspace{7\baselineskip}\hypertarget{ref:eisenstein1850}{}
{\sffamily\bfseries\color{atlasgreen}[eisenstein1850] Über die Irreductibilität und einige andere Eigenschaften der Gleichung, von welcher die Theilung der ganzen Lemniscate abhängt}\par
\color{black} Eisenstein, Gotthold. 1850. Journal für die reine und angewandte Mathematik. 39. pp. 160-179.\par
{\small\sffamily\color{atlasgold}Published work / published}\par
\color{black} Original source for Eisenstein's treatment, motivated by division of the lemniscate and Gaussian integers. The familiar integer-polynomial criterion requires a prime dividing all nonleading coefficients, not the leading coefficient, and not dividing the constant coefficient twice.\par
{\footnotesize\url{https://doi.org/10.1515/crll.1850.39.160}}\par
{\footnotesize\url{https://doi.org/10.1090/proc/14856}}\par
{\footnotesize\url{https://people.math.ethz.ch/~halorenz/4students/Algebra/Schoenemann_Eisenstein.pdf}}\par
\vspace{9pt}
\needspace{7\baselineskip}\hypertarget{ref:capelli1901binomials}{}
{\sffamily\bfseries\color{atlasgreen}[capelli1901binomials] Sulla riduttibilità della funzione x\textasciicircum{}n-A in un campo qualunque di razionalità}\par
\color{black} Capelli, Alfredo. 1901. Mathematische Annalen. 54. pp. 602-603.\par
{\small\sffamily\color{atlasgold}Published work / published}\par
\color{black} Classical reference for the binomial irreducibility criterion, including the exceptional fourth-power obstruction when the exponent is divisible by four. Distinguish this binomial criterion from the related Capelli composition lemma; the note itself points back to Capelli's 1897 and 1898 work.\par
{\footnotesize\url{https://link.springer.com/article/10.1007/BF01450724}}\par
{\footnotesize\url{https://doi.org/10.1007/BF01450724}}\par
{\footnotesize\url{https://www.numdam.org/item/AST_1975__24-25__69_0/}}\par
{\footnotesize\url{https://mathe2.uni-bayreuth.de/stoll/papers/SaShSt.pdf}}\par
\vspace{9pt}
\needspace{7\baselineskip}\hypertarget{ref:dumas1906}{}
{\sffamily\bfseries\color{atlasgreen}[dumas1906] Sur quelques cas d'irréductibilité des polynômes à coefficients rationnels}\par
\color{black} Dumas, Gustave. 1906. Journal de Mathématiques Pures et Appliquées, 6e série. 2. pp. 191-258. 6.\par
{\small\sffamily\color{atlasgold}Published work / published}\par
\color{black} Foundational source for using Newton polygons to restrict possible factor degrees. In the familiar one-edge case, coprime horizontal and vertical increments give an irreducibility criterion; the polygon-product rule is useful far beyond this special case.\par
{\footnotesize\url{https://www.numdam.org/item/JMPA_1906_6_2__191_0/}}\par
{\footnotesize\url{https://www.numdam.org/item/JMPA_1906_6_2__191_0.pdf}}\par
{\footnotesize\url{https://www.tandfonline.com/doi/abs/10.1080/00927872.2014.910800}}\par
\vspace{9pt}
\needspace{7\baselineskip}\hypertarget{ref:perron1907}{}
{\sffamily\bfseries\color{atlasgreen}[perron1907] Neue Kriterien für die Irreduzibilität algebraischer Gleichungen}\par
\color{black} Perron, Oskar. 1907. Journal für die reine und angewandte Mathematik. 132. pp. 288-307.\par
{\small\sffamily\color{atlasgold}Published work / published}\par
\color{black} Classical criteria based on coefficient dominance and the location of complex roots. A standard application is a monic integer polynomial with nonzero constant term and next-to-leading coefficient larger in absolute value than one plus the sum of the lower coefficient magnitudes.\par
{\footnotesize\url{https://www.degruyterbrill.com/document/doi/10.1515/crll.1907.132.288/html}}\par
{\footnotesize\url{https://doi.org/10.1515/crll.1907.132.288}}\par
{\footnotesize\url{https://mathe2.uni-bayreuth.de/stoll/papers/SaShSt.pdf}}\par
\vspace{9pt}
\needspace{7\baselineskip}\hypertarget{ref:schur1908problem226}{}
{\sffamily\bfseries\color{atlasgreen}[schur1908problem226] Problem 226}\par
\color{black} Schur, Issai. 1908. Archiv der Mathematik und Physik. 13. pp. 367. 3.\par
{\small\sffamily\color{atlasgold}Published work / published; citation verified indirectly}\par
\color{black} Original question on products of distinct integral linear factors perturbed by \ensuremath{\pm}1. The link is to the later historical account that verifies this citation, rather than to a scan of the original problem; that account also records Schur's Problem275 (1909) on power compositions.\par
{\small The main link identifies Győry-Hajdu-Tijdeman (2011), the later source used to verify the historical citation, rather than an original scan.}\par
{\footnotesize\url{https://link.springer.com/article/10.1007/s00605-010-0241-9}}\par
{\footnotesize\url{https://math.unideb.hu/sites/default/files/inline-files/ght20.pdf}}\par
\vspace{9pt}
\needspace{7\baselineskip}\hypertarget{ref:polya1919}{}
{\sffamily\bfseries\color{atlasgreen}[polya1919] Verschiedene Bemerkungen zur Zahlentheorie}\par
\color{black} Pólya, George. 1919. Jahresbericht der Deutschen Mathematiker-Vereinigung. 28. pp. 31-40.\par
{\small\sffamily\color{atlasgold}Published work / published}\par
\color{black} Original source for irreducibility criteria obtained from sufficiently small nonzero values at many distinct integers. The interpolation argument and later improvements are developed in Levit1968, Tverberg1973 and Győry-Hajdu-Tijdeman2011.\par
{\footnotesize\url{https://eudml.org/doc/145560}}\par
{\footnotesize\url{https://link.springer.com/article/10.1007/s00605-010-0241-9}}\par
{\footnotesize\url{https://people.math.sc.edu/filaseta/gradcourses/Math788F/latexbook/bibliography.pdf}}\par
\vspace{9pt}
\needspace{7\baselineskip}\hypertarget{ref:ore1923criteria}{}
{\sffamily\bfseries\color{atlasgreen}[ore1923criteria] Zur Theorie der Irreduzibilitätskriterien}\par
\color{black} Ore, Øystein. 1923. Mathematische Zeitschrift. 18. pp. 278-288.\par
{\small\sffamily\color{atlasgold}Published work / published}\par
\color{black} An early systematic extension of valuation-based irreducibility criteria. This belongs in the chain from Eisenstein and Dumas to augmented valuations and higher-order Newton polygons.\par
{\small DOI, volume 18, pages 278-288, and December 1923 publication date checked against the publisher. The original 1922 date is the received date, not the publication year.}\par
{\footnotesize\url{https://link.springer.com/article/10.1007/BF01192407}}\par
{\footnotesize\url{https://doi.org/10.1007/BF01192407}}\par
{\footnotesize\url{https://geodesic.mathdoc.fr/item/MZ_1923__18_167756/}}\par
\vspace{9pt}
\needspace{7\baselineskip}\hypertarget{ref:dorwartore1933}{}
{\sffamily\bfseries\color{atlasgreen}[dorwartore1933] Criteria for the irreducibility of polynomials}\par
\color{black} Dorwart, H. L. and Ore, Øystein. 1933. Annals of Mathematics. 34(1). pp. 81-94.\par
{\small\sffamily\color{atlasgold}Published work / published}\par
\color{black} Studies irreducibility using repeated small or prime values at integer arguments, and extensions of Schur-type constructions. Particularly useful when the polynomial is naturally specified through its values or through products of linear factors.\par
{\footnotesize\url{https://www.jstor.org/stable/1968341}}\par
{\footnotesize\url{https://doi.org/10.2307/1968341}}\par
{\footnotesize\url{https://link.springer.com/article/10.1007/s00605-010-0241-9}}\par
\vspace{9pt}
\needspace{7\baselineskip}\hypertarget{ref:dorwart1935survey}{}
{\sffamily\bfseries\color{atlasgreen}[dorwart1935survey] Irreducibility of polynomials}\par
\color{black} Dorwart, H. L. 1935. The American Mathematical Monthly. 42(6). pp. 369-381.\par
{\small\sffamily\color{atlasgold}Published work / published}\par
\color{black} A compact historical survey organizing criteria through divisibility of coefficients, sizes of coefficients and values at integers. It remains a useful map of the classical literature, supplemented by later work for present-day coverage.\par
{\footnotesize\url{https://www.tandfonline.com/doi/abs/10.1080/00029890.1935.11987732}}\par
{\footnotesize\url{https://doi.org/10.1080/00029890.1935.11987732}}\par
{\footnotesize\url{https://people.math.ethz.ch/~halorenz/4students/Algebra/Schoenemann_Eisenstein.pdf}}\par
\vspace{9pt}
\needspace{7\baselineskip}\hypertarget{ref:maclane1936values}{}
{\sffamily\bfseries\color{atlasgreen}[maclane1936values] A construction for absolute values in polynomial rings}\par
\color{black} MacLane, Saunders. 1936. Transactions of the American Mathematical Society. 40(3). pp. 363-395.\par
{\small\sffamily\color{atlasgold}Published work / published}\par
\color{black} Constructs polynomial valuations through successive augmentations and key polynomials. It supplies a conceptual framework for irreducibility tests that go beyond a first Newton polygon, and underlies modern local factorization methods.\par
{\footnotesize\url{https://users.encs.concordia.ca/~ford/UNCG/smvals.pdf}}\par
{\footnotesize\url{https://doi.org/10.1090/S0002-9947-1936-1501879-8}}\par
{\footnotesize\url{https://celebratio.org/MacLane_S/article/213/}}\par
\vspace{9pt}
\needspace{7\baselineskip}\hypertarget{ref:schinzel1965trinomials}{}
{\sffamily\bfseries\color{atlasgreen}[schinzel1965trinomials] On the reducibility of polynomials and in particular of trinomials}\par
\color{black} Schinzel, Andrzej. 1965. Acta Arithmetica. 11(1). pp. 1-34.\par
{\small\sffamily\color{atlasgold}Published work / published; errata available}\par
\color{black} A major classical source on structural reducibility questions, especially sparse polynomials and trinomials. Consult the 1966 errata together with the original paper before quoting a specific statement.\par
{\footnotesize\url{https://www.impan.pl/en/publishing-house/journals-and-series/acta-arithmetica/all/11/1/95825/on-the-reducibility-of-polynomials-and-in-particular-of-trinomials}}\par
{\footnotesize\url{https://doi.org/10.4064/aa-11-1-1-34}}\par
{\footnotesize\url{https://www.impan.pl/~schinzel/}}\par
\vspace{9pt}
\needspace{7\baselineskip}\hypertarget{ref:schinzel1966errata}{}
{\sffamily\bfseries\color{atlasgreen}[schinzel1966errata] Errata to the paper "On the reducibility of polynomials and in particular of trinomials"}\par
\color{black} Schinzel, Andrzej. 1966. Acta Arithmetica. 11(4). pp. 491.\par
{\small\sffamily\color{atlasgold}Published work / published erratum}\par
\color{black} Corrections to Schinzel's 1965 paper. Keep this linked to the original record in any atlas so readers do not use an uncorrected statement.\par
{\footnotesize\url{https://doi.org/10.4064/aa-11-4-491}}\par
{\footnotesize\url{https://www.impan.pl/en/publishing-house/journals-and-series/acta-arithmetica/all/en/publishing-house/journals-and-series/acta-arithmetica/all/11/4/96020/errata-to-the-paper-on-the-reducibility-of-polynomials-and-in-particular-of-trinomials}}\par
\vspace{9pt}
\needspace{7\baselineskip}\hypertarget{ref:levit1968}{}
{\sffamily\bfseries\color{atlasgreen}[levit1968] Irreducibility of polynomials with low absolute values}\par
\color{black} Levit, R. J. 1968. Transactions of the American Mathematical Society. 132(2). pp. 297-305.\par
{\small\sffamily\color{atlasgold}Published work / published}\par
\color{black} Sharpens the interpolation bounds behind Pólya's small-value criterion. It also excludes factors of specified degrees, which is often more useful than an all-or-nothing test when combined with local information.\par
{\footnotesize\url{https://www.jstor.org/stable/1994841}}\par
{\footnotesize\url{https://doi.org/10.2307/1994841}}\par
{\footnotesize\url{https://www.jstor.org/stable/i334275}}\par
{\footnotesize\url{https://link.springer.com/article/10.1007/s00605-010-0241-9}}\par
\vspace{9pt}
\needspace{7\baselineskip}\hypertarget{ref:tverberg1973smallvalues}{}
{\sffamily\bfseries\color{atlasgreen}[tverberg1973smallvalues] On the irreducibility of polynomials taking small values}\par
\color{black} Tverberg, Helge. 1973. Mathematica Scandinavica. 32. pp. 5-21.\par
{\small\sffamily\color{atlasgold}Published work / published}\par
\color{black} Further development of the small-value method, including asymptotic refinements. Read after Pólya and Levit when integer evaluations give stronger information than coefficient divisibility.\par
{\footnotesize\url{https://journals.msp.org/mscand/article/view/2165}}\par
{\footnotesize\url{https://doi.org/10.7146/math.scand.a-11442}}\par
{\footnotesize\url{https://link.springer.com/article/10.1007/s00605-010-0241-9}}\par
\vspace{9pt}
\needspace{7\baselineskip}\hypertarget{ref:brillhartfilasetaodlyzko1981}{}
{\sffamily\bfseries\color{atlasgreen}[brillhartfilasetaodlyzko1981] On an irreducibility theorem of A. Cohn}\par
\color{black} Brillhart, John and Filaseta, Michael and Odlyzko, Andrew. 1981. Canadian Journal of Mathematics. 33(5). pp. 1055-1059.\par
{\small\sffamily\color{atlasgold}Published work / published}\par
\color{black} Generalizes the criterion that a prime's digit expansion yields an irreducible polynomial when the digits are used as coefficients. Useful for prime-value methods and for understanding precisely which coefficient restrictions make a single prime evaluation decisive.\par
{\footnotesize\url{https://www.cambridge.org/core/journals/canadian-journal-of-mathematics/article/on-an-irreducibility-theorem-of-a-cohn/F1644BDFD66CFC166841810C29CBA91A}}\par
{\footnotesize\url{https://doi.org/10.4153/CJM-1981-080-0}}\par
\vspace{9pt}
\needspace{7\baselineskip}\hypertarget{ref:khandujasaha1997}{}
{\sffamily\bfseries\color{atlasgreen}[khandujasaha1997] On a generalization of Eisenstein's irreducibility criterion}\par
\color{black} Khanduja, Sudesh K. and Saha, Jayanti. 1997. Mathematika. 44(1). pp. 37-41.\par
{\small\sffamily\color{atlasgold}Published work / published}\par
\color{black} A modern extension of the Eisenstein method in the setting of valued fields. It is a useful entry to generalizations where the ordinary divisibility pattern of the coefficients does not suffice.\par
{\footnotesize\url{https://londmathsoc.onlinelibrary.wiley.com/doi/abs/10.1112/S0025579300011931}}\par
{\footnotesize\url{https://doi.org/10.1112/S0025579300011931}}\par
{\footnotesize\url{https://doi.org/10.1090/proc/14856}}\par
\vspace{9pt}
\needspace{7\baselineskip}\hypertarget{ref:filaseta1999smallnorm}{}
{\sffamily\bfseries\color{atlasgreen}[filaseta1999smallnorm] On the factorization of polynomials with small Euclidean norm}\par
\color{black} Filaseta, Michael. 1999. Number Theory in Progress, Vol. 1. pp. 143-163. de Gruyter.\par
{\small\sffamily\color{atlasgold}Published work / published}\par
\color{black} Uses restrictions on the Euclidean norm of the coefficient vector to control factorization. A useful methodological source for sparse polynomials and for results about their nonreciprocal parts, which must be distinguished from irreducibility of the whole polynomial.\par
{\small Pages143-163 follow the author's bibliography and MathSciNet citation reproduced in Sawin-Shusterman-Stoll; the publisher chapter wrapper gives143-164, apparently including a trailing page.}\par
{\footnotesize\url{https://doi.org/10.1515/9783110285581.143}}\par
{\footnotesize\url{https://people.math.sc.edu/filaseta/paperindex.html}}\par
{\footnotesize\url{https://mathe2.uni-bayreuth.de/stoll/papers/SaShSt.pdf}}\par
{\footnotesize\url{https://www.degruyterbrill.com/document/doi/10.1515/9783110285581.143/html}}\par
\vspace{9pt}
\needspace{7\baselineskip}\hypertarget{ref:filasetafordkonyagin2000}{}
{\sffamily\bfseries\color{atlasgreen}[filasetafordkonyagin2000] On an irreducibility theorem of A. Schinzel associated with coverings of the integers}\par
\color{black} Filaseta, Michael and Ford, Kevin and Konyagin, Sergei. 2000. Illinois Journal of Mathematics. 44(3). pp. 633-643.\par
{\small\sffamily\color{atlasgold}Published work / published}\par
\color{black} For large exponents, controls the nonreciprocal part of f(x)x\textasciicircum{}n+g(x), assuming coprimality, nonzero constant terms of both f and g, and explicit Capelli-type exclusions. The method improves the threshold in Schinzel's theorem and connects sparse factorization to covering systems.\par
{\footnotesize\url{https://www.ford126.web.illinois.edu/wwwpapers/ffk.pdf}}\par
{\footnotesize\url{https://doi.org/10.1215/ijm/1256060421}}\par
{\footnotesize\url{https://experts.illinois.edu/en/publications/on-an-irreducibility-theorem-of-a-schinzel-associated-with-coveri/}}\par
\vspace{9pt}
\needspace{7\baselineskip}\hypertarget{ref:murty2002primevalues}{}
{\sffamily\bfseries\color{atlasgreen}[murty2002primevalues] Prime numbers and irreducible polynomials}\par
\color{black} Murty, M. Ram. 2002. The American Mathematical Monthly. 109(5). pp. 452-458.\par
{\small\sffamily\color{atlasgold}Published work / published}\par
\color{black} Readable account linking prime values, coefficient bounds and irreducibility, with Cohn's theorem as a central example. A single prime value needs additional size or coefficient hypotheses; the paper provides such hypotheses rather than asserting a general converse to prime-producing conjectures.\par
{\footnotesize\url{https://doi.org/10.1080/00029890.2002.11919872}}\par
{\footnotesize\url{https://cms.dm.uba.ar/academico/materias/2docuat2011/teoria_de_numeros/Irreducible.pdf}}\par
\vspace{9pt}
\needspace{7\baselineskip}\hypertarget{ref:thangadurai2007survey}{}
{\sffamily\bfseries\color{atlasgreen}[thangadurai2007survey] Irreducibility of polynomials whose coefficients are integers}\par
\color{black} Thangadurai, R. 2007. RMS Mathematical Newsletter. 17(2). pp. 29-37.\par
{\small\sffamily\color{atlasgold}Published work / published}\par
\color{black} Accessible survey of rational irreducibility for integer polynomials, covering several classical tests and their limitations. A useful first reading before the more specialized Newton-polygon, small-value and prime-value papers in this section.\par
{\small The author's bibliography confirms RMS Mathematical Newsletter 17(2) (2007), 29-37. A distinct 2009 chapter with the same title is coauthored with N. Saradha.}\par
{\footnotesize\url{https://www.hri.res.in/~thanga/papers/mnl.pdf}}\par
{\footnotesize\url{https://www.hri.res.in/~thanga/expository.html}}\par
\vspace{9pt}
\needspace{7\baselineskip}\hypertarget{ref:cox2011eisenstein}{}
{\sffamily\bfseries\color{atlasgreen}[cox2011eisenstein] Why Eisenstein Proved the Eisenstein Criterion and Why Schönemann Discovered It First}\par
\color{black} Cox, David A. 2011. The American Mathematical Monthly. 118(1). pp. 3-21.\par
{\small\sffamily\color{atlasgold}Published work / published}\par
\color{black} A particularly useful historical and mathematical exposition linking Gauss's lemma, finite fields, Hensel lifting, Gaussian integers and lemniscatic division. It gives careful attribution and explains the distinct routes by which Schönemann and Eisenstein arrived at their criteria.\par
{\small The cited Monthly version (2011), pp. 3-21, is a reprint with minor changes of the article in Normat 57(2) (2009), pp. 49-73; the Monthly first-page footnote explicitly acknowledges the earlier publication. Publisher online migration date 13 December 2017 is not the publication year.}\par
{\footnotesize\url{https://people.math.ethz.ch/~halorenz/4students/Algebra/Schoenemann_Eisenstein.pdf}}\par
{\footnotesize\url{https://doi.org/10.4169/amer.math.monthly.118.01.003}}\par
{\footnotesize\url{https://www.tandfonline.com/doi/abs/10.4169/amer.math.monthly.118.01.003}}\par
{\footnotesize\url{https://ncm.gu.se/pdf/normat/Normat_2009-2.pdf}}\par
{\footnotesize\url{https://dacox.people.amherst.edu/normat.pdf}}\par
\vspace{9pt}
\needspace{7\baselineskip}\hypertarget{ref:gyoryhajdutijdeman2011}{}
{\sffamily\bfseries\color{atlasgreen}[gyoryhajdutijdeman2011] Irreducibility criteria of Schur-type and Pólya-type}\par
\color{black} Győry, Kálmán and Hajdu, Lajos and Tijdeman, Robert. 2011. Monatshefte für Mathematik. 163(4). pp. 415-443.\par
{\small\sffamily\color{atlasgold}Published work / published}\par
\color{black} Research paper with a substantial historical introduction and bibliography on small-value criteria, perturbed products and compositions. Its extensions use interpolation, algebraic number theory and Diophantine approximation, including results over imaginary quadratic rings and compositions of CM type.\par
{\footnotesize\url{https://link.springer.com/article/10.1007/s00605-010-0241-9}}\par
{\footnotesize\url{https://doi.org/10.1007/s00605-010-0241-9}}\par
{\footnotesize\url{https://math.unideb.hu/sites/default/files/inline-files/ght20.pdf}}\par
\vspace{9pt}
\needspace{7\baselineskip}\hypertarget{ref:bonciocat2015multiprime}{}
{\sffamily\bfseries\color{atlasgreen}[bonciocat2015multiprime] Schönemann-Eisenstein-Dumas-type irreducibility conditions that use arbitrarily many prime numbers}\par
\color{black} Bonciocat, Nicolae Ciprian. 2015. Communications in Algebra. 43(8). pp. 3102-3122.\par
{\small\sffamily\color{atlasgold}Published work / published}\par
\color{black} Combines coefficient divisibility information from several primes through Newton polygons. This is directly useful when each prime leaves possible factor degrees but their combined restrictions exclude every proper factor.\par
{\small Published in 2015; open author preprint v2 dated 21 March 2014. Theorem A gives the several-prime factor-degree intersection criterion. The former secondary verification link DOI 10.1090/proc/14856 is replaced by this paper's own primary text.}\par
{\footnotesize\url{https://www.tandfonline.com/doi/abs/10.1080/00927872.2014.910800}}\par
{\footnotesize\url{https://arxiv.org/pdf/1304.0874}}\par
{\footnotesize\url{https://doi.org/10.1080/00927872.2014.910800}}\par
{\footnotesize\url{https://arxiv.org/abs/1304.0874}}\par
{\footnotesize\url{https://www.tandfonline.com/doi/full/10.1080/00927872.2014.910800}}\par
{\footnotesize\url{https://arxiv.org/abs/1304.0874v2}}\par
\vspace{9pt}
\needspace{7\baselineskip}\hypertarget{ref:coledunnfilaseta2016}{}
{\sffamily\bfseries\color{atlasgreen}[coledunnfilaseta2016] Further irreducibility criteria for polynomials with non-negative coefficients}\par
\color{black} Cole, Morgan and Dunn, Scott and Filaseta, Michael. 2016. Acta Arithmetica. 175(2). pp. 137-181.\par
{\small\sffamily\color{atlasgold}Published work / published}\par
\color{black} Continues the Cohn prime-value approach for nonnegative integer coefficients. It determines sharp coefficient bounds for bases from 3 through 20, allowing coefficients larger than ordinary base digits while retaining an irreducibility conclusion from a prime evaluation.\par
{\footnotesize\url{https://www.impan.pl/en/publishing-house/journals-and-series/acta-arithmetica/all/175/2/91744/further-irreducibility-criteria-for-polynomials-with-non-negative-coefficients}}\par
{\footnotesize\url{https://doi.org/10.4064/aa8376-5-2016}}\par
{\footnotesize\url{https://www.impan.pl/shop/en/publication/transaction/download/product/91744}}\par
\vspace{9pt}
\needspace{7\baselineskip}\hypertarget{ref:filasetali2020}{}
{\sffamily\bfseries\color{atlasgreen}[filasetali2020] On the irreducibility of the non-reciprocal part of polynomials of the form f(x)x\textasciicircum{}n+g(x)}\par
\color{black} Filaseta, Michael and Li, Huixi and Patane, Frank and Skabelund, Dane. 2020. Acta Arithmetica. 196(2). pp. 187-201.\par
{\small\sffamily\color{atlasgold}Published work / published}\par
\color{black} Improves the sufficient exponent threshold in the robust-pair method of Sawin, Shusterman and Stoll. The conclusion concerns the nonreciprocal part, so reciprocal factors still require separate analysis when proving full irreducibility.\par
{\footnotesize\url{https://www.impan.pl/en/publishing-house/journals-and-series/acta-arithmetica/all/196/2/113768/on-the-irreducibility-of-the-non-reciprocal-part-of-polynomials-of-the-form-f-x-x-n-g-x}}\par
{\footnotesize\url{https://doi.org/10.4064/aa190907-11-2}}\par
{\footnotesize\url{https://huixil.wordpress.com/wp-content/uploads/2020/06/on-the-irreducibility-of-the-non-reciprocal-part-of-polynomials-of-the-form-fx-x5en-gx.pdf}}\par
\vspace{9pt}
\needspace{7\baselineskip}\hypertarget{ref:sawinshustermanstoll2020}{}
{\sffamily\bfseries\color{atlasgreen}[sawinshustermanstoll2020] Irreducibility of polynomials with a large gap}\par
\color{black} Sawin, Will and Shusterman, Mark and Stoll, Michael. 2020. Acta Arithmetica. 192(2). pp. 111-139.\par
{\small\sffamily\color{atlasgold}Published work / published}\par
\color{black} Develops Ljunggren's factor-reversal method into practical irreducibility algorithms for x\textasciicircum{}N c(x\textasciicircum{}-1)+d(x). It treats cyclotomic obstructions separately and improves effective bounds under additional robustness hypotheses; its introduction is a useful guide to earlier large-gap work.\par
{\footnotesize\url{https://arxiv.org/abs/1803.10811}}\par
{\footnotesize\url{https://doi.org/10.4064/aa180526-12-6}}\par
{\footnotesize\url{https://www.impan.pl/en/publishing-house/journals-and-series/acta-arithmetica/all/192/2/113301/irreducibility-of-polynomials-with-a-large-gap}}\par
{\footnotesize\url{https://mathe2.uni-bayreuth.de/stoll/papers/SaShSt.pdf}}\par
\vspace{9pt}
\needspace{7\baselineskip}\hypertarget{ref:GuardiaMontesNart2012Newton}{}
{\sffamily\bfseries\color{atlasgreen}[GuardiaMontesNart2012Newton] Newton polygons of higher order in algebraic number theory}\par
\color{black} Guàrdia, Jordi and Montes, Jesús and Nart, Enric. 2012. Transactions of the American Mathematical Society. 364(1). pp. 361-416.\par
{\small\sffamily\color{atlasgold}Published work / Published article; open author preprint v2 (2008)}\par
\color{black} Develops Ore's program of higher Newton polygons for local factorization. Product and residual-polynomial theorems sharpen slope restrictions to irreducibility certificates; higher-order indices control termination. Useful when reduction has repeated factors and a first Newton polygon is inconclusive. Theorem numbers cited in this atlas are explicitly those of arXiv v2 where indicated.\par
{\footnotesize\url{https://doi.org/10.1090/S0002-9947-2011-05442-5}}\par
{\footnotesize\url{https://arxiv.org/abs/0807.2620v2}}\par
{\footnotesize\url{https://arxiv.org/abs/0807.2620}}\par
{\footnotesize\url{https://arxiv.org/pdf/0807.2620}}\par
{\footnotesize\url{https://portalrecerca.uab.cat/en/publications/newton-polygons-of-higher-order-in-algebraic-number-theory/}}\par
{\footnotesize\url{https://www.numdam.org/articles/10.5802/jtnb.782/}}\par
\vspace{9pt}
\needspace{12\baselineskip}\section{Galois group detection and effective prime methods}
\needspace{7\baselineskip}\hypertarget{ref:Jordan1871}{}
{\sffamily\bfseries\color{atlasgreen}[Jordan1871] Théorèmes sur les groupes primitifs}\par
\color{black} Jordan, Camille. 1871. Journal de Mathématiques Pures et Appliquées (2). 16. pp. 383-408.\par
{\small\sffamily\color{atlasgold}Published work / Published article}\par
\color{black} Historical source for the primitive-group results underlying many proofs that a polynomial's Galois group contains A\_\allowbreak{}n. Use modern expositions for the precise prime-cycle criterion and for checking primitivity before applying it.\par
{\footnotesize\url{https://www.numdam.org/item/JMPA_1871_2_16__383_0/}}\par
{\footnotesize\url{https://arxiv.org/abs/1209.5169}}\par
\vspace{9pt}
\needspace{7\baselineskip}\hypertarget{ref:DixonMortimer1996}{}
{\sffamily\bfseries\color{atlasgreen}[DixonMortimer1996] Permutation Groups}\par
\color{black} Dixon, John D. and Mortimer, Brian. 1996. Graduate Texts in Mathematics 163. Springer.\par
{\small\sffamily\color{atlasgold}Book / Published monograph}\par
\color{black} Reference for transitivity, block systems, primitive groups, and the Jordan-type results used to force large Galois groups. It supports the group-theoretic step after reduction modulo primes or local ramification has supplied cycle information.\par
{\footnotesize\url{https://link.springer.com/book/10.1007/978-1-4612-0731-3}}\par
{\footnotesize\url{https://doi.org/10.1007/978-1-4612-0731-3}}\par
{\footnotesize\url{https://people.math.carleton.ca/~jdixon/Errata.pdf}}\par
\vspace{9pt}
\needspace{7\baselineskip}\hypertarget{ref:StevenhagenLenstra1996}{}
{\sffamily\bfseries\color{atlasgreen}[StevenhagenLenstra1996] Chebotarëv and his density theorem}\par
\color{black} Stevenhagen, Peter and Lenstra, Jr., Hendrik W. 1996. The Mathematical Intelligencer. 18(2). pp. 26-37.\par
{\small\sffamily\color{atlasgold}Exposition / Published expository article}\par
\color{black} An accessible explanation of how splitting modulo primes records conjugacy classes in a Galois group, together with the history and proof of Chebotarev's theorem. A strong source for explaining both the power and the limitations of modular factorization experiments.\par
{\footnotesize\url{https://pub.math.leidenuniv.nl/~lenstrahw/PUBLICATIONS/1996d/art.pdf}}\par
{\footnotesize\url{https://doi.org/10.1007/BF03027290}}\par
{\footnotesize\url{https://pub.math.leidenuniv.nl/~lenstrahw/PUBLICATIONS/pub.html}}\par
\vspace{9pt}
\needspace{7\baselineskip}\hypertarget{ref:Jones2014Cycles}{}
{\sffamily\bfseries\color{atlasgreen}[Jones2014Cycles] Primitive permutation groups containing a cycle}\par
\color{black} Jones, Gareth A. 2014. Bulletin of the Australian Mathematical Society. 89(1). pp. 159-165.\par
{\small\sffamily\color{atlasgold}Published work / Published article}\par
\color{black} Classifies primitive finite permutation groups containing a cycle and strengthens Jordan's prime-cycle theorem. Theorems 1.1-1.2 and Corollary 1.3 justify the cycle criteria used here. Prefer the published version: Remark 1.9 corrects the older preprint's distinction between polynomial composition and primitivity of a fixed polynomial's root action.\par
{\small Published online 18 July 2013; journal issue February 2014. Use the published text: preprint Comments 1 and 3 became Remarks 1.4 and 1.6, and the inaccurate fixed-polynomial indecomposability claim in arXiv v1 Comment 6 was corrected to a divisibility criterion in published Remark 1.9.}\par
{\footnotesize\url{https://doi.org/10.1017/S000497271300049X}}\par
{\footnotesize\url{https://www.cambridge.org/core/services/aop-cambridge-core/content/view/BBF3021D8434AD2C0BB0D3D6CE637008/S000497271300049Xa.pdf/primitive-permutation-groups-containing-a-cycle.pdf}}\par
{\footnotesize\url{https://arxiv.org/abs/1209.5169}}\par
{\footnotesize\url{https://www.cambridge.org/core/journals/bulletin-of-the-australian-mathematical-society/article/primitive-permutation-groups-containing-a-cycle/BBF3021D8434AD2C0BB0D3D6CE637008}}\par
\vspace{9pt}
\needspace{7\baselineskip}\hypertarget{ref:LagariasOdlyzko1977}{}
{\sffamily\bfseries\color{atlasgreen}[LagariasOdlyzko1977] Effective versions of the Chebotarev density theorem}\par
\color{black} Lagarias, Jeffrey C. and Odlyzko, Andrew M. 1977. Algebraic Number Fields: L-functions and Galois Properties (Proc. Sympos., Univ. Durham, Durham, 1975). pp. 409-464. Academic Press. Edited by Fröhlich, Albrecht.\par
{\small\sffamily\color{atlasgold}Published work / Published proceedings chapter}\par
\color{black} A foundational source for effective forms of Chebotarev, with conditional and unconditional estimates distinguished. It explains the analytic input needed to turn existence of useful Frobenius elements into bounds on the primes where they occur.\par
{\small The main link is Lagarias's bibliography, which verifies the 1977 proceedings chapter (pp. 409-464); it is not a full-text link.}\par
{\footnotesize\url{https://websites.umich.edu/~lagarias/zeta.html}}\par
{\footnotesize\url{https://www-users.cse.umn.edu/~odlyzko/doc/complete.html}}\par
{\footnotesize\url{https://doi.org/10.1090/mcom/3665}}\par
\vspace{9pt}
\needspace{7\baselineskip}\hypertarget{ref:Cohen1981Hilbert}{}
{\sffamily\bfseries\color{atlasgreen}[Cohen1981Hilbert] The Distribution of Galois Groups and Hilbert's Irreducibility Theorem}\par
\color{black} Cohen, Stephen D. 1981. Proceedings of the London Mathematical Society (3). 43(2). pp. 227-250.\par
{\small\sffamily\color{atlasgold}Published work / Published article}\par
\color{black} A foundational quantitative treatment of exceptional specializations and the distribution of Galois groups in polynomial families. It connects Hilbert irreducibility with sieve estimates, complementing the qualitative existence theorem.\par
{\footnotesize\url{https://londmathsoc.onlinelibrary.wiley.com/doi/10.1112/plms/s3-43.2.227}}\par
{\footnotesize\url{https://doi.org/10.1112/plms/s3-43.2.227}}\par
{\footnotesize\url{https://academic.oup.com/plms/article-abstract/s3-43/2/227/1545903}}\par
\vspace{9pt}
\needspace{7\baselineskip}\hypertarget{ref:Serre1981Chebotarev}{}
{\sffamily\bfseries\color{atlasgreen}[Serre1981Chebotarev] Quelques applications du théorème de densité de Chebotarev}\par
\color{black} Serre, Jean-Pierre. 1981. Publications Mathématiques de l'IHÉS. 54. pp. 123-201.\par
{\small\sffamily\color{atlasgold}Published work / Published article}\par
\color{black} Develops effective Chebotarev estimates and arithmetic applications. An important reference when checking exactly how analytic hypotheses, extension degrees, and discriminants enter arguments based on observed factorization types modulo primes.\par
{\footnotesize\url{https://numdam.org/articles/10.1007/BF02698692/}}\par
{\footnotesize\url{https://doi.org/10.1007/BF02698692}}\par
{\footnotesize\url{https://www.college-de-france.fr/media/jean-pierre-serre/UPL7902704499476033975_UPL61373_Serre_Chebotarev.pdf}}\par
\vspace{9pt}
\needspace{7\baselineskip}\hypertarget{ref:ThornerZaman2019Chebotarev}{}
{\sffamily\bfseries\color{atlasgreen}[ThornerZaman2019Chebotarev] A unified and improved Chebotarev density theorem}\par
\color{black} Thorner, Jesse and Zaman, Asif. 2019. Algebra \& Number Theory. 13(5). pp. 1039-1068.\par
{\small\sffamily\color{atlasgold}Published work / Published article}\par
\color{black} Gives an unconditional effective Chebotarev theorem with improved uniformity, including treatment of a possible exceptional zero. Useful for evaluating whether a probabilistic or computational Galois argument can replace a GRH input by an unconditional estimate within its required parameter range.\par
{\footnotesize\url{https://msp.org/ant/2019/13-5/ant-v13-n5-p02-s.pdf}}\par
{\footnotesize\url{https://doi.org/10.2140/ant.2019.13.1039}}\par
{\footnotesize\url{https://arxiv.org/abs/1803.02823}}\par
{\footnotesize\url{https://experts.illinois.edu/en/publications/a-unified-and-improved-chebotarev-density-theorem/}}\par
\vspace{9pt}
\needspace{7\baselineskip}\hypertarget{ref:KappeWarren1989Quartic}{}
{\sffamily\bfseries\color{atlasgreen}[KappeWarren1989Quartic] An Elementary Test for the Galois Group of a Quartic Polynomial}\par
\color{black} Kappe, Luise-Charlotte and Warren, Bette. 1989. The American Mathematical Monthly. 96(2). pp. 133-137.\par
{\small\sffamily\color{atlasgold}Published work / Published article}\par
\color{black} Original source of the elementary quartic-group test that separates cyclic and dihedral cases after the cubic resolvent and discriminant tests. Publisher metadata verified; the full formulation and proof used in this atlas are checked in Conrad's author-hosted exposition, Theorem 4.1. The publisher also assigns DOI 10.1080/00029890.1989.11972158.\par
{\footnotesize\url{https://www.jstor.org/stable/2323198}}\par
{\footnotesize\url{https://doi.org/10.2307/2323198}}\par
{\footnotesize\url{https://www.jstor.org/stable/i315188}}\par
{\footnotesize\url{https://www.tandfonline.com/doi/abs/10.1080/00029890.1989.11972158}}\par
{\footnotesize\url{https://kconrad.math.uconn.edu/blurbs/galoistheory/cubicquartic.pdf}}\par
\vspace{9pt}
\needspace{12\baselineskip}\section{Special polynomial families}
\needspace{7\baselineskip}\hypertarget{ref:Selmer1956}{}
{\sffamily\bfseries\color{atlasgreen}[Selmer1956] On the irreducibility of certain trinomials}\par
\color{black} Selmer, Ernst S. 1956. Mathematica Scandinavica. 4. pp. 287-302.\par
{\small\sffamily\color{atlasgold}Published work / Published}\par
\color{black} The classical source for irreducibility of x\textasciicircum{}n-x-1 over Q for every n at least 2. A starting point for the arithmetic of sparse polynomials; the full symmetric-group conclusion has a separate attribution.\par
{\footnotesize\url{https://doi.org/10.7146/math.scand.a-10478}}\par
\vspace{9pt}
\needspace{7\baselineskip}\hypertarget{ref:Ljunggren1960}{}
{\sffamily\bfseries\color{atlasgreen}[Ljunggren1960] On the irreducibility of certain trinomials and quadrinomials}\par
\color{black} Ljunggren, Wilhelm. 1960. Mathematica Scandinavica. 8. pp. 65-70.\par
{\small\sffamily\color{atlasgold}Published work / Published}\par
\color{black} A foundational source for factorization of sparse polynomials and the reciprocal/nonreciprocal distinction. The factor-reversal method developed here has later applications far beyond these initial families.\par
{\footnotesize\url{https://doi.org/10.7146/math.scand.a-10593}}\par
\vspace{9pt}
\needspace{7\baselineskip}\hypertarget{ref:NartVila1979}{}
{\sffamily\bfseries\color{atlasgreen}[NartVila1979] Equations of the type X\textasciicircum{}n+aX+b with absolute Galois group S\_\allowbreak{}n}\par
\color{black} Nart, Enric and Vila, Núria. 1979. Revista de la Universidad de Santander. 2. pp. 821-825.\par
{\small\sffamily\color{atlasgold}Published work / Published}\par
\color{black} An early source for symmetric Galois groups of trinomials; in particular, it precedes Osada for x\textasciicircum{}n-x-1. The institutional catalogue supplies an English title although the original paper is in Spanish.\par
{\footnotesize\url{https://portalrecerca.uab.cat/en/publications/equations-of-the-type-xnaxb-with-absolute-galois-group-sn/}}\par
{\footnotesize\url{https://kconrad.math.uconn.edu/blurbs/gradnumthy/galoisselmerpoly.pdf}}\par
\vspace{9pt}
\needspace{7\baselineskip}\hypertarget{ref:Osada1987}{}
{\sffamily\bfseries\color{atlasgreen}[Osada1987] The Galois groups of the polynomials X\textasciicircum{}n+aX\textasciicircum{}l+b}\par
\color{black} Osada, Hiroyuki. 1987. Journal of Number Theory. 25(2). pp. 230-238.\par
{\small\sffamily\color{atlasgold}Published work / Published}\par
\color{black} Gives arithmetic conditions forcing the full symmetric group for trinomials. Useful for the ramification route to transpositions and for the independent proof concerning x\textasciicircum{}n-x-1.\par
{\footnotesize\url{https://doi.org/10.1016/0022-314X(87)90029-1}}\par
{\footnotesize\url{https://kconrad.math.uconn.edu/blurbs/gradnumthy/galoisselmerpoly.pdf}}\par
\vspace{9pt}
\needspace{7\baselineskip}\hypertarget{ref:CohenMovahhediSalinier1997}{}
{\sffamily\bfseries\color{atlasgreen}[CohenMovahhediSalinier1997] Double transitivity of Galois groups of trinomials}\par
\color{black} Cohen, Stephen D. and Movahhedi, A. and Salinier, Alain. 1997. Acta Arithmetica. 82(1). pp. 1-15.\par
{\small\sffamily\color{atlasgold}Published work / Published}\par
\color{black} Studies how arithmetic hypotheses on a trinomial force double transitivity. This is an intermediate group-theoretic step that can be combined with cycle or ramification information to obtain a large Galois group.\par
{\footnotesize\url{https://doi.org/10.4064/aa-82-1-1-15}}\par
{\footnotesize\url{https://www.impan.pl/en/publishing-house/journals-and-series/acta-arithmetica/all/en/publishing-house/journals-and-series/acta-arithmetica/all/82/1/109606/double-transitivity-of-galois-groups-of-trinomials}}\par
\vspace{9pt}
\needspace{7\baselineskip}\hypertarget{ref:CohenMovahhediSalinier1999}{}
{\sffamily\bfseries\color{atlasgreen}[CohenMovahhediSalinier1999] Galois groups of trinomials}\par
\color{black} Cohen, Stephen D. and Movahhedi, A. and Salinier, Alain. 1999. Journal of Algebra. 222(2). pp. 561-573.\par
{\small\sffamily\color{atlasgold}Published work / Published}\par
\color{black} Extends earlier results of Osada using primitivity, cycle structure, and arithmetic restrictions. Apply the precise irreducibility and coprimality assumptions in the paper rather than treating sparsity alone as a sufficient condition.\par
{\footnotesize\url{https://doi.org/10.1006/jabr.1999.8033}}\par
{\footnotesize\url{https://www.sciencedirect.com/science/article/pii/S002186939998033X}}\par
{\footnotesize\url{https://www.unilim.fr/pages_perso/alain.salinier/publications.html}}\par
\vspace{9pt}
\needspace{7\baselineskip}\hypertarget{ref:Schur1929I}{}
{\sffamily\bfseries\color{atlasgreen}[Schur1929I] Einige Sätze über Primzahlen mit Anwendungen auf Irreduzibilitätsfragen, I}\par
\color{black} Schur, Issai. 1929. Sitzungsberichte der Preußischen Akademie der Wissenschaften, Physikalisch-Mathematische Klasse. pp. 125-136.\par
{\small\sffamily\color{atlasgold}Published work / Published; citation verified through later literature}\par
\color{black} Classical prime-divisor arguments applied to irreducibility. Read alongside Part II and the 1930/1931 papers to trace the origin of factorial-coefficient and orthogonal-polynomial methods.\par
{\small The linked source is Rabarison-Pazuki-Molin (2024), bibliography entry [20], which verifies the original citation. It is not a scan of Schur's 1929 paper.}\par
{\footnotesize\url{https://www.numdam.org/articles/10.5802/pmb.57/}}\par
\vspace{9pt}
\needspace{7\baselineskip}\hypertarget{ref:Schur1930}{}
{\sffamily\bfseries\color{atlasgreen}[Schur1930] Gleichungen ohne Affekt}\par
\color{black} Schur, Issai. 1930. Sitzungsberichte der Preußischen Akademie der Wissenschaften, Physikalisch-Mathematische Klasse. pp. 443-449.\par
{\small\sffamily\color{atlasgold}Published work / Published; citation verified through later literature}\par
\color{black} Classical source for the Galois-group theorem for exponential Taylor polynomials. Coleman gives a particularly accessible later proof through Newton polygons.\par
{\small The linked source is Rabarison-Pazuki-Molin (2024), bibliography entry [22], which verifies the original citation. It is not a scan of Schur's 1930 paper.}\par
{\footnotesize\url{https://www.numdam.org/articles/10.5802/pmb.57/}}\par
{\footnotesize\url{https://web.math.princeton.edu/~nmk/kGalois_11bis.pdf}}\par
\vspace{9pt}
\needspace{7\baselineskip}\hypertarget{ref:Schur1931}{}
{\sffamily\bfseries\color{atlasgreen}[Schur1931] Affektlose Gleichungen in der Theorie der Laguerreschen und Hermiteschen Polynome}\par
\color{black} Schur, Issai. 1931. Journal für die reine und angewandte Mathematik. 165. pp. 52-58.\par
{\small\sffamily\color{atlasgold}Published work / Published}\par
\color{black} Foundational treatment of irreducibility and Galois groups in the Laguerre and Hermite families. For Hermite polynomials, parity and the forced factor at zero must be respected in any statement about irreducibility.\par
{\footnotesize\url{https://doi.org/10.1515/crll.1931.165.52}}\par
{\footnotesize\url{https://gdz.sub.uni-goettingen.de/id/PPN243919689_0165?tify=%7B%22view%22%3A%22info%22%2C%22pages%22%3A%5B57%5D%7D}}\par
{\footnotesize\url{https://www.numdam.org/articles/10.5802/pmb.57/}}\par
\vspace{9pt}
\needspace{7\baselineskip}\hypertarget{ref:Coleman1987}{}
{\sffamily\bfseries\color{atlasgreen}[Coleman1987] On the Galois groups of the exponential Taylor polynomials}\par
\color{black} Coleman, Robert F. 1987. L'Enseignement Mathématique. 33. pp. 183-189.\par
{\small\sffamily\color{atlasgold}Published work / Published}\par
\color{black} A short Newton-polygon proof of Schur's theorem for sum\_\allowbreak{}\{j=0\}\textasciicircum{}n x\textasciicircum{}j/j!: irreducible over Q, with group A\_\allowbreak{}n when 4 divides n and S\_\allowbreak{}n otherwise (n at least 2). A model example of local slopes, long cycles, and the discriminant square test working together.\par
{\footnotesize\url{https://doi.org/10.5169/seals-87891}}\par
\vspace{9pt}
\needspace{7\baselineskip}\hypertarget{ref:Martin2004}{}
{\sffamily\bfseries\color{atlasgreen}[Martin2004] The Galois group of x\textasciicircum{}n-x\textasciicircum{}\{n-1\}-...-x-1}\par
\color{black} Martin, Paulo A. 2004. Journal of Pure and Applied Algebra. 190(1--3). pp. 213-223.\par
{\small\sffamily\color{atlasgold}Published work / Published}\par
\color{black} Proves that the generalized Fibonacci polynomial has group S\_\allowbreak{}n when n is even or prime. The prime-degree argument includes irreducibility of its reduction modulo that prime.\par
{\footnotesize\url{https://doi.org/10.1016/j.jpaa.2003.10.028}}\par
{\footnotesize\url{https://core.ac.uk/download/pdf/82404251.pdf}}\par
{\footnotesize\url{https://web.math.princeton.edu/~nmk/kGalois_11bis.pdf}}\par
\vspace{9pt}
\needspace{7\baselineskip}\hypertarget{ref:FilasetaKumchevPasechnik2007}{}
{\sffamily\bfseries\color{atlasgreen}[FilasetaKumchevPasechnik2007] On the irreducibility of a truncated binomial expansion}\par
\color{black} Filaseta, Michael and Kumchev, Angel and Pasechnik, Dmitrii. 2007. Rocky Mountain Journal of Mathematics. 37(2). pp. 455-464.\par
{\small\sffamily\color{atlasgold}Published work / Published}\par
\color{black} Investigates the first k+1 terms of (1+x)\textasciicircum{}n and related coefficient families. A useful bridge from factorial-coefficient criteria to explicit truncated-binomial questions.\par
{\footnotesize\url{https://doi.org/10.1216/rmjm/1181068761}}\par
{\footnotesize\url{https://ora.ox.ac.uk/objects/uuid%3A28aa8966-f92a-4ae1-818b-661a3ced8861}}\par
{\footnotesize\url{https://people.math.sc.edu/filaseta/papers/FKPpaper120304.pdf}}\par
\vspace{9pt}
\needspace{7\baselineskip}\hypertarget{ref:ChambertLoir2011}{}
{\sffamily\bfseries\color{atlasgreen}[ChambertLoir2011] The theorem of Jentzsch--Szegő on an analytic curve: application to the irreducibility of truncations of power series}\par
\color{black} Chambert-Loir, Antoine. 2011. International Journal of Number Theory. 7(7). pp. 1807-1823.\par
{\small\sffamily\color{atlasgold}Published work / Published}\par
\color{black} Relates zero-distribution on analytic curves, including nonarchimedean curves, to irreducibility information for truncations. Use the corrected published version or arXiv v3; the arXiv history specifically records a correction to the application.\par
{\footnotesize\url{https://arxiv.org/abs/1004.3813}}\par
{\footnotesize\url{https://doi.org/10.1142/S1793042111004691}}\par
\vspace{9pt}
\needspace{7\baselineskip}\hypertarget{ref:FilasetaMoy2018}{}
{\sffamily\bfseries\color{atlasgreen}[FilasetaMoy2018] On the Galois group over Q of a truncated binomial expansion}\par
\color{black} Filaseta, Michael and Moy, Richard. 2018. Colloquium Mathematicum. 154(2). pp. 295-308.\par
{\small\sffamily\color{atlasgold}Published work / Published}\par
\color{black} For a fixed truncation degree r other than 6, proves that sufficiently large n gives S\_\allowbreak{}r for the first r+1 terms of (1+x)\textasciicircum{}n. At r=6 it bounds the number of exceptions n up to N by O(log N).\par
{\footnotesize\url{https://doi.org/10.4064/cm7474-3-2018}}\par
\vspace{9pt}
\needspace{7\baselineskip}\hypertarget{ref:ShokriShaffafTaleb2019}{}
{\sffamily\bfseries\color{atlasgreen}[ShokriShaffafTaleb2019] Galois groups of Taylor polynomials of some elementary functions}\par
\color{black} Shokri, Khosro Monsef and Shaffaf, Jafar and Taleb, Reza. 2019. International Journal of Number Theory. 15(6). pp. 1127-1141.\par
{\small\sffamily\color{atlasgold}Published work / Published}\par
\color{black} Studies truncations associated with logarithmic and cosine functions. The cosine examples naturally involve a signed-permutation group because of their even symmetry.\par
{\footnotesize\url{https://doi.org/10.1142/S1793042119500623}}\par
{\footnotesize\url{https://www.numdam.org/articles/10.5802/pmb.57/}}\par
{\footnotesize\url{https://www.worldscientific.com/doi/10.1142/S1793042119500623}}\par
\vspace{9pt}
\needspace{7\baselineskip}\hypertarget{ref:RabarisonPazukiMolin2024}{}
{\sffamily\bfseries\color{atlasgreen}[RabarisonPazukiMolin2024] Exponentielle tronquée et autres contes galoisiens}\par
\color{black} Rabarison, Patrick and Pazuki, Fabien and Molin, Pascal. 2024. Publications mathématiques de Besançon. Algèbre et théorie des nombres. pp. 105-117.\par
{\small\sffamily\color{atlasgold}Published work / Published}\par
\color{black} A compact survey centred on exponential truncations, with further explicit investigations of Padé approximants. Particularly useful as an entry point connecting classical Schur results with later computational questions.\par
{\footnotesize\url{https://www.numdam.org/articles/10.5802/pmb.57/}}\par
{\footnotesize\url{https://doi.org/10.5802/pmb.57}}\par
\vspace{9pt}
\needspace{7\baselineskip}\hypertarget{ref:KatzRivinSpeculations}{}
{\sffamily\bfseries\color{atlasgreen}[KatzRivinSpeculations] Speculations in Galois Theory}\par
\color{black} Katz, Nicholas M. and Rivin, Igor. Undated. Author-hosted note.\par
{\small\sffamily\color{atlasgold}Exposition / Author-hosted contribution to the Laumon volume; publication year not verified}\par
\color{black} Poses limiting-frequency questions for Galois groups of Taylor truncations and reports computations in many families. The note explicitly separates experimental evidence from theorems; its posted copy has no date on the title page.\par
{\footnotesize\url{https://web.math.princeton.edu/~nmk/kGalois_11bis.pdf}}\par
{\footnotesize\url{https://www.laumonvolume.fr/articles}}\par
\vspace{9pt}
\needspace{7\baselineskip}\hypertarget{ref:Grosswald1978}{}
{\sffamily\bfseries\color{atlasgreen}[Grosswald1978] Bessel Polynomials}\par
\color{black} Grosswald, Emil. 1978. Springer. 698. Lecture Notes in Mathematics.\par
{\small\sffamily\color{atlasgold}Book / Published}\par
\color{black} The standard monograph, including chapters on algebraic irreducibility and the Galois group. Read its conjectural irreducibility discussion together with the later Filaseta-Trifonov theorem.\par
{\footnotesize\url{https://link.springer.com/book/10.1007/BFb0063135}}\par
{\footnotesize\url{https://doi.org/10.1007/BFb0063135}}\par
\vspace{9pt}
\needspace{7\baselineskip}\hypertarget{ref:FilasetaLam2002}{}
{\sffamily\bfseries\color{atlasgreen}[FilasetaLam2002] On the irreducibility of the generalized Laguerre polynomials}\par
\color{black} Filaseta, Michael and Lam, Tsit-Yuen. 2002. Acta Arithmetica. 105(2). pp. 177-182.\par
{\small\sffamily\color{atlasgold}Published work / Published}\par
\color{black} For each fixed rational parameter outside the negative integers, proves irreducibility over Q in every sufficiently large degree. The threshold depends on the parameter, so this is not a uniform theorem when the parameter varies with degree.\par
{\footnotesize\url{https://doi.org/10.4064/aa105-2-4}}\par
{\footnotesize\url{https://people.math.sc.edu/filaseta/papers/laguerrepap0202.pdf}}\par
\vspace{9pt}
\needspace{7\baselineskip}\hypertarget{ref:FilasetaTrifonov2002}{}
{\sffamily\bfseries\color{atlasgreen}[FilasetaTrifonov2002] The irreducibility of the Bessel polynomials}\par
\color{black} Filaseta, Michael and Trifonov, Ognian. 2002. Journal für die reine und angewandte Mathematik. 550. pp. 125-140.\par
{\small\sffamily\color{atlasgold}Published work / Published}\par
\color{black} Resolves Grosswald's irreducibility conjecture for Bessel polynomials in every positive degree. Proves a stronger coefficient-perturbation statement using Newton polygons and prime-divisor estimates.\par
{\footnotesize\url{https://doi.org/10.1515/crll.2002.069}}\par
{\footnotesize\url{https://people.math.sc.edu/filaseta/papers/besselconj0701.pdf}}\par
{\footnotesize\url{https://people.math.sc.edu/filaseta/paperindex.html}}\par
\vspace{9pt}
\needspace{7\baselineskip}\hypertarget{ref:Hajir2005}{}
{\sffamily\bfseries\color{atlasgreen}[Hajir2005] On the Galois group of generalized Laguerre polynomials}\par
\color{black} Hajir, Farshid. 2005. Journal de théorie des nombres de Bordeaux. 17(2). pp. 517-525.\par
{\small\sffamily\color{atlasgold}Published work / Published}\par
\color{black} Introduces a Newton-index criterion for containing A\_\allowbreak{}n and applies it, together with irreducibility results, to generalized Laguerre polynomials with fixed rational parameter outside the negative integers in sufficiently large degree. Irreducibility is an explicit input to the group criterion.\par
{\footnotesize\url{https://jtnb.centre-mersenne.org/item/JTNB_2005__17_2_517_0/}}\par
{\footnotesize\url{https://doi.org/10.5802/jtnb.505}}\par
{\footnotesize\url{https://arxiv.org/abs/math/0406308}}\par
{\footnotesize\url{https://numdam.org/articles/10.5802/jtnb.505/}}\par
\vspace{9pt}
\needspace{7\baselineskip}\hypertarget{ref:Hajir2009}{}
{\sffamily\bfseries\color{atlasgreen}[Hajir2009] Algebraic properties of a family of generalized Laguerre polynomials}\par
\color{black} Hajir, Farshid. 2009. Canadian Journal of Mathematics. 61(3). pp. 583-603.\par
{\small\sffamily\color{atlasgold}Published work / Published}\par
\color{black} Studies negative integral parameters linked to degree, a regime different from the fixed-parameter theorem. Newton polygons organize irreducibility and large-group results, together with precise remaining conjectures.\par
{\footnotesize\url{https://doi.org/10.4153/CJM-2009-031-6}}\par
{\footnotesize\url{https://arxiv.org/abs/math/0406307}}\par
{\footnotesize\url{https://www.cambridge.org/core/journals/canadian-journal-of-mathematics/article/algebraic-properties-of-a-family-of-generalized-laguerre-polynomials/0E2DDA7A33FA3B483679E4693D093D36}}\par
\vspace{9pt}
\needspace{7\baselineskip}\hypertarget{ref:FilasetaKiddTrifonov2012}{}
{\sffamily\bfseries\color{atlasgreen}[FilasetaKiddTrifonov2012] Laguerre polynomials with Galois group A\_\allowbreak{}m for each m}\par
\color{black} Filaseta, Michael and Kidd, Travis and Trifonov, Ognian. 2012. Journal of Number Theory. 132. pp. 776-805.\par
{\small\sffamily\color{atlasgold}Published work / Published}\par
\color{black} Completes Schur's generalized Laguerre realizations of A\_\allowbreak{}m by treating m congruent to 2 modulo 4. The constructed polynomials are irreducible for m > 2, while m = 2 is the split case.\par
{\footnotesize\url{https://doi.org/10.1016/j.jnt.2011.09.012}}\par
{\footnotesize\url{https://www.sciencedirect.com/science/article/pii/S0022314X11002502}}\par
{\footnotesize\url{https://people.math.sc.edu/filaseta/papers/FKTpap07.pdf}}\par
\vspace{9pt}
\needspace{7\baselineskip}\hypertarget{ref:BanerjeeFilasetaFinchLeidy2013}{}
{\sffamily\bfseries\color{atlasgreen}[BanerjeeFilasetaFinchLeidy2013] On classifying Laguerre polynomials which have Galois group the alternating group}\par
\color{black} Banerjee, Pradipto and Filaseta, Michael and Finch, Carrie E. and Leidy, J. Russell. 2013. Journal de théorie des nombres de Bordeaux. 25(1). pp. 1-30.\par
{\small\sffamily\color{atlasgold}Published work / Published}\par
\color{black} Classifies the integer parameter pairs with nonzero square discriminant into explicit infinite families and a finite remainder. A square discriminant constrains the group to A\_\allowbreak{}n but does not by itself prove irreducibility or equality with A\_\allowbreak{}n.\par
{\footnotesize\url{https://www.numdam.org/articles/10.5802/jtnb.822/}}\par
{\footnotesize\url{https://doi.org/10.5802/jtnb.822}}\par
\vspace{9pt}
\needspace{7\baselineskip}\hypertarget{ref:JindalLaishram2022}{}
{\sffamily\bfseries\color{atlasgreen}[JindalLaishram2022] Families of Laguerre polynomials with alternating group as Galois group}\par
\color{black} Jindal, Ankita and Laishram, Shanta. 2022. Journal of Number Theory. 241. pp. 387-429.\par
{\small\sffamily\color{atlasgold}Published work / Published}\par
\color{black} Proves irreducibility and group results for integer parameters in the interval [-2n-4,-2n], including Bessel-related cases. Verifies parts of the alternating-group conjectural classification in that range.\par
{\footnotesize\url{https://doi.org/10.1016/j.jnt.2022.04.001}}\par
{\footnotesize\url{https://digitalcommons.isical.ac.in/journal-articles/2870/}}\par
{\footnotesize\url{https://www.sciencedirect.com/science/article/abs/pii/S0022314X22001068}}\par
\vspace{9pt}
\needspace{7\baselineskip}\hypertarget{ref:CullinanHajirSell2009}{}
{\sffamily\bfseries\color{atlasgreen}[CullinanHajirSell2009] Algebraic properties of a family of Jacobi polynomials}\par
\color{black} Cullinan, John and Hajir, Farshid and Sell, Elizabeth. 2009. Journal de théorie des nombres de Bordeaux. 21(1). pp. 97-108.\par
{\small\sffamily\color{atlasgold}Published work / Published}\par
\color{black} For fixed n at least 6, obtains irreducibility for all but finitely many rational parameters in a particular one-parameter Jacobi subfamily. The Galois-group exceptional set is finite for odd n and thin for even n.\par
{\footnotesize\url{https://www.numdam.org/articles/10.5802/jtnb.659/}}\par
{\footnotesize\url{https://doi.org/10.5802/jtnb.659}}\par
\vspace{9pt}
\needspace{7\baselineskip}\hypertarget{ref:CullinanHajir2014}{}
{\sffamily\bfseries\color{atlasgreen}[CullinanHajir2014] On the Galois groups of Legendre polynomials}\par
\color{black} Cullinan, John and Hajir, Farshid. 2014. Indagationes Mathematicae. 25(3). pp. 534-552.\par
{\small\sffamily\color{atlasgold}Published work / Published}\par
\color{black} Studies the groups of parity-reduced Legendre polynomials through discriminants and Newton polygons. Its results and conjectures must be stated for the correct reduced polynomial, with irreducibility assumptions where required.\par
{\footnotesize\url{https://doi.org/10.1016/j.indag.2014.01.004}}\par
{\footnotesize\url{https://www.numdam.org/articles/10.5802/pmb.57/}}\par
\vspace{9pt}
\needspace{7\baselineskip}\hypertarget{ref:CullinanHajir2017}{}
{\sffamily\bfseries\color{atlasgreen}[CullinanHajir2017] Algebraic properties of Kaneko-Zagier lifts of supersingular polynomials}\par
\color{black} Cullinan, John and Hajir, Farshid. 2017. Proceedings of the American Mathematical Society. 145(6). pp. 2291-2304.\par
{\small\sffamily\color{atlasgold}Published work / Published}\par
\color{black} Connects rational lifts of supersingular polynomials with specialized Jacobi polynomials. Gives new irreducibility cases and investigates the Mahlburg-Ono maximal-Galois-group conjecture after removing trivial factors.\par
{\footnotesize\url{https://doi.org/10.1090/proc/13212}}\par
\vspace{9pt}
\needspace{7\baselineskip}\hypertarget{ref:Cullinan2024}{}
{\sffamily\bfseries\color{atlasgreen}[Cullinan2024] On the irreducibility of the Krawtchouck polynomials}\par
\color{black} Cullinan, John. 2024. Journal of the Ramanujan Mathematical Society. 39(4). pp. 323-335.\par
{\small\sffamily\color{atlasgold}Published work / Published}\par
\color{black} Determines Newton polygons for certain Krawtchouck polynomials and proves new irreducibility cases. Shows how a familiar local method applies to a less understood discrete orthogonal family.\par
{\footnotesize\url{https://www.mathjournals.org/jrms/2024-039-004/2024-039-004-002.html}}\par
{\footnotesize\url{https://arxiv.org/abs/2212.07539}}\par
\vspace{9pt}
\needspace{7\baselineskip}\hypertarget{ref:Carlitz1952BernoulliEuler}{}
{\sffamily\bfseries\color{atlasgreen}[Carlitz1952BernoulliEuler] Note on irreducibility of the Bernoulli and Euler polynomials}\par
\color{black} Carlitz, Leonard. 1952. Duke Mathematical Journal. 19. pp. 475-481.\par
{\small\sffamily\color{atlasgold}Published work / Published}\par
\color{black} An early arithmetic source for irreducibility in Bernoulli and Euler families. These are the Appell polynomials defined by exponential generating functions, and should not be conflated with descent-enumerating Eulerian polynomials.\par
{\footnotesize\url{https://doi.org/10.1215/S0012-7094-52-01948-0}}\par
{\footnotesize\url{https://www.mathstat.dal.ca/~dilcher/bernc.html}}\par
{\footnotesize\url{https://www.impan.pl/shop/publication/transaction/download/product/83168}}\par
{\footnotesize\url{https://pmc.ncbi.nlm.nih.gov/articles/PMC7114963/}}\par
\vspace{9pt}
\needspace{7\baselineskip}\hypertarget{ref:Brillhart1969}{}
{\sffamily\bfseries\color{atlasgreen}[Brillhart1969] On the Euler and Bernoulli polynomials}\par
\color{black} Brillhart, John. 1969. Journal für die reine und angewandte Mathematik. 234. pp. 45-64.\par
{\small\sffamily\color{atlasgold}Published work / Published}\par
\color{black} A central classical paper on arithmetic factorization and roots of Euler and Bernoulli polynomials. Useful both for known irreducibility results and for the normalization needed after removing forced factors.\par
{\footnotesize\url{https://doi.org/10.1515/crll.1969.234.45}}\par
{\footnotesize\url{https://www.mathstat.dal.ca/~dilcher/bernb.html}}\par
\vspace{9pt}
\needspace{7\baselineskip}\hypertarget{ref:FoataSchutzenberger1970}{}
{\sffamily\bfseries\color{atlasgreen}[FoataSchutzenberger1970] Théorie géométrique des polynômes eulériens}\par
\color{black} Foata, Dominique and Schützenberger, Marcel-Paul. 1970. Springer. 138. Lecture Notes in Mathematics.\par
{\small\sffamily\color{atlasgold}Book / Published}\par
\color{black} Foundational combinatorial background for Eulerian polynomials: conventions, generating functions, and symmetry. It supplies the structural background needed before formulating an arithmetic irreducibility or Galois question.\par
{\small The linked December 2005 TeX retypesetting is authorized; it reproduces the 1970 monograph with different pagination.}\par
{\footnotesize\url{https://www.mat.univie.ac.at/~slc/books/foaschuetz.html}}\par
\vspace{9pt}
\needspace{7\baselineskip}\hypertarget{ref:SharmaStraus1974}{}
{\sffamily\bfseries\color{atlasgreen}[SharmaStraus1974] On the irreducibility of a class of Euler Frobenius polynomials}\par
\color{black} Sharma, A. and Straus, Ernst G. 1974. Canadian Mathematical Bulletin. 17(2). pp. 265-273.\par
{\small\sffamily\color{atlasgold}Published work / Published}\par
\color{black} A key reference for Euler-Frobenius irreducibility. Theorem 6 at parameter r=1 includes the classical Eulerian polynomial at an odd prime index, where a shift makes Eisenstein's criterion visible.\par
{\footnotesize\url{https://doi.org/10.4153/CMB-1974-052-x}}\par
{\footnotesize\url{https://www.cambridge.org/core/services/aop-cambridge-core/content/view/C8F1019355B9DAEDF0FC0F60677DEAC7/S0008439500070818a.pdf/on_the_irreducibility_of_a_class_of_euler_frobenius_polynomials.pdf}}\par
\vspace{9pt}
\needspace{7\baselineskip}\hypertarget{ref:MinacNguyenTan2023}{}
{\sffamily\bfseries\color{atlasgreen}[MinacNguyenTan2023] Fekete polynomials, quadratic residues, and arithmetic}\par
\color{black} Mináč, Ján and Nguyen, Tung T. and Tân, Nguyen Duy. 2023. Journal of Number Theory. 242. pp. 532-575.\par
{\small\sffamily\color{atlasgold}Published work / Published}\par
\color{black} Introduces normalized reciprocal and trace polynomials associated with prime-conductor Fekete polynomials, relates values to class numbers, and studies their splitting fields. The proposed full wreath-product group is a conjecture in this source, not an all-prime theorem.\par
{\footnotesize\url{https://arxiv.org/abs/2111.05256}}\par
{\footnotesize\url{https://doi.org/10.1016/j.jnt.2022.05.001}}\par
{\footnotesize\url{https://www.sciencedirect.com/science/article/pii/S0022314X22001251}}\par
\vspace{9pt}
\needspace{7\baselineskip}\hypertarget{ref:MinacNguyenTan2024}{}
{\sffamily\bfseries\color{atlasgreen}[MinacNguyenTan2024] On the arithmetic of generalized Fekete polynomials}\par
\color{black} Mináč, Ján and Nguyen, Tung T. and Tân, Nguyen Duy. 2024. Experimental Mathematics. 33(4). pp. 723-754.\par
{\small\sffamily\color{atlasgold}Published work / Published}\par
\color{black} Extends the arithmetic investigation to quadratic characters with possibly composite conductor. Cyclotomic factors, trace-polynomial reductions, and extra symmetries constrain the possible Galois groups.\par
{\footnotesize\url{https://arxiv.org/abs/2206.11778}}\par
{\footnotesize\url{https://doi.org/10.1080/10586458.2023.2293283}}\par
{\footnotesize\url{https://www.tandfonline.com/doi/abs/10.1080/10586458.2023.2293283}}\par
\vspace{9pt}
\needspace{7\baselineskip}\hypertarget{ref:AndersonBertelliODorney2026}{}
{\sffamily\bfseries\color{atlasgreen}[AndersonBertelliODorney2026] Galois groups of reciprocal polynomials and the van der Waerden-Bhargava theorem}\par
\color{black} Anderson, Theresa C. and Bertelli, Adam and O'Dorney, Evan M. 2026. Algebra \& Number Theory. 20(3). pp. 603-628.\par
{\small\sffamily\color{atlasgold}Published work / Published}\par
\color{black} For separable reciprocal polynomials of degree 2n and growing coefficient height H, gives exceptional proportion of order H\textasciicircum{}(-1) log H for nonmonic n \ensuremath{\ge} 1 and monic n \ensuremath{\ge} 3. The corrected version leaves the monic quartic case open; the expected group is the full signed-permutation group.\par
{\footnotesize\url{https://arxiv.org/abs/2406.18970v2}}\par
{\footnotesize\url{https://doi.org/10.2140/ant.2026.20.603}}\par
{\footnotesize\url{https://arxiv.org/abs/2406.18970}}\par
{\footnotesize\url{https://arxiv.org/html/2406.18970v2}}\par
{\footnotesize\url{https://msp.org/ant/2026/20-3/ant-v20-n3-s.pdf}}\par
{\footnotesize\url{https://msp.org/index/ail.php?jpath=ant&l=O}}\par
\vspace{9pt}
\needspace{7\baselineskip}\hypertarget{ref:DLMFBernoulliEuler}{}
{\sffamily\bfseries\color{atlasgreen}[DLMFBernoulliEuler] Bernoulli and Euler Polynomials (Chapter 24)}\par
\color{black} NIST Digital Library of Mathematical Functions. Undated. NIST Digital Library of Mathematical Functions.\par
{\small\sffamily\color{atlasgold}Resource / Maintained reference resource}\par
\color{black} Authoritative conventions, generating identities, Bernoulli denominator arithmetic, and complete multiple-root classifications: every B\_\allowbreak{}n is squarefree; E\_\allowbreak{}5 is the unique Euler exception. Distinguishes Appell Euler polynomials from Eulerian descent polynomials.\par
{\footnotesize\url{https://dlmf.nist.gov/24}}\par
{\footnotesize\url{https://dlmf.nist.gov/24.2}}\par
{\footnotesize\url{https://dlmf.nist.gov/24.4}}\par
{\footnotesize\url{https://dlmf.nist.gov/24.12}}\par
{\footnotesize\url{https://dlmf.nist.gov/24.10}}\par
\vspace{9pt}
\needspace{7\baselineskip}\hypertarget{ref:Deng2026Legendre}{}
{\sffamily\bfseries\color{atlasgreen}[Deng2026Legendre] A proof of the irreducibility conjecture for Legendre polynomials}\par
\color{black} Deng, Zikang. 2026. arXiv preprint.\par
{\small\sffamily\color{atlasgold}Preprint / Preprint claim; version 2, 27 September 2026; no peer-reviewed publication verified}\par
\color{black} Preprint claiming a proof of Stieltjes' irreducibility conjecture for parity-normalized Legendre polynomials in every degree, using upper and lower bounds for a resultant. Included to keep the historical conjecture record current; the full proof has not been independently verified in this atlas.\par
{\footnotesize\url{https://arxiv.org/abs/2609.22336v2}}\par
{\footnotesize\url{https://arxiv.org/html/2609.22336v2}}\par
{\footnotesize\url{https://arxiv.org/abs/2609.22336}}\par
\vspace{9pt}
\needspace{7\baselineskip}\hypertarget{ref:Heidrich1984Eulerian}{}
{\sffamily\bfseries\color{atlasgreen}[Heidrich1984Eulerian] On the factorization of Eulerian polynomials}\par
\color{black} Heidrich, Arnošt J. J. 1984. Journal of Number Theory. 18(2). pp. 157-168.\par
{\small\sffamily\color{atlasgold}Published work / Published article; original full text not accessed in this review}\par
\color{black} A central source on arithmetic factorization of Eulerian polynomials, including the lower bound p\ensuremath{-}1 for the degree of a large irreducible factor when p is the largest prime not exceeding the index. The atlas includes an elementary Newton-polygon derivation and the prime-plus-one consequence.\par
{\small Publication metadata cross-checked in the University of Utah journal bibliography; the factor-degree result is attributed to Heidrich in the introduction of Ta Thi Hoai An and Ngo Quoc Hoan, Multiple Eulerian polynomials. No original-paper theorem number is claimed.}\par
{\footnotesize\url{https://doi.org/10.1016/0022-314X(84)90050-7}}\par
{\footnotesize\url{https://ftp.math.utah.edu/public_html/pub/tex/bib/jnumbertheory1980.pdf}}\par
{\footnotesize\url{https://math.ac.vn/uploads/files/IMH20230604.pdf}}\par
\vspace{9pt}
\needspace{7\baselineskip}\hypertarget{ref:ChidambaramMinacNguyenTan2025}{}
{\sffamily\bfseries\color{atlasgreen}[ChidambaramMinacNguyenTan2025] Fekete polynomials of principal Dirichlet characters}\par
\color{black} Chidambaram, Shiva and Mináč, Ján and Nguyen, Tung T. and Tân, Nguyen Duy. 2025. Journal of Experimental Mathematics. 1(1). pp. 51-93.\par
{\small\sffamily\color{atlasgold}Published work / Published article}\par
\color{black} Studies the 0/1 polynomials associated with principal Dirichlet characters: cyclotomic factors, the remaining polynomial, separability, and Galois groups. Distinct from the prime-conductor quadratic-character family; records conjectures suggested by calculations.\par
{\small Use publisher year2025, volume1(1), pp51-93. Some discovery indexes display2026, but the journal page and PDF state2025.}\par
{\footnotesize\url{https://jexpmath.org/index.php/jem/article/view/Vol-1Issue-1Paper-4}}\par
{\footnotesize\url{https://arxiv.org/abs/2307.14896}}\par
{\footnotesize\url{https://doi.org/10.56994/JXM.001.001.004}}\par
{\footnotesize\url{https://jexpmath.org/index.php/jem/article/download/Vol-1Issue-1Paper-4/Vol-1Issue-1Paper-4/54417}}\par
\vspace{9pt}
\needspace{7\baselineskip}\hypertarget{ref:Dilcher2008MultipleZeros}{}
{\sffamily\bfseries\color{atlasgreen}[Dilcher2008MultipleZeros] On multiple zeros of Bernoulli polynomials}\par
\color{black} Dilcher, Karl. 2008. Acta Arithmetica. 134(2). pp. 149-155.\par
{\small\sffamily\color{atlasgold}Published work / Published work}\par
\color{black} Completes the proof that every Bernoulli polynomial is squarefree, using Brillhart's restriction on possible multiple roots and congruences for centered expansions. Section 4 recalls that E\_\allowbreak{}5 is the unique Euler polynomial with multiple roots.\par
{\footnotesize\url{https://doi.org/10.4064/aa134-2-6}}\par
{\footnotesize\url{https://www.impan.pl/shop/publication/transaction/download/product/83519}}\par
{\footnotesize\url{https://dlmf.nist.gov/24.12}}\par
\vspace{9pt}
\needspace{7\baselineskip}\hypertarget{ref:Kimura2003Bernoulli}{}
{\sffamily\bfseries\color{atlasgreen}[Kimura2003Bernoulli] On universal higher order Bernoulli numbers and polynomials}\par
\color{black} Kimura, Noriaki. 2003. Report of the Research Institute of Industrial Technology, Nihon University. pp. 1-22. Report no. 70.\par
{\small\sffamily\color{atlasgold}Published work / Published work}\par
\color{black} Research report on Bernoulli families and generalizations. Section 2 states Carlitz's even-index irreducibility criterion and its relation to McCarthy's base-p digit-sum condition; later sections discuss higher-order and universal versions.\par
{\small Report number 70 (2003), ISSN 0386-1678. The report's definition of p-Eisenstein indexes coefficients in descending order; the atlas uses the equivalent usual Eisenstein criterion on the reciprocal polynomial.}\par
{\footnotesize\url{https://www.cit.nihon-u.ac.jp/laboratorydata/kenkyu/publication/report/70.pdf}}\par
\vspace{9pt}
\needspace{7\baselineskip}\hypertarget{ref:PatelSiksek2017Powers}{}
{\sffamily\bfseries\color{atlasgreen}[PatelSiksek2017Powers] On powers that are sums of consecutive like powers}\par
\color{black} Patel, Vandita and Siksek, Samir. 2017. Research in Number Theory. 3. Article 2.\par
{\small\sffamily\color{atlasgold}Published work / Published work}\par
\color{black} Proves an unconditional derangement theorem for even Bernoulli Galois groups, using 2-adic Newton polygons and decomposition groups. Applies Chebotarev to equations involving sums of powers; records the even-index irreducibility conjecture.\par
{\footnotesize\url{https://doi.org/10.1007/s40993-016-0068-0}}\par
{\footnotesize\url{https://link.springer.com/article/10.1007/s40993-016-0068-0}}\par
{\footnotesize\url{https://pmc.ncbi.nlm.nih.gov/articles/PMC7114963/}}\par
\vspace{9pt}
\needspace{7\baselineskip}\hypertarget{ref:ODorney2026Twisted}{}
{\sffamily\bfseries\color{atlasgreen}[ODorney2026Twisted] Galois groups of twisted reciprocal polynomials: a uniform asymptotic}\par
\color{black} O'Dorney, Evan M. 2026. arXiv preprint.\par
{\small\sffamily\color{atlasgold}Preprint / Preprint; v1 (25 September 2026)}\par
\color{black} For fixed twist b\ensuremath{\ne}0 and half-degree m\ensuremath{\ge}4, reports an asymptotic for exceptional hyperoctahedral Galois groups in an adapted Euclidean height. Primary arXiv indexed abstract checked; theorem details and constants require the full version.\par
{\small The primary arXiv abstract and v1 endpoint could not be reopened on 12 November 2026. The earlier abstract-only catalogue entry is retained; full theorem statements, constants and any later version remain unverified in this pass.}\par
{\footnotesize\url{https://arxiv.org/abs/2609.31613v1}}\par
{\footnotesize\url{https://arxiv.org/abs/2609.31613}}\par
\vspace{9pt}
\needspace{7\baselineskip}\hypertarget{ref:LaishramShorey2016III}{}
{\sffamily\bfseries\color{atlasgreen}[LaishramShorey2016III] Irreducibility of generalized Hermite-Laguerre polynomials III}\par
\color{black} Laishram, Shanta and Shorey, Tarlok N. 2016. Journal of Number Theory. 164. pp. 303-322.\par
{\small\sffamily\color{atlasgold}Published work / Published}\par
\color{black} All-degree irreducibility for eight third- and quarter-integer Laguerre parameters, with (q,n)=(1/4,2) exceptional. Only the ordinary assertion is used: the printed power-substitution wording needs qualification because L\_\allowbreak{}1\textasciicircum{}(-3/4)(x\ensuremath{{}^4})=1/4\ensuremath{-}x\ensuremath{{}^4} is reducible.\par
{\footnotesize\url{https://doi.org/10.1016/j.jnt.2016.01.003}}\par
{\footnotesize\url{https://arxiv.org/abs/1306.0736v2}}\par
{\footnotesize\url{https://www.isid.ac.in/~statmath/index.php?Action=ViewAbs&PreprintId=311&module=Preprint}}\par
\vspace{9pt}
\needspace{7\baselineskip}\hypertarget{ref:GrosuGrosu2021}{}
{\sffamily\bfseries\color{atlasgreen}[GrosuGrosu2021] The irreducibility of some Wronskian Hermite polynomials}\par
\color{black} Grosu, Codruţ and Grosu, Corina. 2021. Indagationes Mathematicae. 32(2). pp. 456-497.\par
{\small\sffamily\color{atlasgold}Published work / Published}\par
\color{black} Two-part Wronskian Hermite families, Newton polygons, and Schur congruences link arithmetic to zeros. We exclude the unit remainder R\_\allowbreak{}(2,1)=1 when restating Theorem 2 under this atlas's nonconstant irreducibility convention.\par
{\footnotesize\url{https://doi.org/10.1016/j.indag.2020.12.004}}\par
{\footnotesize\url{https://arxiv.org/abs/2007.00065}}\par
{\footnotesize\url{https://arxiv.org/html/2007.00065v1}}\par
{\footnotesize\url{https://www.sciencedirect.com/science/article/pii/S0019357720301233}}\par
\vspace{9pt}
\needspace{7\baselineskip}\hypertarget{ref:BanerjeeBera2022}{}
{\sffamily\bfseries\color{atlasgreen}[BanerjeeBera2022] Classifying families of orthogonal polynomials having Galois group the alternating group}\par
\color{black} Banerjee, Pradipto and Bera, Ranjan. 2022. Journal of Number Theory. 240. pp. 31-49.\par
{\small\sffamily\color{atlasgold}Published work / Published}\par
\color{black} Extends the classification program from square discriminants to alternating Galois groups of generalized Laguerre polynomials, for integer \ensuremath{\alpha}\ensuremath{\ne}n, up to finitely many exceptions. The diagonal \ensuremath{\alpha}=n remains outside the stated scope.\par
{\footnotesize\url{https://doi.org/10.1016/j.jnt.2021.12.009}}\par
{\footnotesize\url{https://www.sciencedirect.com/science/article/abs/pii/S0022314X22000087}}\par
\vspace{9pt}
\needspace{7\baselineskip}\hypertarget{ref:Melotti2025Eulerian}{}
{\sffamily\bfseries\color{atlasgreen}[Melotti2025Eulerian] Distribution of the roots of Eulerian polynomials}\par
\color{black} Melotti, Paul. 2025. Electronic Communications in Probability (accepted on author page); author manuscript arXiv:2507.15908v2.\par
{\small\sffamily\color{atlasgold}Accepted manuscript / Accepted; author manuscript checked}\par
\color{black} A combinatorial proof of the known log-Cauchy zero distribution, using moments that become exactly constant. The paper uses x times the Eulerian polynomial normalized here. The author records acceptance; final journal metadata was not independently retrieved.\par
{\footnotesize\url{https://arxiv.org/abs/2507.15908v2}}\par
{\footnotesize\url{https://arxiv.org/html/2507.15908v2}}\par
{\footnotesize\url{https://www.imo.universite-paris-saclay.fr/~paul.melotti/recherche.html}}\par
\vspace{9pt}
\needspace{7\baselineskip}\hypertarget{ref:Chen2026Dodecic}{}
{\sffamily\bfseries\color{atlasgreen}[Chen2026Dodecic] Elementary characterization of the Galois groups of x\textasciicircum{}12+ax\textasciicircum{}6+b}\par
\color{black} Chen, Malcolm Hoong Wai. 2026. Mathematica Slovaca.\par
{\small\sffamily\color{atlasgold}Published work / Published}\par
\color{black} Classifies the 16 possible groups under irreducibility over Q. Lower-degree power compositions, square and cube tests, and a cubic resolvent give an explicit decision algorithm. Published online 2 September 2026.\par
{\footnotesize\url{https://doi.org/10.1515/ms-2026-0503}}\par
{\footnotesize\url{https://www.degruyterbrill.com/document/doi/10.1515/ms-2026-0503/html}}\par
\vspace{9pt}
\needspace{7\baselineskip}\hypertarget{ref:AwtreyPataneToone2025}{}
{\sffamily\bfseries\color{atlasgreen}[AwtreyPataneToone2025] On Galois groups of power compositional nonic polynomials}\par
\color{black} Awtrey, Chad and Patane, Frank and Toone, Brian. 2025. Bulletin of the Australian Mathematical Society. 112(3). pp. 458-472.\par
{\small\sffamily\color{atlasgold}Published work / Published}\par
\color{black} Classifies the 11 possible transitive groups for irreducible g(x³), with g a rational monic cubic, using two discriminants and two resolvents. Also gives an algorithm to recognize number fields with such a defining polynomial. Online 10 February 2025; issue dated December 2025.\par
{\footnotesize\url{https://doi.org/10.1017/S0004972725000048}}\par
{\footnotesize\url{https://www.cambridge.org/core/journals/bulletin-of-the-australian-mathematical-society/article/abs/on-galois-groups-of-power-compositional-nonic-polynomials/D88BEDEF37A18AA15F6EAF653E4B6200}}\par
{\footnotesize\url{https://math.samford.edu/~cawtrey/power-comp-nonics.pdf}}\par
\vspace{9pt}
\needspace{12\baselineskip}\section{Counting, random polynomials, and random matrices}
\needspace{7\baselineskip}\hypertarget{ref:vanderWaerden1936}{}
{\sffamily\bfseries\color{atlasgreen}[vanderWaerden1936] Die Seltenheit der reduziblen Gleichungen und der Gleichungen mit Affekt}\par
\color{black} van der Waerden, B. L. 1936. Monatshefte für Mathematik und Physik. 43. pp. 133-147.\par
{\small\sffamily\color{atlasgold}Published work / Published}\par
\color{black} Classical starting point for the rarity of reducible polynomials and exceptional Galois groups in coefficient boxes. The asymptotic variable is coefficient height, with degree fixed.\par
{\footnotesize\url{https://link.springer.com/article/10.1007/BF01707594}}\par
{\footnotesize\url{https://doi.org/10.1007/BF01707594}}\par
\vspace{9pt}
\needspace{7\baselineskip}\hypertarget{ref:Chela1963}{}
{\sffamily\bfseries\color{atlasgreen}[Chela1963] Reducible polynomials}\par
\color{black} Chela, R. 1963. Journal of the London Mathematical Society. 38(1). pp. 183-188.\par
{\small\sffamily\color{atlasgold}Published work / Published}\par
\color{black} Classical asymptotic counting of reducible integer polynomials in fixed-degree coefficient boxes. Useful for the main contribution from linear factors and for separating reducibility from exceptional transitive Galois groups.\par
{\footnotesize\url{https://londmathsoc.onlinelibrary.wiley.com/doi/10.1112/jlms/s1-38.1.183}}\par
{\footnotesize\url{https://doi.org/10.1112/jlms/s1-38.1.183}}\par
{\footnotesize\url{https://arxiv.org/pdf/2111.06507}}\par
\vspace{9pt}
\needspace{7\baselineskip}\hypertarget{ref:Gallagher1973}{}
{\sffamily\bfseries\color{atlasgreen}[Gallagher1973] The large sieve and probabilistic Galois theory}\par
\color{black} Gallagher, P. X. 1973. Analytic Number Theory. 24. pp. 91-101. American Mathematical Society. Proceedings of Symposia in Pure Mathematics.\par
{\small\sffamily\color{atlasgold}Published work / Published}\par
\color{black} Foundational use of the large sieve to combine factorization information at many primes and bound exceptional Galois groups. The method turns finite-field factorization statistics into quantitative results for fixed-degree integer coefficient boxes.\par
{\footnotesize\url{https://doi.org/10.1090/pspum/024/0332694}}\par
{\footnotesize\url{https://doi.org/10.1090/pspum/024}}\par
{\footnotesize\url{https://www.cecm.sfu.ca/organics/papers/odlyzko/paper/html/node7.html}}\par
{\footnotesize\url{https://doi.org/10.1017/S1474748023000154}}\par
\vspace{9pt}
\needspace{7\baselineskip}\hypertarget{ref:Dietmann2012}{}
{\sffamily\bfseries\color{atlasgreen}[Dietmann2012] On the distribution of Galois groups}\par
\color{black} Dietmann, Rainer. 2012. Mathematika. 58(1). pp. 35-44.\par
{\small\sffamily\color{atlasgold}Published work / Published}\par
\color{black} Bounds the number of monic integer polynomials of fixed degree and bounded height with a prescribed Galois group. Useful for the resolvent and integral-point approach to exceptional groups.\par
{\footnotesize\url{https://arxiv.org/abs/1010.5341}}\par
{\footnotesize\url{https://doi.org/10.1112/S0025579311002105}}\par
{\footnotesize\url{https://pure.royalholloway.ac.uk/en/publications/on-the-distribution-of-galois-groups/}}\par
{\footnotesize\url{https://www.cambridge.org/core/journals/mathematika/article/abs/on-the-distribution-of-galois-groups/389955B071D757E6F61BDF5776AC9AE5}}\par
\vspace{9pt}
\needspace{7\baselineskip}\hypertarget{ref:ChowDietmann2020}{}
{\sffamily\bfseries\color{atlasgreen}[ChowDietmann2020] Enumerative Galois theory for cubics and quartics}\par
\color{black} Chow, Sam and Dietmann, Rainer. 2020. Advances in Mathematics. 372. Article 107282.\par
{\small\sffamily\color{atlasgold}Published work / Published}\par
\color{black} Studies counts for individual Galois groups in degrees three and four. A useful source for understanding the stronger asymptotic questions that remain after a generic full-symmetric-group result.\par
{\footnotesize\url{https://arxiv.org/abs/1807.05820}}\par
{\footnotesize\url{https://doi.org/10.1016/j.aim.2020.107282}}\par
{\footnotesize\url{https://wrap.warwick.ac.uk/id/eprint/138618/}}\par
{\footnotesize\url{https://www.sciencedirect.com/science/article/pii/S000187082030308X}}\par
\vspace{9pt}
\needspace{7\baselineskip}\hypertarget{ref:BarySorokerBenPorathMatei2024}{}
{\sffamily\bfseries\color{atlasgreen}[BarySorokerBenPorathMatei2024] Probabilistic Galois Theory: The square discriminant case}\par
\color{black} Bary-Soroker, Lior and Ben-Porath, Or and Matei, Vlad. 2024. Bulletin of the London Mathematical Society. 56(6). pp. 2162-2177.\par
{\small\sffamily\color{atlasgold}Published work / Published}\par
\color{black} Studies the frequency of alternating Galois groups and nonzero square discriminants in the fixed-degree large-box model. Constructions and lower bounds show why a naive random-integer model for the discriminant can be misleading.\par
{\footnotesize\url{https://londmathsoc.onlinelibrary.wiley.com/doi/full/10.1112/blms.13049}}\par
{\footnotesize\url{https://doi.org/10.1112/blms.13049}}\par
{\footnotesize\url{https://arxiv.org/abs/2207.12493}}\par
\vspace{9pt}
\needspace{7\baselineskip}\hypertarget{ref:Bhargava2025}{}
{\sffamily\bfseries\color{atlasgreen}[Bhargava2025] Galois groups of random integer polynomials and van der Waerden's Conjecture}\par
\color{black} Bhargava, Manjul. 2025. Annals of Mathematics. 201(2). pp. 339-377.\par
{\small\sffamily\color{atlasgold}Published work / Published}\par
\color{black} Proves the order-of-magnitude bound O(H\textasciicircum{}(n-1)) for monic integer polynomials of each fixed degree n \ensuremath{\ge} 3 and height at most H whose Galois group is not S\_\allowbreak{}n. Degree is fixed; this does not settle the fixed-coefficient-law, growing-degree problem.\par
{\footnotesize\url{https://annals.math.princeton.edu/2025/201-2/p01}}\par
{\footnotesize\url{https://doi.org/10.4007/annals.2025.201.2.1}}\par
{\footnotesize\url{https://arxiv.org/abs/2111.06507}}\par
{\footnotesize\url{https://arxiv.org/html/2609.40123v1}}\par
\vspace{9pt}
\needspace{7\baselineskip}\hypertarget{ref:OdlyzkoPoonen1993}{}
{\sffamily\bfseries\color{atlasgreen}[OdlyzkoPoonen1993] Zeros of polynomials with 0,1 coefficients}\par
\color{black} Odlyzko, A. M. and Poonen, Bjorn. 1993. L'Enseignement Mathématique. 39(3--4). pp. 317-348.\par
{\small\sffamily\color{atlasgold}Published work / Published}\par
\color{black} Historical source for the asymptotic irreducibility question for monic 0-1 polynomials with constant coefficient 1. Its principal results concern the geometry and topology of their zero sets, which provide useful background for restricted-coefficient families.\par
{\footnotesize\url{https://www.cecm.sfu.ca/organics/papers/odlyzko/paper/html/paper.html}}\par
{\footnotesize\url{https://www-users.cse.umn.edu/~odlyzko/doc/sequence.html}}\par
{\footnotesize\url{https://www.cecm.sfu.ca/organics/papers/odlyzko/paper/html/node1.html}}\par
\vspace{9pt}
\needspace{7\baselineskip}\hypertarget{ref:Konyagin1999}{}
{\sffamily\bfseries\color{atlasgreen}[Konyagin1999] On the number of irreducible polynomials with 0,1 coefficients}\par
\color{black} Konyagin, S. V. 1999. Acta Arithmetica. 88(4). pp. 333-350.\par
{\small\sffamily\color{atlasgold}Published work / Published}\par
\color{black} A foundational unconditional result for the 0-1 model, excluding small-degree factors with high probability and obtaining a lower bound of order 1/log n for irreducibility. The finite-field and exponential-sum techniques are important predecessors of later work.\par
{\footnotesize\url{https://www.impan.pl/en/publishing-house/journals-and-series/acta-arithmetica/all/88/44/110736/on-the-number-of-irreducible-polynomials-with-0-1-coefficients}}\par
{\footnotesize\url{https://doi.org/10.4064/aa-88-4-333-350}}\par
{\footnotesize\url{https://arxiv.org/pdf/1810.13360}}\par
\vspace{9pt}
\needspace{7\baselineskip}\hypertarget{ref:BarySorokerKozma2017Bivariate}{}
{\sffamily\bfseries\color{atlasgreen}[BarySorokerKozma2017Bivariate] Is a bivariate polynomial with \ensuremath{\pm}1 coefficients irreducible? Very likely!}\par
\color{black} Bary-Soroker, Lior and Kozma, Gady. 2017. International Journal of Number Theory. 13(4). pp. 933-936.\par
{\small\sffamily\color{atlasgold}Published work / Published}\par
\color{black} A short treatment of irreducibility for random bivariate sign polynomials. It is useful for comparing the extra flexibility of several variables with the substantially harder univariate fixed-coefficient problem.\par
{\footnotesize\url{https://doi.org/10.1142/S1793042117500488}}\par
{\footnotesize\url{https://weizmann.elsevierpure.com/en/publications/is-a-bivariate-polynomial-with-1-coefficients-irreducible-very-li/}}\par
{\footnotesize\url{https://arxiv.org/pdf/1810.13360}}\par
\vspace{9pt}
\needspace{7\baselineskip}\hypertarget{ref:BorstEtAl2018}{}
{\sffamily\bfseries\color{atlasgreen}[BorstEtAl2018] Irreducibility of random polynomials}\par
\color{black} Borst, Christian and Boyd, Evan and Brekken, Claire and Solberg, Samantha and Wood, Melanie Matchett and Wood, Philip Matchett. 2018. Experimental Mathematics. 27(4). pp. 498-506.\par
{\small\sffamily\color{atlasgold}Published work / Published}\par
\color{black} Computational study of several sampling models and a universality heuristic relating reducibility to low-degree factors. Treat its numerical evidence and conjectures separately from its rigorous statements.\par
{\footnotesize\url{https://arxiv.org/abs/1705.03709}}\par
{\footnotesize\url{https://doi.org/10.1080/10586458.2017.1325790}}\par
{\footnotesize\url{https://www.tandfonline.com/doi/abs/10.1080/10586458.2017.1325790}}\par
\vspace{9pt}
\needspace{7\baselineskip}\hypertarget{ref:BreuillardVarju2019}{}
{\sffamily\bfseries\color{atlasgreen}[BreuillardVarju2019] Irreducibility of random polynomials of large degree}\par
\color{black} Breuillard, Emmanuel and Varjú, Péter P. 2019. Acta Mathematica. 223(2). pp. 195-249.\par
{\small\sffamily\color{atlasgold}Published work / Published; main general results conditional on RH}\par
\color{black} Under Riemann hypotheses for suitable Dedekind zeta functions, proves irreducibility and alternating-group containment for broad independent-coefficient models, with careful treatment of zero and cyclotomic factors. The method connects finite-field random walks, root counts and effective prime-ideal estimates.\par
{\footnotesize\url{https://arxiv.org/abs/1810.13360}}\par
{\footnotesize\url{https://doi.org/10.4310/ACTA.2019.v223.n2.a1}}\par
\vspace{9pt}
\needspace{7\baselineskip}\hypertarget{ref:ORourkeWood2019}{}
{\sffamily\bfseries\color{atlasgreen}[ORourkeWood2019] Low-degree factors of random polynomials}\par
\color{black} O'Rourke, Sean and Wood, Philip Matchett. 2019. Journal of Theoretical Probability. 32(2). pp. 1076-1104.\par
{\small\sffamily\color{atlasgold}Published work / Published}\par
\color{black} Converts pointwise delocalization of roots into exclusion of factors of low algebraic degree. Covers several random-matrix characteristic-polynomial models as well as random polynomials; absence of low-degree factors alone is weaker than irreducibility.\par
{\footnotesize\url{https://arxiv.org/abs/1608.01938}}\par
{\footnotesize\url{https://doi.org/10.1007/s10959-018-0839-8}}\par
{\footnotesize\url{https://math.colorado.edu/~seor3821/research/}}\par
\vspace{9pt}
\needspace{7\baselineskip}\hypertarget{ref:BarySorokerKozma2020}{}
{\sffamily\bfseries\color{atlasgreen}[BarySorokerKozma2020] Irreducible polynomials of bounded height}\par
\color{black} Bary-Soroker, Lior and Kozma, Gady. 2020. Duke Mathematical Journal. 169(4). pp. 579-598.\par
{\small\sffamily\color{atlasgold}Published work / Published}\par
\color{black} Proves high-probability irreducibility and containment of A\_\allowbreak{}n for suitable fixed coefficient distributions, including uniform coefficients in \{1,...,210\}. Combines reductions at several primes with probabilistic permutation-group arguments.\par
{\footnotesize\url{https://arxiv.org/abs/1710.05165}}\par
{\footnotesize\url{https://doi.org/10.1215/00127094-2019-0047}}\par
\vspace{9pt}
\needspace{7\baselineskip}\hypertarget{ref:BarySorokerKoukoulopoulosKozma2023}{}
{\sffamily\bfseries\color{atlasgreen}[BarySorokerKoukoulopoulosKozma2023] Irreducibility of random polynomials: general measures}\par
\color{black} Bary-Soroker, Lior and Koukoulopoulos, Dimitris and Kozma, Gady. 2023. Inventiones Mathematicae. 233(3). pp. 1041-1120.\par
{\small\sffamily\color{atlasgold}Published work / Published}\par
\color{black} Develops Fourier equidistribution, a sieve for factor degrees and permutation-group methods for independent integer coefficients, conditioning on nonzero constant term. Proves high-probability irreducibility and Galois group A\_\allowbreak{}n or S\_\allowbreak{}n under explicit distributional hypotheses, including a uniform interval of at least 35 integers; much more general laws yield exclusion of factors up to a positive fraction of the degree.\par
{\footnotesize\url{https://arxiv.org/abs/2007.14567}}\par
{\footnotesize\url{https://doi.org/10.1007/s00222-023-01193-6}}\par
{\footnotesize\url{https://cris.tau.ac.il/en/publications/irreducibility-of-random-polynomials-general-measures/}}\par
\vspace{9pt}
\needspace{7\baselineskip}\hypertarget{ref:BarySorokerGoldgraber2025}{}
{\sffamily\bfseries\color{atlasgreen}[BarySorokerGoldgraber2025] Full Galois groups of polynomials with slowly growing coefficients}\par
\color{black} Bary-Soroker, Lior and Goldgraber, Noam. 2025. Bulletin of the London Mathematical Society. 57(3). pp. 941-955.\par
{\small\sffamily\color{atlasgold}Published work / Published}\par
\color{black} For monic degree-n polynomials with coefficients uniform in [-L(n),L(n)], proves full symmetric Galois group with high probability whenever L(n) tends to infinity. The interval may grow arbitrarily slowly; the theorem does not assert the same conclusion for a fixed interval.\par
{\footnotesize\url{https://arxiv.org/abs/2404.13559}}\par
{\footnotesize\url{https://doi.org/10.1112/blms.70008}}\par
{\footnotesize\url{https://cris.tau.ac.il/en/publications/full-galois-groups-of-polynomials-with-slowly-growing-coefficient-2/}}\par
\vspace{9pt}
\needspace{7\baselineskip}\hypertarget{ref:BarySorokerHokkenKozmaPoonen2025}{}
{\sffamily\bfseries\color{atlasgreen}[BarySorokerHokkenKozmaPoonen2025] Irreducibility of Littlewood polynomials of special degrees}\par
\color{black} Bary-Soroker, Lior and Hokken, David and Kozma, Gady and Poonen, Bjorn. 2025. International Mathematics Research Notices. 2025(21). pp. rnaf326.\par
{\small\sffamily\color{atlasgold}Published work / Published}\par
\color{black} Proves unconditionally that the limsup of the irreducibility probability for uniformly random Littlewood polynomials is 1. Special-degree factorization constraints are essential; the statement should not be replaced by an all-degree limit.\par
{\footnotesize\url{https://arxiv.org/abs/2308.04878}}\par
{\footnotesize\url{https://doi.org/10.1093/imrn/rnaf326}}\par
{\footnotesize\url{https://academic.oup.com/imrn/article/2025/21/rnaf326/8300532}}\par
\vspace{9pt}
\needspace{7\baselineskip}\hypertarget{ref:Bazin2025}{}
{\sffamily\bfseries\color{atlasgreen}[Bazin2025] Irreducibility of random polynomials of \ensuremath{\mathbb Z}[X]}\par
\color{black} Bazin, Pierre-Alexandre. 2025. International Mathematics Research Notices. 2025(11). pp. rnaf136.\par
{\small\sffamily\color{atlasgold}Published work / Published}\par
\color{black} Improves rates in the Bary-Soroker-Koukoulopoulos-Kozma framework when near-uniformity holds modulo more primes. The zero-constant-term obstruction remains: for example, uniform intervals of length at least 10\textasciicircum{}8 give P(reducible) \ensuremath{\le} P(a\_\allowbreak{}0=0)+O(n\textasciicircum{}(-1/2)), rather than an unconditioned bound tending to zero.\par
{\footnotesize\url{https://arxiv.org/abs/2402.06643}}\par
{\footnotesize\url{https://doi.org/10.1093/imrn/rnaf136}}\par
\vspace{9pt}
\needspace{7\baselineskip}\hypertarget{ref:HokkenKoukoulopoulos2025}{}
{\sffamily\bfseries\color{atlasgreen}[HokkenKoukoulopoulos2025] Irreducibility and Galois groups of random reciprocal polynomials of large degree}\par
\color{black} Hokken, David and Koukoulopoulos, Dimitris. 2025.\par
{\small\sffamily\color{atlasgold}Preprint / Preprint; version 2 dated 23 September 2026}\par
\color{black} For reciprocal polynomials of degree 2m and coefficient laws satisfying a Fourier hypothesis, proves high-probability irreducibility and a Galois group equal to the hyperoctahedral group or one of two index-two subgroups. The paired-root symmetry changes both the natural ambient group and the factorization arguments.\par
{\small Version 2, 23 September 2026}\par
{\footnotesize\url{https://arxiv.org/abs/2510.18857v2}}\par
{\footnotesize\url{https://arxiv.org/abs/2510.18857}}\par
\vspace{9pt}
\needspace{7\baselineskip}\hypertarget{ref:BlacharBreuillardKozma2026}{}
{\sffamily\bfseries\color{atlasgreen}[BlacharBreuillardKozma2026] Galois groups of random polynomials of large degree}\par
\color{black} Blachar, Guy and Breuillard, Emmanuel and Kozma, Gady. 2026.\par
{\small\sffamily\color{atlasgold}Preprint / Preprint; version 1 dated 30 September 2026; scope issue noted}\par
\color{black} Preprint reporting an O(1/log n) square-discriminant probability bound in a high-entropy bounded-coefficient model (Rényi entropy of order 2 greater than 12). The atlas reports the constant-term-one subclass and the conditional uniform-interval full-S\_\allowbreak{}n corollary. The unqualified v1 statement permits a zero constant term and misses the persistent discriminant-zero event when the first two coefficients vanish; the full proof has not been independently verified.\par
{\small Version 1, 30 September 2026}\par
{\small Theorem 1.1 in v1 includes square discriminant without excluding zero and permits an arbitrary independent bounded constant-term law. If the constant and linear coefficients are iid with law \ensuremath{\mu} and \ensuremath{\mu}(0)>0, discriminant zero has probability at least \ensuremath{\mu}(0)\textasciicircum{}2, contradicting the unrestricted bound. Section 8 invokes an O(1/n) repeated-root bound without removing this event. The website records only the a0=1 reported subclass and the explicitly conditioned interval corollary; it does not certify or repair the preprint proof.}\par
{\footnotesize\url{https://arxiv.org/abs/2609.40123v1}}\par
{\footnotesize\url{https://arxiv.org/abs/2609.40123}}\par
{\footnotesize\url{https://arxiv.org/html/2609.40123v1}}\par
\vspace{9pt}
\needspace{7\baselineskip}\hypertarget{ref:Kogler2026}{}
{\sffamily\bfseries\color{atlasgreen}[Kogler2026] The Odlyzko-Poonen Conjecture on Irreducibility of Random Polynomials}\par
\color{black} Kogler, Constantin. 2026.\par
{\small\sffamily\color{atlasgold}Preprint / Preprint; version 2 dated 14 October 2026}\par
\color{black} Preprint presenting an unconditional proof of the irreducibility limit for monic polynomials with constant coefficient 1 and all remaining coefficients independent uniform in \{0,1\}. This precise model and the preprint status should accompany any account of the claimed advance.\par
{\small Version 2, 14 October 2026}\par
{\footnotesize\url{https://arxiv.org/abs/2609.26771v2}}\par
{\footnotesize\url{https://arxiv.org/abs/2609.26771}}\par
\vspace{9pt}
\needspace{7\baselineskip}\hypertarget{ref:VarjuXu2026}{}
{\sffamily\bfseries\color{atlasgreen}[VarjuXu2026] A random polynomial with multiplicative coefficients is almost surely irreducible}\par
\color{black} Varjú, Péter P. and Xu, Max Wenqiang. 2026. International Mathematics Research Notices. 2026(16). pp. rnag178.\par
{\small\sffamily\color{atlasgold}Published work / Published; conditional on RH}\par
\color{black} Under RH for Dedekind zeta functions, proves irreducibility with probability 1-O(d\textasciicircum{}(-1/2+epsilon)) for specified random multiplicative-coefficient models. The theorem concerns P\_\allowbreak{}d(x)/x, after removal of the obligatory factor x; these coefficients are dependent.\par
{\footnotesize\url{https://arxiv.org/abs/2511.04240}}\par
{\footnotesize\url{https://doi.org/10.1093/imrn/rnag178}}\par
{\footnotesize\url{https://academic.oup.com/imrn/article/2026/16/rnag178/8761102}}\par
\vspace{9pt}
\needspace{7\baselineskip}\hypertarget{ref:Rivin2008}{}
{\sffamily\bfseries\color{atlasgreen}[Rivin2008] Walks on groups, counting reducible matrices, polynomials, and surface and free group automorphisms}\par
\color{black} Rivin, Igor. 2008. Duke Mathematical Journal. 142(2). pp. 353-379.\par
{\small\sffamily\color{atlasgold}Published work / Published}\par
\color{black} Relates random walks on arithmetic groups and finite quotients to irreducible characteristic polynomials. Provides a route from arithmetic of polynomials to geometric group theory; its random-walk and height models differ from independent bounded matrix entries.\par
{\footnotesize\url{https://arxiv.org/abs/math/0703532}}\par
{\footnotesize\url{https://doi.org/10.1215/00127094-2008-009}}\par
{\footnotesize\url{https://ww2.newschool.edu/WorkArea/DownloadAsset.aspx?id=54362}}\par
\vspace{9pt}
\needspace{7\baselineskip}\hypertarget{ref:Eberhard2022}{}
{\sffamily\bfseries\color{atlasgreen}[Eberhard2022] The characteristic polynomial of a random matrix}\par
\color{black} Eberhard, Sean. 2022. Combinatorica. 42(4). pp. 491-527.\par
{\small\sffamily\color{atlasgold}Published work / Published; conditional on ERH}\par
\color{black} Assuming ERH, proves high-probability irreducibility and containment of A\_\allowbreak{}n for characteristic polynomials of matrices with independent entries from a fixed nontrivial finitely supported integer law. The dependent polynomial coefficients require matrix-specific input.\par
{\footnotesize\url{https://arxiv.org/abs/2008.01223}}\par
{\footnotesize\url{https://doi.org/10.1007/s00493-020-4657-0}}\par
{\footnotesize\url{https://pure.qub.ac.uk/en/publications/the-characteristic-polynomial-of-a-random-matrix/}}\par
\vspace{9pt}
\needspace{7\baselineskip}\hypertarget{ref:FerberJainSahSawhney2023}{}
{\sffamily\bfseries\color{atlasgreen}[FerberJainSahSawhney2023] Random symmetric matrices: rank distribution and irreducibility of the characteristic polynomial}\par
\color{black} Ferber, Asaf and Jain, Vishesh and Sah, Ashwin and Sawhney, Mehtaab. 2023. Mathematical Proceedings of the Cambridge Philosophical Society. 174(2). pp. 233-246.\par
{\small\sffamily\color{atlasgold}Published work / Published; irreducibility application conditional on ERH}\par
\color{black} Under ERH, proves high-probability irreducibility for characteristic polynomials of random symmetric sign matrices. The key innovation is sharp rank-distribution estimates over finite fields, using inverse Littlewood-Offord ideas.\par
{\footnotesize\url{https://arxiv.org/abs/2106.04049}}\par
{\footnotesize\url{https://doi.org/10.1017/S0305004122000226}}\par
{\footnotesize\url{https://www.cambridge.org/core/journals/mathematical-proceedings-of-the-cambridge-philosophical-society/article/abs/random-symmetric-matrices-rank-distribution-and-irreducibility-of-the-characteristic-polynomial/99A806811892A282D875B22DF443BC19}}\par
\vspace{9pt}
\needspace{7\baselineskip}\hypertarget{ref:AndersonODorney2026Matrices}{}
{\sffamily\bfseries\color{atlasgreen}[AndersonODorney2026Matrices] Galois groups of random integer matrices}\par
\color{black} Anderson, Theresa C. and O'Dorney, Evan M. 2026. Mathematical Proceedings of the Cambridge Philosophical Society. pp. 1-16.\par
{\small\sffamily\color{atlasgold}Published work / Published online 12 August 2026; First View}\par
\color{black} Studies fixed-size integer matrices with entries bounded by a height parameter and counts characteristic polynomials with non-symmetric Galois group. Uses the large sieve, Fourier analysis and geometric sieving; this is a different asymptotic regime from Eberhard's growing-dimension theorem.\par
{\footnotesize\url{https://www.cambridge.org/core/journals/mathematical-proceedings-of-the-cambridge-philosophical-society/article/galois-groups-of-random-integer-matrices/B87E2B0D71138A94EF7BB1179891456A}}\par
{\footnotesize\url{https://doi.org/10.1017/S0305004126102126}}\par
\vspace{9pt}
\needspace{7\baselineskip}\hypertarget{ref:BarySorokerGarzoniSodin2026}{}
{\sffamily\bfseries\color{atlasgreen}[BarySorokerGarzoniSodin2026] Irreducibility of the characteristic polynomials of random tridiagonal matrices}\par
\color{black} Bary-Soroker, Lior and Garzoni, Daniele and Sodin, Sasha. 2026. Journal of Number Theory. 280. pp. 973-998.\par
{\small\sffamily\color{atlasgold}Published work / Published; conditional on RH}\par
\color{black} Under RH for certain Dedekind zeta functions, proves irreducibility with exponentially high probability and alternating-or-symmetric Galois group for a specified tridiagonal model. Also analyzes a related model with a smaller Galois group, illustrating why the matrix ensemble matters.\par
{\footnotesize\url{https://arxiv.org/abs/2502.17218}}\par
{\footnotesize\url{https://doi.org/10.1016/j.jnt.2025.09.008}}\par
\vspace{9pt}
\needspace{7\baselineskip}\hypertarget{ref:LuczakPyber1993}{}
{\sffamily\bfseries\color{atlasgreen}[LuczakPyber1993] On random generation of the symmetric group}\par
\color{black} Łuczak, Tomasz and Pyber, László. 1993. Combinatorics, Probability and Computing. 2(4). pp. 505-512.\par
{\small\sffamily\color{atlasgold}Published work / Published}\par
\color{black} Supplies the permutation-group principle behind several probabilistic Galois arguments: a random permutation rarely belongs to a proper transitive subgroup other than A\_\allowbreak{}n. Invariable generation is relevant because modular factorization determines conjugacy classes, rather than labeled permutations.\par
{\footnotesize\url{https://www.cambridge.org/core/journals/combinatorics-probability-and-computing/article/abs/on-random-generation-of-the-symmetric-group/C6FE584196D2BE56BB85CAE23BEDF59D}}\par
{\footnotesize\url{https://doi.org/10.1017/S0963548300000869}}\par
\vspace{9pt}
\needspace{7\baselineskip}\hypertarget{ref:EberhardFordKoukoulopoulos2016}{}
{\sffamily\bfseries\color{atlasgreen}[EberhardFordKoukoulopoulos2016] Permutations contained in transitive subgroups}\par
\color{black} Eberhard, Sean and Ford, Kevin and Koukoulopoulos, Dimitris. 2016. Discrete Analysis. 2016. Article 12.\par
{\small\sffamily\color{atlasgold}Published work / Published}\par
\color{black} Sharpens quantitative control of permutations lying in transitive subgroups other than S\_\allowbreak{}n or A\_\allowbreak{}n. Separates imprimitive and primitive obstructions, making it useful in arguments that first establish irreducibility and then force a large Galois group.\par
{\small Paper No. 12}\par
{\footnotesize\url{https://discreteanalysisjournal.com/article/849-permutations-contained-in-transitive-subgroups}}\par
{\footnotesize\url{https://doi.org/10.19086/da.849}}\par
{\footnotesize\url{https://arxiv.org/abs/1605.01068}}\par
\vspace{9pt}
\needspace{7\baselineskip}\hypertarget{ref:EberhardFordGreen2017}{}
{\sffamily\bfseries\color{atlasgreen}[EberhardFordGreen2017] Invariable generation of the symmetric group}\par
\color{black} Eberhard, Sean and Ford, Kevin and Green, Ben. 2017. Duke Mathematical Journal. 166(8). pp. 1573-1590.\par
{\small\sffamily\color{atlasgold}Published work / Published}\par
\color{black} Shows that three uniform random permutations fail to invariably generate S\_\allowbreak{}n with probability tending to one, whereas four succeed with probability bounded away from zero. Explains why information from a small fixed collection of modular factorizations has a nontrivial combinatorial threshold.\par
{\footnotesize\url{https://arxiv.org/abs/1508.01870}}\par
{\footnotesize\url{https://doi.org/10.1215/00127094-0000007X}}\par
{\footnotesize\url{https://pure.qub.ac.uk/en/publications/invariable-generation-of-the-symmetric-group/}}\par
\vspace{9pt}
\needspace{7\baselineskip}\hypertarget{ref:PeledSenZeitouni2016}{}
{\sffamily\bfseries\color{atlasgreen}[PeledSenZeitouni2016] Double roots of random Littlewood polynomials}\par
\color{black} Peled, Ron and Sen, Arnab and Zeitouni, Ofer. 2016. Israel Journal of Mathematics. 213(1). pp. 55-77.\par
{\small\sffamily\color{atlasgold}Published work / Published}\par
\color{black} Gives sharp parity-sensitive double-root asymptotics for fair-sign Littlewood polynomials. Separates the easier squarefreeness question from irreducibility and explains the dominant repeated roots at \ensuremath{\pm}1.\par
{\small Theorem 1.2, specialized to fair signs; publication metadata confirmed by the publisher, arXiv and the author's institutional record. Publisher and institutional abstract rendering interchanges formulas, so mathematical statements were checked against the full paper.}\par
{\footnotesize\url{https://arxiv.org/abs/1409.2034}}\par
{\footnotesize\url{https://doi.org/10.1007/s11856-016-1328-3}}\par
{\footnotesize\url{https://arxiv.org/html/1409.2034v2}}\par
{\footnotesize\url{https://experts.umn.edu/en/publications/double-roots-of-random-littlewood-polynomials/}}\par
{\footnotesize\url{https://link.springer.com/article/10.1007/s11856-016-1328-3}}\par
\vspace{9pt}
\needspace{7\baselineskip}\hypertarget{ref:FeldheimSen2017}{}
{\sffamily\bfseries\color{atlasgreen}[FeldheimSen2017] Double roots of random polynomials with integer coefficients}\par
\color{black} Feldheim, Ohad N. and Sen, Arnab. 2017. Electronic Journal of Probability. 22. pp. 1-23. Article 10.\par
{\small\sffamily\color{atlasgold}Published work / Published}\par
\color{black} For bounded iid integer coefficients with largest atom at most 1/2, locates repeated-root probability at 0 and \ensuremath{\pm}1 up to o(n\textasciicircum{}(-2)); when zero is excluded, obtains O(n\textasciicircum{}(-2)). A useful safeguard against forgetting the persistent zero-discriminant event.\par
{\small Theorem 1.1. The model randomizes every coefficient, including the leading one; it is not stated as a fixed-monic theorem.}\par
{\footnotesize\url{https://arxiv.org/abs/1603.03811}}\par
{\footnotesize\url{https://doi.org/10.1214/17-EJP24}}\par
{\footnotesize\url{https://arxiv.org/html/1603.03811v1}}\par
{\footnotesize\url{https://experts.umn.edu/en/publications/double-roots-of-random-polynomials-with-integer-coefficients/}}\par
\vspace{9pt}
\needspace{7\baselineskip}\hypertarget{ref:Kuba2009}{}
{\sffamily\bfseries\color{atlasgreen}[Kuba2009] On the distribution of reducible polynomials}\par
\color{black} Kuba, Gerald. 2009. Mathematica Slovaca. 59(3). pp. 349-356.\par
{\small\sffamily\color{atlasgold}Published work / Published}\par
\color{black} Establishes sharp orders for reducible, not necessarily monic integer polynomials of fixed degree and bounded coefficient height: H\textasciicircum{}n in degree n\ensuremath{\ge}3, and H\textasciicircum{}2 log H in degree two. Complements monic counting theorems and isolates the contribution of linear factors.\par
{\small Published in 2009; the author's arXiv copy was deposited in 2020. Theorems 1-4 in that copy.}\par
{\footnotesize\url{https://www.sav.sk/index.php?doc=journal-list&journal_article_no=5434&lang=en&part=article_response_page}}\par
{\footnotesize\url{https://doi.org/10.2478/s12175-009-0131-6}}\par
{\footnotesize\url{https://arxiv.org/abs/2008.05815}}\par
{\footnotesize\url{https://arxiv.org/pdf/2008.05815}}\par
\vspace{9pt}
\needspace{7\baselineskip}\hypertarget{ref:BenPorath2026RootSubsets}{}
{\sffamily\bfseries\color{atlasgreen}[BenPorath2026RootSubsets] Counting Polynomials via Galois Actions on Root Subsets}\par
\color{black} Ben-Porath, Or. 2026.\par
{\small\sffamily\color{atlasgold}Preprint / Preprint; arXiv version 1}\par
\color{black} Obtains upper bounds for monic bounded-height polynomials with specified permutation representations, including actions on subsets, ordered tuples, and regular actions. The method recovers splitting fields from carefully chosen root-subset stabilizers and controls heights of their generators.\par
{\small arXiv:2603.14617v1, submitted 15 March 2026}\par
{\small Theorems 1.1-1.4; Lemma 3.1 and Proposition 3.3 explain the root-subset mechanism. The HTML title page displays 24 August 2026 although the arXiv record lists only v1 submitted 15 March 2026; version identifier is retained explicitly.}\par
{\footnotesize\url{https://arxiv.org/abs/2603.14617}}\par
{\footnotesize\url{https://arxiv.org/html/2603.14617v1}}\par
\vspace{9pt}
\needspace{7\baselineskip}\hypertarget{ref:EntinPopov2024}{}
{\sffamily\bfseries\color{atlasgreen}[EntinPopov2024] Probabilistic Galois theory in function fields}\par
\color{black} Entin, Alexei and Popov, Alexander. 2024. Finite Fields and Their Applications. 98. Article 102466.\par
{\small\sffamily\color{atlasgold}Published work / Published; full-group conclusion in the fixed-coefficient-degree model conditional}\par
\color{black} Studies monic polynomials in y with independent uniformly bounded-degree coefficients in F\_\allowbreak{}q[x]. The two limiting regimes differ: growing coefficient degree gives S\_\allowbreak{}n with probability tending to one; growing y-degree gives irreducibility limit 1\ensuremath{-}q\textasciicircum{}(\ensuremath{-}d) and conditional A\_\allowbreak{}n-or-S\_\allowbreak{}n. The full-S\_\allowbreak{}n refinement in odd characteristic uses a uniform polynomial Chowla conjecture.\par
{\small Theorems 1.1, 1.2, 1.4 and 1.6 checked in arXiv v3 (4 July 2024); the precise uniformity required for the conditional result is Conjecture 1.5. The publisher article endpoint could not be reopened during the 12 November 2026 second pass; publication metadata was retained from the earlier primary-publisher check.}\par
{\footnotesize\url{https://www.sciencedirect.com/science/article/pii/S1071579724001059}}\par
{\footnotesize\url{https://doi.org/10.1016/j.ffa.2024.102466}}\par
{\footnotesize\url{https://arxiv.org/abs/2311.14862}}\par
{\footnotesize\url{https://arxiv.org/abs/2311.14862v3}}\par
{\footnotesize\url{https://arxiv.org/html/2311.14862v3}}\par
\vspace{9pt}
\needspace{7\baselineskip}\hypertarget{ref:MichelenYakir2025Discriminant}{}
{\sffamily\bfseries\color{atlasgreen}[MichelenYakir2025Discriminant] Law of large numbers for the discriminant of random polynomials}\par
\color{black} Michelen, Marcus and Yakir, Oren. 2025.\par
{\small\sffamily\color{atlasgold}Preprint / Preprint; version 1 dated 13 June 2025}\par
\color{black} Proves a universal first-order law for log absolute discriminant of iid mean-zero, variance-one sub-Gaussian Kac polynomials with no zero coefficient atom. Connects Mahler measure, reciprocal symmetry in distribution and root separation. A size asymptotic does not determine the discriminant's square class.\par
{\small arXiv:2506.12206v1}\par
{\small Theorem 1, equation (10), and §1.3. The abstract omits hypotheses that are explicit in Theorem 1; the atlas preserves them.}\par
{\footnotesize\url{https://arxiv.org/abs/2506.12206}}\par
{\footnotesize\url{https://arxiv.org/html/2506.12206v1}}\par
{\footnotesize\url{https://marcusmichelen.org/research.php}}\par
\vspace{9pt}
\needspace{7\baselineskip}\hypertarget{ref:MichelenYakir2026RootSeparation}{}
{\sffamily\bfseries\color{atlasgreen}[MichelenYakir2026RootSeparation] Limit law for root separation in random polynomials}\par
\color{black} Michelen, Marcus and Yakir, Oren. 2026. Advances in Mathematics. 496. Article 110990.\par
{\small\sffamily\color{atlasgold}Published work / Published}\par
\color{black} Establishes a limiting law for the smallest root separation at scale n\textasciicircum{}(-5/4) for mean-zero sub-Gaussian iid coefficients with no atom at zero. Corollary 1.2 implies squarefreeness with probability tending to one without the older integer-law largest-atom restriction.\par
{\small Theorem 1.1 and Corollary 1.2; published in June 2026. The sharp Littlewood repeated-root asymptotic of Peled-Sen-Zeitouni supplies a quantitative conclusion not contained in the limit statement alone.}\par
{\footnotesize\url{https://www.sciencedirect.com/science/article/pii/S0001870826002124}}\par
{\footnotesize\url{https://doi.org/10.1016/j.aim.2026.110990}}\par
{\footnotesize\url{https://arxiv.org/abs/2505.02723}}\par
{\footnotesize\url{https://arxiv.org/html/2505.02723v1}}\par
{\footnotesize\url{https://www.sciencedirect.com/journal/advances-in-mathematics/vol/496/suppl/C}}\par
\vspace{9pt}
\needspace{12\baselineskip}\section{Hilbert irreducibility and inverse Galois theory}
\needspace{7\baselineskip}\hypertarget{ref:Hilbert1892}{}
{\sffamily\bfseries\color{atlasgreen}[Hilbert1892] Ueber die Irreducibilität ganzer rationaler Functionen mit ganzzahligen Coefficienten}\par
\color{black} Hilbert, David. 1892. Journal für die reine und angewandte Mathematik. 110. pp. 104-129.\par
{\small\sffamily\color{atlasgold}Published work / Published article}\par
\color{black} The original source for Hilbert's irreducibility theorem. It is the historical starting point for passing from irreducible polynomial families over a rational function field to irreducible arithmetic specializations; use a modern treatment for the general language of Hilbertian fields.\par
{\footnotesize\url{https://www.degruyterbrill.com/document/doi/10.1515/crll.1892.110.104/html}}\par
{\footnotesize\url{https://doi.org/10.1515/crll.1892.110.104}}\par
\vspace{9pt}
\needspace{7\baselineskip}\hypertarget{ref:DvornicichZannier1994}{}
{\sffamily\bfseries\color{atlasgreen}[DvornicichZannier1994] Fields containing values of algebraic functions}\par
\color{black} Dvornicich, Roberto and Zannier, Umberto. 1994. Annali della Scuola Normale Superiore di Pisa, Classe di Scienze (4). 21(3). pp. 421-443.\par
{\small\sffamily\color{atlasgold}Published work / Published article}\par
\color{black} Studies the number fields generated by specialized values of algebraic functions. A useful bridge from the existence of irreducible specializations to questions about the diversity of the fields produced.\par
{\footnotesize\url{https://www.numdam.org/item/ASNSP_1994_4_21_3_421_0/}}\par
\vspace{9pt}
\needspace{7\baselineskip}\hypertarget{ref:HajirWong2006}{}
{\sffamily\bfseries\color{atlasgreen}[HajirWong2006] Specializations of one-parameter families of polynomials}\par
\color{black} Hajir, Farshid and Wong, Siman. 2006. Annales de l'Institut Fourier. 56(4). pp. 1127-1163.\par
{\small\sffamily\color{atlasgold}Published work / Published article}\par
\color{black} Develops geometric and group-theoretic conditions that make the exceptional set of reducible specializations over a number field finite. For fixed n \ensuremath{\ge} 10, all but finitely many parameter values in any fixed number field give irreducible generalized Laguerre polynomials with Galois group S\_\allowbreak{}n.\par
{\footnotesize\url{https://www.numdam.org/item/AIF_2006__56_4_1127_0/}}\par
{\footnotesize\url{https://doi.org/10.5802/aif.2208}}\par
{\footnotesize\url{https://arxiv.org/abs/math/0405139}}\par
{\footnotesize\url{https://www.numdam.org/item/10.5802/aif.2208.pdf}}\par
\vspace{9pt}
\needspace{7\baselineskip}\hypertarget{ref:DvornicichZannier2007}{}
{\sffamily\bfseries\color{atlasgreen}[DvornicichZannier2007] Cyclotomic Diophantine problems (Hilbert irreducibility and invariant sets for polynomial maps)}\par
\color{black} Dvornicich, Roberto and Zannier, Umberto. 2007. Duke Mathematical Journal. 139(3). pp. 527-554.\par
{\small\sffamily\color{atlasgold}Published work / Published article}\par
\color{black} Develops Hilbert irreducibility in cyclotomic extensions using the arithmetic of torsion points on tori. Relevant when specialization parameters are roots of unity or when polynomial dynamics is restricted to cyclotomic fields; its structural hypotheses must be retained.\par
{\footnotesize\url{https://arpi.unipi.it/handle/11568/112812}}\par
{\footnotesize\url{https://doi.org/10.1215/S0012-7094-07-13934-6}}\par
{\footnotesize\url{https://ror.circle-u.eu/items/44e0b4f8-aef4-40f1-8da3-2a2f2391946c/full}}\par
{\footnotesize\url{https://www.ae-info.org/ae/Member/Zannier_Umberto/Publications}}\par
\vspace{9pt}
\needspace{7\baselineskip}\hypertarget{ref:Serre2007Topics}{}
{\sffamily\bfseries\color{atlasgreen}[Serre2007Topics] Topics in Galois Theory}\par
\color{black} Serre, Jean-Pierre. 2007. A K Peters/CRC Press, Research Notes in Mathematics 1. Edition 2. A K Peters/CRC Press.\par
{\small\sffamily\color{atlasgold}Book / Published monograph}\par
\color{black} A compact route from Hilbert irreducibility and thin sets to explicit Galois realizations, rigidity, and sieve methods. Chapter 3 and Section 4.6 are particularly useful for explaining why generic Galois groups survive many specializations. Section 4.4 proves the Morse criterion through inertia and Riemann-Hurwitz, providing a geometric route to full symmetric monodromy.\par
{\small Notes written by Henri Darmon. Publisher electronic edition copyright 2007; the physical second edition is also commonly catalogued and cited as 2008.}\par
{\footnotesize\url{https://www.routledge.com/Topics-in-Galois-Theory-2nd-Edition/Serre/p/book/9780429064401}}\par
{\footnotesize\url{https://doi.org/10.1201/b10588}}\par
{\footnotesize\url{https://api.pageplace.de/preview/DT0400.9781439865255_A38251687/preview-9781439865255_A38251687.pdf}}\par
{\footnotesize\url{https://indico.math.cnrs.fr/event/11410/attachments/4760/7351/SerreTopicsGaloisTheory.pdf}}\par
\vspace{9pt}
\needspace{7\baselineskip}\hypertarget{ref:Zannier2010}{}
{\sffamily\bfseries\color{atlasgreen}[Zannier2010] Hilbert irreducibility above algebraic groups}\par
\color{black} Zannier, Umberto. 2010. Duke Mathematical Journal. 153(2). pp. 397-425.\par
{\small\sffamily\color{atlasgold}Published work / Published article}\par
\color{black} Studies irreducible fibers above Zariski-dense cyclic subgroups of algebraic tori and powers of an elliptic curve without complex multiplication, under a pullback irreducibility condition. The method combines torsion specialization with descent using Chebotarev, and yields a Bertini-type result for tori.\par
{\footnotesize\url{https://arxiv.org/abs/1001.3079}}\par
{\footnotesize\url{https://doi.org/10.1215/00127094-2010-027}}\par
{\footnotesize\url{https://ricerca.sns.it/handle/11384/64}}\par
\vspace{9pt}
\needspace{7\baselineskip}\hypertarget{ref:DebesLegrand2013}{}
{\sffamily\bfseries\color{atlasgreen}[DebesLegrand2013] Specialization results in Galois theory}\par
\color{black} Dèbes, Pierre and Legrand, François. 2013. Transactions of the American Mathematical Society. 365(10). pp. 5259-5275.\par
{\small\sffamily\color{atlasgold}Published work / Published article}\par
\color{black} Uses twisting to turn prescribed specialization questions into rational-point questions. One application combines irreducible specialization with prescribed residue-degree partitions at finitely many sufficiently large primes for covers over Q with geometric monodromy S\_\allowbreak{}n.\par
{\footnotesize\url{https://arxiv.org/abs/1106.6151}}\par
{\footnotesize\url{https://doi.org/10.1090/S0002-9947-2013-05800-X}}\par
{\footnotesize\url{https://cris.technion.ac.il/en/publications/specialization-results-in-galois-theory/}}\par
{\footnotesize\url{https://pro.univ-lille.fr/fileadmin/user_upload/pages_pros/pierre_debes/A52-DebesLegrand1.pdf}}\par
\vspace{9pt}
\needspace{7\baselineskip}\hypertarget{ref:FriedJarden2023}{}
{\sffamily\bfseries\color{atlasgreen}[FriedJarden2023] Field Arithmetic}\par
\color{black} Fried, Michael D. and Jarden, Moshe. 2023. Springer, Ergebnisse der Mathematik und ihrer Grenzgebiete, 3. Folge, 11. Edition 4. Springer.\par
{\small\sffamily\color{atlasgold}Book / Published monograph}\par
\color{black} A major reference for Hilbertian fields, absolute Galois groups, specialization, and PAC fields. The fourth edition updates the earlier editions and includes a new chapter on Hilbertian subfields of Galois extensions.\par
{\footnotesize\url{https://link.springer.com/book/10.1007/978-3-031-28020-7}}\par
{\footnotesize\url{https://doi.org/10.1007/978-3-031-28020-7}}\par
\vspace{9pt}
\needspace{7\baselineskip}\hypertarget{ref:Volklein1996}{}
{\sffamily\bfseries\color{atlasgreen}[Volklein1996] Groups as Galois Groups: An Introduction}\par
\color{black} Völklein, Helmut. 1996. Cambridge Studies in Advanced Mathematics 53. Cambridge University Press.\par
{\small\sffamily\color{atlasgold}Book / Published monograph}\par
\color{black} Introduces the geometric approach to inverse Galois theory through covers, Riemann's existence theorem, and rigidity. Especially useful for learning how permutation data can be converted into polynomial families with a desired group.\par
{\footnotesize\url{https://www.cambridge.org/core/books/groups-as-galois-groups/A2233CBEE2CD250DDE2B867502C0A4DE}}\par
{\footnotesize\url{https://doi.org/10.1017/CBO9780511471117}}\par
{\footnotesize\url{https://assets.cambridge.org/97805215/62805/frontmatter/9780521562805_frontmatter.pdf}}\par
\vspace{9pt}
\needspace{7\baselineskip}\hypertarget{ref:JensenLedetYui2002}{}
{\sffamily\bfseries\color{atlasgreen}[JensenLedetYui2002] Generic Polynomials: Constructive Aspects of the Inverse Galois Problem}\par
\color{black} Jensen, Christian U. and Ledet, Arne and Yui, Noriko. 2002. Mathematical Sciences Research Institute Publications 45. Cambridge University Press.\par
{\small\sffamily\color{atlasgold}Book / Published monograph}\par
\color{black} A constructive source for generic polynomials and the invariant-theoretic methods behind them. The freely accessible MSRI/SLMath page supplies downloadable chapters and concrete families for specified finite groups.\par
{\footnotesize\url{https://library.slmath.org/books/Book45/index.html}}\par
{\footnotesize\url{https://library.slmath.org/books/Book45/desc.html}}\par
\vspace{9pt}
\needspace{7\baselineskip}\hypertarget{ref:MalleMatzat2018}{}
{\sffamily\bfseries\color{atlasgreen}[MalleMatzat2018] Inverse Galois Theory}\par
\color{black} Malle, Gunter and Matzat, B. Heinrich. 2018. Springer Monographs in Mathematics. Edition 2. Springer.\par
{\small\sffamily\color{atlasgold}Book / Published monograph}\par
\color{black} A systematic reference for constructing Galois extensions through rigidity, braid actions, embedding problems, additive polynomials, and rigid analytic methods. It also includes tables of example polynomials for transitive and primitive groups of small degree.\par
{\footnotesize\url{https://link.springer.com/book/10.1007/978-3-662-55420-3}}\par
{\footnotesize\url{https://doi.org/10.1007/978-3-662-55420-3}}\par
\vspace{9pt}
\needspace{7\baselineskip}\hypertarget{ref:SiksekStoll2026}{}
{\sffamily\bfseries\color{atlasgreen}[SiksekStoll2026] Hilbert's Irreducibility for \(\mathbb{G}_m\)}\par
\color{black} Siksek, Samir and Stoll, Michael. 2026. arXiv preprint.\par
{\small\sffamily\color{atlasgold}Preprint / Preprint; version 1, 3 September 2026}\par
\color{black} Describes exceptional S-unit fibers of a curve cover as a finite set together with finitely many power cosets. Theorem3 makes the cosets effective but leaves the finite remainder ineffective; Corollary1 characterizes infinitely many reducible f(x)\ensuremath{-}p\_\allowbreak{}1\textasciicircum{}\{e\_\allowbreak{}1\}\ensuremath{\cdots}p\_\allowbreak{}s\textasciicircum{}\{e\_\allowbreak{}s\}.\par
{\small Source title uses the notation G\_\allowbreak{}m for the multiplicative group; author order in the PDF is Samir Siksek and Michael Stoll.}\par
{\footnotesize\url{https://arxiv.org/abs/2609.04551v1}}\par
{\footnotesize\url{https://arxiv.org/html/2609.04551v1}}\par
{\footnotesize\url{https://arxiv.org/abs/2609.04551}}\par
{\footnotesize\url{https://www.mathe2.uni-bayreuth.de/stoll/schrift.html}}\par
\vspace{9pt}
\needspace{7\baselineskip}\hypertarget{ref:BarySorokerGarzoniMatei2025Powers}{}
{\sffamily\bfseries\color{atlasgreen}[BarySorokerGarzoniMatei2025Powers] On the irreducibility of f(2\textasciicircum{}n,3\textasciicircum{}m,X) and other such polynomials}\par
\color{black} Bary-Soroker, Lior and Garzoni, Daniele and Matei, Vlad. 2025. Mathematische Zeitschrift. 310. pp. Article 11.\par
{\small\sffamily\color{atlasgold}Published work / Published article}\par
\color{black} Under GRH and the geometric pullback condition (PB), Theorem 1.1 proves density-one irreducibility along independent power parameters. The power-pullback condition is stronger than generic irreducibility and excludes X\ensuremath{{}^2}\ensuremath{-}t. Theorem 1.1 counts exponent vectors in symmetric boxes, including negative exponents.\par
{\footnotesize\url{https://doi.org/10.1007/s00209-025-03717-7}}\par
{\footnotesize\url{https://link.springer.com/article/10.1007/s00209-025-03717-7}}\par
{\footnotesize\url{https://arxiv.org/abs/2405.04058}}\par
\vspace{9pt}
\needspace{7\baselineskip}\hypertarget{ref:Debes1992Powers}{}
{\sffamily\bfseries\color{atlasgreen}[Debes1992Powers] On the irreducibility of the polynomials P(t\textasciicircum{}m, Y)}\par
\color{black} Dèbes, Pierre. 1992. Journal of Number Theory. 42(2). pp. 141-157.\par
{\small\sffamily\color{atlasgold}Published work / Published article}\par
\color{black} Characterizes infinite reducibility along a geometric progression by a reducible scaled power pullback (Theorem 1.2). Under the geometric pullback condition (PB), one obtains irreducibility for all but finitely many integer exponents, unconditionally. The argument uses Siegel's theorem; finiteness is distinct from effective enumeration of exceptional exponents.\par
{\small Original scanned author-hosted paper inspected visually: Theorem 1.2 is on printed page 144. Publication metadata verified against the publisher and author list.}\par
{\footnotesize\url{https://doi.org/10.1016/0022-314X(92)90018-K}}\par
{\footnotesize\url{https://pro.univ-lille.fr/fileadmin/user_upload/pages_pros/pierre_debes/A12-JNT1.pdf}}\par
{\footnotesize\url{https://pro.univ-lille.fr/pierre-debes/publications}}\par
{\footnotesize\url{https://www.sciencedirect.com/science/article/pii/0022314X9290018K}}\par
{\footnotesize\url{https://link.springer.com/article/10.1007/s00209-025-03717-7}}\par
\vspace{9pt}
\needspace{12\baselineskip}\section{Composition, iteration, and arboreal groups}
\needspace{7\baselineskip}\hypertarget{ref:Odoni1985}{}
{\sffamily\bfseries\color{atlasgreen}[Odoni1985] The Galois Theory of Iterates and Composites of Polynomials}\par
\color{black} Odoni, R. W. K. 1985. Proceedings of the London Mathematical Society (3). 51(3). pp. 385-414.\par
{\small\sffamily\color{atlasgold}Published work / Published article}\par
\color{black} A foundational treatment of Galois groups under composition and iteration. In characteristic zero, the iterates of the generic degree-d polynomial have the full iterated wreath-product group, leading to finite-level specialization results and the later arboreal questions.\par
{\footnotesize\url{https://londmathsoc.onlinelibrary.wiley.com/doi/10.1112/plms/s3-51.3.385}}\par
{\footnotesize\url{https://doi.org/10.1112/plms/s3-51.3.385}}\par
{\footnotesize\url{https://arxiv.org/pdf/1410.3814}}\par
\vspace{9pt}
\needspace{7\baselineskip}\hypertarget{ref:Stoll1992}{}
{\sffamily\bfseries\color{atlasgreen}[Stoll1992] Galois groups over Q of some iterated polynomials}\par
\color{black} Stoll, Michael. 1992. Archiv der Mathematik. 59. pp. 239-244.\par
{\small\sffamily\color{atlasgold}Published work / Published article}\par
\color{black} Determines maximality of iterated quadratic Galois groups through arithmetic square-class conditions. A useful concrete model for how critical-orbit information controls the growth of splitting fields under iteration.\par
{\footnotesize\url{https://link.springer.com/article/10.1007/BF01197321}}\par
{\footnotesize\url{https://doi.org/10.1007/BF01197321}}\par
{\footnotesize\url{https://www.mathe2.uni-bayreuth.de/stoll/schrift.html}}\par
\vspace{9pt}
\needspace{7\baselineskip}\hypertarget{ref:AyadMcQuillan2000}{}
{\sffamily\bfseries\color{atlasgreen}[AyadMcQuillan2000] Irreducibility of the iterates of a quadratic polynomial over a field}\par
\color{black} Ayad, Mohamed and McQuillan, Donald L. 2000. Acta Arithmetica. 93(1). pp. 87-97.\par
{\small\sffamily\color{atlasgold}Published work / Published article; correction available}\par
\color{black} Develops criteria for irreducibility of successive quadratic iterates. Read it together with the authors' 2001 correction before using the stated tests in a proof or computation.\par
{\footnotesize\url{https://doi.org/10.4064/aa-93-1-87-97}}\par
{\footnotesize\url{https://www.impan.pl/en/publishing-house/journals-and-series/acta-arithmetica/all/93/1}}\par
{\footnotesize\url{https://geodesic.mathdoc.fr/articles/10.4064/aa99-1-9/}}\par
\vspace{9pt}
\needspace{7\baselineskip}\hypertarget{ref:AyadMcQuillan2001Correction}{}
{\sffamily\bfseries\color{atlasgreen}[AyadMcQuillan2001Correction] Corrections to "Irreducibility of the iterates of a quadratic polynomial over a field" (Acta Arith. 93 (2000), 87--97)}\par
\color{black} Ayad, Mohamed and McQuillan, Donald L. 2001. Acta Arithmetica. 99(1). pp. 97.\par
{\small\sffamily\color{atlasgold}Published work / Published correction}\par
\color{black} The published correction to Ayad and McQuillan's quadratic-iteration paper. Keep this record linked to the original article when building a methods database.\par
{\footnotesize\url{https://geodesic.mathdoc.fr/articles/10.4064/aa99-1-9/}}\par
{\footnotesize\url{https://doi.org/10.4064/aa99-1-9}}\par
\vspace{9pt}
\needspace{7\baselineskip}\hypertarget{ref:Jones2012EveryPrime}{}
{\sffamily\bfseries\color{atlasgreen}[Jones2012EveryPrime] An iterative construction of irreducible polynomials reducible modulo every prime}\par
\color{black} Jones, Rafe. 2012. Journal of Algebra. 369. pp. 114-128.\par
{\small\sffamily\color{atlasgold}Published work / Published article}\par
\color{black} Constructs polynomials of arbitrarily large degree that are irreducible over Q, and over suitable other global fields, but reducible modulo every prime. Failure to find an irreducible reduction is therefore not evidence of reducibility over the ground field.\par
{\footnotesize\url{https://arxiv.org/abs/1012.2857}}\par
{\footnotesize\url{https://doi.org/10.1016/j.jalgebra.2012.05.020}}\par
{\footnotesize\url{https://www.sciencedirect.com/science/article/pii/S0021869312002918}}\par
{\footnotesize\url{https://doi.org/10.1090/bull/1665}}\par
\vspace{9pt}
\needspace{7\baselineskip}\hypertarget{ref:Jones2013Arboreal}{}
{\sffamily\bfseries\color{atlasgreen}[Jones2013Arboreal] Galois representations from pre-image trees: an arboreal survey}\par
\color{black} Jones, Rafe. 2013. Publications mathématiques de Besançon. Algèbre et théorie des nombres. pp. 107-136.\par
{\small\sffamily\color{atlasgold}Published work / Published survey}\par
\color{black} A readable entry point to Galois actions on trees of iterated preimages, stability, and density questions. Use its bibliography for the foundational literature, then supplement its open-problem discussion with later work.\par
{\footnotesize\url{https://www.numdam.org/articles/10.5802/pmb.a-154/}}\par
{\footnotesize\url{https://doi.org/10.5802/pmb.a-154}}\par
{\footnotesize\url{https://arxiv.org/abs/1402.6018}}\par
\vspace{9pt}
\needspace{7\baselineskip}\hypertarget{ref:Specter2018}{}
{\sffamily\bfseries\color{atlasgreen}[Specter2018] Polynomials with Surjective Arboreal Galois Representations Exist in Every Degree}\par
\color{black} Specter, Joel. 2018. arXiv preprint.\par
{\small\sffamily\color{atlasgold}Preprint / Preprint; version 2, 12 March 2018}\par
\color{black} Gives the number-field form of Odoni's existence conjecture in all degrees. The source checked is the 2018 preprint, cited as such in subsequent literature; this must be distinguished from the false assertion over all Hilbertian fields.\par
{\footnotesize\url{https://arxiv.org/abs/1803.00434v2}}\par
{\footnotesize\url{https://arxiv.org/abs/1803.00434}}\par
{\footnotesize\url{https://math.jhu.edu/~jspecter/Odoni.pdf}}\par
{\footnotesize\url{https://arxiv.org/pdf/2012.03076}}\par
\vspace{9pt}
\needspace{7\baselineskip}\hypertarget{ref:BenedettoJuul2019}{}
{\sffamily\bfseries\color{atlasgreen}[BenedettoJuul2019] Odoni's conjecture for number fields}\par
\color{black} Benedetto, Robert L. and Juul, Jamie. 2019. Bulletin of the London Mathematical Society. 51(2). pp. 237-250.\par
{\small\sffamily\color{atlasgold}Published work / Published article}\par
\color{black} Proves the desired existence of full arboreal Galois groups for even degree over every number field and for odd degree when the number field has odd degree over Q. The title should not be read as asserting every number-field case within this paper alone.\par
{\footnotesize\url{https://londmathsoc.onlinelibrary.wiley.com/doi/10.1112/blms.12225}}\par
{\footnotesize\url{https://doi.org/10.1112/blms.12225}}\par
{\footnotesize\url{https://arxiv.org/abs/1803.01987}}\par
\vspace{9pt}
\needspace{7\baselineskip}\hypertarget{ref:Juul2019}{}
{\sffamily\bfseries\color{atlasgreen}[Juul2019] Iterates of Generic Polynomials and Generic Rational Functions}\par
\color{black} Juul, Jamie. 2019. Transactions of the American Mathematical Society. 371(2). pp. 809-831.\par
{\small\sffamily\color{atlasgold}Published work / Published article}\par
\color{black} Extends the generic full wreath-product result to arbitrary characteristic for rational functions and to polynomials except the degree-2, characteristic-2 case. The paper also records the contributions of Odoni's unpublished work to these results.\par
{\footnotesize\url{https://arxiv.org/abs/1410.3814}}\par
{\footnotesize\url{https://arxiv.org/abs/1410.3814v6}}\par
{\footnotesize\url{https://doi.org/10.1090/tran/7229}}\par
{\footnotesize\url{https://arxiv.org/pdf/1410.3814}}\par
{\footnotesize\url{https://sites.google.com/view/jamiejuul/research}}\par
\vspace{9pt}
\needspace{7\baselineskip}\hypertarget{ref:Looper2019}{}
{\sffamily\bfseries\color{atlasgreen}[Looper2019] Dynamical Galois groups of trinomials and Odoni's conjecture}\par
\color{black} Looper, Nicole R. 2019. Bulletin of the London Mathematical Society. 51(2). pp. 278-292.\par
{\small\sffamily\color{atlasgold}Published work / Published article}\par
\color{black} Constructs integral polynomials with surjective arboreal representations in every prime degree. It also gives conditional results in certain composite degrees using Vojta's conjecture, so the unconditional and conditional statements should be separated.\par
{\footnotesize\url{https://arxiv.org/abs/1609.03398}}\par
{\footnotesize\url{https://doi.org/10.1112/blms.12227}}\par
{\footnotesize\url{https://londmathsoc.onlinelibrary.wiley.com/doi/full/10.1112/blms.12227}}\par
{\footnotesize\url{https://sites.google.com/view/nicole-looper/research}}\par
\vspace{9pt}
\needspace{7\baselineskip}\hypertarget{ref:DittmannKadets2022}{}
{\sffamily\bfseries\color{atlasgreen}[DittmannKadets2022] Odoni's conjecture on arboreal Galois representations is false}\par
\color{black} Dittmann, Philip and Kadets, Borys. 2022. Proceedings of the American Mathematical Society. 150(8). pp. 3335-3343.\par
{\small\sffamily\color{atlasgold}Published work / Published article}\par
\color{black} Disproves Odoni's conjecture in its general formulation over Hilbertian fields of characteristic zero. It does not contradict positive existence results over number fields; the difference in the base-field quantifier is essential.\par
{\footnotesize\url{https://arxiv.org/abs/2012.03076}}\par
{\footnotesize\url{https://arxiv.org/abs/2012.03076v2}}\par
{\footnotesize\url{https://doi.org/10.1090/proc/15920}}\par
{\footnotesize\url{https://arxiv.org/pdf/2012.03076}}\par
{\footnotesize\url{https://cris.huji.ac.il/en/publications/odonis-conjecture-on-arboreal-galois-representations-is-false-2/}}\par
\vspace{9pt}
\needspace{7\baselineskip}\hypertarget{ref:Ritt1922}{}
{\sffamily\bfseries\color{atlasgreen}[Ritt1922] Prime and composite polynomials}\par
\color{black} Ritt, Joseph F. 1922. Transactions of the American Mathematical Society. 23(1). pp. 51-66.\par
{\small\sffamily\color{atlasgold}Published work / Published article with erratum}\par
\color{black} The classical source for functional decomposition of complex polynomials into indecomposable components. Here prime means indecomposable under composition, not irreducible under multiplication; read together with the erratum in the same volume, p. 431.\par
{\footnotesize\url{https://www.jstor.org/stable/1988911}}\par
{\footnotesize\url{https://doi.org/10.1090/S0002-9947-1922-1501189-9}}\par
{\footnotesize\url{https://www.ams.org/journals/tran/1922-023-01/S0002-9947-1922-1501189-9/}}\par
{\footnotesize\url{https://doi.org/10.1090/S0002-9947-1922-1501205-4}}\par
\vspace{9pt}
\needspace{7\baselineskip}\hypertarget{ref:Mueller1995Monodromy}{}
{\sffamily\bfseries\color{atlasgreen}[Mueller1995Monodromy] Primitive monodromy groups of polynomials}\par
\color{black} Müller, Peter. 1995. Recent Developments in the Inverse Galois Problem. 186. pp. 385-401. American Mathematical Society. Contemporary Mathematics.\par
{\small\sffamily\color{atlasgold}Published work / Published book chapter; author-hosted version}\par
\color{black} Classifies geometric monodromy groups of indecomposable complex polynomials and discusses restrictions for polynomials with rational coefficients. The group is that of f(X)\ensuremath{-}t over a function field, rather than the splitting-field group of f(X) over the coefficient field.\par
{\footnotesize\url{https://ypfmde.github.io/Papers/mon.pdf}}\par
{\footnotesize\url{https://ypfmde.github.io/publications.html}}\par
{\footnotesize\url{https://doi.org/10.1090/conm/186}}\par
\vspace{9pt}
\needspace{7\baselineskip}\hypertarget{ref:ZieveMueller2008}{}
{\sffamily\bfseries\color{atlasgreen}[ZieveMueller2008] On Ritt's polynomial decomposition theorems}\par
\color{black} Zieve, Michael E. and Müller, Peter. 2008. arXiv preprint.\par
{\small\sffamily\color{atlasgold}Preprint / Preprint; no journal publication verified}\par
\color{black} Develops a structured description of all complete decompositions of a complex polynomial, refining the classical Ritt framework. Both authors' publication pages still list it as a preprint, so it should be cited with its arXiv identifier rather than an invented journal reference.\par
{\footnotesize\url{https://arxiv.org/abs/0807.3578}}\par
{\footnotesize\url{https://websites.umich.edu/~zieve/papers/peter.html}}\par
{\footnotesize\url{https://ypfmde.github.io/publications.html}}\par
\vspace{9pt}
\needspace{7\baselineskip}\hypertarget{ref:MuzychukPakovich2011}{}
{\sffamily\bfseries\color{atlasgreen}[MuzychukPakovich2011] Jordan-Hölder theorem for imprimitivity systems and maximal decompositions of rational functions}\par
\color{black} Muzychuk, Mikhail and Pakovich, Fedor. 2011. Proceedings of the London Mathematical Society. 102(1). pp. 1-24.\par
{\small\sffamily\color{atlasgold}Published work / Published article}\par
\color{black} Studies imprimitivity lattices for permutation groups with a cyclic subgroup having at most two orbits, with applications to decomposition. It extends aspects of Ritt's first theorem and explains why unrestricted rational-function analogues can fail.\par
{\footnotesize\url{https://doi.org/10.1112/plms/pdq009}}\par
{\footnotesize\url{https://academic.oup.com/plms/article-abstract/102/1/1/1478963}}\par
{\footnotesize\url{https://londmathsoc.onlinelibrary.wiley.com/doi/abs/10.1112/plms/pdq009}}\par
\vspace{9pt}
\needspace{7\baselineskip}\hypertarget{ref:Stoll2026Arboreal}{}
{\sffamily\bfseries\color{atlasgreen}[Stoll2026Arboreal] When do all iterates of x\ensuremath{{}^2}+a have maximal Galois group?}\par
\color{black} Stoll, Michael. 2026. arXiv preprint.\par
{\small\sffamily\color{atlasgold}Preprint / Preprint; version 1, 10 November 2026}\par
\color{black} Reworks quadratic critical-orbit criteria and extends explicit sufficient conditions to rational parameters. Theorem6.1 has region and congruence restrictions; Conjecture8.1 asks for uniform finite-level detection. Computational evidence and the reported formalization are distinguished from a general solution of that conjecture.\par
{\footnotesize\url{https://arxiv.org/abs/2611.14014v1}}\par
{\footnotesize\url{https://arxiv.org/html/2611.14014v1}}\par
{\footnotesize\url{https://arxiv.org/abs/2611.14014}}\par
{\footnotesize\url{https://www.mathe2.uni-bayreuth.de/stoll/papers/arboreal-surjective.pdf}}\par
{\footnotesize\url{https://www.mathe2.uni-bayreuth.de/stoll/schrift.html}}\par
{\footnotesize\url{https://arxiv.org/pdf/2611.14014v1}}\par
\vspace{9pt}
\needspace{7\baselineskip}\hypertarget{ref:BenedettoGhiocaJuulTucker2025Periodic}{}
{\sffamily\bfseries\color{atlasgreen}[BenedettoGhiocaJuulTucker2025Periodic] Arboreal Galois groups of postcritically finite quadratic polynomials: The periodic case}\par
\color{black} Benedetto, Robert L. and Ghioca, Dragos and Juul, Jamie and Tucker, Thomas J. 2025. Journal of the London Mathematical Society. 112(2). pp. e70257.\par
{\small\sffamily\color{atlasgold}Published work / Published article}\par
\color{black} Describes groups for x\ensuremath{{}^2}+c with periodic critical point in characteristic different from two. Theorem 1.5 gives equivalent square-root, finite-level, and all-level conditions for attaining the specified constrained group. This is not surjectivity to the full binary-tree automorphism group; the target point and cyclotomic constants are essential.\par
{\footnotesize\url{https://doi.org/10.1112/jlms.70257}}\par
{\footnotesize\url{https://arxiv.org/html/2411.06745v2}}\par
{\footnotesize\url{https://arxiv.org/abs/2411.06745}}\par
{\footnotesize\url{https://londmathsoc.onlinelibrary.wiley.com/doi/10.1112/jlms.70257}}\par
\vspace{9pt}
\needspace{7\baselineskip}\hypertarget{ref:BenedettoGhiocaJuulTucker2026Preperiodic}{}
{\sffamily\bfseries\color{atlasgreen}[BenedettoGhiocaJuulTucker2026Preperiodic] Arboreal Galois groups of postcritically finite quadratic polynomials: The strictly preperiodic case}\par
\color{black} Benedetto, Robert L. and Ghioca, Dragos and Juul, Jamie and Tucker, Thomas J. 2026. Canadian Journal of Mathematics.\par
{\small\sffamily\color{atlasgold}Published work / Published article}\par
\color{black} Complements the periodic case by treating strictly preperiodic critical orbits of x\ensuremath{{}^2}+c. Separates short-tail and long-tail cases and gives criteria for the resulting constrained groups. Section 1.4 excludes the already-classical Chebyshev map x\ensuremath{{}^2}\ensuremath{-}2, whose towers admit a Kummer description.\par
{\small Published online 30 March 2026; no unverified final volume or page range assigned. Mathematical locators refer to author preprint v1.}\par
{\footnotesize\url{https://doi.org/10.4153/S0008414X2610217X}}\par
{\footnotesize\url{https://arxiv.org/html/2507.08347v1}}\par
{\footnotesize\url{https://arxiv.org/abs/2507.08347}}\par
\vspace{9pt}
\needspace{12\baselineskip}\section{Finite fields: construction and distribution}
\needspace{7\baselineskip}\hypertarget{ref:Shoup1990Irreducible}{}
{\sffamily\bfseries\color{atlasgreen}[Shoup1990Irreducible] New algorithms for finding irreducible polynomials over finite fields}\par
\color{black} Shoup, Victor. 1990. Mathematics of Computation. 54. pp. 435-447.\par
{\small\sffamily\color{atlasgold}Published work / Published article; author-hosted version}\par
\color{black} Constructs irreducible polynomials of a prescribed degree, including deterministic polynomial-time results for small characteristic. Distinguish dependence polynomial in the characteristic from dependence polynomial in its bit length.\par
{\footnotesize\url{https://www.shoup.net/papers/detirred.pdf}}\par
{\footnotesize\url{https://doi.org/10.1090/S0025-5718-1990-0993933-0}}\par
{\footnotesize\url{https://www.shoup.net/papers/}}\par
\vspace{9pt}
\needspace{7\baselineskip}\hypertarget{ref:HansenMullen1992}{}
{\sffamily\bfseries\color{atlasgreen}[HansenMullen1992] Primitive polynomials over finite fields}\par
\color{black} Hansen, Tom and Mullen, Gary L. 1992. Mathematics of Computation. 59(200). pp. 639-643.\par
{\small\sffamily\color{atlasgold}Published work / Published article with supplementary tables}\par
\color{black} Provides explicit primitive-polynomial tables and formulates influential existence conjectures with one prescribed coefficient. The primitivity question and the irreducibility question are distinct; a primitive polynomial has a root generating the multiplicative group of the extension field.\par
{\footnotesize\url{https://doi.org/10.1090/S0025-5718-1992-1134730-7}}\par
{\footnotesize\url{https://www.ams.org/journals/mcom/1992-59-200/S0025-5718-1992-1134730-7/S0025-5718-1992-1134730-7.pdf}}\par
{\footnotesize\url{https://pure.psu.edu/en/publications/primitive-polynomials-over-finite-fields/}}\par
\vspace{9pt}
\needspace{7\baselineskip}\hypertarget{ref:Shoup1994}{}
{\sffamily\bfseries\color{atlasgreen}[Shoup1994] Fast construction of irreducible polynomials over finite fields}\par
\color{black} Shoup, Victor. 1994. Journal of Symbolic Computation. 17. pp. 371-391.\par
{\small\sffamily\color{atlasgold}Published work / Published article}\par
\color{black} Develops faster algorithms for constructing an irreducible polynomial of specified degree over a finite field. Useful when the objective is an explicit extension-field representation rather than testing a single given polynomial.\par
{\footnotesize\url{https://doi.org/10.1006/jsco.1994.1025}}\par
{\footnotesize\url{https://www.shoup.net/papers/fastirred.pdf}}\par
{\footnotesize\url{https://www.shoup.net/papers/}}\par
{\footnotesize\url{https://www.sciencedirect.com/science/article/pii/S074771718471025X}}\par
\vspace{9pt}
\needspace{7\baselineskip}\hypertarget{ref:Wan1997}{}
{\sffamily\bfseries\color{atlasgreen}[Wan1997] Generators and irreducible polynomials over finite fields}\par
\color{black} Wan, Daqing. 1997. Mathematics of Computation. 66(219). pp. 1195-1212.\par
{\small\sffamily\color{atlasgold}Published work / Published article}\par
\color{black} Uses character-sum estimates to study generators and irreducible polynomials with prescribed coefficients. Establishes the Hansen-Mullen irreducibility assertion outside finitely many cases; Ham and Mullen settle the remaining cases.\par
{\footnotesize\url{https://doi.org/10.1090/S0025-5718-97-00835-1}}\par
{\footnotesize\url{https://www.cambridge.org/core/journals/canadian-journal-of-mathematics/article/new-proof-of-the-hansenmullen-irreducibility-conjecture/189A3A4627AC3912AC76E77F6E5221E8}}\par
{\footnotesize\url{https://www.pollack-math.net/prescribed.pdf}}\par
\vspace{9pt}
\needspace{7\baselineskip}\hypertarget{ref:HamMullen1998}{}
{\sffamily\bfseries\color{atlasgreen}[HamMullen1998] Distribution of irreducible polynomials of small degrees over finite fields}\par
\color{black} Ham, Kie H. and Mullen, Gary L. 1998. Mathematics of Computation. 67(221). pp. 337-341.\par
{\small\sffamily\color{atlasgold}Published work / Published article}\par
\color{black} Completes the Hansen-Mullen irreducibility conjecture by verifying the finite set of cases left by Wan's theorem. The statement includes necessary exceptions, notably a prescribed zero constant coefficient and the quadratic zero-linear-coefficient case in characteristic two.\par
{\footnotesize\url{https://doi.org/10.1090/S0025-5718-98-00904-1}}\par
{\footnotesize\url{https://pure.psu.edu/en/publications/distribution-of-irreducible-polynomials-of-small-degrees-over-fin/}}\par
{\footnotesize\url{https://citeseerx.ist.psu.edu/document?doi=a767d4e1ec243bd2a44b517e4a0961a31a90255f&repid=rep1&type=pdf}}\par
\vspace{9pt}
\needspace{7\baselineskip}\hypertarget{ref:Cohen2006Primitive}{}
{\sffamily\bfseries\color{atlasgreen}[Cohen2006Primitive] Primitive polynomials with a prescribed coefficient}\par
\color{black} Cohen, Stephen D. 2006. Finite Fields and Their Applications. 12(3). pp. 425-491.\par
{\small\sffamily\color{atlasgold}Published work / Published article}\par
\color{black} Establishes the Hansen-Mullen primitivity assertion for degrees at least nine by character-sum and sieve methods. This is a stronger condition than irreducibility and this paper alone is not the full treatment of the small-degree cases.\par
{\footnotesize\url{https://doi.org/10.1016/j.ffa.2005.08.001}}\par
{\footnotesize\url{https://eprints.gla.ac.uk/12863/}}\par
{\footnotesize\url{https://www.sciencedirect.com/science/article/pii/S1071579705000663}}\par
\vspace{9pt}
\needspace{7\baselineskip}\hypertarget{ref:MullenPanario2013}{}
{\sffamily\bfseries\color{atlasgreen}[MullenPanario2013] Handbook of Finite Fields}\par
\color{black} Mullen, Gary L. and Panario, Daniel (editors). 2013. CRC Press.\par
{\small\sffamily\color{atlasgold}Book / Published edited handbook}\par
\color{black} A broad reference for irreducible and primitive polynomials, prescribed coefficients, factorization algorithms, and finite-field applications. Its extensive bibliography is a practical starting point for extending this guide in characteristic p.\par
{\footnotesize\url{https://doi.org/10.1201/b15006}}\par
{\footnotesize\url{https://pure.psu.edu/en/publications/handbook-of-finite-fields/}}\par
\vspace{9pt}
\needspace{7\baselineskip}\hypertarget{ref:Pollack2013Prescribed}{}
{\sffamily\bfseries\color{atlasgreen}[Pollack2013Prescribed] Irreducible polynomials with several prescribed coefficients}\par
\color{black} Pollack, Paul. 2013. Finite Fields and Their Applications. 22. pp. 70-78.\par
{\small\sffamily\color{atlasgold}Published work / Published article; author-hosted version}\par
\color{black} Gives an explicit estimate for monic irreducibles with coefficients prescribed at arbitrary positions, with a nonzero constant term when it is prescribed. For sufficiently large degree n, it allows about (1\ensuremath{-}\ensuremath{\varepsilon})\(\sqrt{n}\) prescribed coefficients uniformly over the ground finite field.\par
{\footnotesize\url{https://www.pollack-math.net/prescribed.pdf}}\par
{\footnotesize\url{https://doi.org/10.1016/j.ffa.2013.03.001}}\par
{\footnotesize\url{https://www.sciencedirect.com/science/article/pii/S1071579713000336}}\par
\vspace{9pt}
\needspace{7\baselineskip}\hypertarget{ref:Ha2016}{}
{\sffamily\bfseries\color{atlasgreen}[Ha2016] Irreducible polynomials with several prescribed coefficients}\par
\color{black} Ha, Junsoo. 2016. Finite Fields and Their Applications. 40. pp. 10-25.\par
{\small\sffamily\color{atlasgold}Published work / Published article}\par
\color{black} Improves prescribed-coefficient results to a positive proportion of the degree by adapting Bourgain's method. Its large-field range approaches one quarter of the coefficients; the characteristic, size, and nonzero-constant-term conditions must be retained when applying the theorem.\par
{\footnotesize\url{https://arxiv.org/abs/1601.06867}}\par
{\footnotesize\url{https://arxiv.org/html/1601.06867v1}}\par
{\footnotesize\url{https://doi.org/10.1016/j.ffa.2016.02.006}}\par
{\footnotesize\url{https://www.sciencedirect.com/science/article/pii/S1071579716000277}}\par
{\footnotesize\url{https://arxiv.org/abs/1601.06867v1}}\par
\vspace{9pt}
\needspace{7\baselineskip}\hypertarget{ref:TuxanidyWang2018}{}
{\sffamily\bfseries\color{atlasgreen}[TuxanidyWang2018] A new proof of the Hansen-Mullen irreducibility conjecture}\par
\color{black} Tuxanidy, Aleksandr and Wang, Qiang. 2018. Canadian Journal of Mathematics. 70(6). pp. 1373-1389.\par
{\small\sffamily\color{atlasgold}Published work / Published article}\par
\color{black} Gives a different proof of the one-prescribed-coefficient existence theorem using least periods of discrete Fourier transforms over finite fields. A useful example of a method that differs from the character-sum approach while preserving the genuine exceptional cases.\par
{\footnotesize\url{https://arxiv.org/abs/1604.04023}}\par
{\footnotesize\url{https://arxiv.org/abs/1604.04023v1}}\par
{\footnotesize\url{https://doi.org/10.4153/CJM-2017-022-1}}\par
{\footnotesize\url{https://www.cambridge.org/core/journals/canadian-journal-of-mathematics/article/new-proof-of-the-hansenmullen-irreducibility-conjecture/189A3A4627AC3912AC76E77F6E5221E8}}\par
{\footnotesize\url{https://arxiv.org/html/1604.04023v1}}\par
\vspace{9pt}
\needspace{7\baselineskip}\hypertarget{ref:Swan1962}{}
{\sffamily\bfseries\color{atlasgreen}[Swan1962] Factorization of polynomials over finite fields}\par
\color{black} Swan, Richard G. 1962. Pacific Journal of Mathematics. 12(3). pp. 1099-1106.\par
{\small\sffamily\color{atlasgold}Published work / Published article}\par
\color{black} Presents the Stickelberger discriminant-parity theorem, including characteristic two via lifting, and derives reducibility obstructions for binary trinomials. Swan explicitly credits the basic parity theorem to Stickelberger and notes Dalen's earlier proof.\par
{\footnotesize\url{https://msp.org/pjm/1962/12-3/pjm-v12-n3-p27-p.pdf}}\par
{\footnotesize\url{https://doi.org/10.2140/pjm.1962.12.1099}}\par
{\footnotesize\url{https://www.math.uchicago.edu/~swan/bibliography.html}}\par
\vspace{9pt}
\needspace{7\baselineskip}\hypertarget{ref:Cohen1970}{}
{\sffamily\bfseries\color{atlasgreen}[Cohen1970] The distribution of polynomials over finite fields}\par
\color{black} Cohen, Stephen D. 1970. Acta Arithmetica. 17(3). pp. 255-271.\par
{\small\sffamily\color{atlasgold}Published work / Published article}\par
\color{black} A foundational study of the factorization patterns of f(X)-a g(X) as a varies over a finite field. It connects prescribed factor degrees to permutation cycle types in a function-field Galois group, with the arithmetic hypotheses governing the distribution.\par
{\footnotesize\url{https://doi.org/10.4064/aa-17-3-255-271}}\par
{\footnotesize\url{https://www.impan.pl/en/publishing-house/journals-and-series/acta-arithmetica/all/17/3/97248/the-distribution-of-polynomials-over-finite-fields}}\par
{\footnotesize\url{https://bibliotekanauki.pl/articles/1394036.pdf}}\par
\vspace{9pt}
\needspace{7\baselineskip}\hypertarget{ref:GuralnickZieve2010}{}
{\sffamily\bfseries\color{atlasgreen}[GuralnickZieve2010] Polynomials with PSL(2) monodromy}\par
\color{black} Guralnick, Robert M. and Zieve, Michael E. 2010. Annals of Mathematics. 172(2). pp. 1315-1359.\par
{\small\sffamily\color{atlasgold}Published work / Published article}\par
\color{black} Classifies positive-characteristic polynomials with specified PSL(2) monodromy under a ramification hypothesis, with applications to exceptional polynomials and decomposition after extending constants. This concerns function-field Galois groups; splitting fields of individual polynomials over finite fields have cyclic Galois groups.\par
{\footnotesize\url{https://annals.math.princeton.edu/2010/172-2/p11}}\par
{\footnotesize\url{https://doi.org/10.4007/annals.2010.172.1315}}\par
{\footnotesize\url{https://arxiv.org/abs/0707.1835}}\par
\vspace{9pt}
\needspace{7\baselineskip}\hypertarget{ref:BankBarySorokerRosenzweig2015}{}
{\sffamily\bfseries\color{atlasgreen}[BankBarySorokerRosenzweig2015] Prime polynomials in short intervals and in arithmetic progressions}\par
\color{black} Bank, Efrat and Bary-Soroker, Lior and Rosenzweig, Lior. 2015. Duke Mathematical Journal. 164(2). pp. 277-295.\par
{\small\sffamily\color{atlasgold}Published work / Published article}\par
\color{black} Proves asymptotic counts for irreducibles in short polynomial intervals and arithmetic progressions, and treats other factorization types. The asymptotic is fixed degree with the field size tending to infinity, with explicit small-characteristic and interval-length restrictions.\par
{\footnotesize\url{https://arxiv.org/abs/1302.0625}}\par
{\footnotesize\url{https://arxiv.org/html/1302.0625v3}}\par
{\footnotesize\url{https://doi.org/10.1215/00127094-2856728}}\par
{\footnotesize\url{https://arxiv.org/abs/1302.0625v3}}\par
\vspace{9pt}
\needspace{7\baselineskip}\hypertarget{ref:Entin2021Restricted}{}
{\sffamily\bfseries\color{atlasgreen}[Entin2021Restricted] Factorization statistics of restricted polynomial specializations over large finite fields}\par
\color{black} Entin, Alexei. 2021. Israel Journal of Mathematics. 242(1). pp. 37-53.\par
{\small\sffamily\color{atlasgold}Published work / Published article}\par
\color{black} Studies factorization patterns when specialization parameters are restricted to sufficiently large, regular subsets of a prime field. The degree stays fixed as the characteristic tends to infinity; the result is not a growing-degree statement over one fixed finite field.\par
{\footnotesize\url{https://arxiv.org/abs/1810.07512}}\par
{\footnotesize\url{https://doi.org/10.1007/s11856-021-2101-9}}\par
{\footnotesize\url{https://cris.tau.ac.il/en/publications/factorization-statistics-of-restricted-polynomial-specializations/}}\par
{\footnotesize\url{https://www.math.tau.ac.il/~aentin/pub.html}}\par
\vspace{9pt}
\needspace{7\baselineskip}\hypertarget{ref:Porritt2019Restricted}{}
{\sffamily\bfseries\color{atlasgreen}[Porritt2019Restricted] Irreducible polynomials over a finite field with restricted coefficients}\par
\color{black} Porritt, Sam. 2019. Canadian Mathematical Bulletin. 62. pp. 429-439.\par
{\small\sffamily\color{atlasgold}Published work / Published article}\par
\color{black} Counts monic irreducibles when every nonleading coefficient avoids a common small forbidden set. Theorem 1 gives explicit q,n,|R| conditions and a constant-term correction; Theorem 2 uses additional interval structure. The field-size and degree ranges matter.\par
{\footnotesize\url{https://doi.org/10.4153/CMB-2018-027-x}}\par
{\footnotesize\url{https://arxiv.org/html/1711.06243v1}}\par
{\footnotesize\url{https://arxiv.org/abs/1711.06243}}\par
\vspace{9pt}
\needspace{7\baselineskip}\hypertarget{ref:Cheng2025Restricted}{}
{\sffamily\bfseries\color{atlasgreen}[Cheng2025Restricted] Counting irreducible polynomials with restricted coefficients}\par
\color{black} Cheng, Kaimin. 2025. Finite Fields and Their Applications. 108. pp. Article 102691.\par
{\small\sffamily\color{atlasgold}Published work / Published article}\par
\color{black} Extends the common-forbidden-set model to position-dependent forbidden sets. Theorem 1.1 allows at most \(\sqrt{q}\)/2 exclusions per coefficient for q\ensuremath{\ge}300 and sufficiently large degree; Corollary 1.3 treats avoiding one specified value at every position. Includes the necessary constant-term density factor.\par
{\small Issue date December 2025. The publisher-indexed article preview supplies Theorem 1.1 and Corollary 1.3; direct page fetching intermittently returns HTTP 403.}\par
{\footnotesize\url{https://doi.org/10.1016/j.ffa.2025.102691}}\par
{\footnotesize\url{https://www.sciencedirect.com/science/article/pii/S1071579725001212}}\par
{\footnotesize\url{https://www.sciencedirect.com/science/article/abs/pii/S1071579725001212}}\par
\vspace{9pt}
\needspace{7\baselineskip}\hypertarget{ref:BarySorokerShmueli2026Squares}{}
{\sffamily\bfseries\color{atlasgreen}[BarySorokerShmueli2026Squares] Irreducibility of polynomials with square coefficients over finite fields}\par
\color{black} Bary-Soroker, Lior and Shmueli, Roy. 2026. Finite Fields and Their Applications. 109. pp. Article 102714.\par
{\small\sffamily\color{atlasgold}Published work / Published article}\par
\color{black} For fixed degree n, independent nonleading coefficients uniform on the set of square values give irreducibility probability 1/n+O\_\allowbreak{}n(q\textasciicircum{}(\ensuremath{-}1/2)) as q grows. Theorem 2 covers more general sets of bounded algebraic complexity. Uniform square values differ from squaring uniform field elements in odd characteristic.\par
{\small Published January 2026; DOI contains 2025. Mathematical theorem locators refer to the public author preprint v1, while publication metadata is verified from the author institution and publisher.}\par
{\footnotesize\url{https://doi.org/10.1016/j.ffa.2025.102714}}\par
{\footnotesize\url{https://arxiv.org/html/2410.16814v1}}\par
{\footnotesize\url{https://arxiv.org/abs/2410.16814}}\par
{\footnotesize\url{https://cris.tau.ac.il/en/publications/irreducibility-of-polynomials-with-square-coefficients-over-finit/}}\par
\vspace{9pt}
\needspace{12\baselineskip}\section{Algorithms for factorization and Galois groups}
\needspace{7\baselineskip}\hypertarget{ref:Berlekamp1967}{}
{\sffamily\bfseries\color{atlasgreen}[Berlekamp1967] Factoring polynomials over finite fields}\par
\color{black} Berlekamp, Elwyn R. 1967. Bell System Technical Journal. 46(8). pp. 1853-1859.\par
{\small\sffamily\color{atlasgold}Published work / Published article}\par
\color{black} The foundational linear-algebra method for factoring polynomials over finite fields. Read it to understand the Frobenius-fixed subalgebra that detects the irreducible factors.\par
{\footnotesize\url{https://doi.org/10.1002/j.1538-7305.1967.tb03174.x}}\par
{\footnotesize\url{https://onlinelibrary.wiley.com/doi/abs/10.1002/j.1538-7305.1967.tb03174.x}}\par
{\footnotesize\url{https://www.nokia.com/bell-labs/publications-and-media/publications/factoring-polynomials-over-finite-fields/}}\par
\vspace{9pt}
\needspace{7\baselineskip}\hypertarget{ref:Zassenhaus1969}{}
{\sffamily\bfseries\color{atlasgreen}[Zassenhaus1969] On Hensel factorization, I}\par
\color{black} Zassenhaus, Hans. 1969. Journal of Number Theory. 1(3). pp. 291-311.\par
{\small\sffamily\color{atlasgold}Published work / Published article}\par
\color{black} The classical modular route to integer-polynomial factorization: factor modulo a prime, lift p-adically, then recombine the lifted factors. It explains why local factors must be grouped before recovering rational factors.\par
{\footnotesize\url{https://doi.org/10.1016/0022-314X(69)90047-X}}\par
{\footnotesize\url{https://www.sciencedirect.com/science/article/pii/0022314X6990047X}}\par
{\footnotesize\url{https://www.numdam.org/item/JTNB_2009__21_1_15_0/}}\par
\vspace{9pt}
\needspace{7\baselineskip}\hypertarget{ref:Rabin1980}{}
{\sffamily\bfseries\color{atlasgreen}[Rabin1980] Probabilistic algorithms in finite fields}\par
\color{black} Rabin, Michael O. 1980. SIAM Journal on Computing. 9(2). pp. 273-280.\par
{\small\sffamily\color{atlasgold}Published work / Published article}\par
\color{black} A central source for finite-field algorithms for irreducible-polynomial construction, root finding, and factorization. Useful alongside the exact Frobenius-power criterion for testing a monic polynomial for irreducibility.\par
{\small The MIT repository link describes the earlier January 1979 report Probabilistic Algorithm in Finite Fields (MIT-LCS-TR-213). This entry cites the 1980 SIAM journal article; the repository report and journal version have different publication dates.}\par
{\footnotesize\url{https://doi.org/10.1137/0209024}}\par
{\footnotesize\url{https://epubs.siam.org/doi/10.1137/0209024}}\par
{\footnotesize\url{https://dspace.mit.edu/entities/publication/0c9f4e53-8b19-4a3c-8d9e-0a53238d2eb3}}\par
\vspace{9pt}
\needspace{7\baselineskip}\hypertarget{ref:CantorZassenhaus1981}{}
{\sffamily\bfseries\color{atlasgreen}[CantorZassenhaus1981] A new algorithm for factoring polynomials over finite fields}\par
\color{black} Cantor, David G. and Zassenhaus, Hans. 1981. Mathematics of Computation. 36(154). pp. 587-592.\par
{\small\sffamily\color{atlasgold}Published work / Published article}\par
\color{black} A basic randomized finite-field factorization algorithm, especially important for splitting products whose irreducible factors have equal degree. Its randomness concerns the search procedure; successful factors can be checked exactly.\par
{\footnotesize\url{https://doi.org/10.1090/S0025-5718-1981-0606517-5}}\par
{\footnotesize\url{https://flintlib.org/doc/gr_poly/gr_poly_factor_equal_deg.html}}\par
{\footnotesize\url{https://www.ams.org/journals/mcom/1981-36-154/S0025-5718-1981-0606517-5/}}\par
\vspace{9pt}
\needspace{7\baselineskip}\hypertarget{ref:LLL1982}{}
{\sffamily\bfseries\color{atlasgreen}[LLL1982] Factoring polynomials with rational coefficients}\par
\color{black} Lenstra, A. K. and Lenstra, Jr., H. W. and Lovász, L. 1982. Mathematische Annalen. 261(4). pp. 515-534.\par
{\small\sffamily\color{atlasgold}Published work / Published article}\par
\color{black} Proves polynomial-time factorization over the rationals using lattice basis reduction. The LLL algorithm introduced here is also a central tool for finding algebraic relations and constructing certificates.\par
{\small The CWI record is the 1982 preprint report IW 195/82; the UC Davis link supplies scanned full text. The Numdam link is the 2009 Belabas-van Hoeij-Klüners-Steel follow-up, not another edition of this 1982 article.}\par
{\footnotesize\url{https://doi.org/10.1007/BF01457454}}\par
{\footnotesize\url{https://www.math.ucdavis.edu/~deloera/MISC/LA-BIBLIO/trunk/Lovasz/LovaszLenstraLenstrafactor.pdf}}\par
{\footnotesize\url{https://ir.cwi.nl/pub/9304}}\par
{\footnotesize\url{https://www.numdam.org/item/JTNB_2009__21_1_15_0/}}\par
\vspace{9pt}
\needspace{7\baselineskip}\hypertarget{ref:Shoup1995}{}
{\sffamily\bfseries\color{atlasgreen}[Shoup1995] A new polynomial factorization algorithm and its implementation}\par
\color{black} Shoup, Victor. 1995. Journal of Symbolic Computation. 20(4). pp. 363-397.\par
{\small\sffamily\color{atlasgold}Published work / Published article}\par
\color{black} Develops a practical implementation of baby-step/giant-step finite-field factorization. A useful bridge between asymptotic algorithms and actual computational work.\par
{\small The open-text link is the author's version. The published article appeared in Journal of Symbolic Computation 20(4), October 1995, pages 363-397.}\par
{\footnotesize\url{https://doi.org/10.1006/jsco.1995.1055}}\par
{\footnotesize\url{https://www.shoup.net/papers/factorimpl.pdf}}\par
{\footnotesize\url{https://www.sciencedirect.com/science/article/pii/S0747717185710553}}\par
{\footnotesize\url{https://www.shoup.net/papers/}}\par
\vspace{9pt}
\needspace{7\baselineskip}\hypertarget{ref:KaltofenShoup1998}{}
{\sffamily\bfseries\color{atlasgreen}[KaltofenShoup1998] Subquadratic-time factoring of polynomials over finite fields}\par
\color{black} Kaltofen, Erich and Shoup, Victor. 1998. Mathematics of Computation. 67(223). pp. 1179-1197.\par
{\small\sffamily\color{atlasgold}Published work / Published article}\par
\color{black} A major complexity advance in randomized polynomial factorization over finite fields. Its baby-step/giant-step framework is part of the background to modern finite-field factorization libraries.\par
{\small The open-text link is the authors' version of the journal work. The 1998 Mathematics of Computation article expands an extended abstract presented at STOC 1995; the two publication dates identify different versions.}\par
{\footnotesize\url{https://doi.org/10.1090/S0025-5718-98-00944-2}}\par
{\footnotesize\url{https://www.shoup.net/papers/subquadratic.pdf}}\par
{\footnotesize\url{https://www.shoup.net/papers/}}\par
{\footnotesize\url{https://kaltofen.math.ncsu.edu/bibliography/index.html}}\par
{\footnotesize\url{https://flintlib.org/doc/algorithms/polynomials/factoring-finite-fields.html}}\par
\vspace{9pt}
\needspace{7\baselineskip}\hypertarget{ref:VanHoeij2002}{}
{\sffamily\bfseries\color{atlasgreen}[VanHoeij2002] Factoring polynomials and the knapsack problem}\par
\color{black} van Hoeij, Mark. 2002. Journal of Number Theory. 95(2). pp. 167-189.\par
{\small\sffamily\color{atlasgold}Published work / Published article}\par
\color{black} Transforms the expensive recombination stage of modular factorization into a structured knapsack problem handled by lattice reduction. This is one of the main sources behind practical rational-polynomial factorization.\par
{\small The open-text link is the author's version. The journal publication is August 2002; the year 2001 embedded in its DOI is not the publication year.}\par
{\footnotesize\url{https://doi.org/10.1006/jnth.2001.2763}}\par
{\footnotesize\url{https://www.math.fsu.edu/~hoeij/knapsack/paper/knapsack.pdf}}\par
{\footnotesize\url{https://www.math.fsu.edu/~hoeij/papers.html}}\par
{\footnotesize\url{https://www.sciencedirect.com/science/article/pii/S0022314X01927635}}\par
{\footnotesize\url{https://flintlib.org/doc/algorithms/polynomials/factoring-integers.html}}\par
\vspace{9pt}
\needspace{7\baselineskip}\hypertarget{ref:BHKS2009}{}
{\sffamily\bfseries\color{atlasgreen}[BHKS2009] Factoring polynomials over global fields}\par
\color{black} Belabas, Karim and van Hoeij, Mark and Klüners, Jürgen and Steel, Allan. 2009. Journal de théorie des nombres de Bordeaux. 21(1). pp. 15-39.\par
{\small\sffamily\color{atlasgold}Published work / Published article}\par
\color{black} Establishes polynomial-time complexity for van Hoeij's original rational factorization algorithm and natural variants. It also treats global-field extensions of the approach and bivariate factorization over finite fields.\par
{\footnotesize\url{https://www.numdam.org/articles/10.5802/jtnb.655/}}\par
{\footnotesize\url{https://doi.org/10.5802/jtnb.655}}\par
{\footnotesize\url{https://www.numdam.org/item/JTNB_2009__21_1_15_0/}}\par
{\footnotesize\url{https://ris.uni-paderborn.de/record/34889}}\par
\vspace{9pt}
\needspace{7\baselineskip}\hypertarget{ref:Stauduhar1973}{}
{\sffamily\bfseries\color{atlasgreen}[Stauduhar1973] The determination of Galois groups}\par
\color{black} Stauduhar, R. P. 1973. Mathematics of Computation. 27(124). pp. 981-996.\par
{\small\sffamily\color{atlasgold}Published work / Published article}\par
\color{black} Introduces the subgroup-descent strategy based on relative invariants and resolvents. This remains the conceptual starting point for many practical Galois-group computations.\par
{\footnotesize\url{https://doi.org/10.1090/S0025-5718-1973-0327712-4}}\par
{\footnotesize\url{https://docs.magma-maths.org/GlobalFields/GaloisGroupsAndAutomorphisms/bibliography.html}}\par
{\footnotesize\url{https://www.cambridge.org/core/journals/lms-journal-of-computation-and-mathematics/article/computation-of-galois-groups-of-rational-polynomials/70FFFC0B1BBA91C846342618C93D1E08}}\par
\vspace{9pt}
\needspace{7\baselineskip}\hypertarget{ref:LandauMiller1985}{}
{\sffamily\bfseries\color{atlasgreen}[LandauMiller1985] Solvability by radicals is in polynomial time}\par
\color{black} Landau, Susan and Miller, Gary Lee. 1985. Journal of Computer and System Sciences. 30(2). pp. 179-208.\par
{\small\sffamily\color{atlasgold}Published work / Published article}\par
\color{black} Shows that deciding solvability by radicals for rational polynomials is possible in polynomial time. This decision problem should be distinguished from constructing an entire splitting field or listing every element of a large Galois group.\par
{\footnotesize\url{https://doi.org/10.1016/0022-0000(85)90013-3}}\par
{\footnotesize\url{https://www.cs.cmu.edu/~glmiller/Publications/Papers/LaMi85.pdf}}\par
{\footnotesize\url{https://www.cs.cmu.edu/~glmiller/Publications/b2hd-LaMi85.html}}\par
{\footnotesize\url{https://www.sciencedirect.com/science/article/pii/0022000085900133}}\par
\vspace{9pt}
\needspace{7\baselineskip}\hypertarget{ref:SoicherMcKay1985}{}
{\sffamily\bfseries\color{atlasgreen}[SoicherMcKay1985] Computing Galois groups over the rationals}\par
\color{black} Soicher, Leonard and McKay, John. 1985. Journal of Number Theory. 20(3). pp. 273-281.\par
{\small\sffamily\color{atlasgold}Published work / Published article}\par
\color{black} Develops practical computational methods using resolvents and low-degree transitive groups. The paper also gives rational realizations of the transitive permutation groups in degrees three through seven.\par
{\footnotesize\url{https://doi.org/10.1016/0022-314X(85)90022-8}}\par
{\footnotesize\url{https://www.sciencedirect.com/science/article/pii/0022314X85900228}}\par
{\footnotesize\url{https://webspace.maths.qmul.ac.uk/l.h.soicher/pubs.html}}\par
{\footnotesize\url{https://docs.magma-maths.org/GlobalFields/GaloisGroupsAndAutomorphisms/bibliography.html}}\par
\vspace{9pt}
\needspace{7\baselineskip}\hypertarget{ref:Dummit1991}{}
{\sffamily\bfseries\color{atlasgreen}[Dummit1991] Solving solvable quintics}\par
\color{black} Dummit, David S. 1991. Mathematics of Computation. 57(195). pp. 387-401.\par
{\small\sffamily\color{atlasgold}Published work / Published article; corrections available}\par
\color{black} A concrete sextic-resolvent treatment of solvable irreducible quintics and radical expressions. Use it with Robin Chapman's 1992 correction to the proof of Theorem 1 and the numerical errata on the author's webpage.\par
{\footnotesize\url{https://site.uvm.edu/ddummit/}}\par
{\footnotesize\url{https://site.uvm.edu/ddummit/files/2021/04/Solving_Solvable_Quintics__Math_Comp_57_no195_1991__pp_387_401.pdf}}\par
{\footnotesize\url{https://doi.org/10.1090/S0025-5718-1991-1079014-X}}\par
{\footnotesize\url{https://site.uvm.edu/ddummit/files/2021/04/Solving_Solvable_Quintics_note_by_Robin_Chapman.pdf}}\par
\vspace{9pt}
\needspace{7\baselineskip}\hypertarget{ref:GeisslerKlueners2000}{}
{\sffamily\bfseries\color{atlasgreen}[GeisslerKlueners2000] Galois group computation for rational polynomials}\par
\color{black} Geissler, Katharina and Klüners, Jürgen. 2000. Journal of Symbolic Computation. 30(6). pp. 653-674.\par
{\small\sffamily\color{atlasgold}Published work / Published article}\par
\color{black} A practical development of Galois-group computation using p-adic roots and refinements of Stauduhar's method. Read it before the later degree-independent approach of Fieker and Klüners.\par
{\footnotesize\url{https://doi.org/10.1006/jsco.2000.0377}}\par
{\footnotesize\url{https://www.sciencedirect.com/science/article/pii/S0747717100903778}}\par
{\footnotesize\url{https://www.cambridge.org/core/journals/lms-journal-of-computation-and-mathematics/article/computation-of-galois-groups-of-rational-polynomials/70FFFC0B1BBA91C846342618C93D1E08}}\par
\vspace{9pt}
\needspace{7\baselineskip}\hypertarget{ref:KluenersMalle2000}{}
{\sffamily\bfseries\color{atlasgreen}[KluenersMalle2000] Explicit Galois realization of transitive groups of degree up to 15}\par
\color{black} Klüners, Jürgen and Malle, Gunter. 2000. Journal of Symbolic Computation. 30(6). pp. 675-716.\par
{\small\sffamily\color{atlasgold}Published work / Published article}\par
\color{black} Constructs explicit polynomials realizing all transitive groups of degree at most fifteen over the rationals. A valuable source of test examples and methods for organizing computations by permutation group.\par
{\footnotesize\url{https://doi.org/10.1006/jsco.2000.0378}}\par
{\footnotesize\url{https://ris.uni-paderborn.de/record/34899}}\par
{\footnotesize\url{https://www.sciencedirect.com/science/article/pii/S074771710090378X}}\par
\vspace{9pt}
\needspace{7\baselineskip}\hypertarget{ref:DokchitserDokchitser2013}{}
{\sffamily\bfseries\color{atlasgreen}[DokchitserDokchitser2013] Identifying Frobenius elements in Galois groups}\par
\color{black} Dokchitser, Tim and Dokchitser, Vladimir. 2013. Algebra \& Number Theory. 7(6). pp. 1325-1352.\par
{\small\sffamily\color{atlasgold}Published work / Published article}\par
\color{black} Identifies Frobenius conjugacy classes more precisely than the factor-degree data of reduction modulo a prime. Useful when distinct conjugacy classes have the same cycle type in the chosen permutation representation.\par
{\footnotesize\url{https://doi.org/10.2140/ant.2013.7.1325}}\par
{\footnotesize\url{https://msp.org/ant/2013/7-6/ant-v7-n6-p02-s.pdf}}\par
{\footnotesize\url{https://arxiv.org/abs/1009.5388}}\par
\vspace{9pt}
\needspace{7\baselineskip}\hypertarget{ref:FiekerKlueners2014}{}
{\sffamily\bfseries\color{atlasgreen}[FiekerKlueners2014] Computation of Galois groups of rational polynomials}\par
\color{black} Fieker, Claus and Klüners, Jürgen. 2014. LMS Journal of Computation and Mathematics. 17(1). pp. 141-158.\par
{\small\sffamily\color{atlasgold}Published work / Published article}\par
\color{black} Constructs relative invariants for practical Galois-group computation without an a priori degree bound. It is a central reference for understanding current Magma computations; practical time and memory limits still matter.\par
{\footnotesize\url{https://doi.org/10.1112/S1461157013000302}}\par
{\footnotesize\url{https://arxiv.org/abs/1211.3588}}\par
{\footnotesize\url{https://www.cambridge.org/core/journals/lms-journal-of-computation-and-mathematics/article/computation-of-galois-groups-of-rational-polynomials/70FFFC0B1BBA91C846342618C93D1E08}}\par
\vspace{9pt}
\needspace{7\baselineskip}\hypertarget{ref:Cohen1993}{}
{\sffamily\bfseries\color{atlasgreen}[Cohen1993] A Course in Computational Algebraic Number Theory}\par
\color{black} Cohen, Henri. 1993. Springer. 138. Graduate Texts in Mathematics.\par
{\small\sffamily\color{atlasgold}Book / Published book}\par
\color{black} A systematic computational number-theory reference, including polynomial factorization, number fields, and Galois-group computations. Particularly helpful for translating a mathematical method into a reproducible algorithm.\par
{\footnotesize\url{https://doi.org/10.1007/978-3-662-02945-9}}\par
{\footnotesize\url{https://link.springer.com/book/10.1007/978-3-662-02945-9}}\par
\vspace{9pt}
\needspace{7\baselineskip}\hypertarget{ref:GathenGerhard2013}{}
{\sffamily\bfseries\color{atlasgreen}[GathenGerhard2013] Modern Computer Algebra}\par
\color{black} von zur Gathen, Joachim and Gerhard, Jürgen. 2013. Cambridge University Press. Edition 3.\par
{\small\sffamily\color{atlasgold}Book / Published book}\par
\color{black} A comprehensive source on polynomial arithmetic, finite-field factorization, Hensel lifting, and factorization over the rationals. Use it for complexity models and complete algorithmic proofs rather than isolated software recipes.\par
{\footnotesize\url{https://doi.org/10.1017/CBO9781139856065}}\par
{\footnotesize\url{https://www.cambridge.org/core/books/modern-computer-algebra/DB3563D4013401734851CF683D2F03F0}}\par
{\footnotesize\url{https://cosec.bit.uni-bonn.de/science/mca/}}\par
\vspace{9pt}
\needspace{7\baselineskip}\hypertarget{ref:GuardiaMontesNart2011Montes}{}
{\sffamily\bfseries\color{atlasgreen}[GuardiaMontesNart2011Montes] Higher Newton polygons in the computation of discriminants and prime ideal decomposition in number fields}\par
\color{black} Guàrdia, Jordi and Montes, Jesús and Nart, Enric. 2011. Journal de théorie des nombres de Bordeaux. 23(3). pp. 667-696.\par
{\small\sffamily\color{atlasgold}Published work / Published article}\par
\color{black} Presents an optimized Montes method using types and their representatives. Complete types identify irreducible local factors and reveal ramification indices and residue degrees; the algorithm computes prime ideal decomposition and local index information without requiring a global integral basis first.\par
{\footnotesize\url{https://www.numdam.org/articles/10.5802/jtnb.782/}}\par
{\footnotesize\url{https://www.numdam.org/item/10.5802/jtnb.782.pdf}}\par
{\footnotesize\url{https://doi.org/10.5802/jtnb.782}}\par
\vspace{9pt}
\needspace{7\baselineskip}\hypertarget{ref:PoteauxWeimann2022Local}{}
{\sffamily\bfseries\color{atlasgreen}[PoteauxWeimann2022Local] Local Polynomial Factorisation: Improving the Montes Algorithm}\par
\color{black} Poteaux, Adrien and Weimann, Martin. 2022. Proceedings of the 2022 International Symposium on Symbolic and Algebraic Computation (ISSAC '22). pp. 149-157.\par
{\small\sffamily\color{atlasgold}Published work / Published conference paper (5 July 2022); related 2026 arXiv deposit}\par
\color{black} Improves local irreducibility and factorization using generalized Hensel lifting, approximate roots, and divide-and-conquer. The sharper complexity bounds require residue characteristic zero or sufficiently large; the source separately includes residual-field factorization costs. A same-title arXiv record, 2607.02153 (2 July 2026), is linked as a related deposit, not a separate discovery; its full text could not be retrieved in this audit.\par
{\footnotesize\url{https://doi.org/10.1145/3476446.3535487}}\par
{\footnotesize\url{https://weimann.users.lmno.cnrs.fr/issac.pdf}}\par
{\footnotesize\url{https://www.issac-conference.org/2022/papers.php}}\par
{\footnotesize\url{https://arxiv.org/abs/2607.02153}}\par
\vspace{9pt}
\needspace{7\baselineskip}\hypertarget{ref:BrentKung1978}{}
{\sffamily\bfseries\color{atlasgreen}[BrentKung1978] Fast algorithms for manipulating formal power series}\par
\color{black} Brent, Richard P. and Kung, Hsiang-Tsung. 1978. Journal of the ACM. 25(4). pp. 581-595.\par
{\small\sffamily\color{atlasgold}Published work / Published article}\par
\color{black} Classical fast composition and reversion algorithms. The baby-step/giant-step matrix method also underlies modular composition; distinguish the arbitrary-modulus version from techniques that use the special modulus x\textasciicircum{}n. The author page records an erratum on page 590.\par
{\footnotesize\url{https://doi.org/10.1145/322092.322099}}\par
{\footnotesize\url{https://maths-people.anu.edu.au/~brent/pub/pub045.html}}\par
{\footnotesize\url{https://www.eecs.harvard.edu/~htk/publication/1978-jacm-brent-kung}}\par
\vspace{9pt}
\needspace{7\baselineskip}\hypertarget{ref:KedlayaUmans2011}{}
{\sffamily\bfseries\color{atlasgreen}[KedlayaUmans2011] Fast polynomial factorization and modular composition}\par
\color{black} Kedlaya, Kiran S. and Umans, Christopher. 2011. SIAM Journal on Computing. 40(6). pp. 1767-1802.\par
{\small\sffamily\color{atlasgold}Published work / Published article}\par
\color{black} Near-linear bit complexity for modular composition over finite fields leads to the degree-3/2 randomized factorization benchmark. Essential for understanding Frobenius iteration, multivariate evaluation, and the difference between bit and field-operation models. Corollary 7.2 states the modular-composition bound for each fixed positive epsilon.\par
{\footnotesize\url{https://doi.org/10.1137/08073408X}}\par
{\footnotesize\url{https://epubs.siam.org/doi/10.1137/08073408X}}\par
{\footnotesize\url{https://users.cms.caltech.edu/~umans/papers/KU08-final.pdf}}\par
{\footnotesize\url{https://dspace.mit.edu/entities/publication/fd268b79-5301-4239-8ce3-4b964b6b72d6}}\par
\vspace{9pt}
\needspace{7\baselineskip}\hypertarget{ref:GuoNarayananUmans2016}{}
{\sffamily\bfseries\color{atlasgreen}[GuoNarayananUmans2016] Algebraic problems equivalent to beating exponent 3/2 for polynomial factorization over finite fields}\par
\color{black} Guo, Zeyu and Narayanan, Anand Kumar and Umans, Chris. 2016. 41st International Symposium on Mathematical Foundations of Computer Science (MFCS 2016), LIPIcs. 58. pp. 47:1-47:14.\par
{\small\sffamily\color{atlasgold}Published work / Published conference paper}\par
\color{black} Relates finite-field factorization to several algebraic problems by reductions. Gives concrete routes toward improving the degree exponent 3/2; this is a collection of equivalences, not an impossibility result or a proof that 3/2 is optimal.\par
{\footnotesize\url{https://doi.org/10.4230/LIPIcs.MFCS.2016.47}}\par
{\footnotesize\url{https://arxiv.org/abs/1606.04592}}\par
{\footnotesize\url{https://drops.dagstuhl.de/entities/document/10.4230/LIPIcs.MFCS.2016.47}}\par
\vspace{9pt}
\needspace{7\baselineskip}\hypertarget{ref:NeigerSalvySchostVillard2024}{}
{\sffamily\bfseries\color{atlasgreen}[NeigerSalvySchostVillard2024] Faster modular composition}\par
\color{black} Neiger, Vincent and Salvy, Bruno and Schost, Éric and Villard, Gilles. 2024. Journal of the ACM. 71(2). pp. Article 11.\par
{\small\sffamily\color{atlasgold}Published work / Published article}\par
\color{black} A Las Vegas algorithm for modular composition over arbitrary fields of known characteristic, with expected arithmetic complexity O(n\textasciicircum{}1.43). Polynomial matrices encode algebraic relations. Theorem 1.1 detects failed trials; the discussion handles small finite fields by extension. This counts field operations, unlike finite-field bit-complexity bounds.\par
{\footnotesize\url{https://doi.org/10.1145/3638349}}\par
{\footnotesize\url{https://arxiv.org/abs/2110.08354}}\par
{\footnotesize\url{https://cs.uwaterloo.ca/~eschost/publications/NSSV24.pdf}}\par
{\footnotesize\url{https://perso.ens-lyon.fr/gilles.villard/BIBLIOGRAPHIE/biblio.html}}\par
\vspace{9pt}
\needspace{7\baselineskip}\hypertarget{ref:NeigerSalvySchostVillard2026}{}
{\sffamily\bfseries\color{atlasgreen}[NeigerSalvySchostVillard2026] Faster modular composition using two relation matrices}\par
\color{black} Neiger, Vincent and Salvy, Bruno and Schost, Éric and Villard, Gilles. 2026. Proceedings of the 51st International Symposium on Symbolic and Algebraic Computation (ISSAC 2026).\par
{\small\sffamily\color{atlasgold}Published work / Published conference paper}\par
\color{black} Theorem 1.3 improves algebraic modular composition to soft-O(n\textasciicircum{}((omega+3)/4)) for generic inner polynomials relative to a fixed modulus. Genericity failure is detected; an arbitrary-input Las Vegas reduction is left as a research question. Published at ISSAC 2026, with the published version on arXiv v2.\par
{\footnotesize\url{https://doi.org/10.1145/3815436.3815451}}\par
{\footnotesize\url{https://arxiv.org/abs/2601.17422}}\par
{\footnotesize\url{https://arxiv.org/html/2601.17422v2}}\par
{\footnotesize\url{https://perso.ens-lyon.fr/gilles.villard/BIBLIOGRAPHIE/biblio.html}}\par
\vspace{9pt}
\needspace{7\baselineskip}\hypertarget{ref:HoevenLecerf2026}{}
{\sffamily\bfseries\color{atlasgreen}[HoevenLecerf2026] Fast modular composition using spiroids}\par
\color{black} van der Hoeven, Joris and Lecerf, Grégoire. 2026. Journal of Complexity. 96. pp. 102073.\par
{\small\sffamily\color{atlasgold}Published work / Published article}\par
\color{black} Gives a sharper nearly linear bit bound for modular composition over Z/rZ using structured complex evaluation points and controlled-precision Fourier transforms. Prime finite fields are included. The authors distinguish asymptotic improvement from practical competitiveness. Journal publication is October 2026.\par
{\footnotesize\url{https://doi.org/10.1016/j.jco.2026.102073}}\par
{\footnotesize\url{https://researchportal.ip-paris.fr/en/publications/fast-modular-composition-using-spiroids/}}\par
{\footnotesize\url{https://www.texmacs.org/joris/fftcomp/fftcomp.pdf}}\par
{\footnotesize\url{https://www.texmacs.org/joris/fftcomp/fftcomp.html}}\par
\vspace{9pt}
\needspace{7\baselineskip}\hypertarget{ref:UmansWang2025}{}
{\sffamily\bfseries\color{atlasgreen}[UmansWang2025] A number-theoretic conjecture implying faster algorithms for polynomial factorization and integer factorization}\par
\color{black} Umans, Chris and Wang, Siki. 2025. arXiv:2511.10851; listed for ISAAC 2026.\par
{\small\sffamily\color{atlasgold}Preprint / Preprint; listed for ISAAC 2026}\par
\color{black} Proposes a structured divisibility-covering conjecture which, under the paper's hypotheses, would improve the degree exponent for finite-field factorization below 3/2. The proposed algorithmic improvement is conditional. Available as a 2025 preprint and listed in the December 2026 ISAAC program at the verification date; do not cite it as an unconditional breakthrough.\par
{\footnotesize\url{https://arxiv.org/abs/2511.10851}}\par
{\footnotesize\url{https://doi.org/10.48550/arXiv.2511.10851}}\par
{\footnotesize\url{https://www.algo-door.com/isaac2026/program.html}}\par
\vspace{9pt}
\needspace{7\baselineskip}\hypertarget{ref:Doris2020}{}
{\sffamily\bfseries\color{atlasgreen}[Doris2020] Computing the Galois group of a polynomial over a p-adic field}\par
\color{black} Doris, Christopher. 2020. International Journal of Number Theory.\par
{\small\sffamily\color{atlasgold}Published work / Published article}\par
\color{black} Local Galois-group algorithms based on resolvents, ramification, and global models. Useful alongside global computations because a local group is a decomposition subgroup. Remark 3.5 records a heuristic precision bound in the original implementation; the linked package separately documents an exact p-adic input mode for proved results.\par
{\footnotesize\url{https://arxiv.org/abs/2003.05834}}\par
{\footnotesize\url{https://doi.org/10.1142/S179304212050092X}}\par
{\footnotesize\url{https://arxiv.org/html/2003.05834v1}}\par
{\footnotesize\url{https://cjdoris.github.io/}}\par
{\footnotesize\url{https://github.com/cjdoris/pAdicGaloisGroup}}\par
{\footnotesize\url{https://research-information.bris.ac.uk/en/publications/computing-the-galois-group-of-a-polynomial-over-a-p-adic-field/}}\par
\vspace{9pt}
\needspace{12\baselineskip}\section{Databases and software}
\needspace{7\baselineskip}\hypertarget{ref:KluenersMalle2001}{}
{\sffamily\bfseries\color{atlasgreen}[KluenersMalle2001] A database for field extensions of the rationals}\par
\color{black} Klüners, Jürgen and Malle, Gunter. 2001. LMS Journal of Computation and Mathematics. 4. pp. 182-196.\par
{\small\sffamily\color{atlasgold}Published work / Published article}\par
\color{black} Explains the arithmetic database, its construction methods, and its emphasis on Galois groups and signatures. Consult the live database separately, since its coverage has grown beyond the original paper.\par
{\footnotesize\url{https://doi.org/10.1112/S1461157000000851}}\par
{\footnotesize\url{https://arxiv.org/abs/math/0102232}}\par
{\footnotesize\url{https://ris.uni-paderborn.de/record/34897}}\par
{\footnotesize\url{https://www.cambridge.org/core/journals/lms-journal-of-computation-and-mathematics/article/database-for-field-extensions-of-the-rationals/75F226702791F692154D9BE954D08F0C}}\par
\vspace{9pt}
\needspace{7\baselineskip}\hypertarget{ref:FLINTFactorizationDocumentation}{}
{\sffamily\bfseries\color{atlasgreen}[FLINTFactorizationDocumentation] FLINT documentation: polynomial factorization algorithms}\par
\color{black} The FLINT Team. Undated.\par
{\small\sffamily\color{atlasgold}Resource / Maintained software documentation; revision in progress}\par
\color{black} Official explanations of finite-field and integer-polynomial factorization, connecting Berlekamp, Cantor-Zassenhaus, Hensel lifting, and lattice recombination with library routines. These are implementation guides; use the cited original papers for theorem statements and complexity proofs.\par
{\small The documentation was undergoing a major revision when checked. Its October 2026 banner states that parts of the text were generated by an LLM and asks readers to report inaccuracies. Mathematical claims in this atlas are independently supported by the original research papers.}\par
{\footnotesize\url{https://flintlib.org/doc/algorithms/polynomials/factoring-finite-fields.html}}\par
{\footnotesize\url{https://flintlib.org/doc/algorithms/polynomials/factoring-integers.html}}\par
{\footnotesize\url{https://flintlib.org/doc/algorithms/polynomials/irreducibility.html}}\par
\vspace{9pt}
\needspace{7\baselineskip}\hypertarget{ref:KluenersMalleDatabase}{}
{\sffamily\bfseries\color{atlasgreen}[KluenersMalleDatabase] A database for number fields}\par
\color{black} Klüners, Jürgen and Malle, Gunter. Undated.\par
{\small\sffamily\color{atlasgold}Resource / Maintained mathematical database}\par
\color{black} A source of defining polynomials indexed by Galois group and signature. At access, the site reports coverage through degree twenty-three for all transitive groups except M23, with stated limitations on signatures and discriminant minima.\par
{\footnotesize\url{https://galoisdb.math.upb.de/}}\par
\vspace{9pt}
\needspace{7\baselineskip}\hypertarget{ref:LMFDBFieldsGroups}{}
{\sffamily\bfseries\color{atlasgreen}[LMFDBFieldsGroups] The L-functions and Modular Forms Database: number fields and Galois groups}\par
\color{black} The LMFDB Collaboration. Undated.\par
{\small\sffamily\color{atlasgold}Resource / Maintained mathematical database}\par
\color{black} Searchable number-field records and transitive-group data, with labels, invariants, and links between objects. Consult the completeness and reliability pages for each collection; a group appearing in a transitive-group catalogue does not by itself prove its realization over the rationals.\par
{\small Number-field completeness is bounded by degree, signature, discriminant, or ramification conditions. The reliability page marks certain class-group and unit data that depend on GRH. Transitive-group catalogue completeness is a separate statement and does not establish realization over the rationals.}\par
{\footnotesize\url{https://www.lmfdb.org/NumberField/}}\par
{\footnotesize\url{https://www.lmfdb.org/GaloisGroup/}}\par
{\footnotesize\url{https://www.lmfdb.org/NumberField/Completeness}}\par
{\footnotesize\url{https://www.lmfdb.org/NumberField/Reliability}}\par
{\footnotesize\url{https://www.lmfdb.org/GaloisGroup/Completeness}}\par
{\footnotesize\url{https://www.lmfdb.org/GaloisGroup/Reliability}}\par
\vspace{9pt}
\needspace{7\baselineskip}\hypertarget{ref:MagmaGaloisDocumentation}{}
{\sffamily\bfseries\color{atlasgreen}[MagmaGaloisDocumentation] Magma Handbook: Galois groups}\par
\color{black} Magma Computational Algebra System. Undated.\par
{\small\sffamily\color{atlasgold}Resource / Maintained software documentation}\par
\color{black} Official documentation for the Stauduhar-based GaloisGroup implementation, with no fixed degree cap for the documented integer, rational and number-field polynomial inputs. GaloisGroup does not itself certify its conditional steps; the documented GaloisProof overloads attempt certification for integer polynomials or absolute number fields. Record a successful result and the exact input/options.\par
{\footnotesize\url{https://docs.magma-maths.org/GlobalFields/GaloisGroupsAndAutomorphisms/galois.html}}\par
\vspace{9pt}
\needspace{7\baselineskip}\hypertarget{ref:PARIDocumentation}{}
{\sffamily\bfseries\color{atlasgreen}[PARIDocumentation] PARI/GP reference manual: general number fields and polgalois}\par
\color{black} The PARI Group. Undated.\par
{\small\sffamily\color{atlasgold}Resource / Maintained software documentation}\par
\color{black} Official reference for polgalois and number-field computations. PARI/GP 2.19 requires irreducible rational input of degree at most seven, or eleven with galdata. The manual explicitly warns that insufficient initial precision can rarely produce a wrong result; record precision, output convention, and version.\par
{\footnotesize\url{https://pari.math.u-bordeaux.fr/dochtml/ref/General_number_fields.html}}\par
\vspace{9pt}
\needspace{7\baselineskip}\hypertarget{ref:SageGaloisDocumentation}{}
{\sffamily\bfseries\color{atlasgreen}[SageGaloisDocumentation] SageMath reference manual: polynomial and number-field Galois groups}\par
\color{black} The Sage Developers. Undated.\par
{\small\sffamily\color{atlasgold}Resource / Maintained software documentation}\par
\color{black} Official interfaces for polynomial and number-field Galois groups. SageMath 10.10 documents PARI through degree 11, GAP fallback with a documented range through degree 15 and some higher examples, and a Magma wrapper limited by group identification through degree 30. The number-field interface distinguishes original-root and Galois-closure actions and warns about absolute versus relative ground fields. A wrapper is not an independent backend proof.\par
{\footnotesize\url{https://doc.sagemath.org/html/en/reference/polynomial_rings/sage/rings/polynomial/polynomial_rational_flint.html}}\par
{\footnotesize\url{https://doc.sagemath.org/html/en/reference/number_fields/sage/rings/number_field/number_field.html}}\par
{\footnotesize\url{https://doc.sagemath.org/html/en/reference/number_fields/sage/rings/number_field/galois_group.html}}\par
\vspace{9pt}
\needspace{12\baselineskip}\section{Multivariate absolute irreducibility}
\needspace{7\baselineskip}\hypertarget{ref:Ruppert1986}{}
{\sffamily\bfseries\color{atlasgreen}[Ruppert1986] Reduzibilität ebener Kurven}\par
\color{black} Ruppert, Wolfgang M. 1986. Journal für die reine und angewandte Mathematik. 369. pp. 167-191.\par
{\small\sffamily\color{atlasgold}Published work / Published article}\par
\color{black} A classical source on reducibility of plane curves and algebraic criteria for absolute irreducibility. This is a multivariate question: a univariate polynomial of degree greater than one always splits over an algebraic closure.\par
{\small Bibliographic details checked in Lecerf's survey and an AMS bibliography; the publisher full text was not retrieved.}\par
{\footnotesize\url{https://doi.org/10.1515/crll.1986.369.167}}\par
{\footnotesize\url{https://www.numdam.org/item/CCIRM_2013__3_1_A2_0/}}\par
{\footnotesize\url{https://doi.org/10.1090/MCOM/3007}}\par
\vspace{9pt}
\needspace{7\baselineskip}\hypertarget{ref:Kaltofen1995}{}
{\sffamily\bfseries\color{atlasgreen}[Kaltofen1995] Effective Noether irreducibility forms and applications}\par
\color{black} Kaltofen, Erich. 1995. Journal of Computer and System Sciences. 50(2). pp. 274-295.\par
{\small\sffamily\color{atlasgold}Published work / Published article}\par
\color{black} Gives effective bounds for algebraic irreducibility criteria and applications to specialization and factorization complexity. Useful when one wants quantitative control of exceptional specializations in several variables.\par
{\footnotesize\url{https://www.sciencedirect.com/science/article/pii/S0022000085710239}}\par
{\footnotesize\url{https://users.cs.duke.edu/~elk27/bibliography/95/Ka95_jcss.pdf}}\par
{\footnotesize\url{https://doi.org/10.1006/jcss.1995.1023}}\par
\vspace{9pt}
\needspace{7\baselineskip}\hypertarget{ref:Gao2001Polytopes}{}
{\sffamily\bfseries\color{atlasgreen}[Gao2001Polytopes] Absolute irreducibility of polynomials via Newton polytopes}\par
\color{black} Gao, Shuhong. 2001. Journal of Algebra. 237(2). pp. 501-520.\par
{\small\sffamily\color{atlasgold}Published work / Published article}\par
\color{black} Uses integral indecomposability of Newton polytopes to obtain absolute irreducibility criteria for sparse multivariate polynomials. Monomial factors must be excluded when applying a polytope criterion to a polynomial.\par
{\footnotesize\url{https://www.sciencedirect.com/science/article/pii/S0021869300985867}}\par
{\footnotesize\url{https://doi.org/10.1006/jabr.2000.8586}}\par
{\footnotesize\url{https://www.math.clemson.edu/~sgao/pub.html}}\par
\vspace{9pt}
\needspace{7\baselineskip}\hypertarget{ref:Gao2003PDE}{}
{\sffamily\bfseries\color{atlasgreen}[Gao2003PDE] Factoring multivariate polynomials via partial differential equations}\par
\color{black} Gao, Shuhong. 2003. Mathematics of Computation. 72(242). pp. 801-822.\par
{\small\sffamily\color{atlasgold}Published work / Published article}\par
\color{black} A factorization method based on a differential equation and linear algebra, in characteristic zero or sufficiently large characteristic. Connects rational and absolute factorization and provides an effective reduction from multivariate to bivariate problems.\par
{\footnotesize\url{https://doi.org/10.1090/S0025-5718-02-01428-X}}\par
{\footnotesize\url{https://www.math.clemson.edu/~sgao/papers/fac_bipoly.pdf}}\par
{\footnotesize\url{https://www.math.clemson.edu/~sgao/WEB/research.html}}\par
{\footnotesize\url{https://dblp.org/rec/journals/moc/Gao03.html}}\par
\vspace{9pt}
\needspace{7\baselineskip}\hypertarget{ref:Lecerf2013}{}
{\sffamily\bfseries\color{atlasgreen}[Lecerf2013] Factorisation des polynômes à plusieurs variables}\par
\color{black} Lecerf, Grégoire. 2013. Journées Nationales de Calcul Formel (13-17 May 2013). pp. 1-85. CIRM. Les cours du CIRM.\par
{\small\sffamily\color{atlasgold}Published work / Published lecture notes}\par
\color{black} An extensive French-language course on multivariate factorization, with a substantial bibliography. A good continuation when the atlas expands from univariate arithmetic to absolute factorization and computational algebraic geometry.\par
{\small Journées Nationales de Calcul Formel, 13-17 May 2013, talk 2.}\par
{\footnotesize\url{https://www.numdam.org/articles/10.5802/ccirm.18/}}\par
{\footnotesize\url{https://doi.org/10.5802/ccirm.18}}\par
{\footnotesize\url{https://www.numdam.org/item/CCIRM_2013__3_1_A2_0/}}\par
\vspace{9pt}
\needspace{7\baselineskip}\hypertarget{ref:GaoLauder2001}{}
{\sffamily\bfseries\color{atlasgreen}[GaoLauder2001] Decomposition of Polytopes and Polynomials}\par
\color{black} Gao, Shuhong and Lauder, Alan G. B. 2001. Discrete \& Computational Geometry. 26(1). pp. 89-104.\par
{\small\sffamily\color{atlasgold}Published work / Published article}\par
\color{black} Connects Minkowski decompositions of lattice polytopes with possible polynomial factors, gives algorithms for polygon decomposition, and develops projection-based tests in higher dimension. Corollary2 states the absolute-irreducibility certificate with the necessary exclusion of monomial factors. A decomposable polytope is inconclusive.\par
{\footnotesize\url{https://link.springer.com/article/10.1007/s00454-001-0024-0}}\par
{\footnotesize\url{https://arxiv.org/html/math/0012099v1}}\par
{\footnotesize\url{https://doi.org/10.1007/s00454-001-0024-0}}\par
{\footnotesize\url{https://arxiv.org/abs/math/0012099}}\par
{\footnotesize\url{https://arxiv.org/pdf/math/0012099}}\par
\vspace{9pt}
\needspace{12\baselineskip}\section{Free expositions and course materials}
\needspace{7\baselineskip}\hypertarget{ref:FilasetaCourse2016}{}
{\sffamily\bfseries\color{atlasgreen}[FilasetaCourse2016] Math 788F: The Theory of Irreducible Polynomials}\par
\color{black} Filaseta, Michael. 2016. University of South Carolina course materials.\par
{\small\sffamily\color{atlasgold}Exposition / Expository notes}\par
\color{black} A focused course hub with a syllabus, problems, and a Newton polygon introduction. It is a useful model for learning by proof techniques; the page does not currently provide every chapter of the proposed book as an open link.\par
{\footnotesize\url{https://people.math.sc.edu/filaseta/gradcourses/Math788F/Math788F.html}}\par
\vspace{9pt}
\needspace{7\baselineskip}\hypertarget{ref:ConradAllPrimes}{}
{\sffamily\bfseries\color{atlasgreen}[ConradAllPrimes] An irreducible that factors modulo all primes}\par
\color{black} Conrad, Keith. Undated. Author-hosted expository note.\par
{\small\sffamily\color{atlasgold}Exposition / Expository notes}\par
\color{black} Explains why irreducibility over the rationals need not be detectable by one irreducible reduction. The obstruction is visible in the cycle types available in the Galois group.\par
{\small Undated; accessed 11 November 2026.}\par
{\footnotesize\url{https://kconrad.math.uconn.edu/blurbs/ringtheory/reducibleallp.pdf}}\par
{\footnotesize\url{https://kconrad.math.uconn.edu/blurbs/}}\par
\vspace{9pt}
\needspace{7\baselineskip}\hypertarget{ref:ConradCubicQuartic}{}
{\sffamily\bfseries\color{atlasgreen}[ConradCubicQuartic] Galois groups of cubics and quartics (not in characteristic 2)}\par
\color{black} Conrad, Keith. Undated. Author-hosted expository note.\par
{\small\sffamily\color{atlasgold}Exposition / Expository notes}\par
\color{black} Worked low-degree classifications using discriminants and resolvents. The characteristic restriction matters; the author's collection also contains a separate all-characteristics treatment.\par
{\small Undated; accessed 11 November 2026.}\par
{\footnotesize\url{https://kconrad.math.uconn.edu/blurbs/galoistheory/cubicquartic.pdf}}\par
{\footnotesize\url{https://kconrad.math.uconn.edu/blurbs/}}\par
\vspace{9pt}
\needspace{7\baselineskip}\hypertarget{ref:ConradFiniteTests}{}
{\sffamily\bfseries\color{atlasgreen}[ConradFiniteTests] Irreducibility tests in F\_\allowbreak{}p[T]}\par
\color{black} Conrad, Keith. Undated. Author-hosted expository note.\par
{\small\sffamily\color{atlasgold}Exposition / Expository notes}\par
\color{black} An accessible account of irreducibility tests over finite fields. Useful for implementing transparent modular certificates rather than reporting an unexplained software answer.\par
{\small Undated; accessed 11 November 2026.}\par
{\footnotesize\url{https://kconrad.math.uconn.edu/blurbs/ugradnumthy/irredtestFpT.pdf}}\par
{\footnotesize\url{https://kconrad.math.uconn.edu/blurbs/}}\par
\vspace{9pt}
\needspace{7\baselineskip}\hypertarget{ref:ConradIrredQ}{}
{\sffamily\bfseries\color{atlasgreen}[ConradIrredQ] Irreducibility tests in Q[T]}\par
\color{black} Conrad, Keith. Undated. Author-hosted expository note.\par
{\small\sffamily\color{atlasgold}Exposition / Expository notes}\par
\color{black} Clear proofs of Gauss's lemma, reduction modulo a prime, and Eisenstein's criterion. Good first reading before more specialized irreducibility criteria.\par
{\small Undated; accessed 11 November 2026.}\par
{\footnotesize\url{https://kconrad.math.uconn.edu/blurbs/ringtheory/irredtestsoverQ.pdf}}\par
{\footnotesize\url{https://kconrad.math.uconn.edu/blurbs/}}\par
\vspace{9pt}
\needspace{7\baselineskip}\hypertarget{ref:ConradModP}{}
{\sffamily\bfseries\color{atlasgreen}[ConradModP] Galois groups over Q and factorizations mod p}\par
\color{black} Conrad, Keith. Undated. Author-hosted expository note.\par
{\small\sffamily\color{atlasgold}Exposition / Expository notes}\par
\color{black} A short explanation of how squarefree factorizations modulo good primes produce cycle types in a Galois group. Read the discriminant condition carefully.\par
{\small Undated; accessed 11 November 2026.}\par
{\footnotesize\url{https://kconrad.math.uconn.edu/blurbs/gradnumthy/galois-Q-factor-mod-p.pdf}}\par
{\footnotesize\url{https://kconrad.math.uconn.edu/blurbs/}}\par
\vspace{9pt}
\needspace{7\baselineskip}\hypertarget{ref:ConradSnAn}{}
{\sffamily\bfseries\color{atlasgreen}[ConradSnAn] Recognizing Galois groups S\_\allowbreak{}n and A\_\allowbreak{}n}\par
\color{black} Conrad, Keith. Undated. Author-hosted expository note.\par
{\small\sffamily\color{atlasgold}Exposition / Expository notes}\par
\color{black} Useful group-theoretic recognition criteria and examples. A practical bridge from modular factorization and real roots to identifying the full permutation group.\par
{\small Undated; accessed 11 November 2026.}\par
{\footnotesize\url{https://kconrad.math.uconn.edu/blurbs/galoistheory/galoisSnAn.pdf}}\par
{\footnotesize\url{https://kconrad.math.uconn.edu/blurbs/}}\par
\vspace{9pt}
\end{document}
